% FINAL SINGLE-FILE SUBMISSION (2026-07-18).
% CROSS-REFERENCE FIX v5: typed labels + post-aliascnt cleveref loading
% + explicit reference text for the two audited shared-counter labels.
% The aliascnt declarations below are required so cleveref preserves the
% displayed environment type when theorem-like environments share numbering.
\documentclass[11pt]{article}

\usepackage[T1]{fontenc}
\usepackage[utf8]{inputenc}
\usepackage{lmodern}
\usepackage[nopatch=eqnum]{microtype}
\usepackage[a4paper,margin=1in]{geometry}
\usepackage{amsmath,amssymb,amsthm,mathtools,mathrsfs}
\usepackage{booktabs,array}
\usepackage{enumitem}
\usepackage{xurl}
\usepackage{authblk}
\usepackage{xcolor}
\usepackage{aliascnt}
\usepackage[
  colorlinks=true,
  linkcolor=blue!55!black,
  citecolor=blue!55!black,
  urlcolor=blue!55!black,
  unicode=true,
  pdftitle={A Proposed Solution to Erdős Problem 390},
  pdfauthor={Shouqiao Wang}
]{hyperref}

\numberwithin{equation}{section}
\makeatletter
\expandafter\let\expandafter\MT@orig@extended@maketag
  \csname MT_extended_maketag:n\endcsname
\expandafter\let\expandafter\MT@orig@extended@tagform
  \csname MT_extended_tagform:n\endcsname
\expandafter\def\csname MT_extended_tagform:n\endcsname#1{\let\MT@saved@tag@label\df@label
  \ifx\df@label\@empty\else
    \edef\df@label{\expandafter\@firstofone\df@label}\fi
  \MT@orig@extended@tagform{#1}\let\df@label\MT@saved@tag@label
}
\expandafter\def\csname MT_extended_maketag:n\endcsname#1{\let\MT@saved@make@label\df@label
  \ifx\df@label\@empty\else
    \edef\df@label{\expandafter\@firstofone\df@label}\fi
  \MT@orig@extended@maketag{#1}\let\df@label\MT@saved@make@label
}
\makeatother
\mathtoolsset{showonlyrefs}

\newtheorem{theorem}{Theorem}[section]
\newaliascnt{lemma}{theorem}
\newtheorem{lemma}[lemma]{Lemma}
\aliascntresetthe{lemma}
\newaliascnt{proposition}{theorem}
\newtheorem{proposition}[proposition]{Proposition}
\aliascntresetthe{proposition}
\newaliascnt{corollary}{theorem}
\newtheorem{corollary}[corollary]{Corollary}
\aliascntresetthe{corollary}
\newaliascnt{claim}{theorem}
\newtheorem{claim}[claim]{Claim}
\aliascntresetthe{claim}
\theoremstyle{definition}
\newaliascnt{definition}{theorem}
\newtheorem{definition}[definition]{Definition}
\aliascntresetthe{definition}
\theoremstyle{remark}
\newaliascnt{remark}{theorem}
\newtheorem{remark}[remark]{Remark}
\aliascntresetthe{remark}

% Load cleveref only after all alias counters have been declared.  This is
% the order recommended by aliascnt and keeps the reference type stable on
% both older and current TeX Live releases.
\usepackage[nameinlink,capitalise,noabbrev]{cleveref}

\crefname{theorem}{theorem}{theorems}
\Crefname{theorem}{Theorem}{Theorems}
\crefname{lemma}{lemma}{lemmas}
\Crefname{lemma}{Lemma}{Lemmas}
\crefname{proposition}{proposition}{propositions}
\Crefname{proposition}{Proposition}{Propositions}
\crefname{corollary}{corollary}{corollaries}
\Crefname{corollary}{Corollary}{Corollaries}
\crefname{claim}{claim}{claims}
\Crefname{claim}{Claim}{Claims}
\crefname{definition}{definition}{definitions}
\Crefname{definition}{Definition}{Definitions}
\crefname{remark}{remark}{remarks}
\Crefname{remark}{Remark}{Remarks}
% Do not rely on cleveref inferring an alias-counter's type.  In particular,
% these declarations prevent references to Proposition 7.1 or Lemma 5.2 from
% being rendered as ``Theorem 7.1'' or ``Theorem 5.2'' on older installations.
\crefalias{lemma}{lemma}
\crefalias{proposition}{proposition}
\crefalias{corollary}{corollary}
\crefalias{claim}{claim}
\crefalias{definition}{definition}
\crefalias{remark}{remark}

\newcommand{\N}{\mathbb N}
\newcommand{\R}{\mathbb R}
\newcommand{\1}{\mathbf 1}
\newcommand{\vp}[1]{v_{#1}}
\newcommand{\Pplus}{P^{+}}
\newcommand{\Pminus}{P^{-}}
\newcommand{\e}{\mathrm e}
\newcommand{\dd}{\,\mathrm d}
\newcommand{\eps}{\varepsilon}
\DeclareMathOperator{\supp}{supp}
\DeclareMathOperator{\Var}{Var}
\DeclareMathOperator{\Cov}{Cov}
\DeclareMathOperator{\li}{li}

\title{A Proposed Solution to Erd\H{o}s Problem 390}
\author{Shouqiao Wang}
\affil{Columbia University\hspace{2.5em}Multiscalar Intelligence}
\date{}

\begin{document}
\allowdisplaybreaks
\maketitle

\begin{abstract}
Let \(f(n)\) be the least possible largest factor in a representation of
\(n!\) as a product of distinct integers, all larger than \(n\).  We
propose the exact asymptotic
\[
 \lim_{n\to\infty}\frac{(f(n)-2n)\log n}{n}
 =\frac{4029639598}{25970038185}.
\]
The lower bound is a thirteen-layer valuation cut.  The upper bound uses an
exact cofactor-allocation certificate, central-binomial anchors, a guarded
rough-signature selector, a marked friable-number covariance bridge, a
finite-band tangent correction, and deterministic column-sparse rounding
with a precharged universal switch bank.
The complete finite allocation certificate and its exact checker are
supplied in the companion file \texttt{numerical\_verifier.py}.
This proposed solution was found by GPT-5.6.
\end{abstract}

\tableofcontents

\section{Introduction}
\label{sec:introduction}

For an integer \(n\geq3\), let \(f(n)\) be the least possible value of
the largest factor in a factorization
\[
        n!=a_1a_2\cdots a_k,
        \qquad n<a_1<a_2<\cdots<a_k.
\]
Thus the factors are required both to be distinct and to lie strictly above
\(n\).  Erd\H{o}s, Guy and Selfridge proved that there are absolute constants
\(0<c_1<c_2\) for which
\[
  2n+c_1\frac{n}{\log n}<f(n)<
  2n+c_2\frac{n}{\log n}
\]
for all sufficiently large \(n\), and they asked whether the second-order
term has an asymptotic constant \cite[Theorem~3]{ErdosGuySelfridge1982}.

Subsequent discussion isolated two features that are useful here.  Tao
emphasized the complementary product formulation below and observed that
the original lower-bound count misses further carry intervals
\cite{Tao390Comment}.  Mausberg made the first thirteen of these intervals
quantitative and obtained the lower-bound constant displayed in
\eqref{eq:C0-definition} \cite{Mausberg2026}.  The lower-bound argument in
this paper is based on that thirteen-layer obstruction, with an additional
reduction that treats endpoints at or below \(2n\).  We make no claim here
about priority beyond these cited records.
For a recent treatment of the related, but different, problem of
decomposing a factorial into a prescribed number of large factors, see
Alexeev et al.\ \cite{AlexeevEtAl2026}; none of its quantitative estimates
is used below.

Define
\begin{equation}
 \label{eq:C0-definition}
 C_0
 :=
 \frac{\displaystyle\sum_{r=1}^{13}
          \frac{1}{(r+1)(2r+1)}}
      {\displaystyle\sum_{\substack{p\leq 23\\p\ {\rm prime}}}
          \frac{1}{p-1}}
 =\frac{4029639598}{25970038185}
 =0.15516494697830188\ldots.
\end{equation}
Our main result determines the conjectured constant.

\begin{theorem}
\label[theorem]{thm:main}
As \(n\to\infty\),
\[
       f(n)=2n+C_0\frac{n}{\log n}
                 +o\!\left(\frac{n}{\log n}\right).
\]
Equivalently,
\[
       \lim_{n\to\infty}
       \frac{(f(n)-2n)\log n}{n}=C_0.
\]
\end{theorem}

The proof is most naturally expressed after taking complements.  For an
integer \(M>n\), put
\[
              Q(n,M):=\frac{M!}{(n!)^2},
       \qquad \mathcal I(n,M):=(n,M]\cap\mathbb Z.
\]

\begin{lemma}[Complement formulation]
\label[lemma]{lem:complement-formulation}
There is a factorization of \(n!\) into distinct members of
\(\mathcal I(n,M)\) if and only if \(Q(n,M)\) is the product of a subset of
the distinct integers in \(\mathcal I(n,M)\).
\end{lemma}

\begin{proof}
The product of every integer in \(\mathcal I(n,M)\) is \(M!/n!\).  If one
subset has product \(n!\), its complement therefore has product
\((M!/n!)/n!=Q(n,M)\).  The same calculation in reverse proves the
converse.
\end{proof}

We briefly describe the proof of Theorem~\ref{thm:main}.  The lower bound
uses thirteen disjoint layers of primes \(P\) for which
\(v_P(Q(n,M))=1\).  The unique selected multiple of such a prime has the
form \(Pq\), with \(r+1\leq q\leq 2r+1\) in layer \(r\).  Every cofactor
arising in the first thirteen layers has a prime divisor at most \(23\).
Comparing these forced incidences with the available small-prime valuations
gives exactly the quotient in \eqref{eq:C0-definition}.

For the upper bound, fix \(c>C_0\), set
\[
       L=\log n,\qquad N=\frac nL,\qquad
       M=2n+\lceil cN\rceil,
\]
and construct a distinct subset of \((n,M]\) with product \(Q(n,M)\).
The construction has five interacting parts.
First, an exact rational allocation of the large-prime carry layers produces
central anchors; its finite portion is certified by rational arithmetic and
its infinite tail by Nagura's prime-interval theorem.  Second, all factors
needed by the integral repair are reserved in advance, so that their
valuation cost is present in the target ledger before any equation is
solved.  Third, with \(y=n^{2/9}\), a signed selector matches every complete
\(y\)-rough signature exactly and leaves a residual
\(O(N/(pL))\) at each prime \(p\in(W,y]\).  Fourth, marked smooth-number
estimates and the Poisson--Dickman covariance operator fit the smooth row,
the fixed head valuations, the physical logarithm and finitely many exponent
bands.  Finally, local same-row switches cancel the remaining medium-prime
residual, and a precharged universal bank converts the fractional point to a
\(0\)-\(1\) point without changing any exact row quota.  The resulting exact
valuation identity is the product identity required by
Lemma~\ref{lem:complement-formulation}.

The role of the computer-assisted component is deliberately narrow.  It
checks a finite list of rational equalities and inequalities for the anchor
allocation.  The complete finite certificate and its exact checker are the
single companion file \texttt{numerical\_verifier.py};
Appendix~\ref{app:numerical-verification} records its precise scope.  The
analytic tail argument is proved in Section~\ref{sec:infinite-cofactor-allocation},
and no floating-point output is used as a premise of the proof.
 \section{Analytic and probabilistic inputs}
\label{sec:analytic-tools}

This section records the external results used later, together with the
specializations needed in the proof.  All limiting parameters appearing
below are fixed before \(n\to\infty\).  In particular, constants may depend
on fixed compact intervals, the fixed head cutoff and a fixed finite cell
partition, but never on \(n\) or on a marked prime.

\subsection{Smooth numbers and the de Bruijn--Saias normalization}
\label{subsec:smooth-input}

Write \(P^+(m)\) and \(P^-(m)\) for the largest and smallest prime factors
of \(m\), with \(P^+(1)=1\), and put
\[
       \Psi(x,y):=\#\{m\leq x:P^+(m)\leq y\},
       \qquad u:=\frac{\log x}{\log y}.
\]
The Dickman function is defined by
\[
\begin{aligned}
 \rho(u)&=0 &&(u<0),\\
 \rho(u)&=1 &&(0\leq u\leq1),\\
 u\rho'(u)+\rho(u-1)&=0 &&(u>1).
\end{aligned}
\]
At transition points it is important not to replace the smooth-number count
prematurely by \(x\rho(u)\).  We use the de Bruijn--Saias normalization
\begin{equation}
 \label{eq:Lambda-definition}
 \Lambda(x,y)
 :=x\int_{-\infty}^{\infty}
       \rho(u-v)\,d\!\left(\lfloor y^v\rfloor y^{-v}\right),
\end{equation}
where at integral \(x\) the value means the right limit
\(\Lambda(x+0,y)\).  This one-sided convention is part of the definition in
the cited source.

\begin{theorem}[Hildebrand--Tenenbaum--Saias]
\label[theorem]{thm:HT-Saias}
For every \(\varepsilon>0\), uniformly in
\[
 y\geq y_0(\varepsilon),
 \qquad
 1\leq u\leq
 \exp\{(\log y)^{3/5-\varepsilon}\},
\]
one has
\begin{equation}
 \label{eq:HT-Saias}
 \Psi(x,y)=\Lambda(x,y)
 \left\{1+O_\varepsilon\!\left(
       e^{-(\log y)^{3/5-\varepsilon}}
                         \right)\right\}.
\end{equation}
On every fixed compact \(u\)-range contained in \([1,\infty)\),
\begin{equation}
 \label{eq:Lambda-Dickman}
 \Lambda(x,y)=x\rho(u)
       \left\{1+O\!\left(\frac{\log(u+1)}{\log y}\right)\right\}.
\end{equation}
For \(0<u\leq1\), the exact identity
\(\Psi(x,y)=\lfloor x\rfloor\) is used instead.
\end{theorem}

The formulation is Theorem~1.8, (1.29) and Lemma~3.1 of
Hildebrand--Tenenbaum \cite{HildebrandTenenbaum1993}; the underlying sharp
estimate is due to Saias \cite{Saias1989}.  In our application
\(y=n^{2/9}\to\infty\), while all values of \(u\) remain in a fixed compact
set.  Thus both range conditions in Theorem~\ref{thm:HT-Saias} hold with
room to spare, and the exponential error is smaller than every fixed
negative power of \(\log n\).  Formula \eqref{eq:HT-Saias}, rather than only
\eqref{eq:Lambda-Dickman}, will be retained at terminal and Dickman
transition points.

We shall repeatedly use the following derived marked form.  It is included
here to specify exactly what is, and what is not, obtained from the cited
smooth-number theorem.

\begin{lemma}[Uniform marked smooth cells]
\label[lemma]{lem:marked-smooth-cells}
Let \(y=n^{2/9}\), so that
\(U=\log n/\log y=9/2\), and put \(L=\log n\).  Fix a finite set \(H\) of
primes, a vector \({\bf e}=(e_\ell)_{\ell\in H}\) of nonnegative integers,
and fixed reals \(0<A<B<\infty\).  Define
\[
 h_{\bf e}:=\prod_{\ell\in H}\ell^{e_\ell},
 \qquad M_H:=\prod_{\ell\in H}\ell,
 \qquad
 \delta_{\bf e}:=\frac1{h_{\bf e}}
                  \frac{\varphi(M_H)}{M_H}.
\]
If \(d\leq y^4\) is \(y\)-smooth and \((d,M_H)=1\), then, uniformly in
\(d\),
\begin{align}
 &\#\{m:An<m\leq Bn,\ P^+(m)\leq y,\ d\mid m,
             \ v_\ell(m)=e_\ell\ (\ell\in H)\}
 \notag\\
 &\hspace{25mm}
 =\delta_{\bf e}\frac{(B-A)n}{d}
   \left\{\rho\!\left(U-\frac{\log d}{\log y}\right)
                  +O_{H,A,B}\!\left(\frac1L\right)\right\}.
 \label{eq:marked-smooth-count}
\end{align}
If \(\phi\) has bounded variation on \([A,B]\), the corresponding
Stieltjes-weighted form is
\begin{align}
 &\sum_{\substack{An<m\leq Bn,\ P^+(m)\leq y,\ d\mid m\\
                   v_\ell(m)=e_\ell\ (\ell\in H)}}
        \phi(m/n)
 \notag\\
 &\quad=
 \delta_{\bf e}\frac nd
 \rho\!\left(U-\frac{\log d}{\log y}\right)
 \int_A^B\phi(t)\,dt
 +O_{H,A,B}\!\left(
       \frac{n}{dL}\{\|\phi\|_\infty+\operatorname{Var}\phi\}
                    \right).
 \label{eq:marked-smooth-Stieltjes}
\end{align}
Consequently, under normalized counting measure on the unmarked cell,
\begin{equation}
 \label{eq:marked-smooth-probability}
 \mathbb P(d\mid m)
 =\frac1d
   \frac{\rho(U-\log d/\log y)}{\rho(U)}
   +O_{H,A,B}\!\left(\frac1{dL}\right).
\end{equation}
The conclusions persist under a fixed finite disjoint union, or a
compactly weighted fixed finite mixture, of such cells.  An overlapping
finite family is always first disjointified and then interpreted through
the induced convex mixture.
\end{lemma}

\begin{proof}
The exact valuation conditions are finite divisibility
inclusion--exclusion:
\[
 1_{\{v_\ell(m)=e_\ell\ (\ell\in H)\}}
 =1_{h_{\bf e}\mid m}
   \sum_{a\mid M_H}\mu(a)
       1_{a\mid m/h_{\bf e}}.
\]
After also dividing by \(d\), each summand is a difference of two values of
\(\Psi(\,\cdot\,,y)\).  The fixed factors \(A,B,h_{\bf e},a\) alter the
smoothness parameter by \(O_H(1/L)\), and \(\rho\) is Lipschitz on the
resulting compact range.  Theorem~\ref{thm:HT-Saias}, with the exact formula
below parameter one, therefore gives \eqref{eq:marked-smooth-count} after
summing
\[
       \frac1{h_{\bf e}}
       \sum_{a\mid M_H}\frac{\mu(a)}a
       =\delta_{\bf e}.
\]
The endpoint floor error is \(O(1)\).  Since
\(n/d\geq n/y^4=n^{1/9}\), it is absorbed by \(O(n/(dL))\).
Stieltjes partial summation proves
\eqref{eq:marked-smooth-Stieltjes}; division by the case \(d=1\) proves
\eqref{eq:marked-smooth-probability}.
\end{proof}

Only exact valuation and coprimality patterns, and finite disjoint unions
of them (after disjointification when necessary), will be called
\emph{head cells}.  Lemma~\ref{lem:marked-smooth-cells} makes
no assertion about arbitrary nonzero residue classes modulo a fixed integer;
such a statement would require a smooth-number theorem in arithmetic
progressions.  The bound \(d\leq y^4\) covers the marks \(p\), \(pq\),
\(p^k\), and \(p^kq\) used below and leaves the genuine cofactor margin
\(n/y^4=n^{1/9}\).

\subsection{Prime counts and Mertens estimates}
\label{subsec:pnt-mertens}

We use the classical zero-free-region prime number theorem in the form
\begin{equation}
 \label{eq:cumulative-PNT}
 \pi(x)=\operatorname{li}(x)
       +O\!\left(xe^{-c\sqrt{\log x}}\right)
       \qquad(x\geq2)
\end{equation}
for an absolute \(c>0\); see
\cite[Theorem~6.9]{MontgomeryVaughan2007}.  In particular, if
\(0<a<b\) are fixed and
\(n^a\leq X\leq Z\leq n^b\), subtraction at the two exact endpoints gives
\begin{equation}
 \label{eq:PNT-difference}
 \pi(Z)-\pi(X)
 =\int_X^Z\frac{dt}{\log t}
   +O_{a,b}\!\left(Ze^{-c_{a,b}\sqrt{\log n}}\right).
\end{equation}
The error in \eqref{eq:PNT-difference} is absolute.  We do not invoke a
relative prime number theorem for arbitrary short intervals.  Whenever a
relative conclusion is used later, the displayed main term is separately
checked to dominate this absolute error.

Mertens' estimates, in the forms used here, are
\begin{align}
 \sum_{p\leq x}\frac1p
   &=\log\log x+B_1+O\!\left(\frac1{\log x}\right),
 \label{eq:Mertens-harmonic}\\
 \prod_{p\leq x}\left(1-\frac1p\right)^{-1}
   &=e^\gamma\log x+O(1).
 \label{eq:Mertens-product}
\end{align}
See \cite{Mertens1874} or
\cite[Theorem~2.7(d),(e)]{MontgomeryVaughan2007}.  Thus, for fixed
\(0<a<b\),
\begin{equation}
 \label{eq:Mertens-fixed-range}
 \sum_{y^a<p\leq y^b}\frac1p
 =\log\frac ba+O_{a,b}\!\left(\frac1{\log y}\right),
\end{equation}
and
\begin{equation}
 \label{eq:Mertens-sieve-density}
 \prod_{p\leq y}\left(1-\frac1p\right)
 =\frac{e^{-\gamma}}{\log y}
  +O\!\left(\frac1{\log^2y}\right).
\end{equation}
These formulas justify every bounded reciprocal-prime mass used in a
fixed-depth expansion.  They also give, for \(X\geq y\),
\[
 \sum_{\substack{R\leq X\\P^-(R)>y}}\frac1R
 \leq\prod_{y<p\leq X}\left(1-\frac1p\right)^{-1}
 \ll\frac{\log X}{\log y}.
\]
For \(y=n^{2/9}\) and \(X\leq n^C\), this is \(O_C(1)\).

\subsection{The interval Selberg sieve}
\label{subsec:selberg-sieve}

For a positive integer \(P\), define
\[
 S(x,H;P):=\#\{m:x<m\leq x+H,\ (m,P)=1\},
 \qquad
 L_P(R):=\sum_{\substack{d\leq R\\d\mid P}}
             \frac{\mu(d)^2}{\varphi(d)}.
\]
The one-dimensional Selberg upper-bound sieve gives, uniformly in the real
translate \(x\),
\begin{equation}
 \label{eq:Selberg-general}
 S(x,H;P)
 \leq \frac{H}{L_P(R)}
      +O\!\left(\frac{R^2}{L_P(R)^2}\right)
 \qquad(H>0,\ R\geq1).
\end{equation}
This is \cite[Theorem~3.2]{MontgomeryVaughan2007}.  Indeed, for weights
supported on \(d,e\leq R\), the elementary interval remainder is
\[
 \#\{m:x<m\leq x+H,\ [d,e]\mid m\}
       =\frac H{[d,e]}+O(1).
\]

Take \(P=\prod_{p\leq y}p\) and choose Selberg weights supported on
\(d\leq R=y^2\).  The elementary lower bound
\[
 L_P(y^2)\geq\sum_{d\leq y}\frac{\mu(d)^2}{\varphi(d)}
 \geq\sum_{d\leq y}\frac{\mu(d)^2}{d}
 =\frac1{\zeta(2)}\log y+O(1)
 \gg\log y
\]
follows, for example, by inserting
\(\mu(d)^2=\sum_{a^2\mid d}\mu(a)\) in the last sum.  Consequently,
\begin{equation}
 \label{eq:Selberg-specialized}
 \#\{m:x<m\leq x+H:P^-(m)>y\}
 \ll\frac H{\log y}+\frac{y^4}{(\log y)^2}.
\end{equation}
The weight support is \(y^2\), whereas the largest lcm remainder modulus is
\([d,e]\leq y^4\); we keep these two levels distinct.  In the exceptional
row application, \(y=n^{2/9}\) and
\(H\gg n^{1-\delta_*}/\log n\), where \(\delta_*<1/18\).  Hence
\(y^4=n^{8/9}=o(H)\), even after logarithmic factors, and
\eqref{eq:Selberg-specialized} is \(O(H/\log y)\) uniformly in the
translated interval.

\subsection{The Poisson--Dickman bridge}
\label{subsec:Poisson-Dickman}

Let \(\Pi\) be a Poisson point process on \((0,1]\) with intensity
\(dt/t\), and set
\[
                         T:=\sum_{t\in\Pi}t.
\]
Arratia, Barbour and Tavar\'e proved that \(T\) has density
\begin{equation}
 \label{eq:Poisson-Dickman-density}
                         g(t)=e^{-\gamma}\rho(t)
                         \qquad(t>0)
\end{equation}
\cite[(2.3)--(2.6)]{ArratiaBarbourTavare1999}.  In particular,
\(\int_0^\infty\rho(t)\,dt=e^\gamma\).

We fix a canonical conditional version at every exact total.  For
\(0<\varepsilon<1\), write
\[
 \Pi_{>\varepsilon}:=\Pi\cap(\varepsilon,1],\qquad
 T_{>\varepsilon}:=\sum_{x\in\Pi_{>\varepsilon}}x,
\]
and let \(T_{\leq\varepsilon}\) be the independent mass of the remaining
atoms.  After scaling by \(\varepsilon\), the process on
\((0,\varepsilon]\) is again the scale-invariant process on \((0,1]\).
Consequently \(T_{\leq\varepsilon}\) has the continuous density
\[
                         g_\varepsilon(v)
          =\varepsilon^{-1}g(v/\varepsilon)\qquad(v>0).
\]
If \(G\) is a bounded cylinder function depending only on
\(\Pi_{>\varepsilon}\), define
\begin{equation}
 \mathbb E_uG
 :=\frac1{g(u)}\mathbb E\!\left[
   G(\Pi_{>\varepsilon})
   g_\varepsilon(u-T_{>\varepsilon})
   \mathbf1_{\{T_{>\varepsilon}<u\}}\right].
 \label{eq:canonical-PD-bridge}
\end{equation}
The convolution identity for the density of
\(T_{>\varepsilon}+T_{\leq\varepsilon}\) shows that this is a probability
law.  Here is the consistency calculation.  If
\(0<\varepsilon'<\varepsilon\) and \(G\) depends only on
\(\Pi_{>\varepsilon}\), condition first on \(\Pi_{>\varepsilon}\) in
the \(\varepsilon'\)-formula.  With \(V\) the independent mass on
\((\varepsilon',\varepsilon]\), the convolution identity is
\begin{equation}
 \mathbb E\!\left[g_{\varepsilon'}(v-V)
                   \1_{\{V<v\}}\right]=g_\varepsilon(v)
 \qquad(v>0).
 \label{eq:PD-bridge-consistency-convolution}
\end{equation}
Substitution of \(v=u-T_{>\varepsilon}\) shows that the two definitions
of \(\mathbb E_uG\) coincide.  The configuration space of locally finite
point measures on \((0,1]\), with the sigma-field generated by the
restrictions to \((\varepsilon,1]\), is standard Borel.  The consistent
finite restrictions therefore define a unique probability law
\(\mathbb P_u\) for every \(u>0\).
For each bounded cylinder function, the right-hand side of
\eqref{eq:canonical-PD-bridge} is a Borel function of \(u\).  A monotone
class argument, starting from the cylinder events, therefore shows that
\(u\mapsto\mathbb P_u(A)\) is Borel for every configuration event \(A\).
Thus \((\mathbb P_u)_{u>0}\) is a probability kernel, as required for the
disintegrations below.

This projective law is supported on configurations of total mass exactly
\(u\), rather than merely at most \(u\).  Indeed its restrictions give
\(T_{>\varepsilon}\leq u\), while the finite-\(\varepsilon\)
conditional Campbell calculation gives directly
\begin{equation}
 \mathbb E_u\{u-T_{>\varepsilon}\}
 =\int_0^\varepsilon\frac{\rho(u-s)}{\rho(u)}\,ds
 \leq C_u\varepsilon.
 \label{eq:PD-small-mass-bound}
\end{equation}
Thus \(u-T_{>\varepsilon}\to0\) in \(L^1(\mathbb P_u)\), and monotone
convergence yields \(\sum_{x\in\Pi}x=u\) almost surely.  This selects the
continuous regular conditional version used below, rather than leaving
its value at a probability-zero event arbitrary.

Let \(\mathbb E_U\) denote expectation under \(\mathbb P_U\).
Campbell--Mecke and its multivariate form
\cite[Theorems~4.1 and~4.4]{LastPenrose2017}, applied first to cylinder functions in
\eqref{eq:canonical-PD-bridge}, and then monotone convergence, gives after
division by the density in \eqref{eq:Poisson-Dickman-density}, for bounded
measurable test functions for which the integrals converge,
\begin{align}
 \mathbb E_U\sum_{t\in\Pi}F(t)
 &=\int_0^1 F(s)
     \frac{\rho(U-s)}{\rho(U)}\,\frac{ds}{s},
 \label{eq:Palm-one-point}\\
 \mathbb E_U\sum_{\substack{s,t\in\Pi\\s\ne t}}F(s)G(t)
 &=\int_0^1\!\int_0^1 F(s)G(t)
     \frac{\rho(U-s-t)}{\rho(U)}
                  \frac{ds}{s}\frac{dt}{t}.
 \label{eq:Palm-two-point}
\end{align}
Here \(\rho(v)=0\) for \(v<0\).  Since \(g(U)>0\) for every \(U>0\),
these density ratios define the required conditional version at the fixed
value \(U=9/2\).  In particular, with
\[
 h(s):=\frac{\rho(U-s)}{\rho(U)},
 \qquad K(s,t):=\frac{\rho(U-s-t)}{\rho(U)}-h(s)h(t),
\]
equations \eqref{eq:Palm-one-point}--\eqref{eq:Palm-two-point} give exactly
the multiplication and covariance kernels used by the smooth-row bridge.
The formulas hold first for nonnegative tests; a signed test follows by
positive and negative parts whenever the right side is absolutely
integrable.  In every later use \(F(t)=O(t)\), or a product of two such
functions, so \eqref{eq:PD-small-mass-bound} and
\(|K(s,t)|\ll st\) verify the required first and second integrability.

\subsection{An asymmetric local-lemma criterion}
\label{subsec:LLL}

We use the asymmetric Lov\'asz local lemma in the following standard form
\cite[Lemma~5.1.1]{AlonSpencer2016}.  If events \((A_i)\) have a dependency
digraph, and numbers \(0\leq x_i<1\) satisfy
\begin{equation}
 \label{eq:asymmetric-LLL}
 \mathbb P(A_i)\leq x_i
       \prod_{j\in\Gamma(i)}(1-x_j)
\end{equation}
for every \(i\), then
\(\mathbb P(\bigcap_i\overline{A_i})>0\).

The collision events later in the proof depend on at most two independently
sampled request variables.  Suppose that, for every one request, the sum of
the probabilities of collision events involving it is at most \(B\).  A
collision event \(E\) then has
\[
       \sum_{E'\sim E}\mathbb P(E')\leq2B.
\]
Set \(x_E=2\mathbb P(E)\).  If \(B\leq1/8\), then
\[
       \sum_{E'\sim E}x_{E'}\leq4B\leq\frac12
\]
and, using \(\prod_i(1-z_i)\geq1-\sum_i z_i\),
\[
 x_E\prod_{E'\sim E}(1-x_{E'})
 \geq2\mathbb P(E)
       \left(1-\sum_{E'\sim E}x_{E'}\right)
 \geq\mathbb P(E).
\]
Thus \eqref{eq:asymmetric-LLL} holds.  The factor \(4\) between the
per-request probability bound and the dependency-neighborhood \(x\)-mass is
included explicitly here.

\subsection{Nagura's prime interval}
\label{subsec:Nagura}

Nagura proved that for every real \(x>25\) there is a prime \(p\) with
\begin{equation}
 \label{eq:Nagura}
                         x<p<\frac65x
\end{equation}
\cite{Nagura1952}.  Consequently, if \(p_-<p\) are consecutive primes and
\(p>401\), then \(p<(6/5)p_-\), and
\begin{equation}
 \label{eq:Nagura-tail}
 (p-1)\sum_{r=p_-}^{p-1}
       \frac1{(r+1)(2r+1)}
 <\frac{p(p-p_-)}{2p_-^2}<\frac3{25}.
\end{equation}
The finite allocation certificate is checked through \(401\); hence
\eqref{eq:Nagura-tail} is invoked only where the hypothesis of
\eqref{eq:Nagura} is automatic.
 \section{The thirteen-layer lower bound}
\label{sec:lower-bound}

We begin with the lower bound.  This is Mausberg's thirteen-layer
valuation argument~\cite{Mausberg2026}, with a direct preliminary lemma
that rules out all endpoints at or below \(2n\).  Put
\[
        N:=\frac{n}{\log n},
        \qquad
        \mathcal P_{\mathrm{sm}}
        :=\{2,3,5,7,11,13,17,19,23\},
\]
and define
\begin{equation}
 \begin{split}
 A_{13}
   &:=\sum_{r=1}^{13}\frac{1}{(r+1)(2r+1)}
     =\frac{2014819799}{5736673800},\\
 S_{23}
   &:=\sum_{\ell\in\mathcal P_{\mathrm{sm}}}
               \frac{1}{\ell-1}
     =\frac{17927}{7920},\\
 C_0&:=\frac{A_{13}}{S_{23}}
     =\frac{4029639598}{25970038185}.
 \end{split}
 \label{eq:lower-constant}
\end{equation}

For an integer \(M>n\), set
\begin{equation}
        Q(n,M):=\frac{M!}{(n!)^2}.
        \label{eq:complement-quotient}
\end{equation}
We call \(M\) \emph{admissible} if \(Q(n,M)\) is the product of a
subset of the distinct integers in \((n,M]\).  This terminology is
equivalent to \(f(n)\le M\).  Indeed,
\[
       \prod_{n<a\le M}a=\frac{M!}{n!}.
\]
Thus, if one subset has product \(n!\), its complement has product
\(Q(n,M)\); conversely, the complement of a subset with product
\(Q(n,M)\) has product \(n!\).

\begin{lemma}
\label[lemma]{lem:no-admissible-endpoint-below-two-n}
For all sufficiently large \(n\), no integer \(M\le 2n\) is admissible.
\end{lemma}

\begin{proof}
Suppose that \(M\le2n\) is admissible, and let
\(\mathcal B\subset(n,M]\cap\mathbb Z\) have product \(Q(n,M)\).
Adjoin every integer in \((M,2n]\) and call the resulting set
\(\mathcal B'\).  The factors remain distinct and
\begin{equation}
 \prod_{b\in\mathcal B'}b
   =\frac{M!}{(n!)^2}\frac{(2n)!}{M!}
   =\binom{2n}{n}.
 \label{eq:two-n-extension-product}
\end{equation}

For \(1\le r\le13\), let
\begin{equation}
 \mathcal L_r^{(0)}
 :=\left\{P\ \text{prime}:
       \frac{n}{r+1}<P\le\frac{2n}{2r+1}\right\}.
 \label{eq:stationary-layers}
\end{equation}
If \(P\in\mathcal L_r^{(0)}\), then
\[
 \left\lfloor\frac nP\right\rfloor=r,
 \qquad
 \left\lfloor\frac{2n}{P}\right\rfloor=2r+1.
\]
Moreover \(P>n/14\), so \(P^2>2n\) when \(n\) is sufficiently
large.  Legendre's formula therefore gives
\begin{equation}
       v_P\binom{2n}{n}=1.
       \label{eq:stationary-carry-one}
\end{equation}
The layers are pairwise disjoint, since membership determines
\(r=\lfloor n/P\rfloor\).

Equation~\eqref{eq:stationary-carry-one} forces exactly one factor of
\(\mathcal B'\) to be divisible by \(P\).  Write this factor as \(Pq(P)\).
Because it belongs to \((n,2n]\), the endpoint inequalities in
\eqref{eq:stationary-layers} imply
\begin{equation}
             r+1\le q(P)\le2r+1.
             \label{eq:stationary-cofactor-range}
\end{equation}
No factor at most \(2n\) contains two primes from the union of the
thirteen layers: each such prime is larger than \(n/14\), and the product
of two of them is larger than \(2n\) for large \(n\).  Consequently the
forced factors \(Pq(P)\) are distinct.

Every integer between \(2\) and \(27\) has a prime divisor in
\(\mathcal P_{\mathrm{sm}}\).  For each \(P\), choose one such divisor
\(\ell(P)\mid q(P)\).  Distinctness of the forced factors gives the
incidence inequality
\begin{equation}
 \begin{split}
 \sum_{r=1}^{13}\lvert\mathcal L_r^{(0)}\rvert
  &\le
    \sum_{\ell\in\mathcal P_{\mathrm{sm}}}
       \sum_{b\in\mathcal B'}v_\ell(b)\\
  &=\sum_{\ell\in\mathcal P_{\mathrm{sm}}}
       v_\ell\binom{2n}{n}.
 \end{split}
 \label{eq:stationary-incidence}
\end{equation}
For each fixed \(r\), the prime number theorem yields
\begin{equation}
 \lvert\mathcal L_r^{(0)}\rvert
  =\left(\frac{2}{2r+1}-\frac{1}{r+1}+o(1)\right)N
  =\left(\frac{1}{(r+1)(2r+1)}+o(1)\right)N.
 \label{eq:stationary-layer-count}
\end{equation}
On the other hand, for each fixed prime \(\ell\),
\begin{equation}
 \begin{split}
 v_\ell\binom{2n}{n}
  &=\sum_{j\ge1}
    \left(
      \left\lfloor\frac{2n}{\ell^j}\right\rfloor
      -2\left\lfloor\frac n{\ell^j}\right\rfloor
    \right)
   =O_\ell(\log n).
 \end{split}
 \label{eq:central-small-prime-valuation}
\end{equation}
Indeed, every summand is either \(0\) or \(1\), and only
\(O_\ell(\log n)\) summands are nonzero.  The left side of
\eqref{eq:stationary-incidence} is
\((A_{13}+o(1))N\) by \eqref{eq:stationary-layer-count}, whereas its
right side is \(O(\log n)=o(N)\).  This contradiction proves the lemma.
\end{proof}

\begin{lemma}
\label[lemma]{lem:thirteen-layer-lower-bound}
As \(n\to\infty\),
\[
       f(n)\ge 2n+(C_0-o(1))\frac{n}{\log n}.
\]
Equivalently,
\[
 \liminf_{n\to\infty}
   \frac{(f(n)-2n)\log n}{n}\ge C_0.
\]
\end{lemma}

\begin{proof}
Let \(M\) be admissible.
Lemma~\ref{lem:no-admissible-endpoint-below-two-n} allows us, for large \(n\),
to write
\begin{equation}
                    M=2n+h,\qquad h>0.
                    \label{eq:lower-h-definition}
\end{equation}
If \(h\ge N\), the desired inequality follows from \(C_0<1\).  We may
therefore assume
\begin{equation}
                    0<h<N.
                    \label{eq:lower-h-range}
\end{equation}

For \(1\le r\le13\), define the moving layer
\begin{equation}
 \mathcal L_r:=
 \left\{P\ \text{prime}:
   \frac{M}{2r+2}<P\le\frac{M}{2r+1},
   \quad
   \frac{n}{r+1}<P\le\frac nr
 \right\}.
 \label{eq:moving-layers}
\end{equation}
The second pair of inequalities is redundant for all sufficiently large
\(n\), uniformly for \(h\) in \eqref{eq:lower-h-range}.  The lower
inequality follows from \(M>2n\), and the upper inequality follows from
\[
       \frac{M}{2r+1}<\frac nr
       \quad\Longleftrightarrow\quad rh<n,
\]
which holds for every fixed \(r\le13\) because \(h<N\).
The prime number theorem at the two endpoints now gives, uniformly for
\(0<h<N\),
\begin{equation}
 \begin{split}
 \lvert\mathcal L_r\rvert
  &=\left(\frac{2}{2r+1}-\frac{1}{r+1}+o(1)\right)N\\
  &=\left(\frac{1}{(r+1)(2r+1)}+o(1)\right)N.
 \end{split}
 \label{eq:moving-layer-count}
\end{equation}
For clarity, this uniformity is an immediate consequence of
\[
  \pi(an+O(N))=\frac{an}{\log n}+o(N)
\]
for fixed \(a>0\), uniformly when the \(O(N)\)-shift is in a fixed
bounded range.  There are only thirteen endpoint pairs.

If \(P\in\mathcal L_r\), then
\[
 \left\lfloor\frac MP\right\rfloor=2r+1,
 \qquad
 \left\lfloor\frac nP\right\rfloor=r.
\]
Also \(P>M/28\), and hence \(P^2>M\) for large \(n\).  Thus Legendre's
formula has no higher-power terms and gives the exact identity
\begin{equation}
          v_P(Q(n,M))
          =\left\lfloor\frac MP\right\rfloor
            -2\left\lfloor\frac nP\right\rfloor
          =1.
 \label{eq:moving-carry-one}
\end{equation}
Let \(\mathcal B\subset(n,M]\) be a distinct-factor representation of
\(Q(n,M)\).  For every \(P\) in the thirteen layers, exactly one member
of \(\mathcal B\) is divisible by \(P\); write it as \(Pq(P)\).
The inequalities \(P\le n/r\), \(Pq(P)>n\),
\(P>M/(2r+2)\), and \(Pq(P)\le M\) imply
\begin{equation}
              r+1\le q(P)\le2r+1.
              \label{eq:moving-cofactor-range}
\end{equation}
The layers are disjoint because \(r=\lfloor n/P\rfloor\).  Moreover, no
factor at most \(M\) can contain two layer primes: both would exceed
\(M/28\), whose square exceeds \(M\) for large \(M\).  Thus the forced
factors are distinct.

Choose a prime
\(\ell(P)\in\mathcal P_{\mathrm{sm}}\) dividing \(q(P)\), which is
possible by \eqref{eq:moving-cofactor-range}.  The same incidence count
as before gives
\begin{equation}
 \sum_{r=1}^{13}\lvert\mathcal L_r\rvert
 \le
 \sum_{\ell\in\mathcal P_{\mathrm{sm}}}v_\ell(Q(n,M)).
 \label{eq:moving-incidence}
\end{equation}
Factor the quotient as
\[
 Q(n,M)
   =\binom{2n}{n}\prod_{2n<a\le2n+h}a.
\]
For a fixed prime \(\ell\), Legendre's formula and
\eqref{eq:central-small-prime-valuation} yield
\begin{equation}
 \begin{split}
 v_\ell(Q(n,M))
  &=v_\ell\binom{2n}{n}
    +\sum_{j\ge1}
      \left(
       \left\lfloor\frac{2n+h}{\ell^j}\right\rfloor
       -\left\lfloor\frac{2n}{\ell^j}\right\rfloor
      \right)\\
  &\le \frac{h}{\ell-1}+O_\ell(\log n).
 \end{split}
 \label{eq:tail-small-prime-capacity}
\end{equation}
Here each floor difference is at most \(h/\ell^j+1\), and there are
only \(O_\ell(\log n)\) relevant powers.

Summing \eqref{eq:moving-layer-count} and using
\eqref{eq:moving-incidence}--\eqref{eq:tail-small-prime-capacity}, we
obtain
\[
       (A_{13}+o(1))N
       \le hS_{23}+O(\log n).
\]
Since \(\log n=o(N)\), this is
\[
              h\ge(C_0-o(1))N
\]
by \eqref{eq:lower-constant}.  Together with the already treated case
\(h\ge N\), this proves the assertion for every admissible endpoint and
hence for \(M=f(n)\).
\end{proof}
 \section{The certified allocation and the central anchors}
\label{sec:allocation-anchors}

For the upper bound, fix a constant
\[
             c>C_0=\frac{4029639598}{25970038185},
\]
and write
\begin{equation}
 \varepsilon:=c-C_0,\qquad
 L:=\log n,\qquad N:=\frac nL,\qquad
 h:=\left\lceil cN\right\rceil,\qquad M:=2n+h.
 \label{eq:upper-basic-parameters}
\end{equation}
We split the complement quotient as
\begin{equation}
 C_n:=\binom{2n}{n},\qquad
 T_n:=\prod_{1\le j\le h}(2n+j),\qquad
 Q(n,M)=C_nT_n.
 \label{eq:central-tail-split}
\end{equation}
This section constructs distinct factors in \((n,2n]\) whose product is
\(C_n\) times a controlled divisor of \(T_n\).  In particular, it removes
the central binomial coefficient exactly, not merely up to an asymptotic
valuation error.

\subsection{An exact infinite cofactor allocation}
\label{sec:infinite-cofactor-allocation}

Set
\begin{equation}
             \alpha_r:=\frac{1}{(r+1)(2r+1)}
             \qquad(r\ge1).
             \label{eq:allocation-row-mass}
\end{equation}
The following finite statement is verified entirely in rational
arithmetic.

\begin{lemma}[Finite allocation certificate]
\label[lemma]{lem:finite-allocation-certificate}
There are \(211\) positive rational numbers
\[
 x^{\mathrm{fin}}_{r,q},
 \qquad
 1\le r\le200,\quad r+1\le q\le2r+1,
\]
all omitted coordinates being zero, with the following properties.
\begin{enumerate}
 \item For every \(1\le r\le200\),
 \[
       \sum_{q=r+1}^{2r+1}x^{\mathrm{fin}}_{r,q}
       =\alpha_r.
 \]
 \item If
 \[
       \lambda_\ell^{\mathrm{fin}}
       :=\sum_{r=1}^{200}\sum_q
             v_\ell(q)x^{\mathrm{fin}}_{r,q},
 \]
 then
 \[
       \lambda_\ell^{\mathrm{fin}}
       \le\frac{C_0}{\ell-1}
 \]
 for every prime \(\ell\le401\).
 \item If \(201<p\le401\) is prime and \(p_-\) denotes the prime
 immediately preceding \(p\), then
 \[
 \lambda_p^{\mathrm{fin}}
  +\sum_{r=\max(201,p_-)}^{p-1}\alpha_r
       \le\frac{C_0}{p-1}.
 \]
\end{enumerate}
\end{lemma}

The complete 211-entry rational array and its exact checker are contained
in the single companion file \texttt{numerical\_verifier.py}, distributed
with this paper.  The checker verifies positivity and the permitted cofactor
range, adds every row in exact rational arithmetic, factors every cofactor,
and checks the two families of capacity inequalities in
Lemma~\ref{lem:finite-allocation-certificate}.  Thus the file is both an
explicit certificate for the existential statement above and a reproducible
verification of every finite calculation used here.  No floating-point
calculation enters the verification; see
Appendix~\ref{app:numerical-verification} for its precise scope.

\begin{lemma}[Infinite cofactor allocation]
\label[lemma]{lem:nagura-allocation-tail}
There are nonnegative rational numbers \(x_{r,q}\), indexed by
\[
       r\ge1,\qquad r+1\le q\le2r+1,
\]
such that
\begin{align}
 \sum_{q=r+1}^{2r+1}x_{r,q}
    &=\alpha_r
       &&(r\ge1),                                     
 \label{eq:allocation-row-sums}\\
 \sum_{r\ge1}\sum_{q=r+1}^{2r+1}
       v_\ell(q)x_{r,q}
    &\le\frac{C_0}{\ell-1}
       &&(\ell\ \text{prime}).
 \label{eq:allocation-capacities}
\end{align}
\end{lemma}

\begin{proof}
For \(r\le200\), use the array in
Lemma~\ref{lem:finite-allocation-certificate}.  For \(r\ge201\), let \(p(r)\)
be the least prime strictly larger than \(r\), and put
\[
          x_{r,p(r)}:=\alpha_r,
          \qquad x_{r,q}:=0\quad(q\ne p(r)).
\]
Bertrand's postulate gives \(r<p(r)<2r\), so \(p(r)\) is an allowed
cofactor.  The row identities \eqref{eq:allocation-row-sums} follow.

It remains to check the prime capacities.  A tail row contributes only
to the prime \(p(r)\).  If \(p\) and \(p_-\) are consecutive primes, the
indices for which the least-prime rule routes to \(p\) are precisely
\begin{equation}
                   p_-\le r\le p-1.
                   \label{eq:least-prime-tail-block}
\end{equation}
The actual tail contribution is the intersection of this block with
\(r\ge201\).
For \(p\le200\), there is no tail contribution, while for
\(200<p\le401\), the finite-plus-tail inequality is exactly the third
check in Lemma~\ref{lem:finite-allocation-certificate}.

Now let \(p>401\).  Every finite cofactor is at most \(401\), so its
\(p\)-load is zero.  Nagura's theorem~\cite{Nagura1952}, applied to
\(p_->25\), gives \(p<6p_-/5\).  Hence \(p-p_-<p_-/5\), and
\begin{equation}
 \begin{split}
 (p-1)\sum_{r=p_-}^{p-1}\alpha_r
  &<
   p\,\frac{p-p_-}{2p_-^2}\\
  &<\frac{3}{25}
   <C_0.
 \end{split}
 \label{eq:nagura-tail-capacity}
\end{equation}
The first strict inequality uses
\((r+1)(2r+1)>2p_-^2\) throughout the block, and the last one follows
directly from the rational value of \(C_0\).  Dividing
\eqref{eq:nagura-tail-capacity} by \(p-1\) proves
\eqref{eq:allocation-capacities} for every remaining prime.
\end{proof}

\subsection{Exact realization of a fixed prefix}
\label{sec:fixed-prefix-anchors}

For \(r\ge1\), define the central carry row
\begin{equation}
 \mathcal C_r(n):=
 \left\{P\ \text{prime}:
      \frac{n}{r+1}<P\le\frac{2n}{2r+1}\right\}.
 \label{eq:central-carry-row}
\end{equation}
If \(P\in\mathcal C_r(n)\), then
\[
 \left\lfloor\frac nP\right\rfloor=r,
 \qquad
 \left\lfloor\frac{2n}{P}\right\rfloor=2r+1.
\]
For fixed \(r\) and sufficiently large \(n\), one also has \(P^2>2n\),
and therefore
\begin{equation}
                    v_P(C_n)=1.
                    \label{eq:central-row-valuation}
\end{equation}
Every integer \(q\) in the allocation range obeys
\begin{equation}
       r+1\le q\le2r+1
       \quad\Longrightarrow\quad
       n<Pq\le2n.
 \label{eq:central-anchor-location}
\end{equation}
The prime number theorem gives, for each fixed \(r\),
\begin{equation}
       \lvert\mathcal C_r(n)\rvert
       =(\alpha_r+o(1))N.
       \label{eq:central-row-count}
\end{equation}

Fix an integer \(R\ge200\), independently of \(n\).  For each \(r\le R\),
partition \(\mathcal C_r(n)\) among the positive coordinates
\(x_{r,q}\).  More precisely, round
\(x_{r,q}N\) to integers for all but one positive coordinate and assign
all remaining primes in the row to the last coordinate.  Because the
number of coordinates is fixed and \eqref{eq:central-row-count} holds,
the resulting parts \(\mathcal C_{r,q}(n)\) satisfy
\begin{equation}
 \lvert\mathcal C_{r,q}(n)\rvert=x_{r,q}N+o(N)
 \label{eq:integer-prefix-parts}
\end{equation}
for every positive coordinate; in particular, the last part is
nonnegative for all sufficiently large \(n\).  The parts exhaust the row
exactly.

Use the factors
\begin{equation}
 \mathcal H_R^{\mathrm{pre}}
 :=\{Pq:P\in\mathcal C_{r,q}(n),\ 1\le r\le R\}.
 \label{eq:prefix-anchor-set}
\end{equation}
They lie in \((n,2n]\) by \eqref{eq:central-anchor-location}.  They are
also distinct.  Indeed, for fixed \(R\) and large \(n\),
\[
        P>\frac{n}{R+1}>2R+1\ge q.
\]
Thus \(P\) is the unique prime divisor of \(Pq\) larger than \(2R+1\);
equality of two prefix anchors forces equality of both \(P\) and \(q\).
This also shows that repeated use of the same cofactor in different rows
cannot create a collision.

Define the cofactor product
\begin{equation}
 D_R^{\mathrm{pre}}
 :=\prod_{r=1}^{R}\prod_q
        q^{\lvert\mathcal C_{r,q}(n)\rvert}.
 \label{eq:prefix-cofactor-product}
\end{equation}
For every prime \(\ell\), the finite realization
\eqref{eq:integer-prefix-parts} and the full capacity inequality
\eqref{eq:allocation-capacities} give
\begin{equation}
 v_\ell(D_R^{\mathrm{pre}})
 \le\left(\frac{C_0}{\ell-1}+o_R(1)\right)N.
 \label{eq:prefix-cofactor-load}
\end{equation}
All prime divisors of \(D_R^{\mathrm{pre}}\) are at most \(2R+1\).

There is also a row \(r=0\).  Every prime \(P\) with
\begin{equation}
                         n<P\le2n
                         \label{eq:row-zero-primes}
\end{equation}
occurs once in \(C_n\), and we use \(P\) itself as an anchor.  These
singletons have cofactor \(1\), consume no tail valuation, and cannot
collide with the composite factors in \eqref{eq:prefix-anchor-set}.

\subsection{Promotion of all residual central valuations}
\label{sec:residual-central-promotion}

Put
\begin{equation}
                       X_R^{\rm anc}:=\frac{n}{R+1}.
                       \label{eq:anchor-signature-cutoff}
\end{equation}
Once \(R\) is fixed, take \(n\) sufficiently large that
\begin{equation}
       X_R^{\rm anc}>2R+1,\qquad (X_R^{\rm anc})^2>2n.
       \label{eq:anchor-scale-separation}
\end{equation}
Every prime \(P>X_R^{\rm anc}\) occurring in \(C_n\) then occurs to the first power
and is accounted for by exactly one of the rows
\eqref{eq:central-carry-row} with \(1\le r\le R\), or by the row-zero
set \eqref{eq:row-zero-primes}.  Indeed,
\(\lfloor n/P\rfloor\le R\), and the equality
\[
 \left\lfloor\frac{2n}{P}\right\rfloor
 -2\left\lfloor\frac nP\right\rfloor=1
\]
is precisely the membership condition for the corresponding carry row.

It remains to account for primes \(p\le X_R^{\rm anc}\).  Put
\[
             e_p:=v_p(C_n),\qquad B_p:=p^{e_p}.
\]
Every summand in Legendre's formula
\begin{equation}
 e_p=\sum_{j\ge1}
 \left(
    \left\lfloor\frac{2n}{p^j}\right\rfloor
    -2\left\lfloor\frac n{p^j}\right\rfloor
 \right)
 \label{eq:residual-central-exponent}
\end{equation}
is \(0\) or \(1\).  Consequently
\begin{equation}
       e_p\le\lfloor\log_p(2n)\rfloor,
       \qquad B_p\le2n.
       \label{eq:central-prime-power-bound}
\end{equation}
For every \(p\le X_R^{\rm anc}\) with \(e_p>0\), let \(k_p\) be the least
nonnegative integer for which
\begin{equation}
                         n<2^{k_p}B_p.
                         \label{eq:promotion-exponent}
\end{equation}
Minimality and \eqref{eq:central-prime-power-bound} imply
\[
                         n<2^{k_p}B_p\le2n.
\]
We therefore use the single promoted factor
\begin{equation}
                         A_p:=2^{k_p}p^{e_p}.
                         \label{eq:promoted-central-factor}
\end{equation}

The promoted factors are mutually distinct.  Factors belonging to two
different odd base primes are distinguished by unique factorization, and
the factor with base prime \(2\) is a pure power of two.  They also cannot
collide with an earlier anchor: every prime divisor of a promoted factor
is at most \(X_R^{\rm anc}\), whereas a prefix anchor has a unique prime divisor
above \(X_R^{\rm anc}\), and a row-zero singleton is itself a prime above \(n\).

The additional power of two has small total cost.  Define
\begin{equation}
 K_R(n):=\sum_{\substack{p\le X_R^{\rm anc}\\e_p>0}}k_p.
 \label{eq:total-promotion-cost}
\end{equation}
Since \(B_p\ge p\), minimality in
\eqref{eq:promotion-exponent} gives
\begin{equation}
                     k_p\le1+\log_2\frac np.
                     \label{eq:pointwise-promotion-cost}
\end{equation}
For primes \(p\le\sqrt n\), the total contribution is
\[
        O\bigl(\pi(\sqrt n)\log n\bigr)=O(\sqrt n)=o(N).
\]
For \(p>\sqrt n\), split the primes into dyadic blocks
\[
       \frac{n}{2^{j+1}}<p\le\frac{n}{2^j}.
\]
Only indices
\[
 j\ge\lfloor\log_2(R+1)\rfloor-O(1),
 \qquad 2^j\le\sqrt n,
\]
occur.  The standard Chebyshev upper bound for primes, uniformly in this
range, gives
\[
 \pi(n/2^j)\ll
 \frac{n}{2^j\log(n/2^j)}
 \ll\frac{N}{2^j}.
\]
Together with \eqref{eq:pointwise-promotion-cost}, this yields
\begin{equation}
 \begin{split}
 K_R(n)
  &\ll N\sum_{j\ge\log_2(R+1)-O(1)}
          \frac{j+1}{2^j}+o(N)\\
  &\ll N\frac{1+\log R}{R}+o(N),
 \end{split}
 \label{eq:dyadic-promotion-bound}
\end{equation}
with an absolute implied constant.  In particular,
\begin{equation}
 \lim_{R\to\infty}\limsup_{n\to\infty}\frac{K_R(n)}N=0.
 \label{eq:promotion-double-limit}
\end{equation}

\subsection{The exact anchor product and its tail reserve}
\label{sec:full-central-anchor}

Let \(\mathcal H_R\) be the union of the prefix anchors
\eqref{eq:prefix-anchor-set}, the row-zero singletons
\eqref{eq:row-zero-primes}, and the promoted factors
\eqref{eq:promoted-central-factor}.  The preceding arguments show that
these are distinct integers in \((n,2n]\).  Every prime valuation of
\(C_n\) appears exactly once in their product: the valuations above
\(X_R^{\rm anc}\) occur in their carry or row-zero anchors, and each valuation at
or below \(X_R^{\rm anc}\) occurs in its complete block \(p^{e_p}\).  Therefore,
if
\begin{equation}
 D_R:=D_R^{\mathrm{pre}}\,2^{K_R(n)},
 \label{eq:full-anchor-divisor}
\end{equation}
then the following is an exact integer identity:
\begin{equation}
                  \prod_{a\in\mathcal H_R}a=C_nD_R.
                  \label{eq:exact-anchor-product}
\end{equation}
Moreover \(D_R\) is supported on the fixed prime set
\begin{equation}
        F_R:=\{\ell\ \text{prime}:\ell\le2R+1\}.
        \label{eq:anchor-prime-set}
\end{equation}

We now choose \(R=R(c)\).  By
\eqref{eq:promotion-double-limit}, we may fix \(R\ge201\) so large that
\begin{equation}
                       K_R(n)\le\frac{\varepsilon}{4}N
                       \label{eq:chosen-promotion-budget}
\end{equation}
for every sufficiently large \(n\).  This choice is made before \(n\)
tends to infinity.  For each fixed prime \(\ell\), Legendre's formula
over the tail interval gives
\begin{equation}
 \begin{split}
 v_\ell(T_n)
  &=\sum_{j\ge1}
     \left(
       \left\lfloor\frac{2n+h}{\ell^j}\right\rfloor
       -\left\lfloor\frac{2n}{\ell^j}\right\rfloor
     \right)\\
  &=\frac{h}{\ell-1}+O_\ell(\log n)
   =\left(\frac{c}{\ell-1}+o(1)\right)N.
 \end{split}
 \label{eq:tail-prime-valuation-asymptotic}
\end{equation}
For an odd prime \(\ell\in F_R\),
\eqref{eq:prefix-cofactor-load} and \eqref{eq:full-anchor-divisor} give
\[
       v_\ell(D_R)
       \le\left(\frac{C_0}{\ell-1}+o(1)\right)N.
\]
At \(\ell=2\), the promotion cost adds at most
\(\varepsilon N/4\), so
\[
       v_2(D_R)
       \le\left(C_0+\frac{\varepsilon}{4}+o(1)\right)N.
\]
Because \(F_R\) is fixed, these inequalities and
\eqref{eq:tail-prime-valuation-asymptotic} hold simultaneously at all
of its primes.  After increasing \(n\), if necessary, they imply
\begin{equation}
       v_\ell(T_n)-v_\ell(D_R)
       \ge\frac{\varepsilon}{3(\ell-1)}N
       \qquad(\ell\in F_R).
 \label{eq:strict-anchor-reserve}
\end{equation}
Outside \(F_R\), the divisor \(D_R\) has zero valuation.  Thus
\begin{equation}
                       D_R\mid T_n.
                       \label{eq:anchor-divides-tail}
\end{equation}

We record the conclusion in the form used later.

\begin{lemma}[Full central-anchor lemma]
\label[lemma]{lem:full-central-anchor}
Fix \(c>C_0\), and let \(h,M,C_n,T_n\) be as in
\eqref{eq:upper-basic-parameters}--\eqref{eq:central-tail-split}.  There
is a finite prime set \(F_{\mathrm{anc}}=F_{\mathrm{anc}}(c)\) such
that, for all sufficiently large \(n\), one can find an
\(F_{\mathrm{anc}}\)-supported divisor \(D\mid T_n\) and a set
\(\mathcal H\) of distinct integers in \((n,2n]\) satisfying
\begin{equation}
                \prod_{a\in\mathcal H}a=C_nD.
                \label{eq:central-anchor-lemma-product}
\end{equation}
In addition,
\begin{equation}
       v_\ell(T_n)-v_\ell(D)
       \ge\frac{c-C_0}{3(\ell-1)}N
       \qquad(\ell\in F_{\mathrm{anc}}).
       \label{eq:central-anchor-lemma-reserve}
\end{equation}
Consequently, removal of these anchors leaves the exact quotient
\[
            \frac{Q(n,M)}
                 {\prod_{a\in\mathcal H}a}
            =\frac{T_n}{D}.
\]
\end{lemma}

\begin{proof}
Take \(R=R(c)\) as above, set \(F_{\mathrm{anc}}:=F_R\),
\(D:=D_R\), and \(\mathcal H:=\mathcal H_R\).  Equations
\eqref{eq:exact-anchor-product}, \eqref{eq:strict-anchor-reserve}, and
\eqref{eq:anchor-divides-tail} give all assertions.
\end{proof}

We finally note an endpoint detail that will be useful when the residual
tail is realized.  A signature prime \(P>X_R^{\rm anc}\) used in a carry or
row-zero anchor can occur once in \(C_n\) and also occur in one or more
tail terms.  The anchor construction uses exactly its central occurrence.
Since \(D\) is supported on primes at most \(2R+1<X_R^{\rm anc}\), all of these
additional \(P\)-occurrences remain in \(T_n/D\).  This causes no
numerical collision: every anchor is at most \(2n\), whereas every actual
tail term is larger than \(2n\).  In particular, if
\(n<P\le n+h/2\), the row-zero anchor is the singleton \(P\), while
\(2P\), if it is used later, is a distinct tail factor.  Primes in
\((2n,M]\) occur only in the tail and are untouched by the anchor
construction.
 \section{A precharged two-sided universal bank}
\label{sec:rounding-bank}

We reserve now the switches that will absorb the discrepancy produced by
integral rounding.  This is done before exceptional tail factors are
retained and before a fractional selector is solved: the chosen base state
of every switch is part of the charged residual product.

Put
\begin{equation}
 \theta:=\frac29,\qquad y:=n^\theta,\qquad
 d_n:=\lceil\log_2(3n)\rceil,\qquad
 \beta_p:=4d_n\left\lceil\frac{\log(3n)}{\log p}\right\rceil
 \quad(p\le y).
 \label{eq:bank-discrepancy-box}
\end{equation}
Lemma~\ref{lem:floating-rounding} will show that retaining every integer
rough-row sum while rounding creates an integral error \(e\), supported on
\(p\le y\), with \(|e_p|\le\beta_p\).  The elementary PNT estimates
\[
 \sum_{p\le y}\frac1{\log p}\ll\frac{y}{(\log y)^2},
 \qquad \pi(y)\ll\frac{y}{\log y}
\]
give
\begin{align}
 \sum_{p\le y}\beta_p&=O(y),
 \label{eq:bank-number-of-paths}\\
 \sum_{p\le y}\beta_p\log p&=O(yL).
 \label{eq:bank-total-path-length}
\end{align}

\subsection{Complete signatures and path states}
\label{subsec:bank-signatures}

For every legal lower or upper factor define its complete rough signature
\begin{equation}
 S(a):=(v_P(a))_{P>y}\qquad(n<a\le M).
 \label{eq:complete-rough-signature}
\end{equation}
Equivalently, the integer label used in
Section~\ref{sec:rough-selector} is
\(R_y(a)=\prod_{P>y}P^{S_P(a)}\); the vector and integer conventions
partition the factors into exactly the same rows.
For a finite set \(G\) of legal factors, put
\[
 m_S(G):=\#\{a\in G:S(a)=S\},\qquad
 v(G):=v\!\left(\prod_{a\in G}a\right).
\]
A switch \(g\) has two states \(G_g^0,G_g^1\) and change
\(\Delta_g=v(G_g^1)-v(G_g^0)\).  We require
\begin{equation}
 m_S(G_g^0)=m_S(G_g^1)
 \qquad\text{for every complete signature }S.
 \label{eq:bank-all-factor-invariance}
\end{equation}
Our switches are paths of component edges.  Each component edge has one
factor in each state, and those factors have the same complete signature;
different components use different marker primes.  Hence a complete path
state has one token in each of its component marker rows, and its two full
states have identical signature counts row by row.  In particular, a full
path has one token per component row, not one token in total.

Denote the eventual collection of path switches by \(\mathcal B\).  It is
\emph{two-sided \(\beta\)-universal} if every integer vector
\begin{equation}
 z=(z_p)_{p\le y},\qquad |z_p|\le\beta_p,
 \label{eq:universal-box-vector}
\end{equation}
is the exact valuation change of a collision-free collection of toggles.
We now construct such a bank while backing every one of its rough tokens by
a distinct actual tail occurrence.

\subsection{The literal geometric descent}
\label{subsec:bank-ordinary-descent}

Use the grid
\begin{equation}
                         Q_j:=4(4/3)^j\qquad(j\ge0).
 \label{eq:bank-literal-grid}
\end{equation}
Every integer core \(q>4\) lies in a unique cell \((Q,4Q/3]\).  At a
nonbottom cell choose
\begin{equation}
 q\in(Q,4Q/3],\qquad 3Q/4<b<q,\qquad
 P\in I_Q:=\left(\frac{4n}{3Q},\frac{3n}{2Q}\right].
 \label{eq:bank-ordinary-cell}
\end{equation}
Then
\begin{equation}
                         n<Pb<Pq\le2n,
 \label{eq:bank-ordinary-endpoints}
\end{equation}
so \(Pq\leftrightarrow Pb\) is a legal lower switch of ratio \(b/q\).

Here is a literal descent which avoids the power-of-two cores of the
promoted anchors.  If \(q\ge6\) is not a power of two and \(Q>20\), set
\(b_0=\lceil4q/5\rceil\) and
\begin{equation}
 b:=\begin{cases}
       b_0,&b_0\text{ is not a power of two},\\
       b_0-1,&b_0\text{ is a power of two}.
    \end{cases}
 \label{eq:bank-large-core-descent}
\end{equation}
Then \(b\) is a non-power and
\[
                  b>\frac45Q-1>\frac34Q,\qquad 5\le b<q.
\]
For the non-power cores in cells with \(Q\le20\), use
\begin{equation}
\begin{array}{c|rrrrrrrrrrrrrrr}
q&6&7&9&10&11&12&13&14&15&17&18&19&20&21&22\\ \hline
b&5&6&7&9&9&10&11&12&12&14&15&15&15&17&18.
\end{array}
\label{eq:bank-finite-descent-table}
\end{equation}
Every displayed pair satisfies the strict cell inequalities.  A source
prime \(p\ge5\) is a non-power, so induction keeps every intermediate core
a non-power.  The core decreases, and above the table it is at most
\(4q/5+1\); the path therefore reaches \(5\) in \(O(\log p)\) steps.  We
stop there.  Moreover, when \(Q>20\) one has
\(b<17q/20\); hence two consecutive moves whose source cores remained in
the same cell would reduce a core at most \(4Q/3\) below
\((17/20)^2(4Q/3)<Q\).  Thus a path uses at most two edges at every
large scale, while the finite table gives an absolute bound at the five
small scales.  In particular it uses \(O(1)\) edges at each geometric
scale.  The strict ordinary descent is not asserted to reach \(2\).

For a nonbottom scale \(Q\le y\), define the donor multiplicity
\begin{equation}
 s_P:=\#\left\{u:\frac{7Q}{5}\le u\le\frac{29Q}{20},\
       P^+(u)\le y,\ 2n<Pu\le M\right\}.
 \label{eq:bank-donor-multiplicity}
\end{equation}
When \(Q>20\), there are \(\asymp Q\) eligible \(y\)-smooth \(u\)'s.
Indeed, \(u<2y\), so an integer here having a prime factor above \(y\)
must itself be such a prime; this discards only \(O(Q/\log Q)\) values.
For fixed eligible \(u\), the interval
\[
                         2n/u<P\le M/u
\]
lies in a fixed shrunken part of \(I_Q\).  The uniform PNT
\eqref{eq:PNT-difference} gives
\begin{equation}
 \#\{P:2n/u<P\le M/u\}\asymp_c\frac{N}{QL}.
 \label{eq:bank-one-donor-count}
\end{equation}
Indeed its endpoints have size \(n/Q\), with
\(\log(n/Q)\ge(1-\theta)L\), and its length is
\(h/u\asymp_c N/Q\).  The cumulative-PNT error at the two moving
endpoints is
\(O((n/Q)e^{-c_0\sqrt L})=o(N/(QL))\), uniformly for \(Q\le y\).
Thus \eqref{eq:bank-one-donor-count} does not invoke a prime theorem for
arbitrary short intervals.
At the five nonbottom scales \(Q_j\le20\), use respectively the fixed
donor cofactors
\begin{equation}
\begin{array}{c|ccccc}
j&1&2&3&4&5\\ \hline
u&8&10&13&17&23.
\end{array}
\label{eq:bank-small-scale-donors}
\end{equation}
Each lies between \(4Q_j/3\) and \(3Q_j/2\), and the same PNT estimate
applies.

Double-counting \((P,u)\) gives
\begin{equation}
 C_Q:=\sum_{P\in I_Q}s_P\asymp_c\frac NL.
 \label{eq:bank-marker-occurrences}
\end{equation}
A fixed \(P\in I_Q\) admits only
\(O(h/P+1)=O_c(Q/L+1)\) donor cofactors.  Hence
\begin{equation}
 K_Q:=\#\{P\in I_Q:s_P>0\}
      \asymp_c\frac{N}{\max(Q,L)}.
 \label{eq:bank-marker-signatures}
\end{equation}
The upper bound follows from either the length of \(I_Q\) or
\eqref{eq:bank-marker-occurrences}; division by the maximum multiplicity
gives the lower bound.

Choose one donor for each of the \(K_Q\) eligible marker primes, and split
this set of marker--donor pairs into two parts whose sizes differ by at
most one.  These are the two ordinary orientation pools.  Thus each
orientation separately retains \(\asymp K_Q\) distinct signatures; the
occurrence bound \eqref{eq:bank-marker-occurrences} remains available as a
separate verification of donor load.

Successive marker intervals are disjoint: if \(Q'=4Q/3\), then
\[
 I_{Q'}=\left(\frac nQ,\frac{9n}{8Q}\right],\qquad
 I_Q=\left(\frac{4n}{3Q},\frac{3n}{2Q}\right],
\]
with a fixed gap.  Moreover
\begin{equation}
                         P\gg n/y>y,\qquad P^2>M.
 \label{eq:bank-marker-is-rough}
\end{equation}
Thus the donor and both endpoints have the same complete signature,
namely the single marker \(P\) with exponent one.

\subsection{Four bottom pools and the signed unit lattice}
\label{subsec:bank-bottom-lattice}

Continue the descent from \(5\) through the four pairwise disjoint pools
\begin{equation}
\begin{array}{c|c|c|c}
\text{move}&\text{marker interval}&\text{states}&\text{tail donor}\\ \hline
5\to4&(n/3,M/6]&5P\leftrightarrow4P&6P\\
4\to3&(2n/5,M/5]&4P\leftrightarrow3P&5P\\
3\to2&(2n/3,M/3]&3P\leftrightarrow2P&3P\\
2\to1&(n/2,M/4]&4P\leftrightarrow2P&4P.
\end{array}
\label{eq:bank-four-bottom-pools}
\end{equation}
All state factors are in \((n,M]\).  In the last two lines the upper state
is itself the tail donor; it is one occurrence, not a retained donor plus
a second bank factor.  The unsplit prime capacities are
\begin{equation}
 \left(\frac c6+o(1)\right)\frac NL,\quad
 \left(\frac c5+o(1)\right)\frac NL,\quad
 \left(\frac c3+o(1)\right)\frac NL,\quad
 \left(\frac c4+o(1)\right)\frac NL.
 \label{eq:bank-bottom-capacities}
\end{equation}
These are differences of the cumulative PNT at endpoints of size
\(\asymp n\): the four interval lengths are
\(h/6,h/5,h/3,h/4\), and the absolute PNT remainder
\(O(ne^{-c_0\sqrt L})\) is \(o(N/L)\).
Split each interval into two fixed positive subintervals for the two
orientations.

The word ``split'' means a split into two subintervals whose lengths are
fixed positive proportions of the original interval.  Consequently, for
constants \(c_i=c_i(c)>0\), each of the eight oriented bottom pools
contains
\begin{equation}
             (c_i+o(1))\frac{N}{L}
             \quad\hbox{marker primes}.
 \label{eq:bank-oriented-bottom-capacities}
\end{equation}
This is much larger than the \(O(y)\) requests made to any one pool.

For \(p\ge5\), concatenate the ordinary \(p\to5\) path and the four
bottom moves.  In the downward orientation its ratio and valuation change
are
\begin{align}
 \frac4p\frac34\frac23\frac12&=\frac1p,
 \label{eq:bank-unit-ratio}\\
 (2\mathbf e_2-\mathbf e_p)+(\mathbf e_3-2\mathbf e_2)
 +(\mathbf e_2-\mathbf e_3)-\mathbf e_2&=-\mathbf e_p.
 \label{eq:bank-unit-vector}
\end{align}
For \(p=3\), use only \(3\to2\to1\), and for \(p=2\), only the terminal
move.  These yield \(-\mathbf e_3\) and \(-\mathbf e_2\).  Reverse every
base orientation in a disjoint copy to obtain the positive unit vector.
Reserve \(\beta_p\) copies of each orientation for every \(p\le y\).

There are \(O(y)\) paths by \eqref{eq:bank-number-of-paths}.  A path uses
only \(O(1)\) edges at any geometric scale.  At the worst scale \(Q=y\),
\begin{equation}
 \frac{O(y)}{K_Q}\ll\frac{y^2}{N}=o(1),\qquad
 \frac{O(y)}{C_Q}\ll\frac{yL}{N}=o(1).
 \label{eq:bank-ordinary-congestion}
\end{equation}
Thus every ordinary component can receive its own marker prime and donor.
Each bottom pool receives \(O(y)\) requests, and its separate utilization
is
\begin{equation}
                         O(y)/(N/L)=O(yL/N)=o(1).
 \label{eq:bank-bottom-congestion}
\end{equation}
The total number of component factors is
\begin{equation}
 O\left(\sum_{p\le y}\beta_p\log p\right)=O(yL)=o(N).
 \label{eq:bank-total-size}
\end{equation}

\begin{lemma}[Disjoint donor assignment]
\label[lemma]{lem:bank-donor-injection}
For all sufficiently large \(n\), every component in the fully reserved
bank can be assigned a pair \((P,Pu)\), where \(P\) is its marker and
\(Pu\in E\) is its backing donor, so that the following hold
simultaneously.
\begin{enumerate}
 \item Distinct components have distinct markers and distinct donor
       occurrences.
 \item The two orientations at a fixed scale use disjoint pools; pools at
       different ordinary scales, and the four bottom pools, are mutually
       disjoint.
 \item The donor and both states of the component have the same complete
       rough signature.
 \item If a bottom donor is itself one of the two states, it denotes that
       one occurrence, not an additional copy of the same integer.
\end{enumerate}
Thus assignment of the donors is an injection, rather than only a
comparison of their total number with the number of requests.
\end{lemma}

\begin{proof}
At an ordinary scale \(Q\), first choose one donor for every eligible
marker, as above, and split the resulting marker--donor pairs into the two
orientation pools.  Let \(m_Q^\pm\) be the numbers of component requests
of the two orientations.  A path has \(O(1)\) components at a given
scale, by the literal descent estimate preceding
\eqref{eq:bank-ordinary-congestion}, and
\eqref{eq:bank-number-of-paths} gives
\(m_Q^\pm=O(y)\).  On the other hand, each orientation pool has
\(\asymp K_Q\) pairs.  Uniformly for \(Q\le y\),
\[
 \frac{m_Q^\pm}{K_Q}
 \ll \frac{y\max(Q,L)}N
 \le \frac{y^2+yL}{N}=o(1)
\]
by \eqref{eq:bank-marker-signatures}.  A greedy assignment therefore
uses a fresh pair for every request.  Successive ordinary marker intervals
are disjoint by the calculation following
\eqref{eq:bank-marker-signatures}, so assignments made at different
scales cannot meet.

For a bottom pool, \eqref{eq:bank-oriented-bottom-capacities} supplies
\((c_i+o(1))N/L\) marker--donor occurrences, whereas the request count is
\(O(y)=o(N/L)\).  The same greedy construction works separately in each
oriented subpool.  The four displayed marker intervals are disjoint for
large \(n\), and the last ordinary marker interval lies below the first
bottom interval, so these choices do not meet any ordinary choice.

Finally, at an ordinary scale the three cofactors \(q,b,u\) are
\(y\)-smooth and \(P>y\), \(P^2>M\), by
\eqref{eq:bank-marker-is-rough}; the same statement is immediate from the
four fixed bottom rows.  Hence every factor in one component has complete
signature \(P\) with exponent one.  In the last two bottom rows the donor
and one state coincide by definition, and the construction counts that
single occurrence only once.  This proves all assertions.
\end{proof}

Given \(z\) in \eqref{eq:universal-box-vector}, toggle \(|z_p|\) copies
having its sign.  Equation~\eqref{eq:bank-unit-vector}, with its
\(p=2,3\) truncations, gives total change exactly \(z\).  Every such plan
is a subplan of the fully reserved plan, so
\eqref{eq:bank-ordinary-congestion}--\eqref{eq:bank-bottom-congestion}
give uniform \(o(1)\) congestion throughout the box.

\subsection{Anchor modifications and guards}
\label{subsec:bank-anchor-guards}

Choose all bank markers before freezing the integer anchor partitions of
Section~\ref{sec:allocation-anchors}.  The following lemma records every
point at which the resulting modification is used later.

\begin{lemma}[Collision-free modification of the anchors]
\label[lemma]{lem:bank-anchor-modification}
The integer partitions in the fixed prefix can be chosen so that no bank
state is an anchor.  More precisely, there are modified anchors
\(\mathcal H'\) and an \(F_{\mathrm{anc}}\)-supported integer \(D'\) such
that
\begin{enumerate}
 \item every forced bottom marker belongs to the asserted central carry
       row, and every replacement cofactor is legal in that row;
 \item the forced row-\(1\) and row-\(2\) markers use only
       \(o(N)\) of positive certificate cells of size \(\asymp N\);
 \item the anchors in \(\mathcal H'\) are mutually distinct and avoid both
       incident states of every bank component;
 \item
       \[
          \prod_{a\in\mathcal H'}a=C_nD',\qquad D'\mid T_n;
       \]
 \item for every \(f\in F_{\mathrm{anc}}\), the change from the original
       cofactor product is \(o(N)\), so a fixed positive linear part of
       \eqref{eq:central-anchor-lemma-reserve} remains.
\end{enumerate}
These statements remain valid when a legal replacement cofactor was a
zero coordinate of the rational allocation: only \(O(y)=o(N)\) such
occurrences are inserted.
\end{lemma}

\begin{proof}
For a \(5\to4\) marker \(P\in(n/3,M/6]\), one has, for all sufficiently
large \(n\),
\[
 \left\lfloor n/P\right\rfloor=2,\qquad
 \left\lfloor2n/P\right\rfloor=5,\qquad v_P(C_n)=1.
\]
Thus it lies in the row-\(2\) carry interval, and we force its anchor to
be \(3P\).  The relevant certificate coordinate is strictly positive:
\begin{equation}
                  x_{2,3}=\frac{3432239399}{51940076370}>0.
 \label{eq:bank-positive-row-two-cell}
\end{equation}
Likewise a terminal marker \(P\in(n/2,M/4]\) belongs to the row-\(1\)
carry interval and is assigned the anchor \(3P\), where
\begin{equation}
                  x_{1,3}=\frac{597400199}{51940076370}>0.
 \label{eq:bank-positive-row-one-cell}
\end{equation}
Each of these two certificate cells contains
\(x_{r,3}N+o(N)\asymp N\) marker primes by
\eqref{eq:integer-prefix-parts}.  The number forced into either cell is
\(O(y)=o(N)\), by \eqref{eq:bank-number-of-paths}; hence the integer
partition can accommodate every forced marker.  The terminal bank switch
is still only \(4P\leftrightarrow2P\): the anchor \(3P\) uses the central
\(P\)-occurrence, while either bank state uses its separately backed tail
occurrence.

We record explicitly why the fixed prefix meets only \(O(y)\) bank
components, although the entire bank has \(O(yL)\) components.  If an
ordinary component at scale \(Q\) meets a fixed-prefix anchor, equality of
the two integers forces equality of their complete rough signatures and
hence forces the component marker \(P\) to be the anchor prime: for fixed
\(R\) and large \(n\), its prime is \(>n/(R+1)>y\), whereas its cofactor
is at most \(2R+1<y\).  A fixed-prefix anchor has
\(P>n/(R+1)\), while
\(P\in I_Q\subset(4n/(3Q),3n/(2Q)]\).  Consequently
\[
                         Q<\frac{3(R+1)}2.
\]
There are only \(O_R(1)\) geometric scales satisfying this inequality.
By the literal descent estimate preceding
\eqref{eq:bank-ordinary-congestion}, each of the \(O(y)\) paths has
\(O(1)\) components at any one scale.  (The stronger assertion ``at most
one'' is neither needed nor generally true for the finite descent table.)
Hence ordinary components meeting the prefix number \(O_R(y)\).  The
four bottom moves contribute at most another \(O(y)\).  Thus
\begin{equation}
 \#\{\text{bank components meeting a fixed-prefix anchor}\}
                         =O_R(y)=O(y).
 \label{eq:bank-prefix-component-count}
\end{equation}

At any other fixed prefix row \(r\ge2\), the legal cofactors
\(r+1,\ldots,2r+1\) contain at least three integers.  At most two are the
incident component cores, so choose a third legal cofactor.  It need not
have had positive limiting allocation weight: by
\eqref{eq:bank-prefix-component-count}, adding only \(O(y)=o(N)\)
occurrences changes every fixed capacity by \(o(N)\).  Beyond the fixed
prefix, a promoted nonsmooth anchor has power-of-two smooth core, whereas
the ordinary descent uses non-power cores.  The only unavoidable core
\(4\) is covered by the forced row-\(2\) assignment; the \(4\to3\) and
\(3\to2\) bottom pools have no competing central occurrence.

Distinctness is now literal.  A prefix anchor has its unique marker
prime \(P>n/(R+1)>2R+1\); different components use different markers,
and for a fixed marker we selected a cofactor different from both state
cores.  Promoted anchors were already mutually distinct, and the
power-of-two-core observation separates them from ordinary states.

For each changed occurrence let \(a_i\) and \(b_i\) be its old and new
cofactors.  If \(D\) is written as the product of its individual cofactor
occurrences, deletion of the selected occurrences and insertion of their
replacements gives the integer
\begin{equation}
                         D':=D\prod_i\frac{b_i}{a_i}.
 \label{eq:bank-modified-divisor}
\end{equation}
This notation does not assert \(a_i\mid b_i\): integrality follows because
the displayed denominators cancel the corresponding factors in that
occurrence-level product.  Every \(b_i\) lies in a fixed-prefix legal
range, so all its prime divisors belong to \(F_{\mathrm{anc}}\).  Moreover,
by \eqref{eq:bank-prefix-component-count}, for each
\(f\in F_{\mathrm{anc}}\),
\begin{equation}
 |v_f(D')-v_f(D)|=O_{F_{\mathrm{anc}}}(y)=o(N).
 \label{eq:bank-anchor-change-cost}
\end{equation}
Replacing the corresponding anchors changes their product by exactly the
same ratio, and hence
\begin{equation}
 D'\mid T_n,\qquad
 \prod_{a\in\mathcal H'}a=C_nD'.
 \label{eq:bank-modified-anchor-identity}
\end{equation}
Here divisibility and the surviving positive linear reserve follow from
\eqref{eq:central-anchor-lemma-reserve} and the \(o(N)\) loss in
\eqref{eq:bank-anchor-change-cost}.  This proves the lemma.
\end{proof}

Use a different marker prime at every component, across paths, scales, and
orientations.  Because every cofactor is \(O(y)<P\), equality of two bank
factors forces equality of their marker and then of their core.  Donors
have \(y\)-smooth cofactor and the same unique marker, and each actual
donor is used once.  Donors exceed \(2n\), so they cannot equal a lower
endpoint or anchor.  Remove both possible states from later flexible
lists, and guard them from every fixed residual factor.  These choices
make anchors, bank states, donors, and all later fixed or flexible factors
jointly collision-free.

\subsection{Donor-backed row tokens and the base charge}
\label{subsec:bank-precharge}

For each ordinary component choose one donor counted by
\eqref{eq:bank-donor-multiplicity}; for each bottom component use the
donor in \eqref{eq:bank-four-bottom-pools}.  Remove it before a retained
list or row quota is frozen and replace its signature token by the chosen
base-state factor.  If the donor is an upper state, the upper orientation
keeps that occurrence and the lower orientation replaces it; it is never
used twice.

For each component the donor and either endpoint have the same complete
signature and contribute one factor to that row.  The adjusted quota is
therefore nonnegative and invariant under a toggle.  For a full path this
replacement is made separately in every component marker row, which proves
\eqref{eq:bank-all-factor-invariance} for the full states.

Fix every zero orientation and put
\begin{equation}
 G_{\mathrm{bank}}^0:=\bigcup_{g\in\mathcal B}G_g^0,
 \qquad
 B_{\mathrm{bank}}^0:=\prod_{a\in G_{\mathrm{bank}}^0}a.
 \label{eq:bank-base-product}
\end{equation}
Unused copies remain in this charged state.  Uniformly for \(p\le y\),
\eqref{eq:bank-total-size} gives
\begin{equation}
 v_p(B_{\mathrm{bank}}^0)
 \ll yL\frac{L}{\log p}
 =o\left(\frac{N}{pL}\right).
 \label{eq:bank-small-prime-profile}
\end{equation}
Also, uniformly in this range,
\begin{equation}
 v_p(T_n)=\frac{h}{p-1}+O\left(\frac{L}{\log p}\right)
         =\left(\frac{c}{p-1}+o(1)\right)N.
 \label{eq:bank-tail-small-prime-supply}
\end{equation}
For \(p\in F_{\mathrm{anc}}\), use the reserve left by
\eqref{eq:bank-anchor-change-cost}; outside that set, \(v_p(D')=0\).
For \(P>y\), the one-to-one donor assignment injects the complete rough
valuation of every base token into distinct factors of \(T_n\).  These
coordinatewise statements prove the simultaneous preliminary charge
\begin{equation}
                         D'B_{\mathrm{bank}}^0\mid T_n.
 \label{eq:preliminary-bank-divisibility}
\end{equation}
This concerns only the bank.  Section~\ref{sec:rough-selector} checks the
combined product with the retained exceptional factors before solving the
fractional equations; those and the bank states exhaust the fixed residual
token classes.

\begin{lemma}[Precharged universal bank]
\label[lemma]{lem:precharged-universal-bank}
There are modified anchors \(\mathcal H'\), a divisor \(D'\), and a
jointly guarded bank such that:
\begin{enumerate}
 \item \(\prod_{a\in\mathcal H'}a=C_nD'\) and
       \(D'B_{\mathrm{bank}}^0\mid T_n\);
 \item both states of every path have one token in each component marker
       row and satisfy \eqref{eq:bank-all-factor-invariance};
 \item every integer \(z\) supported on \(p\le y\), with
       \(|z_p|\le\beta_p\), is an exact collision-free bank change; and
 \item the fully reserved bank has \(o(1)\) utilization at each ordinary
       scale and in each of the four oriented bottom pools, and has
       \(O(yL)=o(N)\) component factors.
\end{enumerate}
Every donor and both states have been removed or guarded before later
retained lists and row quotas are defined.
\end{lemma}
 \section{The guarded rough-signature selector}
\label{sec:rough-selector}

We now construct the fractional point which fixes every nontrivial large-prime
signature.  Throughout this section
\[
 L=\log n,\qquad N=\frac nL,\qquad \theta=\frac29,
 \qquad y=n^\theta,
\]
and \(h,M,T=T_n\) are as in \eqref{eq:upper-basic-parameters} and
\eqref{eq:central-tail-split}.  The divisor supplied by the anchor and bank
construction will be denoted by \(D'\).  It is supported on the fixed set
\(F_{\rm anc}\), and the replacement of \(D\) by \(D'\) preserves a fixed
positive fraction of the reserve in
\eqref{eq:central-anchor-lemma-reserve}.

We fix the order of all choices used in this and the later correction
stages.  After \(c\), the prefix length \(R=R(c)\), the set
\(F_{\rm anc}\), and the modified anchor construction are fixed.  Next
choose constants
\begin{equation}
              1<r_0<\frac32,\qquad 1<\rho,\qquad \rho^3<r_0,
 \label{eq:global-ratio-parameter-order}
\end{equation}
for the common-list construction of \Cref{sec:tangent-absorber}.  Then
choose a sufficiently large nonprime real number
\(W>\max F_{\rm anc}\), before any tilt box or fine mesh is chosen, and
put
\[
 F_{\rm hd}:=\{p:p\le W\},\qquad
 P_{\rm hd}:=\prod_{p\in F_{\rm hd}}p,\qquad
 \delta_{\rm hd}:=\frac{\varphi(P_{\rm hd})}{P_{\rm hd}}.
\]
Thus \(F_{\rm anc}\subset F_{\rm hd}\), but these two sets have different
roles: only the former supports \(D'\), while the primes in
\(F_{\rm hd}\setminus F_{\rm anc}\) retain their full tail supply.  All
constants depending on \(W\) are fixed before \(n\) tends to infinity.

For precision, choose constants \(\gamma>0\), \(A_{\rm exc}<\infty\), and
\(C_{\rm sv}<\infty\), depending only on parameters already fixed, with
the following meanings.  For all sufficiently large \(n\), the valuation
available after the anchor charge is at least \(\gamma N/p\), uniformly
for every \(p\le y\): at \(p\in F_{\rm anc}\) this is the surviving
reserve from Lemma~\ref{lem:bank-anchor-modification}, and otherwise it follows
from \(v_p(T_n)=(c+o(1))N/(p-1)\).  The constant \(A_{\rm exc}\) is an
admissible constant in the nonnegative estimate
\eqref{eq:uniform-exceptional-charge} below, and \(C_{\rm sv}\) is an
admissible constant in the exceptional-list estimate
\eqref{eq:tangent-exceptional-loss}.  Their existence is proved without
using \(\delta_*\) except through the displayed linear factor.

We now choose \(\delta_*\) once, before defining any cutoff-dependent
set, so that
\begin{equation}
 0<\delta_*<
 \min\left\{
   \frac1{18},
   \frac{\theta\gamma}{4A_{\rm exc}},
   \frac{\theta\delta_{\rm hd}(2-r_0)}{4C_{\rm sv}}
 \right\},
 \qquad X_0=n^{\delta_*}.
 \label{eq:rough-cutoff}
\end{equation}
It is never decreased later.  After the remaining finite-dimensional
parameters (including a permitted fixed mesh and its tilt box) have been
chosen in the order specified in \Cref{sec:smooth-bridge}, there is a
threshold \(n_0\) depending on all those fixed choices.  No uniformity of
\(n_0\) over a varying family of meshes is asserted.  Throughout this
section, \(o_n(1)\) means that all parameters just listed are fixed first;
an \(O_W(1)\) constant may depend on the already fixed head cutoff but not
on \(n\), the later tilt-box radius, or a later mesh.
Equivalently, if \(\mathfrak P(\eta_0)\) denotes the permitted fixed
meshes of \Cref{sec:smooth-bridge}, including their subsequently fixed
finite-dimensional data, the quantifiers used for the upper construction
have the form
\begin{equation}
 \begin{gathered}
 \forall c>C_0\ \exists
 (R,F_{\rm anc},r_0,\rho,W,\delta_*)\ \exists(\eta_0,C_{\rm tan}),\\
 \forall\mathcal P\in\mathfrak P(\eta_0)\ \exists n_0(\mathcal P)\
 \forall n\ge n_0(\mathcal P):\\
 \text{all stated construction conclusions hold}.
 \end{gathered}
 \label{eq:upper-construction-nested-quantifiers}
\end{equation}
The bridge proves that \(\eta_0\) and \(C_{\rm tan}\) are independent of
the particular permitted fine mesh; only \(n_0\) is allowed to depend on
that mesh.  Display \eqref{eq:upper-construction-nested-quantifiers} is a
dependency convention for the assertions which follow, not an assertion
that one common threshold works for arbitrarily varying meshes.

For every positive integer \(a\le M\), write uniquely
\begin{equation}
 a=R_y(a)S_y(a),\qquad P^-(R_y(a))>y,\qquad
 P^+(S_y(a))\le y.
 \label{eq:complete-rough-decomposition}
\end{equation}
The integer \(R_y(a)\), including all its prime-power multiplicities, is
called the \emph{complete rough signature} of \(a\).  A signature \(R\) is
exceptional when
\begin{equation}
 X_R:=\frac{2n}{R}<X_0.
 \label{eq:exceptional-signature}
\end{equation}
This complete-signature convention is important: equality of row masses
then gives equality of every valuation at every prime greater than \(y\),
not merely equality of their squarefree supports.

\subsection{A head-compatible balanced point}
\label{subsec:head-compatible-point}

Let \(K>2\) be a fixed integer, to be chosen in a moment, and set
\[
 E=(2n,2n+h],\qquad H=(2n-Kh,2n],\qquad
 J=(n,2n-Kh].
\]
All endpoints are integers, so the number of integers in each half-open
interval equals its Lebesgue length.  Choose a small fixed \(\beta>0\),
then \(K\) sufficiently large, so that
\begin{equation}
 \alpha:=\frac{h/\delta_{\rm hd}-(\beta/L)|J|}{|H|}
 \label{eq:rough-alpha}
\end{equation}
lies, for all sufficiently large \(n\), in a fixed compact subinterval of
\((0,1/2)\).  This is possible because \(h=(c+o(1))n/L\): first take
\(\beta<c/(3\delta_{\rm hd})\), and then take a sufficiently large fixed
\(K\).  By definition,
\begin{equation}
 \delta_{\rm hd}\left(\alpha |H|+\frac\beta L|J|\right)=h
 \label{eq:head-balanced-normalization}
\end{equation}
exactly, not only asymptotically.  Define on \((n,2n]\)
\begin{equation}
 x_a^{\rm raw}:=\1_{(a,P_{\rm hd})=1}
   \left(\alpha\1_H(a)+\frac\beta L\1_J(a)\right).
 \label{eq:head-compatible-raw-point}
\end{equation}
Every flexible nonsmooth factor used below is therefore
\(F_{\rm hd}\)-free.  Split
\(
 \beta=\beta_{\rm prot}+\beta_{\rm act}
\)
with both summands fixed and positive.  On the smooth row, the two pieces
of \eqref{eq:head-compatible-raw-point} are recorded separately as
\begin{align}
 x_a^{\rm top}&=\alpha\1_H(a)\1_{(a,P_{\rm hd})=1}
                    \1_{P^+(a)\le y},\notag\\
 f_a^{\rm prot}&=\frac{\beta_{\rm prot}}L\1_J(a)
                    \1_{(a,P_{\rm hd})=1}\1_{P^+(a)\le y},\notag\\
 z_a^{\rm base}&=\frac{\beta_{\rm act}}L\1_J(a)
                    \1_{(a,P_{\rm hd})=1}\1_{P^+(a)\le y}.
 \label{eq:rough-smooth-splitting}
\end{align}
The first two summands will be frozen.  Only the last is redistributed by
the smooth bridge; \eqref{eq:rough-smooth-splitting} is an additive
decomposition of an already normalized point, not extra mass.  The
de Bruijn--Saias estimate with finite head inclusion--exclusion gives
\begin{equation}
 \sum_a z_a^{\rm base}
 =\{\delta_{\rm hd}\beta_{\rm act}\rho(1/\theta)+o_W(1)\}N
 \asymp_W N.
 \label{eq:rough-active-smooth-mass}
\end{equation}
Thus this summand supplies a macroscopic active measure despite its
pointwise \(1/L\) scale.

There is already strong cancellation before rows are separated.  If
\(p>W\), finite inclusion--exclusion gives, for every physical interval
\(I\) and every \(k\ge1\),
\begin{equation}
 \#\{a\in I:p^k\mid a,\ (a,P_{\rm hd})=1\}
 =\delta_{\rm hd}\frac{|I|}{p^k}+O_W(1).
 \label{eq:rough-head-crt}
\end{equation}
The main terms for \(\1_E-x^{\rm raw}\) cancel by
\eqref{eq:head-balanced-normalization}.  Hence
\begin{equation}
 \left|\sum_a(\1_E(a)-x_a^{\rm raw})v_p(a)\right|
 \ll_W 1+\frac L{\log p}
 =o\!\left(\frac{N}{pL}\right)
 \quad(W<p\le y).
 \label{eq:rough-global-crt}
\end{equation}
The last estimate is uniform up to \(p=y\).  The work below shows that
charging exceptional rows and making all other row quotas exact costs no
more than the right side of \eqref{eq:rough-global-crt} at its natural
scale.

\subsection{The transition-complete balanced-block estimate}
\label{subsec:transition-balanced-block}

We retain the de Bruijn--Saias normalization from
\Cref{subsec:smooth-input}, in the form proved by
Hildebrand--Tenenbaum and Saias
\cite{HildebrandTenenbaum1993,Saias1989}.  For nonintegral \(z>0\),
integration by
parts in \eqref{eq:Lambda-definition} gives the exact identity
\begin{align}
 \Lambda(z,y)&=zG_y(u)-\{z\},\qquad
 u=\frac{\log z}{\log y},\notag\\
 G_y(u)&=\rho(u)-\int_0^{(u-1)_+}
       \rho'(u-v)\{y^v\}y^{-v}\,dv.
 \label{eq:rough-G-normalization}
\end{align}
At integral \(z\) we use the right-limit convention of
\Cref{thm:HT-Saias}.  For \(0<u\le1\), the formula reduces to
\(\Lambda(z,y)=\lfloor z\rfloor\).
Every later invocation in this section has \(y=n^\theta\) and
\(z=t x/d\), where \(t\) ranges over a fixed compact subset of
\((0,\infty)\), \(d\mid P_{\rm hd}\) is fixed, and
\(n^{\delta_*}/C\le x\le Cn\).  Hence
\(\log z/\log y\) remains in a fixed compact subset of \((0,\infty)\),
and the compact-\(u\) error in \Cref{thm:HT-Saias} applies uniformly.
The stated right-limit convention covers all integer physical endpoints.

\begin{lemma}[Uniform transition normalization]
\label[lemma]{lem:rough-G-Lipschitz}
On every fixed compact interval \(0<u_-\le u\le u_+\),
\[
 |G_y(u)-G_y(v)|\ll_{u_-,u_+}|u-v|
\]
uniformly for \(y\ge2\).
\end{lemma}

\begin{proof}
The Dickman function is Lipschitz on compact intervals.  Extend
\(\rho'\) by zero to the left of \(1\).  On a fixed compact interval this
extension has bounded variation (including its single jump at \(1\)).
For a compactly supported function \(q\) of bounded variation,
\[
 \int_{\mathbb R}|q(t+s)-q(t)|\,dt
 \le \operatorname{Var}(q)|s|.
\]
Apply this to the translated copies of \(\rho'\) in
\eqref{eq:rough-G-normalization}, using
\(0\le\{y^v\}y^{-v}\le1\).  The moving endpoint is included by the
zero extension.  This proves the assertion also at the transition
\(u=1\).
\end{proof}

\begin{lemma}[Balanced smooth blocks]
\label[lemma]{lem:balanced-smooth-blocks}
Let \(x\ge n^{\delta_*}/C\).  Let a fixed number of \(t_i\)'s range in a
fixed compact subset of \((0,\infty)\), and let \(a_i=O(1)\).  If
\begin{equation}
 \sum_i a_it_i=0,
 \qquad
 \sum_i|a_i|t_i|\log t_i|\ll\frac1L,
 \label{eq:balanced-block-hypotheses}
\end{equation}
then
\begin{equation}
 \sum_i a_i\Psi(t_ix,y)=O(x/L^2+1).
 \label{eq:balanced-block-conclusion}
\end{equation}
The same estimate holds with the smooth integers restricted to be
coprime to \(P_{\rm hd}\).
\end{lemma}

\begin{proof}
Put \(u=\log x/\log y\) and
\(d_i=\log t_i/\log y\).  Theorem~\ref{thm:HT-Saias},
\eqref{eq:rough-G-normalization}, and the exact formula below \(y\) give,
uniformly across every Dickman transition,
\[
 \Psi(t_ix,y)=t_ixG_y(u+d_i)
   +O\!\left(1+x\exp\{-cL^{1/2}\}\right).
\]
Subtract \(t_ixG_y(u)\), sum, and use
\Cref{lem:rough-G-Lipschitz} and
\eqref{eq:balanced-block-hypotheses}.  This gives
\eqref{eq:balanced-block-conclusion}; the exponentially small term is
absorbed because \(x\ge n^{\delta_*}/C\).  For the coprime version use
the exact finite identity
\[
 \1_{(m,P_{\rm hd})=1}
 =\sum_{d\mid P_{\rm hd}}\mu(d)\1_{d\mid m}
\]
and apply the first assertion at \(x/d\).  This proof retains every
endpoint \(O(1)\), and differentiates neither \(\rho\) nor \(\Psi\) at a
transition.
\end{proof}

For later use isolate the effect of a fixed head divisor on the entire
physical block.  Put
\begin{equation}
 B_{\rm ph}(x):=
 \alpha\{\Psi(x,y)-\Psi((1-K\kappa)x,y)\}
 +\frac{\beta}{L}\{\Psi((1-K\kappa)x,y)-\Psi(x/2,y)\}.
 \label{eq:rough-physical-block}
\end{equation}

\begin{lemma}[Fixed-head shift of a physical block]
\label[lemma]{lem:fixed-head-physical-shift}
Uniformly for \(n^{\delta_*}/C\leq x\leq Cn\) and every fixed
\(d\mid P_{\rm hd}\),
\begin{equation}
 B_{\rm ph}(x/d)=\frac1dB_{\rm ph}(x)
                 +O_W(x/L^2+1).
 \label{eq:fixed-head-physical-shift}
\end{equation}
\end{lemma}

\begin{proof}
Write \(b=1-K\kappa\), \(u=\log x/\log y\),
\(e=\log d/\log y\), and
\[
 \Delta_e(s):=G_y(s-e)-G_y(s).
\]
The fixedness of \(d\), \Cref{lem:rough-G-Lipschitz}, and the triangle
inequality give, on the relevant compact parameter interval,
\begin{equation}
 |\Delta_e(s)|\ll_W L^{-1},\qquad
 |\Delta_e(s)-\Delta_e(t)|\ll_W |s-t|.
 \label{eq:rough-head-shift-modulus}
\end{equation}
The transition normalization used in the proof of
\Cref{lem:balanced-smooth-blocks}, with every endpoint floor retained,
therefore gives
\begin{align}
 B_{\rm ph}(x/d)-d^{-1}B_{\rm ph}(x)
 &=\frac{x}{d}\left[
   \alpha\{\Delta_e(u)-b\Delta_e(u+\log b/\log y)\}\right.\notag\\
 &\hspace{26mm}\left.
  +\frac\beta L\{b\Delta_e(u+\log b/\log y)
       -\tfrac12\Delta_e(u-\log2/\log y)\}\right]\notag\\
 &\quad+O_W(xe^{-cL^{1/2}}+1).
 \label{eq:rough-head-shift-expansion}
\end{align}
Here \(1-b=K\kappa=O(L^{-1})\) and
\(|\log b|/\log y=O(L^{-2})\).  Rewrite the first brace as
\[
 (1-b)\Delta_e(u)
 +b\{\Delta_e(u)-\Delta_e(u+\log b/\log y)\}.
\]
Both of its terms are \(O_W(L^{-2})\) by
\eqref{eq:rough-head-shift-modulus}; the second brace in
\eqref{eq:rough-head-shift-expansion} is \(O_W(L^{-1})\) and already
has the prefactor \(L^{-1}\).  The exponential error is
\(O(x/L^2)\) in the stated range.  This proves
\eqref{eq:fixed-head-physical-shift}, including the \(O(1)\) endpoint
terms.
\end{proof}

Fix a nonexceptional signature \(R\), put \(x=X_R\), and write
\(\kappa=h/(2n)\).  The exact upper-row count is
\[
 t_R=\Psi((1+\kappa)x,y)-\Psi(x,y).
\]
The mass of \eqref{eq:head-compatible-raw-point} in this row is the
corresponding \(P_{\rm hd}\)-coprime smooth count in
\(((1-K\kappa)x,x]\) with coefficient \(\alpha\), plus that in
\((x/2,(1-K\kappa)x]\) with coefficient \(\beta/L\).  Dividing
\eqref{eq:head-balanced-normalization} by \(R\) gives the exact real-length
identity
\begin{equation}
 \delta_{\rm hd}\left\{\alpha K\kappa x
   +\frac\beta L(1/2-K\kappa)x\right\}=\kappa x.
 \label{eq:row-real-length-balance}
\end{equation}
The head inclusion--exclusion is handled before invoking the balanced
lemma.  Exactly from the definition of \(B_{\rm ph}\),
\[
 m_R^{\rm raw}
 =\sum_{d\mid P_{\rm hd}}\mu(d)B_{\rm ph}(x/d).
\]
By \Cref{lem:fixed-head-physical-shift},
\begin{equation}
 m_R^{\rm raw}=\delta_{\rm hd}B_{\rm ph}(x)
                   +O_W(x/L^2+1).
 \label{eq:rough-head-shift-summed}
\end{equation}
It remains to compare the unscaled upper block with
\(\delta_{\rm hd}B_{\rm ph}(x)\).  Its endpoint pairs are
\[
 \Psi((1+\kappa)x,y)-\Psi(x,y)-\delta_{\rm hd}B_{\rm ph}(x).
\]
Their coefficients \(a_i\) and relative endpoints \(t_i\) have
\(\sum_i a_it_i=0\) by \eqref{eq:row-real-length-balance}.  Moreover,
the endpoints \(1+\kappa\) and \(1-K\kappa\) contribute
\(O(\kappa)=O(L^{-1})\) to
\(\sum_i|a_i|t_i|\log t_i|\), the endpoint \(1\) contributes zero,
and the fixed endpoint \(1/2\) has coefficient
\(\delta_{\rm hd}\beta/L\).  Thus the absolute-log hypothesis
\eqref{eq:balanced-block-hypotheses} really does hold for this unscaled
comparison.  Applying \Cref{lem:balanced-smooth-blocks} and then
\eqref{eq:rough-head-shift-summed}, with
\(m_R^{\rm raw}=\sum_{R_y(a)=R}x_a^{\rm raw}\),
\begin{equation}
 \boxed{\ \varepsilon_R:=t_R-m_R^{\rm raw}
      =O_W(X_R/L^2+1)\ }
 \qquad(X_R\ge X_0).
 \label{eq:active-row-error}
\end{equation}
The \(+1\) in \eqref{eq:active-row-error} includes the integer endpoint
error of every one of the finitely many head inclusion--exclusion terms.
The same argument over the fixed-ratio interval \(J/R\), together with
the positivity of \(G_y(u)=\rho(u)+O(1/L)\) on the relevant compact
range, gives
\begin{equation}
 \#\{a\in J:R_y(a)=R,\ (a,P_{\rm hd})=1\}\gg_W X_R.
 \label{eq:active-clean-pool}
\end{equation}

\subsection{The fixed-depth exceptional expansion}
\label{subsec:exceptional-expansion}

The estimate needed for exceptional rows lies strictly between the fourth
and fifth births of a rough integer.  We give the short proof because it
is what permits the signed sum below; a rowwise sawtooth estimate would not
be uniform at the one-prime boundary.

For an interval \(I\), write
\[
 \Phi(I;y):=\#\{r\in I:P^-(r)>y\}.
\]

\begin{lemma}[Four-to-five rough chamber]
\label[lemma]{lem:four-five-rough-chamber}
For each fixed \(C\geq1\), let \(Z\) satisfy
\[
 4.2\le \frac{\log Z}{\log y}\le4.6,
\]
and let \(I\subset[Z/C,CZ]\cap(0,3n]\) have arbitrary real endpoints.
There is a \(C^1\) function \(\mathcal K\) on \([4.1,4.7]\), with bounded
derivative, such that
\begin{equation}
 \Phi(I;y)=\int_I\frac{\mathcal K(\log t/\log y)}{\log y}\,dt
       +O_C\!\left(Ze^{-c\sqrt L}+\sqrt Z\right).
 \label{eq:four-five-rough-asymptotic}
\end{equation}
\end{lemma}

\begin{proof}
Changing the inclusion convention at either endpoint of \(I\) changes
\(\Phi(I;y)\) by at most two, which is absorbed by the claimed error.
We may therefore write \(I=(A,B]\).  For \(1\le j\le4\), define the
ordered count
\begin{equation}
 \mathcal N_j(I;y):=
 \#\{(p_1,\ldots,p_j):p_i>y\ {\rm prime},\
                         p_1\cdots p_j\in I\}.
 \label{eq:rough-ordered-count}
\end{equation}

We first record the reciprocal-prime discrepancy with all constants
uniform on the compact range used below.  On \([1,4.8]\), put
\[
 d\mu_y(s):=\sum_{p>y}\frac1p
      \delta_{\log p/\log y}(ds),
 \qquad d\lambda(s):=\frac{ds}{s}.
\]
For \(1\le T\le4.8\), partial summation gives
\begin{equation}
 \mu_y([1,T])
 =\frac{\pi(y^T)}{y^T}-\frac{\pi(y)}y
    +\int_y^{y^T}\frac{\pi(v)}{v^2}\,dv.
 \label{eq:rough-reciprocal-partial-summation}
\end{equation}
Write \(E(v)=\pi(v)-\operatorname{li}(v)\).  The
\(\operatorname{li}\)-part of
\eqref{eq:rough-reciprocal-partial-summation} is exactly
\(\int_y^{y^T}dv/(v\log v)=\log T\).  By
\eqref{eq:cumulative-PNT}, uniformly in \(T\),
\begin{align}
 &\left|\frac{E(y^T)}{y^T}\right|
  +\left|\frac{E(y)}y\right|
  +\int_y^{y^{4.8}}\frac{|E(v)|}{v^2}\,dv\notag\\
 &\quad\ll e^{-c_0\sqrt{\log y}}
   +\int_y^{y^{4.8}}e^{-c_0\sqrt{\log v}}\frac{dv}{v}
 \ll L e^{-c_0\sqrt{(2/9)L}}
 \ll e^{-c_1\sqrt L}.
 \label{eq:rough-reciprocal-error-sum}
\end{align}
Consequently
\begin{equation}
 D_y:=\sup_{1\le T\le4.8}
 \left|\mu_y([1,T])-\lambda([1,T])\right|
 \ll e^{-c_1\sqrt L}.
 \label{eq:exceptional-Stieltjes-PNT}
\end{equation}
The same bound holds for either convention at the moving endpoint:
changing one atom costs at most \(1/y\ll e^{-c_1\sqrt L}\).
In particular, if \(1\le a\le4.8\) and \(f\) has bounded variation on
\([1,a]\), Stieltjes integration by parts gives
\begin{equation}
 \left|\int_{[1,a]}f\,d(\mu_y-\lambda)\right|
 \le2(D_y+y^{-1})
       \{\|f\|_\infty+\operatorname{Var}_{[1,a]}f\}.
 \label{eq:rough-one-dimensional-BV}
\end{equation}

We next expose every moving endpoint in the ordered count.  Fix \(j\),
put \(m=j-1\), and write \(q=p_1\cdots p_m\), with the empty product
interpreted as one.  If \(B/q\le y\), this tuple contributes nothing.
Otherwise put \(a_q=\max(y,A/q)\).  The two nonempty endpoints
\(a_q,B/q\) lie between \(y=n^{2/9}\) and \(CZ\le C y^{4.6}\).
Subtracting \eqref{eq:cumulative-PNT} at these exact endpoints gives
\begin{align}
 \#\{p_j>y:qp_j\in(A,B]\}
 &=\pi(B/q)-\pi(a_q)\notag\\
 &=\int_{a_q}^{B/q}\frac{dv}{\log v}
   +O_C\!\left(\frac Zq e^{-c_2\sqrt L}\right)\notag\\
 &=\int_A^B\frac{\1_{qy<t}}{q\log(t/q)}\,dt
   +O_C\!\left(\frac Zq e^{-c_2\sqrt L}\right).
 \label{eq:exceptional-last-prime}
\end{align}
This is an absolute cumulative-PNT error, not a relative assertion for a
short interval.  Moreover,
\begin{equation}
 \sum_{\substack{p_1,\ldots,p_m>y\\q<B/y}}\frac1q
 \le\left(\sum_{y<p\le CZ/y}\frac1p\right)^m
 \ll_{C,m}1
 \label{eq:rough-reciprocal-tuple-mass}
\end{equation}
by \eqref{eq:Mertens-fixed-range}.  Thus the total error in
\eqref{eq:exceptional-last-prime} is
\(O_{C,j}(Ze^{-c_2\sqrt L})\), uniformly in both endpoints.

Summing the main term in \eqref{eq:exceptional-last-prime} and
interchanging the finite sum with the \(t\)-integral gives
\begin{equation}
 \mathcal N_j(I;y)
 =\int_I S_{j,y}(t)\,dt
   +O_{C,j}(Ze^{-c_2\sqrt L}).
 \label{eq:rough-ordered-after-last-prime}
\end{equation}
If \(u=\log t/\log y\), then
\begin{equation}
 S_{j,y}(t)=\frac1{\log y}
 \int_{\substack{s_1,\ldots,s_m\ge1\\
                  s_1+\cdots+s_m<u-1}}
 \frac{d\mu_y(s_1)\cdots d\mu_y(s_m)}
      {u-s_1-\cdots-s_m};
 \label{eq:rough-exact-product-Stieltjes}
\end{equation}
for \(m=0\), the last integral means \(1/u\).  The strict face causes no
outer-integral error: for each fixed prime tuple it is a single value of
\(t\).

For \(t\in[Z/C,CZ]\), the hypotheses on \(Z\) imply, for all sufficiently
large \(n\) depending only on \(C\),
\begin{equation}
                         4.1\le u\le4.7.
 \label{eq:rough-padded-u-range}
\end{equation}
Replace the \(m\le3\) copies of \(d\mu_y\) in
\eqref{eq:rough-exact-product-Stieltjes} one at a time by \(d\lambda\).
Conditioned on the other variables, the current variable ranges over
\([1,a]\), where
\(
 a=u-1-\sum_{\ell\ne i}s_\ell
\), or over the empty set.  On a nonempty interval the conditional
kernel
\[
 f(s)=\frac1{u-s-\sum_{\ell\ne i}s_\ell}
\]
has denominator at least one, and hence
\begin{equation}
 \|f\|_\infty\le1,\qquad \operatorname{Var}f\le1.
 \label{eq:rough-conditional-BV}
\end{equation}
For large \(n\), both \(\mu_y\) and \(\lambda\) have mass at most
\(M_0:=1+\log4.8\) on the relevant compact interval.  Therefore the
literal \(m\)-term telescope and
\eqref{eq:rough-one-dimensional-BV} give
\begin{align}
 &\left|
 \int_{\sum s_i<u-1}
       \frac{d\mu_y(s_1)\cdots d\mu_y(s_m)}{u-\sum s_i}
 -\int_{\sum s_i\le u-1}
       \frac{ds_1\cdots ds_m}
        {s_1\cdots s_m(u-\sum s_i)}
 \right|\notag\\
 &\hspace{35mm}
 \le4mM_0^{m-1}(D_y+y^{-1})
 \ll_j e^{-c_1\sqrt L}.
 \label{eq:rough-product-measure-telescope}
\end{align}
This includes a moving or degenerate simplex face with one constant.

Define
\begin{align}
 \mathcal K_1(u)&=\frac1u,\notag\\
 \mathcal K_j(u)&=
 \int_{\substack{s_1,\ldots,s_{j-1}\ge1\\
                  s_1+\cdots+s_{j-1}\le u-1}}
 \frac{ds_1\cdots ds_{j-1}}
 {s_1\cdots s_{j-1}(u-s_1-\cdots-s_{j-1})}
 \quad(2\le j\le4).
 \label{eq:exceptional-Kj}
\end{align}
Equations \eqref{eq:rough-ordered-after-last-prime} and
\eqref{eq:rough-product-measure-telescope} prove
\begin{equation}
 \mathcal N_j(I;y)
 =\frac1{\log y}\int_I
       \mathcal K_j(\log t/\log y)\,dt
   +O_{C,j}(Ze^{-c\sqrt L}).
 \label{eq:rough-ordered-asymptotic}
\end{equation}

For completeness, the regularity of the kernels is also uniform.  For
\(j\ge2\), put \(m=j-1\), \(T=u-j\), and
\(\Delta_m=\{z_i\ge0:\sum z_i\le1\}\).  The substitutions
\(s_i=1+Tz_i\) give
\begin{equation}
 \mathcal K_j(u)=T^m\int_{\Delta_m}
 \frac{dz_1\cdots dz_m}
 {\prod_{i=1}^m(1+Tz_i)
  \{1+T(1-\sum_i z_i)\}}.
 \label{eq:rough-kernel-fixed-simplex}
\end{equation}
On \(4.1\le u\le4.7\) and \(2\le j\le4\), one has \(T\ge0.1\)
and \(T\ll1\).  The integrand and its first \(T\)-derivative are bounded
on the fixed simplex, so differentiation under the integral proves
\begin{equation}
 \sup_{4.1\le u\le4.7}
 \{|\mathcal K_j(u)|+|\mathcal K_j'(u)|\}\ll_j1.
 \label{eq:rough-kernel-C1}
\end{equation}
The same assertion is immediate for \(\mathcal K_1\).

It remains to pass from ordered tuples to rough integers.  Put
\(\Omega(r)=\sum_pv_p(r)\), and let \(A_j(I;y)\) count the \(r\in I\)
with \(P^-(r)>y\) and \(\Omega(r)=j\).  Since
\(y^5=n^{10/9}>3n\), every integer counted by \(\Phi(I;y)\) has
\(1\le\Omega(r)\le4\).  If
\[
                  r=\prod_{\nu=1}^kq_\nu^{e_\nu},
 \qquad            \sum_{\nu=1}^ke_\nu=j,
\]
then its exact multiplicity in \(\mathcal N_j\) is
\begin{equation}
                         \frac{j!}{e_1!\cdots e_k!}.
 \label{eq:rough-exact-ordered-multiplicity}
\end{equation}
Thus division by \(j!\) gives weight one for a squarefree \(r\), and
weight \(1/(e_1!\cdots e_k!)\in(0,1]\) otherwise.  Consequently
\begin{equation}
 0\le\Phi(I;y)-\sum_{j=1}^4\frac{\mathcal N_j(I;y)}{j!}
 \le\#\{r\in I:q^2\mid r\text{ for some prime }q>y\}.
 \label{eq:rough-ordered-unordered-discrepancy}
\end{equation}
The last quantity is at most
\begin{align}
 \sum_{\substack{y<q\le\sqrt{CZ}\\q\ {\rm prime}}}
       \left(\frac{|I|}{q^2}+1\right)
 &\ll_C\frac{|I|}{y\log y}+\frac{\sqrt Z}{\log Z}\notag\\
 &\ll_CZe^{-c\sqrt L}+\sqrt Z.
 \label{eq:rough-repeated-prime-total}
\end{align}
Here the first bound follows by partial summation from
\(\pi(t)\ll t/\log t\), and the last uses
\(y^{-1}\ll e^{-c\sqrt L}\).  Combining
\eqref{eq:rough-ordered-asymptotic} and
\eqref{eq:rough-ordered-unordered-discrepancy} proves
\eqref{eq:four-five-rough-asymptotic} with
\[
                 \mathcal K(u)=\sum_{j=1}^4\frac{\mathcal K_j(u)}{j!}.
\]
\end{proof}

For later use we record explicitly the simultaneous form just proved.
For fixed \(J_0,G\), let \(J\le J_0\), let
\(I_\nu\subset[Z/C,CZ]\cap(0,3n]\) be intervals, possibly empty or
clipped at arbitrary moving endpoints, and let \(|\gamma_\nu|\le G\).
Then the same kernel and one common implied constant give
\begin{equation}
 \sum_{\nu=1}^J\gamma_\nu\Phi(I_\nu;y)
 =\frac1{\log y}\sum_{\nu=1}^J\gamma_\nu
       \int_{I_\nu}\mathcal K(\log t/\log y)\,dt
 +O_{C,J_0,G}(Ze^{-c\sqrt L}+\sqrt Z).
 \label{eq:four-five-rough-finite-family}
\end{equation}

Put \(R_0=2n/X_0\), and let
\[
 w_{\rm hd}(a):=\1_E(a)-x_a^{\rm raw}.
\]
For \(W<p\le y\), the exact signed contribution of the exceptional rows
to the uncharged comparison is
\begin{equation}
 \mathcal E_p=
 \sum_{a:R_y(a)>R_0}w_{\rm hd}(a)v_p(a)
 =\sum_{b<2X_0}v_p(b)
    \sum_{\substack{P^-(R)>y\\R>R_0}}w_{\rm hd}(bR).
 \label{eq:exceptional-core-first}
\end{equation}
Indeed, \(a=Rb\) in this range has \(b<2X_0<y\), and so
\(v_p(a)=v_p(b)\).  Moreover, throughout the support in the inner sum,
\begin{equation}
 4.25+o(1)<\frac{\log R}{\log y}<4.51.
 \label{eq:exceptional-four-prime-range}
\end{equation}
The lower bound follows from
\((1-\delta_*)/\theta>17/4\), and the upper bound from \(R<3n\).
Thus \Cref{lem:four-five-rough-chamber} applies with a fixed separation
from both birth faces.

Define the mean-zero periodic function
\begin{equation}
 c_{\rm hd}(b):=1-\delta_{\rm hd}^{-1}
                       \1_{(b,P_{\rm hd})=1}.
 \label{eq:rough-periodic-core}
\end{equation}
For every \(b<2X_0\), put
\[
 Z_b:=\frac{2n}{b},\qquad
 u_b:=\frac{\log Z_b}{\log y},\qquad
 g_b:=\1_{(b,P_{\rm hd})=1},
\]
and introduce the three disjoint physical intervals
\begin{align*}
 I_{+,b}&=(Z_b,(1+\kappa)Z_b],\\
 I_{-,b}&=((1-K\kappa)Z_b,Z_b],\\
 I_{0,b}&=(Z_b/2,(1-K\kappa)Z_b].
\end{align*}
Since \(b<2X_0<y\), the complete rough factor is coprime to
\(P_{\rm hd}\).  The definitions of \(E,H,J\) and
\eqref{eq:head-compatible-raw-point} therefore give the exact identity
\begin{equation}
 \sum_{P^-(R)>y}w_{\rm hd}(bR)
 =\Phi(I_{+,b};y)-g_b\alpha\Phi(I_{-,b};y)
       -g_b\frac\beta L\Phi(I_{0,b};y).
 \label{eq:rough-three-physical-intervals}
\end{equation}

Suppose first that \(b\le X_0/2\).  The cutoff \(R>R_0\) in
\eqref{eq:exceptional-core-first} is then automatic.  All three intervals
lie in \([Z_b/2,2Z_b]\cap(0,3n]\), and the range in
\eqref{eq:exceptional-four-prime-range} puts \(u_b\) in the compact
chamber of \Cref{lem:four-five-rough-chamber}.  With
\[
 \mathfrak E_b:=Z_be^{-c\sqrt L}+\sqrt{Z_b},
\]
the simultaneous estimate
\eqref{eq:four-five-rough-finite-family}, applied once to the three terms
of \eqref{eq:rough-three-physical-intervals}, yields
\begin{align}
 \sum_{P^-(R)>y}w_{\rm hd}(bR)
 &=\frac1{\log y}\left\{
   \int_{I_{+,b}}\mathcal K\!\left(\frac{\log t}{\log y}\right)dt
  -g_b\alpha\int_{I_{-,b}}
       \mathcal K\!\left(\frac{\log t}{\log y}\right)dt\right.\notag\\
 &\hspace{31mm}\left.
  -g_b\frac\beta L\int_{I_{0,b}}
       \mathcal K\!\left(\frac{\log t}{\log y}\right)dt
  \right\}+O_W(\mathfrak E_b).
 \label{eq:rough-simultaneous-physical-expansion}
\end{align}
This is one endpoint-uniform invocation, rather than three asymptotics
with potentially different moving-endpoint constants.

Put \(K_b=\mathcal K(u_b)\).  By the mean-value theorem and the uniform
bound for \(\mathcal K'\), the error made by replacing every kernel in
\eqref{eq:rough-simultaneous-physical-expansion} by \(K_b\) is at most
\begin{align}
 &\frac1{\log y}\left\{
   \int_{I_{+,b}}\left|\mathcal K\!\left(\frac{\log t}{\log y}\right)-K_b\right|dt
  +\alpha\int_{I_{-,b}}
       \left|\mathcal K\!\left(\frac{\log t}{\log y}\right)-K_b\right|dt\right.\notag\\
 &\hspace{29mm}\left.
  +\frac\beta L\int_{I_{0,b}}
       \left|\mathcal K\!\left(\frac{\log t}{\log y}\right)-K_b\right|dt
  \right\}
 \ll_W\frac{Z_b}{(\log y)^2}
       \left(\kappa^2+K^2\kappa^2+\frac1L\right)
 \ll_W\frac{Z_b}{L^3}.
 \label{eq:rough-physical-kernel-variation}
\end{align}
Indeed, the two short intervals have respectively relative lengths
\(\kappa\) and \(K\kappa\), while
\(|\log(t/Z_b)|\le\log2\) on the broad interval; also
\(\kappa\ll L^{-1}\).

The constant-kernel coefficient has no asymptotic error.  Dividing
\eqref{eq:head-balanced-normalization} by \(b\), or equivalently using
\eqref{eq:row-real-length-balance} with \(x=Z_b\), gives
\begin{align}
 &\kappa Z_b-g_b\left\{\alpha K\kappa Z_b
       +\frac\beta L(1/2-K\kappa)Z_b\right\}\notag\\
 &\hspace{25mm}=\kappa Z_b\left(1-\frac{g_b}{\delta_{\rm hd}}\right)
 =\frac hb\,c_{\rm hd}(b).
 \label{eq:rough-exact-physical-coefficient}
\end{align}
Finally, \(Z_b\ge4n/X_0\) in the present range, so
\begin{equation}
 \mathfrak E_b\ll Z_b/L^3.
 \label{eq:rough-chamber-error-absorbed}
\end{equation}
Combining
\eqref{eq:rough-simultaneous-physical-expansion}--
\eqref{eq:rough-chamber-error-absorbed}, and retaining a harmless unit
endpoint allowance, proves
\begin{align}
 \sum_{P^-(R)>y}w_{\rm hd}(bR)
 &=\frac{h}{b\log y}
   \mathcal K\!\left(\frac{\log(2n/b)}{\log y}\right)c_{\rm hd}(b)
   +O_W\!\left(\frac{n}{bL^3}+1\right).
 \label{eq:exceptional-deep-core}
\end{align}

For \(p>W\), multiplication by \(p^k\) permutes the residue classes
modulo \(P_{\rm hd}\).  Hence
\[
 \sup_X\left|\sum_{m\le X}c_{\rm hd}(p^km)\right|\ll_W1.
\]
Partial summation, followed by the expansion
\(v_p(b)=\sum_{k\ge1}\1_{p^k\mid b}\), shows that for every bounded
\(C^1\) function \(V\) of \(\log b/\log y\),
\begin{equation}
 \sum_{b\le X}\frac{c_{\rm hd}(b)v_p(b)}b
       V\!\left(\frac{\log b}{\log y}\right)
 \ll_W\frac1p.
 \label{eq:periodic-prime-power-sum}
\end{equation}
For the single cutoff band \(X_0/2<b<2X_0\), set
\[
 I_{\nu,b}^{\circ}:=I_{\nu,b}\cap(R_0,\infty)
 \qquad(\nu\in\{+,-,0\}).
\]
The three physical intervals are disjoint, so taking the absolute value
of each coefficient is exact:
\begin{align}
 &\sum_{\substack{P^-(R)>y\\R>R_0}}|w_{\rm hd}(bR)|\notag\\
 &\quad=\Phi(I_{+,b}^{\circ};y)
      +g_b\alpha\Phi(I_{-,b}^{\circ};y)
      +g_b\frac\beta L\Phi(I_{0,b}^{\circ};y).
 \label{eq:rough-cutoff-three-intervals}
\end{align}
Every clipped interval is still contained in
\([Z_b/2,2Z_b]\cap(0,3n]\).  Applying
\eqref{eq:four-five-rough-finite-family} simultaneously, now with
positive coefficients, gives its main term bounded by
\begin{align}
 \frac{\|\mathcal K\|_\infty}{\log y}
 \left\{|I_{+,b}^{\circ}|+\alpha|I_{-,b}^{\circ}|
       +\frac\beta L|I_{0,b}^{\circ}|\right\}
 &\le\frac{\|\mathcal K\|_\infty}{\log y}
 \left\{\kappa Z_b+\alpha K\kappa Z_b
       +\frac{\beta Z_b}{2L}\right\}\notag\\
 &\ll_W\frac{Z_b}{L^2}.
 \label{eq:rough-cutoff-weighted-length}
\end{align}
Uniformly in this band, \(Z_b\asymp n/X_0\), and hence the common
analytic error \(Z_be^{-c\sqrt L}+\sqrt{Z_b}\) is
\(o(Z_b/L^2)\).  Since \(Z_b=2n/b\), we obtain, with the existing
harmless endpoint allowance,
\begin{equation}
 \sum_{\substack{P^-(R)>y\\R>R_0}}|w_{\rm hd}(bR)|
 \ll_W\frac{n}{bL^2}+1.
 \label{eq:exceptional-cutoff-band}
\end{equation}

Finally, the elementary prime-power identities
\begin{align}
 \sum_{b\le X}\frac{v_p(b)}b&\ll\frac{\log X+1}{p},
 &\sum_{b\le X}v_p(b)&\ll\frac Xp,\notag\\
 \sum_{X/2<b<2X}\frac{v_p(b)}b&\ll\frac1p
 \label{eq:exceptional-prime-power-sums}
\end{align}
make the last summation completely explicit.  The function
\(b\mapsto\mathcal K(\log(2n/b)/\log y)\) is a bounded \(C^1\)
function of \(\log b/\log y\).  Hence
\eqref{eq:periodic-prime-power-sum} and \(h\ll n/L\) give
\begin{equation}
 \frac h{\log y}\left|
  \sum_{b\le X_0/2}\frac{c_{\rm hd}(b)v_p(b)}b
   \mathcal K\!\left(\frac{\log(2n/b)}{\log y}\right)\right|
 \ll_W\frac{n}{pL^2}.
 \label{eq:exceptional-leading-ledger}
\end{equation}
The accumulated error in the deep range is, with no suppressed logarithm,
\begin{align}
 \sum_{b\le X_0/2}v_p(b)
      \left(\frac n{bL^3}+1\right)
 &\ll \frac n{L^3}\frac{\log X_0+1}{p}+\frac{X_0}{p}
 \ll\frac{n}{pL^2}+\frac{X_0}{p}.
 \label{eq:exceptional-deep-error-ledger}
\end{align}
Likewise, the entire cutoff band contributes
\begin{align}
 \sum_{X_0/2<b<2X_0}v_p(b)
      \left(\frac n{bL^2}+1\right)
 &\ll\frac{n}{pL^2}+\frac{X_0}{p}.
 \label{eq:exceptional-cutoff-error-ledger}
\end{align}
Finally,
\begin{equation}
 \frac{X_0}{p}=o\!\left(\frac{n}{pL^2}\right)
 \label{eq:exceptional-X0-absorbed}
\end{equation}
uniformly in \(p\), because \(X_0=n^{\delta_*}\) and
\(\delta_*<1\).  Applying these three ledgers to
\eqref{eq:exceptional-core-first}, we obtain
\begin{equation}
 \boxed{\ |\mathcal E_p|\ll_W\frac{n}{pL^2}
       =\frac{N}{pL}\ }
 \qquad(W<p\le y).
 \label{eq:signed-exceptional-bound}
\end{equation}
The order of summation in \eqref{eq:exceptional-core-first} is essential:
absolute values are taken only after the balanced physical pieces have
been summed for a fixed smooth core \(b\).  Thus every one-sided rough
birth remains in the displayed signed sum.

\subsection{Exact row correction, including all endpoint terms}
\label{subsec:exact-rough-row-correction}

We first describe the correction without guards.  In each nonexceptional,
nontrivial row choose a clean subpool
\(
 \mathcal P_R\subset J
\)
as in \eqref{eq:active-clean-pool}, and let \(\mu_R\) be its uniform
probability measure.  Adding
\begin{equation}
 \nu_R(a)=\varepsilon_R\mu_R(a)
 \label{eq:rough-row-correction}
\end{equation}
makes the row mass exactly \(t_R\).  Equations
\eqref{eq:active-row-error} and \eqref{eq:active-clean-pool} give
\begin{equation}
 \|\nu_R\|_\infty\ll_W L^{-2}+X_R^{-1}=o(L^{-1}),
 \label{eq:rough-row-correction-slack}
\end{equation}
uniformly because \(X_R\ge X_0\).  Thus the broad floor remains between
fixed positive multiples of \(1/L\), and every weight stays strictly
between zero and one.

For \(W<p\le y\), the number of core candidates in \(\mathcal P_R\)
divisible by \(p^k\) is at most \(O_W(X_R/p^k+1)\).  If it is nonzero,
then \(p^k\ll_W X_R\), so the endpoint \(+1\) is itself
\(O_W(X_R/p^k)\).  Consequently
\begin{equation}
 \sum_a\mu_R(a)v_p(a)\ll_W\frac1p.
 \label{eq:rough-pool-valuation}
\end{equation}
The rough harmonic and counting bounds
\begin{equation}
 \sum_{\substack{R\le2n/X_0\\P^-(R)>y}}\frac1R\ll_\theta1,
 \qquad
 \#\{R\le2n/X_0:P^-(R)>y\}\ll_\theta\frac{N}{X_0}
 \label{eq:rough-row-harmonic-count}
\end{equation}
follow respectively from the truncated Euler product and the elementary
upper-bound sieve.  Therefore
\begin{align}
 \sum_R\left|\sum_a\nu_R(a)v_p(a)\right|
 &\ll_W\frac1p\sum_R(X_R/L^2+1)\notag\\
 &\ll_W\frac1p\left(\frac n{L^2}+\frac{N}{X_0}\right)
 \ll_W\frac{N}{pL}.
 \label{eq:rough-total-row-correction}
\end{align}
This computation includes all prime powers and all \(+1\)'s.  For an
independent check, group rows dyadically by \(X_R\asymp X\).  There are
\(O_\theta(N/X)\) such signatures.  If an endpoint \(+1\) is retained
instead of being absorbed into \(X_R/p^k\), only scales
\(X\gg\max(X_0,p)\) occur, and its sum is
\[
 \ll\frac L{\log p}\sum_{X\gg\max(X_0,p)}
       \frac NX\left(\frac1{L^2}+\frac1X\right)
 \ll\frac{N}{pL};
\]
here either \(X_0\) or \(p\) dominates every power of \(L\).

We now insert the charged guard system.  Define the exceptional upper
list and the global bank-donor list by
\begin{equation}
 E_{\rm exc}:=\{a\in E:X_{R_y(a)}<X_0\},\qquad
 E_{\rm donor}:=\{a\in E:a\hbox{ is a designated bank donor}\}.
 \label{eq:rough-exceptional-donor-lists}
\end{equation}
Let \(G_{\rm bank}^0\) be the chosen state-zero bank factors and define,
exhaustively,
\begin{equation}
                         G_{\rm fix}:=E_{\rm exc}\setminus E_{\rm donor}.
 \label{eq:rough-fixed-list-closed}
\end{equation}
There is no further class of dedicated fixed residual tokens: anchors are
external to the residual product, bank states belong to
\(G_{\rm bank}^0\), and all later bridge and tangent factors remain
flexible.  Put
\begin{equation}
 B_0:=\prod_{a\in G_{\rm fix}}a
       \prod_{a\in G_{\rm bank}^0}a.
 \label{eq:rough-charged-base}
\end{equation}
Also write \(G_{\rm bank}^1:=\bigcup_{g\in\mathcal B}G_g^1\).  The
following ownership ledger distinguishes selection classes from source or
guard roles; this matters because a bottom donor can be the same
occurrence as one of its two bank states.
The symbol \(\mathcal A=\coprod_R\mathcal A_R\) will denote the lower
coordinates left flexible after these exclusions.
To avoid conflict with the quotient space denoted by \(\mathcal G\) in
\Cref{sec:smooth-bridge}, denote the exhaustive set of numerical guards by
\begin{equation}
 \Gamma_{\rm num}:=\mathcal H'\cup G_{\rm fix}\cup E_{\rm donor}
       \cup G_{\rm bank}^0\cup G_{\rm bank}^1.
 \label{eq:numerical-guard-set}
\end{equation}

\begin{center}
\small
\begin{tabular}{@{}>{\raggedright\arraybackslash}p{2.35cm}
 >{\raggedright\arraybackslash}p{2.25cm}
 >{\raggedright\arraybackslash}p{3.05cm}
 >{\raggedright\arraybackslash}p{3.0cm}
 >{\raggedright\arraybackslash}p{3.25cm}@{}}
\toprule
Object & Location & Selection status & Row-quota status & Later operations \\
\midrule
\(\mathcal H'\) & \((n,2n]\) & fixed external anchors & outside the residual quotas & never varied \\
\(E_{\rm exc}\) & \((2n,M]\) & source list, not itself a selection class & supplies exceptional upper tokens & split into \(G_{\rm fix}\) and donor occurrences \\
\(E_{\rm donor}\) & \((2n,M]\) & withheld backing token; not inserted separately & one token is replaced by its bank state & never bridged or rounded \\
\(G_{\rm fix}\) & \((2n,M]\) & fixed in \(B_0\) and in the final residual set & charged before \(q_R\) is frozen & never varied \\
\(G_{\rm bank}^0,G_{\rm bank}^1\) & \((n,M]\) & alternative full path states; state \(0\) is in \(B_0\) & equal complete-signature counts & toggled only after rounding \\
\(\mathcal A\) & \((n,2n]\) & flexible candidates & supplies the remaining quota \(q_R\) & bridge, tangent, and rounding \\
\(\Gamma_{\rm num}\) & \((n,M]\) & exclusion role, not an additional selection class & no additional quota beyond its constituent roles & used only to enforce disjointness \\
\bottomrule
\end{tabular}
\end{center}

The final selection classes are
\(\mathcal H'\), \(G_{\rm fix}\), one full state of each bank path, and
the rounded subset of \(\mathcal A\).  They are pairwise disjoint.  The
only allowed overlap in the table is a role overlap: in the last two
bottom pools an element of \(E_{\rm donor}\) may be a bank-state
occurrence, and it is then counted once, as stipulated in
\Cref{lem:bank-donor-injection}.

\begin{lemma}[Guard census]
\label[lemma]{lem:rough-guard-census}
Before the quotas \(q_R\) are frozen, the numerical exclusions are
exhausted by the external anchors, the exceptional source and donor
decisions, and both states of the reserved bank.  They have the following
two, logically different, size properties.
\begin{enumerate}
 \item Globally,
 \[
  |\mathcal H'|=O_c(N),\quad |G_{\rm fix}|=O_c(N),\quad
  |E_{\rm donor}|+|G_{\rm bank}^0|+|G_{\rm bank}^1|=O(yL).
 \]
 Only \(O(y)\) anchors are modified, and the number of promoted anchors
 having trivial rough signature is \(O(\pi(y))\).
 \item In any active nonexceptional row \(R\ne1\), at most one anchor and
 at most one bank component occur; \(G_{\rm fix}\) contributes nothing.
 Hence deletion of all guarded lower coordinates removes \(O(1)\)
 candidates from that row.  For a fixed medium-prime pair \(u,v\), the
 only anchors which can equal \(ua\) or \(va\) are the
 \(O(\pi(y))\) promoted smooth anchors; together with the bank endpoints,
 they give \(O(yL)\) potentially relevant lower guards.
\end{enumerate}
In particular, the proof uses a local rowwise census and a
common-list-relevant census.  It does not require the stronger statement
that all anchors and all fixed exceptional factors together number
\(O(yL^A)\), which need not hold.
\end{lemma}

\begin{proof}
The first global bounds follow from the construction of the central
anchors, \(|E|=h=O(N)\), and \eqref{eq:bank-total-size}.  Anchor
modification affects \(O(y)\) occurrences by the fixed-prefix scale count
\eqref{eq:bank-prefix-component-count}.  A promoted anchor has trivial
rough signature only when its
base prime is at most \(y\), which gives \(O(\pi(y))\) possibilities.

A nonsmooth prefix or row-zero anchor has one unique rough marker
\(P>y\); a promoted nonsmooth anchor has one complete signature
\(p^{e_p}>1\).  There is at most one such anchor in a given nontrivial
row.  Distinct bank components have distinct marker primes by
\Cref{lem:bank-donor-injection}, so a nontrivial marker row contains at
most one component, its two possible endpoint occurrences, and its one
donor occurrence.  The set \(G_{\rm fix}\) is supported only on
exceptional rows.  This proves the rowwise assertion.

Finally, a prefix anchor has smooth cofactor supported on
\(F_{\rm anc}\subset[2,W]\), a row-zero anchor is prime, and a promoted
nonsmooth anchor has a power of two as smooth cofactor.  None is divisible
by a label in \((W,y]\).  Only promoted smooth anchors remain possible in
a lower common-multiplier list.  The fixed exceptional factors and all
withheld donors lie above \(2n\); bank endpoints contribute \(O(yL)\)
numbers.  This proves the last assertion.
\end{proof}

All factors in this product are mutually distinct and are guarded from
the anchors and from the flexible candidates.  If \(m_R(G)\) denotes the
number of factors of \(G\) having complete signature \(R\), define, for
every nontrivial row,
\begin{equation}
 q_R:=t_R-m_R(G_{\rm fix})-m_R(G_{\rm bank}^0)\in\mathbb Z_{\ge0}.
 \label{eq:postcharge-rough-quota}
\end{equation}
The backing map is literal:
\begin{equation}
 \iota(g)=g\quad(g\in G_{\rm fix}),\qquad
 \iota(g)=\text{the designated donor of its bank component}
                   \quad(g\in G_{\rm bank}^0).
 \label{eq:rough-explicit-backing-injection}
\end{equation}
Its image lies in \(E\), preserves the complete signature, and is
injective.  Indeed, different bank components have different donors, and
every such donor has been removed from \(G_{\rm fix}\) (a nonexceptional
donor never belonged to \(G_{\rm fix}\) in the first place).  This proves
the nonnegativity of \eqref{eq:postcharge-rough-quota}.  In an exceptional
row, the partition into retained factors and designated donors, together
with the one base token replacing each donor, exhausts all \(t_R\) upper
tokens; hence \(q_R=0\) there.

We also make the whole-row convention explicit.  Let
\(\mathscr B_{\rm row}\) be the set of nontrivial complete signatures of
bank component marker rows.  Every such signature is the component's
unique marker prime \(P>y\).  Its designated donor is \(Pu\in E\), with
\(u\ll y\), so
\[
                   X_P=\frac{2n}{P}\ll y.
\]
The component count \eqref{eq:bank-total-size} and the use of a different
marker for every component give
\begin{equation}
 |\mathscr B_{\rm row}|=O(yL),\qquad
 X_R\ll y\quad(R\in\mathscr B_{\rm row}).
 \label{eq:rough-bank-touched-row-data}
\end{equation}
In the endpoint-deletion convention used in this paper, every
nonexceptional bank-touched row remains active and is corrected below; no
such row is dedicated in its entirety.  Thus the set of fully dedicated
nonexceptional rows is
\begin{equation}
 \mathscr D_{\rm row}:=\varnothing\subset\mathscr B_{\rm row},\qquad
 |\mathscr D_{\rm row}|=O(yL),\qquad
 X_R\ll y\quad(R\in\mathscr D_{\rm row}).
 \label{eq:rough-dedicated-row-data}
\end{equation}
The last two bounds would remain true if any subset of the bank-touched
rows were instead dedicated.  Exceptional rows \(X_R<X_0\) are fixed
separately by \eqref{eq:rough-fixed-list-closed} and are not put in
\(\mathscr D_{\rm row}\).

\begin{lemma}[Uniform clean pool after all guards]
\label[lemma]{lem:rough-uniform-clean-pool}
For every active nonexceptional row \(R\ne1\), let \(\mathcal A_R\) be
the lower candidates remaining after the exhaustive guard census in
\Cref{lem:rough-guard-census}, and let \(\mathcal P_R\subset\mathcal A_R\)
be the broad correction pool.  Uniformly in these rows,
\begin{equation}
 \left|q_R-\sum_{a\in\mathcal A_R}x_a^{\rm raw}\right|
       \ll_W X_R/L^2+1,\qquad
 \#\mathcal P_R\gg_W X_R.
 \label{eq:guard-local-conditions}
\end{equation}
Consequently the constant correction on \(\mathcal P_R\) has density
\begin{equation}
 O_W(L^{-2}+X_R^{-1})=o(L^{-1}),
 \label{eq:guard-clean-pool-density}
\end{equation}
and the two-sided raw slack survives.
\end{lemma}

\begin{proof}
Before guards, \eqref{eq:active-clean-pool} supplies at least
\(c_WX_R\) broad candidates.  By
\Cref{lem:rough-guard-census}, at most a fixed number are deleted in an
active nontrivial row: at most one anchor, and the endpoints and donor
role of at most one bank component.  Fixed exceptional factors lie above
\(2n\) and occur only in exceptional rows.  Since \(X_R\ge X_0\to\infty\),
at least \(c_WX_R/2\) clean candidates remain for all sufficiently large
\(n\), proving the second assertion in
\eqref{eq:guard-local-conditions}.

For the first assertion, let \(\mathcal D_R\) be exactly the lower
coordinates removed from the raw row.  Subtracting the definitions of
\(q_R\) and \(\mathcal A_R\) gives the identity
\begin{align}
 &\left(q_R-\sum_{a\in\mathcal A_R}x_a^{\rm raw}\right)
  -\left(t_R-\sum_{R_y(a)=R}x_a^{\rm raw}\right)\notag\\
 &\qquad
 =-m_R(G_{\rm fix})-m_R(G_{\rm bank}^0)
                         +\sum_{a\in\mathcal D_R}x_a^{\rm raw}.
 \label{eq:guard-exact-local-ledger}
\end{align}
Here \(m_R(G_{\rm fix})=0\).  If a bank component occurs in the row, its
contribution to the right side is
\[
             -1+\sum_{\substack{a\in\mathcal D_R\\
                    a\ {\rm is\ one\ of\ its\ lower\ endpoints}}}
                         x_a^{\rm raw}=O_W(1);
\]
the expression is not claimed to vanish.  The at-most-one anchor
contributes another \(O_W(1)\).  There are no further classes in the
census.  Therefore
\begin{equation}
 \left|q_R-\sum_{a\in\mathcal A_R}x_a^{\rm raw}\right|
 \le |\varepsilon_R|+O_W(1),
 \label{eq:guard-local-ledger-bound}
\end{equation}
and \eqref{eq:active-row-error} proves the first assertion.  Dividing by
\(\#\mathcal P_R\gg_WX_R\) gives
\eqref{eq:guard-clean-pool-density}.  Since \(X_R\ge n^{\delta_*}\), this
is \(o(L^{-1})\), so a fixed fraction of both the lower and upper broad
margin remains.
\end{proof}

The already charged medium-prime
guard profile is \(O_W(N/(pL))\).  Replacing \(\varepsilon_R\) in
\eqref{eq:rough-row-correction} by the left side of
\eqref{eq:guard-local-conditions} therefore repeats
\eqref{eq:rough-row-correction-slack}--\eqref{eq:rough-total-row-correction}
verbatim and proves the exact postcharge identity
\begin{equation}
 \sum_{a\in\mathcal A_R}x_a=q_R
 \qquad(R\ne1).
 \label{eq:exact-complete-signature-row}
\end{equation}
The row \(R=1\) is intentionally left for the integral smooth-row quota
and structured head fit of \Cref{sec:smooth-bridge}; its top component and
protected floor in \eqref{eq:rough-smooth-splitting} remain frozen.

\subsection{Nonnegative exceptional charge and simultaneous feasibility}
\label{subsec:rough-charges}

\begin{remark}[Signed discrepancy is not a charge]
\label[remark]{rem:signed-versus-nonnegative-exceptional}
The quantity \(\mathcal E_p\) in
\eqref{eq:exceptional-core-first} is signed: it compares the actual upper
exceptional rows with their lower raw baseline and is used only to bound
the residual error.  By contrast, the quantity \(C_p^{\rm exc}\) below is
a sum of nonnegative valuations of factors that are actually retained; it
is used to prove divisibility.  Cancellation in \(\mathcal E_p\) gives no
upper bound for \(C_p^{\rm exc}\), whose natural size can be \(N/p\).
\end{remark}

We therefore estimate the nonnegative valuation independently.  Let
\[
 C_p^{\rm exc}:=
 \sum_{\substack{a\in E,\ X_{R_y(a)}<X_0\\
                  a\text{ not a bank donor}}}v_p(a)
 \qquad(p\le y).
\]
If \(a=Rb\) occurs, then \(b<2X_0\) and \(v_p(a)=v_p(b)\).  For fixed
\(b\), discard the lower cutoff on \(R\) and apply the interval Selberg
sieve to
\[
            2n/b<R\le(2n+h)/b.
\]
Its length is \(h/b\gg n^{1-\delta_*}/L\).  We use Selberg weights
supported on \(d\le y^2\); the elementary interval remainders have lcm
moduli
\(
 [d,e]\le y^4=n^{8/9}
\), not merely \(y^2\).  Since \(\delta_*<1/18\), these remainders are
uniformly negligible compared with the interval length after logarithmic
factors.  By \eqref{eq:Selberg-specialized},
\begin{equation}
 \#\{R:2n/b<R\le(2n+h)/b,\ P^-(R)>y\}
 \ll \frac{h}{b\log y}+\frac{y^4}{(\log y)^2}.
 \label{eq:exceptional-selberg-count}
\end{equation}
Now
\begin{equation}
 \sum_{b<2X_0}\frac{v_p(b)}b
 =\sum_{k\ge1}\frac1{p^k}
       \sum_{m<2X_0/p^k}\frac1m
 \ll\frac{\delta_*L+1}{p},
 \label{eq:exceptional-charge-harmonic}
\end{equation}
and \(\sum_{b<2X_0}v_p(b)\ll X_0/p\).  Multiplying
\eqref{eq:exceptional-selberg-count} by \(v_p(b)\) and summing \(b\)
therefore yields, uniformly in \(p\le y\), the following bound.  Here
\(\delta_*>0\) has already been fixed; hence, after increasing the
threshold for \(n\), one has \(\delta_*L+1\leq2\delta_*L\).  The factor
two is absorbed into a fixed \(A_{\rm exc}\), while the displayed
little-oh term below contains only the sieve remainder.  Thus
\begin{equation}
 \boxed{\ C_p^{\rm exc}
   \le A_{\rm exc}\frac{\delta_*}{\theta}\frac Np
       +o\!\left(\frac{N}{pL}\right)\ }.
 \label{eq:uniform-exceptional-charge}
\end{equation}
For \(p>2X_0\), the left side is exactly zero.  The sieve remainder is
indeed harmless, since
\[
 \frac{y^4X_0}{p(\log y)^2}
 =o\!\left(\frac{N}{pL}\right)
\]
under \(8/9+\delta_*<1\).

Let \(B_{\rm exc}\) denote the product of the exceptional factors in
\(G_{\rm fix}\).  The bank and guard construction gives
\begin{equation}
 v_p(B_0/B_{\rm exc})=v_p(B_{\rm bank}^0)
       =o_W\!\left(\frac{N}{pL}\right)
 \qquad(p\le y),
 \label{eq:base-guard-small-prime-profile}
\end{equation}
by \eqref{eq:bank-small-prime-profile}.  There is no unenumerated fixed
class in this estimate.  Legendre's formula gives uniformly in this range
\begin{equation}
 v_p(T)=\frac{h}{p-1}+O\!\left(\frac L{\log p}\right)
      =\frac{(c+o(1))N}{p-1}.
 \label{eq:tail-uniform-small-valuation}
\end{equation}
\begin{lemma}[Combined charge]
\label[lemma]{lem:combined-charge}
With the one-time choice \eqref{eq:rough-cutoff}, the complete fixed base
and the modified anchor divisor satisfy, simultaneously,
\begin{equation}
                   D'B_0\mid T,
 \qquad             B_0\mid Y:=T/D'.
 \label{eq:combined-precharge-divisibility}
\end{equation}
\end{lemma}

\begin{proof}
We verify the valuation inequality in three disjoint prime ranges.

If \(p\in F_{\rm anc}\), the definition of \(\gamma\) above
\eqref{eq:rough-cutoff} gives
\[
             v_p(T)-v_p(D')\ge\gamma N/p.
\]
Equations \eqref{eq:uniform-exceptional-charge} and
\eqref{eq:base-guard-small-prime-profile}, together with
\(A_{\rm exc}\delta_*/\theta<\gamma/4\), show that
\(v_p(B_0)<\gamma N/(2p)\) for all sufficiently large \(n\).
Thus the required inequality has a fixed positive margin.

If \(p\le y\) but \(p\notin F_{\rm anc}\), then \(v_p(D')=0\), while
\eqref{eq:tail-uniform-small-valuation} and the definition of \(\gamma\)
give \(v_p(T)\ge\gamma N/p\).  The same two charge estimates again give
\(v_p(B_0)<\gamma N/(2p)\).  This includes
\(p\in F_{\rm hd}\setminus F_{\rm anc}\): no anchor valuation is charged
at such a prime.

Finally let \(p>y\).  Then \(v_p(D')=0\), because
\(F_{\rm anc}\subset[2,W]\).  The injection
\eqref{eq:rough-explicit-backing-injection} maps every factor of
\(G_{\rm fix}\cup G_{\rm bank}^0\) to a distinct actual factor of \(E\),
preserving the entire complete rough signature and hence every
\(p\)-adic valuation.  A fixed exceptional factor is mapped to itself;
a bank-base factor is mapped to its designated donor.  These two image
classes are disjoint by construction.  Therefore
\(v_p(B_0)\le v_p(T)\) for every \(p>y\).

The three ranges cover every prime and prove
\(v_p(D'B_0)\le v_p(T)\) coordinatewise.  This proves both assertions in
\eqref{eq:combined-precharge-divisibility}.  In particular, the conclusion
is not obtained by multiplying two separately known divisors of \(T\);
it is a simultaneous verification for the product \(D'B_0\).
\end{proof}

\subsection{Output of the rough stage}
\label{subsec:rough-output}

Let \(\mathcal A=\coprod_R\mathcal A_R\) be the remaining flexible lower
candidates, and use the residual convention
\begin{equation}
 r_p:=v_p(Y)-v_p(B_0)-\sum_{a\in\mathcal A}x_av_p(a).
 \label{eq:rough-residual-convention}
\end{equation}
We now give the full sign ledger rather than absorbing guard changes into
words.  Let \(\mathcal D_{\rm ne}\) be the set of raw lower coordinates
deleted in nonexceptional rows, and put
\[
 E_{\rm donor}^{\rm exc}:=E_{\rm donor}\cap E_{\rm exc}.
\]
For \(W<p\le y\), define
\begin{align}
 \mathcal R_p^{\rm raw}
   &:=\sum_a\{\1_E(a)-x_a^{\rm raw}\}v_p(a),\notag\\
 \Delta_p^{\rm guard}
   &:=\sum_{a\in\mathcal D_{\rm ne}}x_a^{\rm raw}v_p(a)
      +\sum_{a\in E_{\rm donor}^{\rm exc}}v_p(a)
      -\sum_{a\in G_{\rm bank}^0}v_p(a).
 \label{eq:rough-guard-difference-definition}
\end{align}
The first sum in \(\Delta_p^{\rm guard}\) restores the raw coordinates
which were removed outside the exceptional rows; the second restores the
exceptional donor terms included in \(\mathcal E_p\); and the last charges
the bank base which replaced those donor tokens.  Since
\(v_p(D')=0\) in this medium range, direct expansion of
\(Y=T/D'\), \(B_0=G_{\rm fix}G_{\rm bank}^0\), and
\(G_{\rm fix}=E_{\rm exc}\setminus E_{\rm donor}\) gives the exact
identity
\begin{equation}
 \boxed{\quad
 r_p=\mathcal R_p^{\rm raw}-\mathcal E_p
      -\sum_{\substack{R\ne1\\X_R\ge X_0}}
          \sum_a\nu_R(a)v_p(a)
      +\Delta_p^{\rm guard}.
 \quad}
 \label{eq:rough-complete-residual-algebra}
\end{equation}
This also verifies the convention ``target minus current'' term by term.

For completeness, the guard census makes the last error quantitative.
Prefix, row-zero, and promoted nonsmooth anchors have zero \(p\)-valuation
when \(W<p\le y\).  For a fixed \(p\), at most one promoted smooth anchor
has nonzero \(p\)-valuation, of size at most \(L/\log p\).  All other
terms in \eqref{eq:rough-guard-difference-definition} come from
\(O(yL)\) bank endpoints or donors, and each has valuation at most
\(O(L/\log p)\).  Hence, uniformly in the medium range,
\begin{align}
 |\Delta_p^{\rm guard}|
 &\ll \frac{yL^2}{\log p}+\frac L{\log p}
   =o\!\left(\frac{N}{pL}\right),
 \label{eq:rough-guard-difference-bound}
\end{align}
because
\(pyL^3/N\le y^2L^3/N=n^{-5/9}L^4\to0\).
Equations \eqref{eq:rough-global-crt},
\eqref{eq:signed-exceptional-bound},
\eqref{eq:rough-total-row-correction}, and
\eqref{eq:rough-guard-difference-bound}, inserted with their displayed
signs in \eqref{eq:rough-complete-residual-algebra}, prove
\begin{equation}
 |r_p|\ll_W\frac{N}{pL}\qquad(W<p\le y).
 \label{eq:strict-rate-rough-residual}
\end{equation}
Equation \eqref{eq:exact-complete-signature-row} proves exact equality in
every nontrivial complete rough-signature row and hence \(r_p=0\) for every
\(p>y\), with all multiplicities.  At this stage the smooth row and the
coordinates \(p\le W\) are provisional.  The smooth-row decomposition has
a frozen \(F_{\rm hd}\)-free protected floor and an \(F_{\rm hd}\)-free
active excess of mass \(\asymp_WN\); all fixed-head valuations will be
inserted by the structured cells of the next two sections.  Thus the
distinction \(F_{\rm anc}\subset F_{\rm hd}\), all fixed charges, and every
guard are already present in the ledger to which the bridge is applied.

For reference, the exact invariants carried through the remaining stages
are summarized below.  ``Exact rows'' always means complete rough
signatures, with prime-power multiplicity; ``mass'' includes integral
fixed tokens with weight one and fractional coordinates with their actual
weights.

\begin{center}
\small
\begin{tabular}{@{}>{\raggedright\arraybackslash}p{1.55cm}
 >{\raggedright\arraybackslash}p{2.45cm}
 >{\raggedright\arraybackslash}p{2.25cm}
 >{\raggedright\arraybackslash}p{2.55cm}
 >{\raggedright\arraybackslash}p{2.7cm}
 >{\raggedright\arraybackslash}p{2.25cm}@{}}
\toprule
Stage & Rough rows and \(p>y\) & Head \(p\le W\) & Medium \(W<p\le y\) & Total mass and ordinary log & Coordinates \\
\midrule
After rough selection & every \(R\ne1\) exact; hence \(p>y\) exact & provisional & residual bounded by \eqref{eq:strict-rate-rough-residual} & provisional & two-sided interior on clean pools \\
After smooth bridge & every row exact & exact & band sums exact; pointwise residual retains the tangent bound & mass and ordinary log exact & interior \\
After tangent & every row exact & exact & every individual valuation exact & mass and ordinary log exact & in \([0,1]\), with endpoint slack where used \\
After floating rounding & every row exact, so no \(p>y\) error & integer error \(e_p\) & integer error \(e_p\) & mass exact; log changes by \(\sum e_p\log p\) & flexible coordinates are \(0\)-\(1\) \\
After bank replacement & every row exact & exact & exact & mass and ordinary log exact & all selected factors are \(0\)-\(1\) and distinct \\
\bottomrule
\end{tabular}
\end{center}

Each row of this table is proved at the indicated stage: the present
section proves the first, the bridge and tangent sections prove the next
two, and \Cref{lem:floating-rounding,prop:guarded-integral-exactification}
prove the last two.  In particular, ordinary logarithm after rounding is
not asserted separately from the valuation error; it becomes exact only
when the bank changes valuation by \(-e\).
 \section{Marked friable counts on the structured cells}
\label{sec:marked-smooth}

The smooth bridge uses only physical intervals and exact finite
valuation and coprimality patterns at the head primes.  This restriction is
essential: ordinary smooth-number estimates do not count an arbitrary
nonzero residue class modulo a fixed integer.  In this section we prove all
one- and two-marked estimates needed for these structured cells, including
their stability under the compact tilts used later.

Throughout,
\begin{equation}
 y=n^{2/9},\qquad U:=\frac{\log n}{\log y}=\frac92,
 \qquad L=\log n.
 \label{eq:marked-parameters}
\end{equation}
The numerical value \(U=9/2\) is used below.  In particular, we make no
terminal prime-power assertion at \(U=5/2\).

\subsection{Exact head cells and a four-mark estimate}
\label{subsec:marked-structured-cells}

Fix a finite set \(H\subset F_{\rm hd}\), a vector of fixed nonnegative
integers \({\bf e}=(e_\ell)_{\ell\in H}\), and fixed
\(0<A<B<\infty\).  Put
\begin{equation}
 h_{\bf e}=\prod_{\ell\in H}\ell^{e_\ell},\qquad
 M_H=\prod_{\ell\in H}\ell,\qquad
 \delta_{\bf e}=\frac1{h_{\bf e}}
                    \frac{\varphi(M_H)}{M_H},
 \label{eq:marked-head-density}
\end{equation}
and define the structured cell
\begin{equation}
 \mathcal C_{\bf e}(A,B)=
 \{m:An<m\le Bn,\ P^+(m)\le y,\
                         v_\ell(m)=e_\ell\ (\ell\in H)\}.
 \label{eq:marked-structured-cell}
\end{equation}
The finite collection of \(A,B,H,{\bf e}\) used in the construction is
fixed before \(n\to\infty\).
For a bounded-variation function we abbreviate
\(\|\phi\|_{BV}:=\|\phi\|_\infty+\operatorname{Var}\phi\).

For a \(y\)-smooth integer \(d\), write
\[
 v(d)=\frac{\log d}{\log y},\qquad
 h_U(v)=\frac{\rho(U-v)}{\rho(U)}.
\]

If \(\nu_i\) is normalized uniform measure on a structured cell
\(C_i\), a \emph{finite mixture of cells} means the probability measure
\begin{equation}
                  \nu=\sum_{i=1}^m\lambda_i\nu_i,
 \qquad \lambda_i\geq0,\qquad \sum_{i=1}^m\lambda_i=1,
 \label{eq:marked-finite-mixture-definition}
\end{equation}
where \(m\) is fixed.  When compact positivity is required later, the
weight vector is restricted to a fixed compact subset of the relative
interior of the relevant simplex.  A piecewise-constant baseline weight
on the disjoint cells is exactly of the form
\eqref{eq:marked-finite-mixture-definition} after normalization.  We
apply count formulas to each \(C_i\); only probabilities, expectations,
and covariances are then averaged under \(\nu\).  We never interpret a
nonintegral convex mixture as a count.

\begin{proposition}[Uniform marked structured cells]
\label[proposition]{prop:marked-structured-cells}
Suppose that \(d\le y^4\) is \(y\)-smooth and
\((d,M_H)=1\).  Uniformly in \(d\),
\begin{align}
 A_{{\bf e},A,B}(d)
 &:=\#\{m\in\mathcal C_{\bf e}(A,B):d\mid m\}\notag\\
 &=\delta_{\bf e}\frac{(B-A)n}{d}
   \left\{\rho(U-v(d))+O_{H,A,B}(L^{-1})\right\}.
 \label{eq:marked-cell-count}
\end{align}
Consequently, for the normalized uniform measure on the cell,
\begin{equation}
 \mathbb P_{{\bf e},A,B}(d\mid m)
 =\frac{h_U(v(d))}{d}+O_{H,A,B}\!\left(\frac1{dL}\right).
 \label{eq:marked-cell-probability}
\end{equation}
If \(\phi:[A,B]\to\mathbb R\) has bounded variation, then
\begin{align}
 &\sum_{\substack{m\in\mathcal C_{\bf e}(A,B)\\d\mid m}}
       \phi(m/n)\notag\\
 &\quad=\delta_{\bf e}\frac nd\rho(U-v(d))
       \int_A^B\phi(t)\,dt
 +O_{H,A,B}\!\left(\frac{n}{dL}
       \{\|\phi\|_\infty+\operatorname{Var}\phi\}\right).
 \label{eq:marked-cell-Stieltjes}
\end{align}
The normalized assertions hold for a fixed finite disjoint union, or a
probability mixture in the sense of
\eqref{eq:marked-finite-mixture-definition}.  Any
overlapping family is first disjointified and then interpreted through its
induced convex mixture.  In particular,
any two allowed normalized cells \(C,C'\) satisfy
\begin{equation}
 \left|\mathbb P_C(d\mid m)-\mathbb P_{C'}(d\mid m)\right|
 \ll_W\frac1{dL}.
 \label{eq:marked-cell-comparison}
\end{equation}
\end{proposition}

\begin{proof}
The exact valuation pattern, unlike a general residue class, has the
finite divisibility identity
\begin{equation}
 \1_{\{v_\ell(m)=e_\ell\ (\ell\in H)\}}
 =\1_{h_{\bf e}\mid m}
    \sum_{a\mid M_H}\mu(a)\1_{a\mid m/h_{\bf e}}.
 \label{eq:marked-head-inclusion-exclusion}
\end{equation}
Since \((d,M_H)=1\), it follows exactly that
\begin{align}
 A_{{\bf e},A,B}(d)
 =\sum_{a\mid M_H}\mu(a)
 \left\{\Psi\!\left(\frac{Bn}{h_{\bf e}ad},y\right)
       -\Psi\!\left(\frac{An}{h_{\bf e}ad},y\right)\right\}.
 \label{eq:marked-cell-Psi-identity}
\end{align}
The smallest unscaled cofactor is
\begin{equation}
             \frac n{y^4}=n^{1/9}=y^{1/2}\longrightarrow\infty.
 \label{eq:marked-four-cofactor-margin}
\end{equation}
Thus every smooth parameter in \eqref{eq:marked-cell-Psi-identity} lies in
the fixed compact interval
\([1/2+o(1),9/2+o(1)]\).  When it is at least one, apply
\Cref{thm:HT-Saias} and \eqref{eq:Lambda-Dickman}
\cite{HildebrandTenenbaum1993,Saias1989}; when it is at most
one, use the exact identity \(\Psi(x,y)=\lfloor x\rfloor\).  The fixed
factors \(A,B,h_{\bf e},a\) shift the parameter by \(O_H(1/L)\), and
\(\rho\) is Lipschitz on this compact range.  Hence, uniformly for
\(C=A,B\),
\begin{equation}
 \Psi\!\left(\frac{Cn}{h_{\bf e}ad},y\right)
 =\frac{Cn}{h_{\bf e}ad}
    \left\{\rho(U-v(d))+O_H(L^{-1})\right\}+O(1).
 \label{eq:marked-Psi-specialization}
\end{equation}
By \eqref{eq:marked-four-cofactor-margin}, the endpoint \(O(1)\) is
\(O(n/(dL))\).  Substitute \eqref{eq:marked-Psi-specialization} in
\eqref{eq:marked-cell-Psi-identity} and use
\[
 \frac1{h_{\bf e}}\sum_{a\mid M_H}\frac{\mu(a)}a
 =\delta_{\bf e}.
\]
This proves \eqref{eq:marked-cell-count}.  With \(d=1\), it also gives
\begin{equation}
 \#\mathcal C_{\bf e}(A,B)
 =\delta_{\bf e}(B-A)n\{\rho(U)+O_W(L^{-1})\}
 \asymp_W n,
 \label{eq:marked-cell-size}
\end{equation}
so division proves \eqref{eq:marked-cell-probability}.

Apply \eqref{eq:marked-cell-count} uniformly to every cumulative
subinterval \((An,tn]\), \(A\le t\le B\).  Its error is
\(O_W(n/(dL))\), independent of \(t\).  Stieltjes partial summation
against \(\phi\) proves \eqref{eq:marked-cell-Stieltjes}.  Every cell has
the same normalized main divisibility profile, so finite disjoint unions
and compact convex mixtures preserve it; this also proves
\eqref{eq:marked-cell-comparison}.
\end{proof}

The cutoff \(y^4\) covers exactly the divisors used later:
\begin{equation}
 p,\quad pq,\quad p^k,\quad p^kq,\quad p^kqr^j,
 \label{eq:marked-covered-divisors}
\end{equation}
where \(W<p,q,r\le y\), \(p^k\le y^2\), and \(r^j\le y\).
In particular, the last product is at most \(y^4\).  This is the genuine
reason for fixing \(U=9/2\): all pointwise marked products retain the
cofactor margin \eqref{eq:marked-four-cofactor-margin}.

\subsection{One and two marks, prime powers, and physical tests}
\label{subsec:marked-consequences}

Let \(C\) be an allowed cell or finite mixture and put
\(t_p=\log p/\log y\).  For distinct \(W<p,q\le y\),
Proposition~\ref{prop:marked-structured-cells} gives
\begin{align}
 \mathbb E_C\1_{p\mid m}
 &=\frac{h_U(t_p)}p+O_W(1/(pL)),\notag\\
 \mathbb E_C\1_{pq\mid m}
 &=\frac{h_U(t_p+t_q)}{pq}+O_W(1/(pqL)).
 \label{eq:marked-one-two}
\end{align}
It follows that
\begin{align}
 \operatorname{Var}_C(\1_{p\mid m})
 &=\frac{h_U(t_p)}p+O_W(1/(pL)+1/p^2),\notag\\
 \operatorname{Cov}_C(\1_{p\mid m},\1_{q\mid m})
 &=\frac{h_U(t_p+t_q)-h_U(t_p)h_U(t_q)}{pq}
       +O_W(1/(pqL)).
 \label{eq:marked-one-two-covariance}
\end{align}
Likewise, if \(p^k\le y^2\) and \(p\ne q\), then
\begin{align}
 \mathbb E_C\1_{p^kq\mid m}
 &=\frac{h_U(kt_p+t_q)}{p^kq}+O_W(1/(p^kqL)),
 \label{eq:marked-power-two}\\
 \operatorname{Cov}_C(\1_{p^k\mid m},\1_{q\mid m})
 &=\frac{h_U(kt_p+t_q)-h_U(kt_p)h_U(t_q)}{p^kq}
       +O_W(1/(p^kqL)).
 \label{eq:marked-power-covariance}
\end{align}

There is also a comparison for the full valuation.  Expand
\(v_p=\sum_{k\ge1}\1_{p^k\mid m}\).  Applying
\eqref{eq:marked-cell-comparison} for \(p^k\le y^4\) gives total error
\(O_W(1/(pL))\).  The remaining powers satisfy, by the elementary
multiple count and \eqref{eq:marked-cell-size},
\[
 \sum_{p^k>y^4}\mathbb P_C(p^k\mid m)
 \ll y^{-4}+L/n=o(1/(pL)).
\]
Consequently
\begin{equation}
 |\mathbb E_Cv_p(m)-\mathbb E_{C'}v_p(m)|
 \ll_W\frac1{pL}\qquad(W<p\le y).
 \label{eq:marked-full-valuation-cell-comparison}
\end{equation}

If \(\phi\) has bounded variation on the common physical compact set,
\eqref{eq:marked-cell-Stieltjes} and
\eqref{eq:marked-cell-probability} give
\begin{equation}
 |\operatorname{Cov}_C(\1_{d\mid m},\phi(m/n))|
 \ll_W\frac{\|\phi\|_{BV}}{dL}.
 \label{eq:marked-physical-covariance}
\end{equation}
The same estimate holds for finite mixtures: within-cell covariances have
this bound, while between-cell covariances are controlled by
\eqref{eq:marked-cell-comparison}.  Summing over \(d=p^k\) yields
\begin{equation}
 |\operatorname{Cov}_C(v_p(m),\phi(m/n))|
 \ll_W\frac{\|\phi\|_{BV}}{pL}.
 \label{eq:marked-physical-full-valuation}
\end{equation}
For example, with \(\phi(t)=\log t\) and arbitrary coefficients \(c_p\),
\begin{equation}
 \left|\operatorname{Cov}_C\!\left(
   \sum_{W<p\le y}c_pv_p(m),\log(m/n)\right)\right|
 \ll_W\frac1L\sum_{W<p\le y}\frac{|c_p|}{p}.
 \label{eq:marked-physical-linear-combination}
\end{equation}

\subsection{Stability under compact bridge tilts}
\label{subsec:marked-compact-tilts}

The bridge changes the measure by bounded band fugacities divided by
\(L\).  We prove that the error scale above survives this homotopy.  Let
\begin{equation}
 S(m)=\sum_{W<p\le y}\eta_pv_p(m),\qquad |\eta_p|\le B,
 \label{eq:marked-tilt-score}
\end{equation}
where \(B\) is fixed, and denote expectation after reweighting by
\(e^{S/L}\) by \(\mathbb E_{C,S}\).

\begin{lemma}[Squarefree marked transfer under tilts]
\label[lemma]{lem:marked-tilted-squarefree}
Uniformly for \(d=p\) or \(d=pq\), with distinct primes in \((W,y]\),
\begin{equation}
 \mathbb P_{C,S}(d\mid m)
 =\mathbb P_C(d\mid m)+O_{B,W}(1/(dL)).
 \label{eq:marked-tilted-squarefree}
\end{equation}
If \(\phi\) is a fixed bounded-variation physical test, then
\begin{equation}
 |\operatorname{Cov}_{C,S}(\1_{d\mid m},\phi(m/n))|
 \ll_{B,W}\frac{\|\phi\|_{BV}}{dL}.
 \label{eq:marked-tilted-physical-squarefree}
\end{equation}
\end{lemma}

\begin{proof}
Put \(\Omega(m)=\sum_{W<p\le y}v_p(m)\).  Since
\(\Omega(m)\log W\le\log(Bn)\), the factor
\(e^{B\Omega/L}\) is uniformly bounded.  Moreover, for \(j=1,2\),
elementary counting gives
\begin{equation}
 \mathbb E_C\{\1_{d\mid m}\Omega(m)^j\}
 \ll_{W,j}\frac{(1+\log L)^j}{d}.
 \label{eq:marked-Omega-moments}
\end{equation}
Indeed write \(m=dk\), discard smoothness and the head restrictions, and
enlarge to \(k\le Bn/d\).  Since \(d\) has at most two prime factors,
\(\Omega(dk)\le2+\Omega(k)\).  Expanding the first two powers into
prime-power divisibility indicators and using
\(\lfloor X/e\rfloor\le X/e\) gives
\[
 \frac1X\sum_{k\le X}\Omega(k)^j
 \ll_j\left(1+
   \sum_{W<r\le y}\sum_{a\ge1}\frac1{r^a}\right)^j
 \ll_j(1+\log L)^j.
\]
Division by the cell size \eqref{eq:marked-cell-size} proves
\eqref{eq:marked-Omega-moments}.

Taylor's formula, with the uniform exponential bound just noted, gives
\begin{align}
 \mathbb E_Ce^{S/L}
 &=1+L^{-1}\mathbb E_CS+O_B((\log L)^2/L^2),\notag\\
 \mathbb E_C(\1_{d\mid m}e^{S/L})
 &=\mathbb P_C(d\mid m)+L^{-1}\mathbb E_C(\1_{d\mid m}S)
       +O_B((\log L)^2/(dL^2)).
 \label{eq:marked-tilt-Taylor}
\end{align}
It remains to use cancellation in the centered first-order term.  Expand
\(S\) into prime-power indicators.  If the score prime \(r\) is not a
factor of \(d\), Proposition~\ref{prop:marked-structured-cells} applies to
\(dr^a\) for \(r^a\le y\); the product is then at most \(y^3\).  Its
main covariance is bounded by
\[
 \frac1{dr^a}
 |h_U(v(d)+at_r)-h_U(v(d))h_U(at_r)|
 \ll_U\frac{at_r}{dr^a},
\]
because the expression vanishes at \(at_r=0\) and \(h_U\) is Lipschitz
on the relevant compact range.  The at most two score primes dividing
\(d\) contribute \(O_B(1/d)\) directly.  The marked errors and the tail
\(r^a>y\) are summable by elementary common-multiple bounds.  Since
\begin{equation}
 \sum_{W<r\le y}\sum_{r^a\le y}\frac{at_r}{r^a}=O_W(1),
 \label{eq:marked-score-summability}
\end{equation}
we obtain
\begin{equation}
 |\operatorname{Cov}_C(\1_{d\mid m},S)|\ll_{B,W}\frac1d.
 \label{eq:marked-first-order-cancellation}
\end{equation}
Divide the two lines of \eqref{eq:marked-tilt-Taylor} and use
\((\log L)^2/L^2=o(1/L)\).  This proves
\eqref{eq:marked-tilted-squarefree}.

Insert \(\phi(m/n)\) into the same calculation.  In every leading marked
term, \eqref{eq:marked-cell-Stieltjes} factors the same Lebesgue average
of \(\phi\).  More explicitly, the first derivative at zero of the
tilted covariance is the third centered combination
\begin{align}
 \kappa_C(\1_d,\phi,S)
 &:={\mathbb E}_C(\1_d\phi S)
   -{\mathbb E}_C(\1_d\phi){\mathbb E}_CS
   -{\mathbb E}_C(\1_dS){\mathbb E}_C\phi\notag\\
 &\quad-{\mathbb E}_C(\phi S){\mathbb P}_C(d\mid m)
   +2{\mathbb P}_C(d\mid m){\mathbb E}_C\phi\,{\mathbb E}_CS.
 \label{eq:marked-physical-three-centering}
\end{align}
For each prime-power term of \(S\), the common Stieltjes main term in
these five summands cancels exactly.  Summing the remaining Stieltjes
errors gives
\[
 \frac1L|\kappa_C(\1_d,\phi,S)|
 \ll_{B,W}\frac{\|\phi\|_{BV}\log L}{dL^2}.
\]
The three-factor Taylor remainder is
\(O_{B,W}(\|\phi\|_{BV}(\log L)^2/(dL^2))\), hence is
\(o(\|\phi\|_{BV}/(dL))\).  Together with the un-tilted covariance
\eqref{eq:marked-physical-covariance}, this proves
\eqref{eq:marked-tilted-physical-squarefree}.
\end{proof}

Individual very high prime powers need not satisfy a relative
\(O(1/L)\) stability statement when the tilt itself contains that prime.
The bridge does not need such a statement.  It needs the following
aggregate transfer, which we now prove.

For a set \(E\) of at most two primes, let \(S_E\) be the score
\eqref{eq:marked-tilt-score} with the primes of \(E\) omitted.  The
extension from squarefree marks to arbitrary powers supported on \(E\)
is recorded separately because it is used in every local restoration.

\begin{lemma}[Omitted-local-score transfer]
\label[lemma]{lem:marked-omitted-local-score}
If \(d\leq y^3\) is supported on \(E\), then, uniformly in \(d\),
\begin{equation}
 \frac{\mathbb E_C(\1_{d\mid m}e^{S_E/L})}
      {\mathbb E_Ce^{S_E/L}}
 =\frac{h_U(v(d))}{d}+O_{B,W}(1/(dL)).
 \label{eq:marked-omitted-local-score}
\end{equation}
For two divisibility indicators the same assertion holds with \(d\)
replaced by their least common multiple, provided that lcm is at most
\(y^3\).  If a bounded-variation physical test \(\phi(m/n)\) is
inserted, its leading Lebesgue mean factors and the remaining error is
\(O_{B,W}(\|\phi\|_{BV}/(dL))\).
\end{lemma}

\begin{proof}
Put
\[
 \Omega_E(m)=\sum_{\substack{W<r\leq y\\r\notin E}}v_r(m).
\]
Writing \(m=dk\), discarding smoothness and the fixed cell restrictions,
and expanding one or two powers of \(\Omega_E\) into divisibility
indicators gives, exactly as an elementary multiple-counting estimate,
\begin{equation}
 \mathbb E_C\{\1_{d\mid m}\Omega_E(m)^j\}
 \ll_{W,j}\frac{(1+\log L)^j}{d}\qquad(j=1,2).
 \label{eq:marked-omitted-local-moments}
\end{equation}
The endpoint contribution is absorbed because \(d\leq y^3\) and
\(n/y^3=n^{1/3}\).  Since \(|S_E|\leq B\Omega_E\) and
\(e^{B\Omega_E/L}\ll_{B,W}1\), Taylor's formula reduces the desired
ratio to bounding the centered first-order term.

Expand \(S_E\) as
\(\sum_{r\notin E}\eta_r\sum_{a\geq1}\1_{r^a\mid m}\).  When
\(r^a\leq y\), the product \(dr^a\leq y^4\), and
Proposition~\ref{prop:marked-structured-cells} gives the main covariance
\[
 \frac{h_U(v(d)+at_r)-h_U(v(d))h_U(at_r)}{dr^a}.
\]
The numerator vanishes at \(at_r=0\) and is
\(O_U(at_r)\) by Lipschitz continuity.  Hence these main terms sum to
\(O_{B,W}(1/d)\), because
\[
 \sum_{W<r\leq y}\sum_{r^a\leq y}\frac{at_r}{r^a}=O_W(1).
\]
Their marked errors sum to
\(O_{B,W}((1+\log L)/(dL))\).  For \(r^a>y\), elementary common-multiple
counting gives
\[
 \mathbb P_C(dr^a\mid m)\ll_W\frac1{dr^a}+\frac1n.
\]
The reciprocal tails are \(O_W(1/d)\); the literal endpoints total
\(O_W(yL/n)=O_W(1/d)\), since \(d\leq y^3\) and
\(y^4L=o(n)\).  We have therefore proved
\begin{equation}
 |\operatorname{Cov}_C(\1_{d\mid m},S_E)|\ll_{B,W}\frac1d.
 \label{eq:marked-omitted-local-first-order}
\end{equation}
Using \eqref{eq:marked-omitted-local-moments}, the second-order Taylor
remainder is
\(O_{B,W}((1+\log L)^2/(dL^2))=o(1/(dL))\).  Division by the analogous
denominator expansion, followed by
\eqref{eq:marked-cell-probability}, proves
\eqref{eq:marked-omitted-local-score}.  A product of two divisibility
indicators is the indicator of their lcm.  Finally, inserting \(\phi\)
and using \eqref{eq:marked-cell-Stieltjes} makes the same Lebesgue
average appear in every leading centered term; it cancels, while the
Stieltjes errors are bounded by
\(O_{B,W}(\|\phi\|_{BV}/(dL))\).  This proves all assertions.
\end{proof}

Restore an omitted prime \(p\).  If
\(\lambda_p=e^{\eta_p/L}\), the exact identity
\begin{equation}
 \lambda_p^{v_p(m)}
 =1+\sum_{a\ge1}(\lambda_p^a-\lambda_p^{a-1})
                         \1_{p^a\mid m}
 \label{eq:marked-local-factor-expansion}
\end{equation}
has
\begin{equation}
 |\lambda_p^a-\lambda_p^{a-1}|
 \le \frac{C_B}{L}e^{Ba/L}.
 \label{eq:marked-local-factor-coefficient}
\end{equation}
Let \(C_*>0\) be a fixed upper bound for every physical endpoint in the
finite cell family.  Fix once and for all a bounded Lipschitz extension
\(F:[0,\infty)\to\mathbb R\) of \(h_U|_{[0,3]}\), with \(F(0)=1\).
For local fugacities \(\lambda_p=e^{\eta_p/L}\), define the finite
coefficients
\begin{equation}
 A_p=\left\lfloor\frac{\log(C_*n)}{\log p}\right\rfloor,\qquad
 c_{p,0}=1,\qquad
 c_{p,a}=
 \begin{cases}
  \lambda_p^a-\lambda_p^{a-1},&1\leq a\leq A_p,\\
 0,&a>A_p,
 \end{cases}
 \label{eq:marked-local-coefficients-definition}
\end{equation}
and similarly for \(q\).  The explicit two-prime model is
\begin{equation}
 Q_{p,q}(r,s):=
 \sum_{a,b\geq0}c_{p,a}c_{q,b}
 \frac{F(\max(r,a)t_p+\max(s,b)t_q)}
 {p^{\max(r,a)}q^{\max(s,b)}}.
 \label{eq:marked-two-prime-model-definition}
\end{equation}
The following lemma supplies the two-local-factor calculation, including
the terms beyond the four-mark range.

\begin{lemma}[Two-local-factor expansion]
\label[lemma]{lem:marked-two-local-expansion}
Let \(p\ne q\) lie in \((W,y]\), put \(I_{p^r}=\1_{p^r\mid m}\),
and allow \(r,s\ge0\), with \(I_{p^0}=I_{q^0}=1\).  If
\(p^rq^s\le y^3\), then the exact restoration of the two omitted local
factors satisfies the explicit formula
\begin{equation}
 {\mathbb E}_{C,S}(I_{p^r}I_{q^s})
 =\frac{Q_{p,q}(r,s)}{Q_{p,q}(0,0)}
 +O_{B,W}\!\left(\frac1{p^rq^sL}+\frac1n\right).
 \label{eq:marked-two-local-total-error}
\end{equation}
Moreover
\(Q_{p,q}(0,0)=1+O_{B,W}(1/(Lp)+1/(Lq)+1/n)\), so the denominator is
bounded away from zero for all sufficiently large \(n\).
In particular, whenever \(p^k\le y^2\),
\begin{equation}
 |\operatorname{Cov}_{C,S}(I_{p^k},I_q)|
 \le C_B\frac{t_q}{p^kq}
       +O_{B,W}\!\left(\frac1{p^kqL}\right).
 \label{eq:marked-tilted-power-covariance}
\end{equation}
The transposed assertion is
\begin{equation}
 |\operatorname{Cov}_{C,S}(I_p,I_{q^l})|
 \le C_B\frac{t_p}{pq^l}
       +O_{B,W}\!\left(\frac1{pq^lL}\right)
 \qquad(q^l\le y^2).
 \label{eq:marked-tilted-power-covariance-transpose}
\end{equation}
\end{lemma}

\begin{proof}
Let \(\nu\) be the probability law obtained by tilting the allowed base
law \(\nu_C\) by \(e^{S_{\{p,q\}}/L}\).  If \(C\) is an actual cell,
\(\nu_C\) is its normalized counting measure; if \(C\) denotes a finite
mixture, \(\nu_C\) is the probability mixture defined in
\eqref{eq:marked-finite-mixture-definition}.  The density of \(\nu\)
relative to \(\nu_C\) is bounded above and below by constants depending
only on \(B,W\).  Elementary multiple counting and
\eqref{eq:marked-cell-size} give the following estimate on every actual
cell.  For a mixture, apply that estimate cellwise to the tilted
numerator and use the uniform two-sided bounds for the (cellwise and
mixture) normalizing constants.  Thus, for every integer \(D\),
\begin{equation}
                  \nu(D\mid m)\ll_{B,W}\frac1D+\frac1n.
 \label{eq:marked-omitted-elementary-bound}
\end{equation}
For \(D\le y^3\) supported on \(p,q\), the stronger formula
\eqref{eq:marked-omitted-local-score} applies.

\paragraph{One local factor.}
Use the coefficients from
\eqref{eq:marked-local-coefficients-definition}.  Since
\(v_p(m)\le A_p\) throughout
the cell family, this truncation leaves
\eqref{eq:marked-local-factor-expansion} exact.  From
\eqref{eq:marked-local-factor-coefficient}, geometric summation gives,
uniformly for \(r\ge0\),
\begin{align}
 \sum_{a\ge1}|c_{p,a}|&\ll_{B,W}1,
 &\sum_{a\ge1}\frac{|c_{p,a}|}{p^a}&\ll_B\frac1{Lp},
 &\sum_{a\ge0}\frac{|c_{p,a}|}{p^{\max(r,a)}}
      &\ll_{B,W}\frac1{p^r}.
 \label{eq:marked-local-coefficient-sums}
\end{align}
The first estimate also follows directly by telescoping separately when
\(\lambda_p\ge1\) and \(\lambda_p<1\); indeed \(A_p=O_W(L)\).

\paragraph{Two local factors.}
Multiplying the two identities
\eqref{eq:marked-local-factor-expansion} gives the exact numerator
\begin{equation}
 \sum_{a,b\ge0}c_{p,a}c_{q,b}\,
 \nu\!\left(p^{\max(r,a)}q^{\max(s,b)}\mid m\right).
 \label{eq:marked-two-local-exact-numerator}
\end{equation}
The denominator is the same expression with \(r=s=0\).  Split the
\((a,b)\)-sum according as
\begin{equation}
       p^{\max(r,a)}q^{\max(s,b)}\le y^3
       \quad\hbox{or}\quad
       p^{\max(r,a)}q^{\max(s,b)}>y^3.
 \label{eq:marked-two-local-tail-split}
\end{equation}
On the first part insert \eqref{eq:marked-omitted-local-score}.  By
\eqref{eq:marked-local-coefficient-sums}, the sum of all its marked
errors is
\begin{equation}
 \ll_{B,W}\frac1L
 \sum_{a,b\ge0}\frac{|c_{p,a}c_{q,b}|}
 {p^{\max(r,a)}q^{\max(s,b)}}
 \ll_{B,W}\frac1{p^rq^sL}.
 \label{eq:marked-two-local-inside-error}
\end{equation}
On the complementary part, at least one of \(a>r\) or \(b>s\) holds.
Using \eqref{eq:marked-local-factor-coefficient} in that coordinate and
\eqref{eq:marked-local-coefficient-sums} in the other gives explicitly
\begin{align}
 &\sum_{\substack{a,b\ge0\\
 p^{\max(r,a)}q^{\max(s,b)}>y^3}}
 \frac{|c_{p,a}c_{q,b}|}
 {p^{\max(r,a)}q^{\max(s,b)}}
 \notag\\
 &\qquad\le
 \left(\sum_{a>r}\frac{|c_{p,a}|}{p^a}\right)
       \left(\sum_{b\ge0}\frac{|c_{q,b}|}{q^{\max(s,b)}}\right)
 +\left(\sum_{a\ge0}\frac{|c_{p,a}|}{p^{\max(r,a)}}\right)
       \left(\sum_{b>s}\frac{|c_{q,b}|}{q^b}\right)
 \ll_{B,W}\frac1{p^rq^sL}.
 \label{eq:marked-two-local-reciprocal-tail}
\end{align}
The endpoint part of \eqref{eq:marked-omitted-elementary-bound} is
separate and is
\begin{equation}
 \frac1n\sum_{a,b\ge0}|c_{p,a}c_{q,b}|\ll_{B,W}\frac1n.
 \label{eq:marked-two-local-endpoint-tail}
\end{equation}
This proves the asserted tail split with every floor contribution
retained.

For completeness, we now identify the covariance in the explicit model
\eqref{eq:marked-two-prime-model-definition}.
Replacing the terms outside \eqref{eq:marked-two-local-tail-split} by
this extension costs no more than
\eqref{eq:marked-two-local-reciprocal-tail}.  Hence
\begin{equation}
 {\mathbb E}_{C,S}(I_{p^r}I_{q^s})
 =\frac{Q_{p,q}(r,s)}{Q_{p,q}(0,0)}
 +O_{B,W}\!\left(\frac1{p^rq^sL}+\frac1n\right),
 \label{eq:marked-two-local-model-approximation}
\end{equation}
The \((a,b)=(0,0)\) term of \(Q_{p,q}(0,0)\) is one, while
\eqref{eq:marked-local-coefficient-sums} bounds all remaining terms by
\(O_{B,W}(1/(Lp)+1/(Lq))\).  This proves the denominator assertion in
the lemma (with harmless room for the literal endpoint), and in
particular \(Q_{p,q}(0,0)\asymp_{B,W}1\).

\paragraph{Rank-one cancellation.}
Write
\[
 K_F(x,z)=F(x+z)-F(x)F(z).
\]
Lipschitz continuity and \(F(0)=1\) imply
\begin{equation}
 |K_F(x,z)|\ll_U z,\qquad K_F(0,z)=K_F(x,0)=0.
 \label{eq:marked-two-local-product-kernel}
\end{equation}
If
\[
 \mathcal P_p(r)=\sum_{a\ge0}c_{p,a}
       \frac{F(\max(r,a)t_p)}{p^{\max(r,a)}},
\]
and \(\mathcal P_q(s)\) is defined similarly, then
\(Q_{p,q}(r,s)=\mathcal P_p(r)\mathcal P_q(s)+\mathcal R(r,s)\).
Equations \eqref{eq:marked-local-coefficient-sums} and
\eqref{eq:marked-two-local-product-kernel} give
\begin{align}
 |\mathcal P_p(r)|&\ll_{B,W}p^{-r},
 &|\mathcal P_q(1)|&\ll_{B,W}q^{-1},\notag\\
 |\mathcal R(r,1)|&\ll_{B,W}\frac{t_q}{p^rq},
 &|\mathcal R(r,0)|&\ll_{B,W}\frac{t_q}{p^rLq},\notag\\
 |\mathcal R(0,1)|&\ll_{B,W}\frac{t_q}{Lpq},
 &|\mathcal R(0,0)|&\ll_{B,W}\frac{t_q}{L^2pq}.
 \label{eq:marked-two-local-remainder-bounds}
\end{align}
For example, the \(q\)-sum in the first remainder is bounded by
\[
 \sum_{b\ge0}|c_{q,b}|
       \frac{\max(1,b)t_q}{q^{\max(1,b)}}
 \le\frac{t_q}{q}
   +\frac{C_Bt_q}{L}\sum_{b\ge1}\frac b{q^b}
 \ll_B\frac{t_q}{q},
\]
whereas for \(\mathcal R(r,0)\) its \(b=0\) term vanishes and the
extra factor \(L^{-1}\) remains.  The two analogous \(p\)-sums give
the last two bounds.

Expand
\(
 Q_{p,q}(r,1)Q_{p,q}(0,0)
 -Q_{p,q}(r,0)Q_{p,q}(0,1)
\).
The rank-one products \(\mathcal P_p\mathcal P_q\) cancel exactly, and
\eqref{eq:marked-two-local-remainder-bounds} bounds what remains by
\begin{equation}
 \left|Q_{p,q}(r,1)Q_{p,q}(0,0)
 -Q_{p,q}(r,0)Q_{p,q}(0,1)\right|
 \ll_{B,W}\frac{t_q}{p^rq}.
 \label{eq:marked-two-local-model-covariance}
\end{equation}
Combining this with the four instances of
\eqref{eq:marked-two-local-model-approximation} gives
\[
 |\operatorname{Cov}_{C,S}(I_{p^r},I_q)|
 \le C_B\frac{t_q}{p^rq}
 +O_{B,W}\!\left(\frac1{p^rqL}+\frac1n\right).
\]
When \(p^r\le y^2\), the endpoint term is absorbed because
\(p^rq\le y^3\) and \(n/y^3=n^{1/3}\).  This proves
\eqref{eq:marked-tilted-power-covariance}.  Interchanging \(p,r\) and
\(q,s\) proves \eqref{eq:marked-tilted-power-covariance-transpose}.
\end{proof}

We now perform all four aggregate prime-power sums.  Let
\begin{equation}
 I_p=\1_{p\mid m},\qquad V_p=v_p(m),\qquad
 J_p:=V_p-I_p=\sum_{k\ge2}I_{p^k},
 \label{eq:marked-IJV-definitions}
\end{equation}
and use the weighted row norm
\begin{equation}
 \|K\|_{\rm row}:=\sup_{W<p\le y}
      p\sum_{W<q\le y}|K(p,q)|.
 \label{eq:marked-row-norm}
\end{equation}

The order in which the local powers are summed matters for the later
moving-low-cell argument.  We therefore record a box-uniform refinement
of the preceding two-local calculation.  Its leading constant is
independent of the radius of the compact tilt box.

\begin{lemma}[Product-weighted aggregate power transfer]
\label[lemma]{lem:marked-product-weighted-power-transfer}
There is a constant \(C_{\rm pow}>0\), depending only on the fixed
Dickman parameter and the fixed physical compact set, with the following
property.  For every fixed \(B,W\) there is a quantity
\(\epsilon_{B,W}(n)\to0\) such that, for distinct primes
\(W<p,q\le y\),
\begin{align}
 \sum_{k\ge2}
   |\operatorname{Cov}_{C,S}(I_{p^k},I_q)|
 &\le C_{\rm pow}\frac{t_pt_q}{p^2q}
      +\frac{\epsilon_{B,W}(n)}{p^2q}
      +O_{B,W}(L/n),
 \label{eq:marked-product-JI}\\
 \sum_{l\ge2}
   |\operatorname{Cov}_{C,S}(I_p,I_{q^l})|
 &\le C_{\rm pow}\frac{t_pt_q}{pq^2}
      +\frac{\epsilon_{B,W}(n)}{pq^2}
      +O_{B,W}(L/n),
 \label{eq:marked-product-IJ}\\
 \sum_{k,l\ge2}
   |\operatorname{Cov}_{C,S}(I_{p^k},I_{q^l})|
 &\le C_{\rm pow}\frac{t_pt_q}{p^2q^2}
      +\frac{\epsilon_{B,W}(n)}{p^2q^2}
      +O_{B,W}(L^2/n).
 \label{eq:marked-product-JJ}
\end{align}
Moreover
\begin{equation}
 \mathbb E_{C,S}J_p^2
 \le \frac{C_{\rm pow}}{p^2}
       +\frac{\epsilon_{B,W}(n)}{p^2}
       +O_{B,W}(L^2/n).
 \label{eq:marked-product-diagonal}
\end{equation}
Consequently
\begin{equation}
 \left\|\operatorname{Cov}_{C,S}(V_p,V_q)
       -\operatorname{Cov}_{C,S}(I_p,I_q)\right\|_{\rm row}
 \le \frac{C_{\rm pow}}W+\epsilon_{B,W}(n),
 \label{eq:marked-box-uniform-power-row}
\end{equation}
after enlarging \(C_{\rm pow}\) by an absolute factor.  In particular,
the coefficient of \(1/W\) in this estimate does not depend on \(B\).
Here \(\epsilon_{B,W}(n)\) denotes one fixed maximum of the finitely
many pair, tail, endpoint, and row remainders appearing in the proof;
the explicit admissible choice is given at the end of the proof.  Thus
the same function is used simultaneously in all four displays.
\end{lemma}

\begin{proof}
\medskip\noindent\emph{Four-mark chamber transfer.}\quad
We first transfer the structured-cell formula through the complete
four-mark range while omitting the forced local primes from the tilt.
For a set \(E\) of one or two primes in \((W,y]\), put
\[
 \Omega_E(m)=\sum_{\substack{W<r\le y\\r\notin E}}v_r(m),
 \qquad
 S_E(m)=\sum_{\substack{W<r\le y\\r\notin E}}\eta_rv_r(m),
\]
and let \(\nu_E\) be the probability law obtained by tilting the allowed
base law \(\nu_C\) by \(e^{S_E/L}\).  If \(D\) is supported on \(E\),
then \(\Omega_E(Da)=\Omega_E(a)\).  For each actual cell \(C_i\), one has
\(\#C_i\asymp_W n\), and elementary multiple counting consequently gives
the following bound on \(C_i\), uniformly for \(D\le y^4\).  Since the
number of cells is fixed, averaging the cellwise expectations with the
weights in \eqref{eq:marked-finite-mixture-definition} gives exactly the
same bound for an allowed finite mixture:
\begin{equation}
 \mathbb E_C\!\left(\1_{D\mid m}\Omega_E(m)^j\right)
 \ll_{W,j}\frac{(1+\log L)^j}{D}\qquad(j=1,2).
 \label{eq:marked-omitted-four-moments}
\end{equation}
Here are the details needed for the second moment.  On writing \(m=Da\)
and discarding smoothness and the cell restrictions, expand
\[
 \Omega_E(a)=\sum_{\substack{W<r\le y\\r\notin E}}
                 \sum_{b\ge1}\1_{r^b\mid a}.
\]
The distinct-prime terms in its square have mean at most \(A_y^2\),
where
\[
 A_y:=\sum_{W<r\le y}\sum_{b\ge1}\frac1{r^b}
      =\sum_{W<r\le y}\frac1{r-1}\ll_W1+\log L.
\]
The same-prime terms have mean at most
\[
 \sum_{W<r\le y}\sum_{a,b\ge1}\frac1{r^{\max(a,b)}}
 =\sum_{W<r\le y}\sum_{c\ge1}\frac{2c-1}{r^c}
 \ll_W1+\log L.
\]
This proves \eqref{eq:marked-omitted-four-moments} literally.

We have \(|S_E|\le B\Omega_E\), while
\(\Omega_E(m)\log W\le\log(C_*n)\) for a fixed physical endpoint
\(C_*\).  Thus \(e^{B\Omega_E/L}\ll_{B,W}1\) pointwise.  Taylor's
formula and \eqref{eq:marked-omitted-four-moments} show that
\begin{align*}
 \mathbb E_C(\1_{D\mid m}e^{S_E/L})
 &=\mathbb P_C(D\mid m)
   +O_{B,W}\!\left(\frac{1+\log L}{DL}
              +\frac{(1+\log L)^2}{DL^2}\right),\\
 \mathbb E_Ce^{S_E/L}
 &=1+O_{B,W}\!\left(\frac{1+\log L}{L}
              +\frac{(1+\log L)^2}{L^2}\right).
\end{align*}
The marked formula \eqref{eq:marked-cell-probability}, valid uniformly
through \(D\le y^4\), may now be inserted before division.  We obtain
\begin{equation}
 \nu_E(D\mid m)=\frac{h_U(v(D))}{D}
       +O_{B,W}\!\left(\frac{1+\log L}{DL}\right)
 \qquad(D\le y^4, \operatorname{supp}D\subset E).
 \label{eq:marked-omitted-four-transfer}
\end{equation}
There is no separate endpoint here: \(n/D\ge n^{1/9}\), so the endpoint
in the structured-cell count is absorbed by \(1/(DL)\).  For arbitrary
\(D\), the bounded density ratio and elementary multiple counting give
the fallback
\begin{equation}
                         \nu_E(D\mid m)\ll_{B,W}D^{-1}+n^{-1}.
 \label{eq:marked-omitted-four-fallback}
\end{equation}

We next record the product estimate for the Dickman profile using only
the regularity available across its kink.  Put \(H=h_U|_{[0,4]}\).
The function \(H\) is bounded and Lipschitz, \(H(0)=1\), and it is
\(C^{1,1}\) on a fixed neighborhood of zero.  Hence
\begin{equation}
 |H(x+z)-H(x)H(z)|\le C_Kxz
 \qquad(x,z\ge0,\ x+z\le4),
 \label{eq:marked-four-product-kernel}
\end{equation}
where \(C_K\) depends only on \(U\).  Indeed, in a sufficiently small
square at the origin the kernel vanishes on both axes and its mixed
derivative is bounded.  If \(x\) is in that small interval and \(z\)
is outside it, Lipschitz continuity and \(H(0)=1\) give
\[
 |H(x+z)-H(x)H(z)|
 \le |H(x+z)-H(z)|+|H(z)|\,|H(x)-1|\ll x\ll xz;
\]
the transposed case is identical.  When both variables are outside the
small interval, boundedness gives \(O(1)\ll xz\).  This proves
\eqref{eq:marked-four-product-kernel}; in particular, no
\(C^{1,1}\) assertion on all of \([0,4]\) is being used.

It remains to restore the primes of \(E\) exactly.  Use the truncated
coefficients \(c_{p,a}\) from the proof of
\Cref{lem:marked-two-local-expansion}, so that
\eqref{eq:marked-local-factor-expansion} is exact on every allowed cell.
Fix \(B,W\) and take \(n\) large enough that \(\lambda_p\le3/2\).
For \(r\ge0\), define
\[
 D_p(r):=\sum_{a\ge1}\frac{|c_{p,a}|}{p^{\max(r,a)}},
 \qquad C_p(r):=p^{-r}+D_p(r).
\]
Splitting at \(a=r\), using
\(|c_{p,a}|\le(2B/L)\lambda_p^{a-1}\) and
\(\lambda_p^{A_p}\ll_{B,W}1\), and summing the geometric tail gives
\begin{align}
 \sum_{a\ge1}|c_{p,a}|&\ll_{B,W}1,
 &D_p(0)&\ll_B\frac1{Lp},\notag\\
 D_p(r)&\ll_{B,W}\frac{r+1}{Lp^r}\quad(r\ge1),
 &C_p(r)&\ll_{B,W}p^{-r}\quad(r\ge1),
 \label{eq:marked-four-local-coefficients}
\end{align}
and therefore
\begin{align}
 \sum_{r\ge2}D_p(r)&\ll_{B,W}\frac1{Lp^2},
 &\sum_{r\ge2}C_p(r)&\ll_{B,W}\frac1{p^2},\notag\\
 \sum_{r\ge2}(2r-3)D_p(r)&\ll_{B,W}\frac1{Lp^2},
 &\sum_{r\ge2}\frac{2r-3}{p^r}&\ll\frac1{p^2}.
 \label{eq:marked-four-local-aggregate}
\end{align}
For example, in \(D_p(r)\) the two ranges satisfy
\[
 p^{-r}\sum_{1\le a\le r}|c_{p,a}|
 \ll_{B,W}\frac r{Lp^r},\qquad
 \sum_{a>r}\frac{|c_{p,a}|}{p^a}
 \ll_{B,W}\frac1{Lp^{r+1}}.
\]
These estimates also prove the aggregate displays.

Take first \(E=\{p,q\}\).  If \(N_{r,s}\) is the numerator for
\(p^rq^s\mid m\) after restoring the two local fugacities, exact
expansion gives
\begin{equation}
 N_{r,s}=\sum_{a,b\ge0}c_{p,a}c_{q,b}\,
 \nu_E\!\left(p^{\max(r,a)}q^{\max(s,b)}\mid m\right),
 \qquad
 \mathbb P_{C,S}(p^rq^s\mid m)=\frac{N_{r,s}}{N_{0,0}}.
 \label{eq:marked-four-exact-restoration}
\end{equation}
Apply \eqref{eq:marked-omitted-four-fallback} to every term other than
\((a,b)=(0,0)\).  The coefficient sums just proved yield
\begin{equation}
 |N_{r,s}-\nu_E(p^rq^s\mid m)|
 \ll_{B,W}D_p(r)C_q(s)+p^{-r}D_q(s)+n^{-1}.
 \label{eq:marked-four-restoration-error}
\end{equation}
The endpoint is really \(O_{B,W}(n^{-1})\), because the unweighted
absolute coefficient sums are bounded.  Taking \(r=s=0\) shows
\begin{equation}
 N_{0,0}=1+O_{B,W}\!\left(\frac1{Lp}+\frac1{Lq}+\frac1n\right).
 \label{eq:marked-four-restoration-denominator}
\end{equation}
Thus the local denominator is \(1+o_{B,W}(1)\), not merely bounded by a
box-dependent constant.

For clarity, we make the summation implicit in the next displays
literal.  Put
\[
 a_n=C_{B,W}\frac{1+\log L}{L},\qquad
 z_{p,q}=C_{B,W}\left(\frac1{Lp}+\frac1{Lq}+\frac1n\right),
\]
where the constants are enlarged once below.  Division by
\eqref{eq:marked-four-restoration-denominator}, together with
\eqref{eq:marked-omitted-four-transfer} and
\eqref{eq:marked-four-restoration-error}, gives, whenever
\(rt_p+st_q\le4\),
\begin{align}
 &\left|\mathbb P_{C,S}(p^rq^s\mid m)
       -\frac{H(rt_p+st_q)}{p^rq^s}\right|\notag\\
 &\quad\ll_{B,W}
 \frac{a_n+z_{p,q}}{p^rq^s}
 +D_p(r)C_q(s)+p^{-r}D_q(s)+\frac1n.
 \label{eq:marked-four-pointwise-error-ledger}
\end{align}
The one-prime restoration gives symmetrically
\begin{equation}
 \left|\mathbb P_{C,S}(p^r\mid m)-\frac{H(rt_p)}{p^r}\right|
 \ll_{B,W}\frac{a_n+z_p}{p^r}+D_p(r)+\frac1n,
 \qquad z_p\ll_{B,W}\frac1{Lp}+\frac1n.
 \label{eq:marked-four-marginal-pointwise-ledger}
\end{equation}
The required geometric sums are
\begin{equation}
 \sum_{r\ge2}\frac1{p^r}=\frac1{p(p-1)},\qquad
 \sum_{r\ge2}\frac r{p^r}=\frac{2p-1}{p(p-1)^2},\qquad
 \sum_{r\ge2}\frac{2r-3}{p^r}=\frac{p+1}{p(p-1)^2};
 \label{eq:marked-four-geometric-ledger}
\end{equation}
in particular, all three are \(O(p^{-2})\).

Combine \eqref{eq:marked-omitted-four-transfer} and
\eqref{eq:marked-four-local-aggregate}--
\eqref{eq:marked-four-geometric-ledger}, and sum these pointwise error
majorants at fixed \(p,q\) before passing to a prime-row norm.  There is
a quantity
\begin{equation}
 \alpha_{B,W}(n)\ll_{B,W}\frac{1+\log L}{L},
 \qquad \alpha_{B,W}(n)\log L\longrightarrow0,
 \label{eq:marked-four-pair-error-rate}
\end{equation}
such that, throughout the indicated four-mark ranges,
\begin{align}
 \sum_{\substack{r\ge2\\rt_p+t_q\le4}}
 \left|\mathbb P_{C,S}(p^rq\mid m)
       -\frac{H(rt_p+t_q)}{p^rq}\right|
 &\le\frac{\alpha_{B,W}(n)}{p^2q}+O_{B,W}(L/n),\notag\\
 \sum_{\substack{s\ge2\\t_p+st_q\le4}}
 \left|\mathbb P_{C,S}(pq^s\mid m)
       -\frac{H(t_p+st_q)}{pq^s}\right|
 &\le\frac{\alpha_{B,W}(n)}{pq^2}+O_{B,W}(L/n),\notag\\
 \sum_{\substack{r,s\ge2\\rt_p+st_q\le4}}
 \left|\mathbb P_{C,S}(p^rq^s\mid m)
       -\frac{H(rt_p+st_q)}{p^rq^s}\right|
 &\le\frac{\alpha_{B,W}(n)}{p^2q^2}+O_{B,W}(L^2/n).
 \label{eq:marked-four-joint-aggregate}
\end{align}
The one-prime restoration gives in exactly the same way
\begin{align}
 \sum_{\substack{r\ge2\\rt_p\le4}}
 \left|\mathbb P_{C,S}(p^r\mid m)-\frac{H(rt_p)}{p^r}\right|
 &\le\frac{\alpha_{B,W}(n)}{p^2}+O_{B,W}(L/n),\notag\\
 \left|\mathbb P_{C,S}(q\mid m)-\frac{H(t_q)}q\right|
 &\le\frac{\alpha_{B,W}(n)}q+O_{B,W}(1/n).
 \label{eq:marked-four-marginal-aggregate}
\end{align}
By symmetry, the high-power marginal needed in the \(IJ\) and \(JJ\)
products is
\begin{equation}
 \sum_{\substack{s\ge2\\st_q\le4}}
 \left|\mathbb P_{C,S}(q^s\mid m)-\frac{H(st_q)}{q^s}\right|
 \le\frac{\alpha_{B,W}(n)}{q^2}+O_{B,W}(L/n);
 \label{eq:marked-four-symmetric-marginal-aggregate}
\end{equation}
the analogous first-power estimate at \(p\) follows by exchanging
\(p,q\) in the second line of
\eqref{eq:marked-four-marginal-aggregate}.
All errors in these displays have been accounted for explicitly:
\eqref{eq:marked-omitted-four-transfer} contributes
\(O_{B,W}((1+\log L)/L)\) times the reciprocal series,
\eqref{eq:marked-four-local-aggregate} contributes \(O_{B,W}(1/L)\),
and there are only \(O_W(L)\), respectively \(O_W(L^2)\), endpoint
terms.

Let
\[
 \Gamma_{r,s}:=\operatorname{Cov}_{C,S}(I_{p^r},I_{q^s}),\qquad
 \mathcal K_{r,s}:=
 \frac{H(rt_p+st_q)-H(rt_p)H(st_q)}{p^rq^s}.
\]
Use
\(
 |ab-a_0b_0|\le |a-a_0||b|+|a_0||b-b_0|
\)
with \eqref{eq:marked-four-joint-aggregate}--
\eqref{eq:marked-four-symmetric-marginal-aggregate}.  For the \(JJ\)
marginal error only, enlarge the coupled chamber to the containing
Cartesian product.  The geometric sums in
\eqref{eq:marked-four-geometric-ledger} then give the literal covariance
error ledger
\begin{align}
 \sum_{\substack{r\ge2\\rt_p+t_q\le4}}
       |\Gamma_{r,1}-\mathcal K_{r,1}|
 &\le\frac{\alpha_{B,W}(n)}{p^2q}+O_{B,W}(L/n),\notag\\
 \sum_{\substack{s\ge2\\t_p+st_q\le4}}
       |\Gamma_{1,s}-\mathcal K_{1,s}|
 &\le\frac{\alpha_{B,W}(n)}{pq^2}+O_{B,W}(L/n),\notag\\
 \sum_{\substack{r,s\ge2\\rt_p+st_q\le4}}
       |\Gamma_{r,s}-\mathcal K_{r,s}|
 &\le\frac{\alpha_{B,W}(n)}{p^2q^2}+O_{B,W}(L^2/n).
 \label{eq:marked-four-covariance-error-ledger}
\end{align}
Finally \eqref{eq:marked-four-product-kernel} and
\(\sum_{r\ge2}r/p^r\ll p^{-2}\) show that the complete four-mark chamber
contributes at most
\begin{align}
 C\frac{t_pt_q}{p^2q}
   +\frac{\alpha_{B,W}(n)}{p^2q}+O_{B,W}(L/n),\qquad
 C\frac{t_pt_q}{pq^2}
   +\frac{\alpha_{B,W}(n)}{pq^2}+O_{B,W}(L/n),\notag\\
 C\frac{t_pt_q}{p^2q^2}
   +\frac{\alpha_{B,W}(n)}{p^2q^2}+O_{B,W}(L^2/n),
 \label{eq:marked-four-covariance-aggregate}
\end{align}
in the \(JI\), \(IJ\), and \(JJ\) orientations.  The constant \(C\)
depends only on \(U\), and in particular is independent of \(B,W\).

\medskip\noindent\emph{Beyond-four-mark reciprocal tails.}\quad
We now treat every forced product beyond \(y^4\).  The full tilted
measure also has bounded density ratio, so
\begin{equation}
 \mathbb P_{C,S}(D\mid m)\ll_{B,W}D^{-1}+n^{-1}.
 \label{eq:marked-full-tilt-tail-count}
\end{equation}
For \(JI\), if \(rt_p+t_q>4\), then
\[
 (r-2)t_p>4-t_q-2t_p\ge1,
\]
so the reciprocal tail is at most
\(O_{B,W}(Ly^{-1}/(p^2q))\).  The transpose is identical.  For \(JJ\),
put \(E_{r,s}=(r-2)t_p+(s-2)t_q\).  The case \(r=s=2\) cannot occur.
If exactly one exponent exceeds two, say \(r\ge3,s=2\), then
\(rt_p>2\) and
\[
 E_{r,2}=(r-2)t_p>\frac{2(r-2)}r\ge\frac23.
\]
If both exceed two, then
\(rt_p+st_q=2(t_p+t_q)+E_{r,s}\le3E_{r,s}\), so
\(E_{r,s}>4/3\).  The \(JJ\) reciprocal tail is consequently
\(O_{B,W}(L^2y^{-2/3}/(p^2q^2))\).  The elementary inequality
\[
 |\operatorname{Cov}(\1_A,\1_B)|
 \le\mathbb P(A\cap B)+\mathbb P(A)\mathbb P(B)
\]
and \eqref{eq:marked-full-tilt-tail-count} also show that all literal
endpoints are \(O_{B,W}(L/n)\) in \(JI,IJ\) and
\(O_{B,W}(L^2/n)\) in \(JJ\).  Thus the complete tail and endpoint
ledger is
\begin{equation}
 \begin{array}{c|c|c}
  &\text{reciprocal tail}&\text{literal endpoint}\\ \hline
 JI&O_{B,W}(Ly^{-1}/(p^2q))&O_{B,W}(L/n)\\
 IJ&O_{B,W}(Ly^{-1}/(pq^2))&O_{B,W}(L/n)\\
 JJ&O_{B,W}(L^2y^{-2/3}/(p^2q^2))&O_{B,W}(L^2/n).
 \end{array}
 \label{eq:marked-four-tail-endpoint-ledger}
\end{equation}
Together with
\eqref{eq:marked-four-covariance-aggregate}, this proves
\eqref{eq:marked-product-JI}--\eqref{eq:marked-product-JJ} with an
intermediate pair error
\begin{equation}
 \epsilon^{(0)}_{B,W}(n)
 \ll_{B,W}\frac{1+\log L}{L}+\frac Ly+\frac{L^2}{y^{2/3}},
 \qquad \epsilon^{(0)}_{B,W}(n)\log L\longrightarrow0.
 \label{eq:marked-pair-error-explicit-rate}
\end{equation}

For the diagonal, repeat the omitted-score and exact-restoration argument
with \(E=\{p\}\).  Since
\[
 J_p^2=\sum_{r\ge2}(2r-3)\1_{p^r\mid m},
\]
\eqref{eq:marked-four-local-aggregate} and the one-prime version of
\eqref{eq:marked-omitted-four-transfer} give, for \(rt_p\le4\), an
absolute leading bound \(C/p^2\), an
\(\epsilon^{(0)}_{B,W}(n)/p^2\) error, and an
\(O_{B,W}(L^2/n)\) endpoint.  If \(rt_p>4\), then
\((r-2)t_p>4-2t_p\ge2\), so the remaining reciprocal tail is
\(O_{B,W}(L^2y^{-2}/p^2)\).  This proves
\eqref{eq:marked-product-diagonal} with a leading constant independent
of \(B,W\).

\medskip\noindent\emph{Aggregate row norm.}\quad
Finally, sum the three off-diagonal bounds in the row norm.  For fixed
\(p\), the contraction uses the following literal finite-sum
factorizations (with \(q=p\) omitted throughout):
\[
 p\sum_{q\ne p}
 \left(\frac{C_{\rm pow}t_pt_q}{p^2q},
       \frac{C_{\rm pow}t_pt_q}{pq^2},
       \frac{C_{\rm pow}t_pt_q}{p^2q^2}\right)
 =C_{\rm pow}\left(\frac{t_p}{p}\sum_{q\ne p}\frac{t_q}{q},
          t_p\sum_{q\ne p}\frac{t_q}{q^2},
          \frac{t_p}{p}\sum_{q\ne p}\frac{t_q}{q^2}\right).
\]
The same identities with the \(t\)'s removed factor the three error
terms.  Thus no interchange of an unscaled \(o(1)\) with the prime sum is
being used.  These exact identities give the explicit orientation ledger
\begin{align}
 \mathcal R_{JI}^{\rm main}
 &\ll \frac{t_p}{p}\sum_q\frac{t_q}{q},
 &\mathcal R_{JI}^{\rm err}
 &\ll \frac{\epsilon^{(0)}_{B,W}(n)}p\sum_q\frac1q,\notag\\
 \mathcal R_{IJ}^{\rm main}
 &\ll t_p\sum_q\frac{t_q}{q^2},
 &\mathcal R_{IJ}^{\rm err}
 &\ll \epsilon^{(0)}_{B,W}(n)\sum_q\frac1{q^2},\notag\\
 \mathcal R_{JJ}^{\rm main}
 &\ll \frac{t_p}{p}\sum_q\frac{t_q}{q^2},
 &\mathcal R_{JJ}^{\rm err}
 &\ll \frac{\epsilon^{(0)}_{B,W}(n)}p\sum_q\frac1{q^2}.
 \label{eq:marked-row-orientation-ledger}
\end{align}
The leading terms use
\[
 \sum_{W<q\le y}\frac{t_q}{q}=O(1),\qquad
 \sum_{q>W}\frac1{q^2}\ll\frac1W.
\]
The only harmonic loss is in the unweighted \(JI\) pair error, and
\eqref{eq:marked-pair-error-explicit-rate} gives
\[
 \epsilon^{(0)}_{B,W}(n)\sum_{q\le y}\frac1q
 \ll\epsilon^{(0)}_{B,W}(n)\log L=o(1).
\]
For the endpoints, \(p\le y\), \(\pi(y)\le y\), and the worst \(JJ\)
term in \eqref{eq:marked-four-tail-endpoint-ledger} give explicitly
\[
 p\,\pi(y)\frac{L^2}{n}\le\frac{y^2L^2}{n}=o(1).
\]
On the
diagonal, \(I_pJ_p=J_p\), \(0\le J_p\le J_p^2\), and hence
\[
 |\operatorname{Var}(V_p)-\operatorname{Var}(I_p)|
 \le2|\operatorname{Cov}(I_p,J_p)|+\operatorname{Var}(J_p)
 \le3\mathbb E J_p^2.
\]
Thus its weighted row contribution is \(O(1/p)\le O(1/W)\), with the
same box-independent leading constant.  With a sufficiently large
constant \(C_{B,W}\), define
\[
 \epsilon_{B,W}(n):=C_{B,W}\left\{
   \epsilon^{(0)}_{B,W}(n)\log L+\frac{y^2L^2}{n}\right\}.
\]
This single function simultaneously dominates the pair and row
remainders and tends to zero.
This proves \eqref{eq:marked-box-uniform-power-row} and completes the
proof.
\end{proof}

We retain the following direct row calculation as an explicit verification
of the truncated marked range, every high-power tail, and every endpoint.
Its crude \(B\)-dependent leading constants are not used to choose \(W\);
that box-uniform choice comes only from
\Cref{lem:marked-product-weighted-power-transfer}.  The calculation also
supplies the physical and coefficient-weighted corollaries used below.

First, \eqref{eq:marked-tilted-power-covariance} and Mertens' estimates
give the \(J_pI_q\) orientation
\begin{align}
 &p\sum_{\substack{W<q\le y\\q\ne p}}
   \sum_{\substack{k\ge2\\p^k\le y^2}}
 |\operatorname{Cov}_{C,S}(I_{p^k},I_q)|\notag\\
 &\qquad\le C_Bp\sum_{k\ge2}\frac1{p^k}
             \sum_{q\le y}\frac{t_q}{q}
 +O_{B,W}\!\left\{\frac pL\sum_{k\ge2}\frac1{p^k}
             \sum_{q\le y}\frac1q\right\}
 \ll_B\frac1p+o_{W,n}(1/p).
 \label{eq:marked-power-row-JI}
\end{align}
For the transpose \(I_pJ_q\),
\begin{align}
 &p\sum_{\substack{W<q\le y\\q\ne p}}
   \sum_{\substack{l\ge2\\q^l\le y^2}}
 |\operatorname{Cov}_{C,S}(I_p,I_{q^l})|\notag\\
 &\qquad\le
 \left\{C_Bt_p+O_{B,W}(L^{-1})\right\}
       \sum_{q>W}\sum_{l\ge2}\frac1{q^l}
 \ll_B\frac1W+o_{W,n}(1).
 \label{eq:marked-power-row-IJ}
\end{align}
Here \(\sum t_q/q=O(1)\), \(L^{-1}\sum1/q=o(1)\), and
\(\sum_{q>W,l\ge2}q^{-l}\ll W^{-1}\).

The bounded tilt has density ratio \(O_B(1)\) relative to \(C\),
independently of \(W\).  Thus Proposition~\ref{prop:marked-structured-cells} gives,
for distinct \(p,q\) and \(p^kq^l\le y^4\),
\begin{equation}
 |\operatorname{Cov}_{C,S}(I_{p^k},I_{q^l})|
 \ll_B\frac1{p^kq^l}
       +O_{B,W}\!\left(\frac1{p^kq^lL}\right).
 \label{eq:marked-common-multiple-covariance}
\end{equation}
For arbitrary \(k,l\), elementary common-multiple counting gives the
fallback
\begin{equation}
 |\operatorname{Cov}_{C,S}(I_{p^k},I_{q^l})|
 \ll_{B,W}\frac1{p^kq^l}+\frac1n.
 \label{eq:marked-common-multiple-covariance-tail}
\end{equation}
There are \(O_W(L)\) relevant powers at each prime.  Consequently the
part with \(p^kq^l\le y^4\) contributes
\[
 \ll_B\frac1{pW}+o_{W,n}(1).
\]
For the reciprocal part beyond \(y^4\), split according as
\(p^k>y^2\) or \(p^k\le y^2\); in the latter case necessarily
\(q^l>y^2\).  Hence
\begin{align}
 p\sum_{\substack{q>W,\ k,l\ge2\\p^kq^l>y^4}}
       \frac1{p^kq^l}
 &\ll_W
 py^{-2}\sum_{q>W,l\ge2}\frac1{q^l}
 +p\sum_{k\ge2}\frac1{p^k}
       \sum_{\substack{q\le y,\ l\ge2\\q^l>y^2}}\frac1{q^l}\notag\\
 &\ll_W\frac1y.
 \label{eq:marked-JJ-reciprocal-tail}
\end{align}
Using \eqref{eq:marked-common-multiple-covariance-tail} on this part,
the complete two-extra-power orientation is therefore
\begin{align}
 p\sum_{\substack{W<q\le y\\q\ne p}}
 |\operatorname{Cov}_{C,S}(J_p,J_q)|
 &\ll_B\frac1{pW}+o_{W,n}(1)
       +O_{B,W}\!\left(\frac{pyL^2}{n}\right).
 \label{eq:marked-power-row-JJ}
\end{align}
The endpoint term is uniform and is \(o_n(1)\), since
\(p\le y=n^{2/9}\).

The portions omitted from \eqref{eq:marked-power-row-JI} and
\eqref{eq:marked-power-row-IJ} are also explicit.  From
\eqref{eq:marked-common-multiple-covariance-tail},
\begin{align}
 p\sum_{q\le y}\sum_{p^k>y^2}
 |\operatorname{Cov}_{C,S}(I_{p^k},I_q)|
 &\ll_{B,W}py^{-2}\sum_{q\le y}\frac1q
       +O_{B,W}\!\left(\frac{pyL}{n}\right)=o_n(1),
 \label{eq:marked-power-tail-JI}\\
 p\sum_{q\le y}\sum_{q^l>y^2}
 |\operatorname{Cov}_{C,S}(I_p,I_{q^l})|
 &\ll_{B,W}y^{-1}
       +O_{B,W}\!\left(\frac{pyL}{n}\right)=o_n(1).
 \label{eq:marked-power-tail-IJ}
\end{align}
Thus every endpoint contribution has the advertised form
\(O(pyL^{O(1)}/n)=o(1)\).

Finally consider the same-prime diagonal.  Since
\[
 J_p^2=\sum_{k\ge2}(2k-3)I_{p^k},
\]
the marked bound through \(p^k\le y^4\), followed by the elementary
tail bound, gives
\begin{equation}
 {\mathbb E}_{C,S}J_p^2
 \ll_B\frac1{p^2}
 +O_{B,W}\!\left(\frac1{p^2L}+\frac{L}{y^4}+\frac{L^2}{n}\right).
 \label{eq:marked-power-diagonal-moment}
\end{equation}
Moreover \(I_pJ_p=J_p\), so
\begin{align}
 p\left|\operatorname{Var}_{C,S}(V_p)
             -\operatorname{Var}_{C,S}(I_p)\right|
 &\le p\{2|\operatorname{Cov}_{C,S}(I_p,J_p)|
                    +\operatorname{Var}_{C,S}(J_p)\}\notag\\
 &\ll_B\frac1p+o_{W,n}(1)+O_{B,W}(pL^2/n).
 \label{eq:marked-power-row-diagonal}
\end{align}
The preceding displays control every high-power and endpoint tail.  The
box-uniform leading terms are supplied by
\Cref{lem:marked-product-weighted-power-transfer}; together they give
\begin{equation}
 \left\|\operatorname{Cov}_{C,S}(V_p,V_q)
       -\operatorname{Cov}_{C,S}(I_p,I_q)\right\|_{\rm row}
 \le \frac{C_{\rm pow}}W+\epsilon_{B,W}(n).
 \label{eq:marked-prime-power-aggregate}
\end{equation}
To record explicitly the relative estimate used later, suppose that a
coefficient vector satisfies
\begin{equation}
 \|b\|_\infty+\sum_p\frac{|b_p|}{p}\le C_{\rm cmp}w,\qquad
                  \sum_p\frac{b_p^2}{p}\le C_{\rm cmp}w^2,
 \label{eq:marked-compensated-coefficient-norms}
\end{equation}
where \(C_{\rm cmp}\) is fixed independently of the compact tilt box and
before \(W\) is enlarged.  Put
\[
 A_1:=\sum_p\frac{t_p|b_p|}{p},\quad
 A_2:=\sum_p\frac{t_p|b_p|}{p^2},\quad
 B_1:=\sum_p\frac{|b_p|}{p},\quad
 B_2:=\sum_p\frac{|b_p|}{p^2}.
\]
Since \(0<t_p\le1\) and \(p>W\),
\begin{equation}
 A_1,B_1\le C_{\rm cmp}w,\qquad
 A_2,B_2\le \frac{C_{\rm cmp}}Ww.
 \label{eq:marked-compensated-product-norm-ledger}
\end{equation}
The product-weighted bounds
\eqref{eq:marked-product-JI}--\eqref{eq:marked-product-JJ} now give,
with their reciprocal tails already included in
\(\epsilon_{B,W}(n)\),
\begin{align}
 &\sum_{p\ne q}|b_pb_q|
 \left\{|\operatorname{Cov}(J_p,I_q)|
             +|\operatorname{Cov}(I_p,J_q)|\right\}\notag\\
 &\quad\le C_{\rm pow}(A_2A_1+A_1A_2)
       +\epsilon_{B,W}(n)(B_2B_1+B_1B_2)\notag\\
 &\quad\le \frac{C(C_{\rm cmp})}{W}w^2
       +C(C_{\rm cmp})\epsilon_{B,W}(n)w^2.
 \label{eq:marked-compensated-JI-IJ}
\end{align}
Similarly, \eqref{eq:marked-product-JJ} and
\eqref{eq:marked-product-diagonal} contribute
\[
 C_{\rm pow}A_2^2+\epsilon_{B,W}(n)B_2^2
 \quad\hbox{and}\quad
 \{C_{\rm pow}+\epsilon_{B,W}(n)\}
       \sum_p\frac{b_p^2}{p^2},
\]
respectively.  By
\eqref{eq:marked-compensated-coefficient-norms} and
\eqref{eq:marked-compensated-product-norm-ledger}, these are
\(O_{C_{\rm cmp}}(w^2/W)\) plus
\(C(C_{\rm cmp})\epsilon_{B,W}(n)w^2\).
Keeping the harmless numerical factors visible, and enlarging
\(C_{\rm pow}\) once, the entire non-endpoint contribution is bounded by
\begin{equation}
 \frac{3C_{\rm pow}(C_{\rm cmp}^2+C_{\rm cmp})}{W}w^2
 +3(C_{\rm cmp}^2+C_{\rm cmp})\epsilon_{B,W}(n)w^2.
 \label{eq:marked-compensated-nonendpoint-contraction}
\end{equation}
Indeed the off-diagonal products use \(A_2,B_2\le
C_{\rm cmp}w/W\), while
\(\sum_pb_p^2/p^2\le W^{-1}\sum_pb_p^2/p\); the last comparison uses
only \(W^{-1}\le1\) and the nonnegativity of
\(C_{\rm cmp}^2+C_{\rm cmp}\).  This is the scalar contraction recorded
in the finite Lean audit, rather than an implicit nonlinear use of the
\(O\)-notation.

It remains only to sum the literal endpoints, rather than fold them into
an unscaled \(o(1)\).  The coefficient bound gives
\[
 \sum_p|b_p|\le\pi(y)\|b\|_\infty\le C_{\rm cmp}wy.
\]
Consequently the \(JI/IJ\), \(JJ\), and diagonal endpoints are,
respectively,
\begin{align}
 \frac Ln\sum_{p\ne q}|b_pb_q|
 &\ll C_{\rm cmp}^2w^2\frac{y^2L}{n},\notag\\
 \frac{L^2}{n}\sum_{p\ne q}|b_pb_q|
 &\ll C_{\rm cmp}^2w^2\frac{y^2L^2}{n},\notag\\
 \frac{L^2}{n}\sum_pb_p^2
 &\ll C_{\rm cmp}^2w^2\frac{yL^2}{n}.
 \label{eq:marked-relative-endpoint-ledger}
\end{align}
All three are \(o_n(w^2)\).  Define, for this fixed comparison constant,
\begin{equation}
 \epsilon^{\rm rel}_{B,W,C_{\rm cmp}}(n)
 :=C(C_{\rm cmp})\left\{\epsilon_{B,W}(n)
             +\frac{y^2L^2}{n}\right\}.
 \label{eq:marked-relative-remainder-definition}
\end{equation}
After increasing the fixed factor \(C(C_{\rm cmp})\), this also contains
the smaller \(y^2L/n\) and \(yL^2/n\) endpoint terms.  Hence
\(\epsilon^{\rm rel}_{B,W,C_{\rm cmp}}(n)\to0\) for every fixed
\(B,W,C_{\rm cmp}\).  The complete reciprocal tails and all coefficient
factors are now recorded in this single relative remainder.
Consequently
\begin{equation}
 \left|\operatorname{Var}_{C,S}\!\left(\sum_pb_pV_p\right)
       -\operatorname{Var}_{C,S}\!\left(\sum_pb_pI_p\right)\right|
 \le \frac{C_{\rm rel}(C_{\rm cmp})}{W}w^2
       +\epsilon^{\rm rel}_{B,W,C_{\rm cmp}}(n)w^2,
 \label{eq:marked-compensated-power-relative}
\end{equation}
where \(C_{\rm rel}(C_{\rm cmp})\) is independent of the tilt-box
radius.  Formula
\eqref{eq:marked-relative-remainder-definition} records explicitly the
fixed coefficient loss and all endpoint and local-fugacity remainders;
the threshold for its convergence may depend on \(B,W,C_{\rm cmp}\).
Thus no unrecorded transpose or diagonal term is being absorbed in the
relative estimate used in \Cref{sec:smooth-bridge}.

The Stieltjes form of \eqref{eq:marked-omitted-local-score}, followed by
the same local expansion, gives for \(p^k\le y^2\)
\[
 |\operatorname{Cov}_{C,S}(\1_{p^k\mid m},\phi(m/n))|
 \ll_{B,W}\frac{\|\phi\|_{BV}}{p^kL}.
\]
For the omitted exponents, comparison of the tilted and untilted
densities and the elementary count of multiples give, uniformly for
\(W<p\le y\),
\[
 \sum_{p^k>y^2}\mathbb P_{C,S}(p^k\mid m)
 \ll_{B,W}y^{-2}+\frac Ln
 =o\!\left(\frac1{pL}\right).
\]
Indeed, after division by the cell size, the geometric multiple-counting
term is \(O(y^{-2})\), while the \(+1\) endpoint is summed over at most
\(O(L)\) exponents.  Relative to \(1/(pL)\), the two displayed terms
have ratios at most \(L/y\) and \(yL^2/n\), respectively, and both tend
to zero.  Since the covariance of a divisibility indicator with the
bounded physical test is at most twice its probability times
\(\|\phi\|_\infty\), summation over every exponent now proves
\begin{equation}
 |\operatorname{Cov}_{C,S}(v_p(m),\phi(m/n))|
 \ll_{B,W}\frac{\|\phi\|_{BV}}{pL}.
 \label{eq:marked-tilted-physical-full}
\end{equation}

A further physical factor \(e^{\zeta\phi(m/n)/L}=1+O_{B}(1/L)\) is
handled by the same expansion.  Reweighting the finitely many exact head
cells in a compact positive box also preserves every estimate, because
their normalized main marked profile in
\eqref{eq:marked-cell-probability} is common.  We have therefore proved
uniform one- and two-prime kernels, the physical Stieltjes column, and the
aggregate \(O(1/W)\) prime-power transfer throughout the complete compact
band, physical, and finite-head homotopy used in the next section.  At no
point was an arbitrary residue class or a \(U=5/2\) terminal estimate
invoked.

For later reference we isolate the preceding stability statement,
including the effect of the guarded deletions.

\begin{corollary}[Compact full tilt and guarded cells]
\label[corollary]{cor:marked-full-tilt-guard-stability}
Let \(\nu=\sum_{i=1}^m\lambda_i\nu_i\) be a fixed finite cell mixture
as in \eqref{eq:marked-finite-mixture-definition}, with its nonzero
weights in a fixed compact subset of the relative interior.  Reweight it
by
\begin{equation}
 \frac{d\nu_\xi}{d\nu}(m)=\frac1{Z_\xi}
 \exp\!\left\{\frac1L\left(S(m)+\zeta\phi(m/n)
                  +\sum_{\ell=1}^e\kappa_\ell H_\ell^\circ(m)
                         \right)\right\},
 \label{eq:marked-full-tilt-density}
\end{equation}
where the effective prime fugacities, \(\zeta\), and the finitely many
\(\kappa_\ell\) lie in a fixed box, \(\phi\) is bounded with bounded
variation, and the centered head indicators are bounded.  Then
\begin{enumerate}[label=\textup{(\roman*)}]
\item there are constants \(0<c_{B,W}<C_{B,W}<\infty\) such that
\(c_{B,W}\leq d\nu_\xi/d\nu\leq C_{B,W}\) pointwise;
\item all squarefree, prime-power, physical, and coefficient-weighted
estimates in this section hold uniformly under \(\nu_\xi\), with the
same box-independent coefficients of \(1/W\);
\item if a set \(\mathcal R_n^{\rm guard}\) of at most
\(O(yL^A)\) numerical guards
is deleted before normalization, the change in every normalized
one- or two-mark prime-row operator, and hence in its harmonically
averaged band operator, is
\begin{equation}
 O_{B,W}\!\left(\frac{y^2L^{A+3}}n\right)=o_n(\alpha_0),
 \qquad \alpha_0\asymp\frac1{\log L}.
 \label{eq:marked-guard-relative-low-error}
\end{equation}
In particular the deletion error remains relative at the moving low
cell, not merely \(o_n(1)\).
\end{enumerate}
\end{corollary}

\begin{proof}
The elementary inequality
\[
 \Omega(m)\log W\leq\log(C_*n)
\]
gives \(|S(m)|/L\leq B/\log W+O_{B,W}(1/L)\).  The physical and head scores
in \eqref{eq:marked-full-tilt-density} are \(O_{B,W}(1/L)\).
Normalization therefore proves (i).  Taylor expansion of the latter
two factors has a uniform \(O_{B,W}(1/L)\) remainder.  The Stieltjes
formula makes their leading physical average cancel in centered marked
moments, while the common marked profile of the head cells makes the
between-cell term cancel.  Combining this with
 \Cref{lem:marked-tilted-squarefree} and
 \Cref{lem:marked-product-weighted-power-transfer}, together with
 \eqref{eq:marked-compensated-power-relative} proves (ii).
We shall also use the following component-level consequence of the same
calculation.  Let \(\nu_{i,\xi}\) be the conditional law on the \(i\)-th
exact component cell before the numerical guards are deleted.  Then,
uniformly over the finitely many nonzero components and the fixed box,
\begin{equation}
 \sup_{i,i'}\left|
  \mathbb E_{\nu_{i,\xi}}V_p-\mathbb E_{\nu_{i',\xi}}V_p
 \right|\ll_{B,W}\frac1{pL}
 \qquad(W<p\leq y).
 \label{eq:marked-full-tilt-component-valuation-comparison}
\end{equation}
Indeed \eqref{eq:marked-full-valuation-cell-comparison} proves this before
tilting.  For \(p^k\leq y^2\), the one-local specialization of the
omitted-local-score expansion and exact fugacity restoration used in
\Cref{lem:marked-product-weighted-power-transfer} writes every component
mean as the same restored local model, with error
\[
 O_{B,W}\!\left(\frac1{p^kL}+\frac1n\right).
\]
Summing \(k\) gives \(O_{B,W}(1/(pL)+L/n)=O_{B,W}(1/(pL))\), since
\(pL^2/n\leq yL^2/n=o(1)\).  The omitted powers are
\(o(1/(pL))\) by the tail calculation immediately preceding
\eqref{eq:marked-tilted-physical-full}.  A physical tilt is
\(1+O_B(1/L)\), and a head tilt is constant on each exact component, so
neither changes this comparison.  Thus
\eqref{eq:marked-full-tilt-component-valuation-comparison} is a displayed
consequence of (ii), rather than an implicit appeal to a common profile.

We first make the component-weight bookkeeping literal.  If the later
construction prescribes weights \(\lambda_i\) for the individually
guarded component laws, put
\[
 s_i:=\nu_i((\mathcal R_n^{\rm guard})^c),\qquad
 \widetilde\lambda_i:=
 \frac{\lambda_i/s_i}{\sum_j\lambda_j/s_j}.
\]
Every component has \(\Theta_W(n)\) candidates, so
\[
 1-s_i\ll_W\frac{yL^A}{n}=o(1)
\]
uniformly in the fixed finite component family.  Hence all \(s_i\) are
positive and the weights \(\widetilde\lambda_i\) remain in a fixed compact
subset of the relative interior whenever the prescribed
\(\lambda_i\)'s do.  The exact measure identity
\begin{equation}
 \left(\sum_i\widetilde\lambda_i\nu_i\right)
       \left(\,\cdot\mid(\mathcal R_n^{\rm guard})^c\right)
 =\sum_i\lambda_i
       \nu_i\left(\,\cdot\mid(\mathcal R_n^{\rm guard})^c\right)
 \label{eq:marked-guard-component-weight-identity}
\end{equation}
follows because the conditioned weight of component \(i\) on the left is
\(\widetilde\lambda_is_i/\sum_j\widetilde\lambda_js_j=\lambda_i\).
The identity remains exact after applying and normalizing the same
positive exponential tilt on both sides.  We may therefore apply the
pre-deletion estimates to the compact mixture with weights
\(\widetilde\lambda_i\), and recover precisely the prescribed guarded
mixture after conditioning.  This also covers the simpler case in which
the guards are deleted directly from a globally weighted mixture.

Let \(\nu_\xi\) denote the resulting pre-deletion tilted law,
\(\widehat\nu_\xi\) the actual guarded tilted law, and put
\[
\delta_\xi:=\nu_\xi(\mathcal R_n^{\rm guard}).
\]
Every component cell has \(\Theta_W(n)\) candidates, and (i) makes the
probability of a single candidate \(O_{B,W}(1/n)\).  Hence, uniformly on
the fixed box,
\begin{equation}
 \delta_\xi\ll_{B,W}\frac{yL^A}{n}=o(1).
 \label{eq:marked-guard-deleted-mass}
\end{equation}
We use the following elementary conditional-deletion identities.  If
\(\mu\) is a probability law, \(G\) has \(\mu(G)=\delta<1\), and
\(\widehat\mu=\mu(\,\cdot\mid G^c)\), then, for every bounded \(F\),
\begin{equation}
 \mathbb E_{\widehat\mu}F-\mathbb E_\mu F
 =\frac{\delta\mathbb E_\mu F-\mathbb E_\mu(F\1_G)}{1-\delta},
 \qquad
 \left|\mathbb E_{\widehat\mu}F-\mathbb E_\mu F\right|
 \leq\frac{2\delta}{1-\delta}\|F\|_\infty.
 \label{eq:marked-conditional-deletion-expectation}
\end{equation}
Writing the covariance difference as
\begin{align*}
 \operatorname{Cov}_{\widehat\mu}(F,H)-\operatorname{Cov}_\mu(F,H)
 &=\{\mathbb E_{\widehat\mu}(FH)-\mathbb E_\mu(FH)\}\notag\\
 &\quad-\{\mathbb E_{\widehat\mu}F-\mathbb E_\mu F\}
       \mathbb E_{\widehat\mu}H\notag\\
 &\quad-\mathbb E_\mu F
       \{\mathbb E_{\widehat\mu}H-\mathbb E_\mu H\}
\end{align*}
and applying \eqref{eq:marked-conditional-deletion-expectation} three
times gives the explicit recentered bound
\begin{equation}
 \left|\operatorname{Cov}_{\widehat\mu}(F,H)
       -\operatorname{Cov}_\mu(F,H)\right|
 \leq\frac{6\delta}{1-\delta}\|F\|_\infty\|H\|_\infty.
 \label{eq:marked-conditional-deletion-covariance}
\end{equation}
Thus both normalization and the two changes of means are included.

Apply this with \(\mu=\nu_\xi\), \(G=\mathcal R_n^{\rm guard}\), and
\(V_p=v_p(m)\).  On every candidate, for all sufficiently large \(n\),
\begin{equation}
 pV_p(m)\ll yL,
 \qquad
 \sum_{W<r\leq y}V_r(m)\leq\frac{\log m}{\log W}\ll_W L.
 \label{eq:marked-guard-score-envelopes}
\end{equation}
Let
\[
 E^{\rm guard}_{pr}:=
 \operatorname{Cov}_{\widehat\nu_\xi}(V_p,V_r)
 -\operatorname{Cov}_{\nu_\xi}(V_p,V_r).
\]
For fixed \(p\), duality with
\(H_b=\sum_r b_rV_r\), \(\|b\|_\infty\leq1\), followed by
\eqref{eq:marked-conditional-deletion-covariance}--
\eqref{eq:marked-guard-score-envelopes}, gives
\begin{align}
 p\sum_{W<r\leq y}|E^{\rm guard}_{pr}|
 &=\sup_{\|b\|_\infty\leq1}
      p\left|\operatorname{Cov}_{\widehat\nu_\xi}(V_p,H_b)
             -\operatorname{Cov}_{\nu_\xi}(V_p,H_b)\right|\notag\\
 &\ll_{B,W}\delta_\xi yL^2
 \ll_{B,W}\frac{y^2L^{A+2}}n
 \leq C_{B,W}\frac{y^2L^{A+3}}n.
 \label{eq:marked-guard-explicit-prime-row}
\end{align}
The expectation bound in
\eqref{eq:marked-conditional-deletion-expectation} proves the analogous
one-mark estimate; replacing either \(V\) by \(I_p=\1_{p\mid m}\), or
the second score by a bounded physical or finite-head score, only makes
the envelope smaller.

It remains to identify the norm after band aggregation.  For
\(\|b\|_\infty\leq1\), the covariance-error contribution on band \(i\)
is bounded exactly by
\begin{align*}
 \frac1{H_i}\left|\sum_{p\in\mathcal P_i}\sum_r
       E^{\rm guard}_{pr}b_{j(r)}\right|
 &\leq\frac1{H_i}\sum_{p\in\mathcal P_i}\frac1p
       \left(p\sum_r|E^{\rm guard}_{pr}|\right)\\
 &\leq\sup_u u\sum_r|E^{\rm guard}_{ur}|.
\end{align*}
Denote this raw normalized band covariance-error operator by
\(E^{\rm guard}_{\rm band}\).  Thus harmonic averaging is a contraction
of the normalized prime-row norm \eqref{eq:marked-row-norm}; the uniformly
bounded arithmetic gauge projections do not alter its order.  For
completeness, the sharp conjugation used later also follows from the
displayed relative rate.  Put
\(S_n=\operatorname{diag}(\alpha_0,\ldots,\alpha_s)\).  If
\(\|q\|_\infty\leq1\), apply the raw row estimate to
\(b_j=\alpha_jq_j\), for which \(\|b\|_\infty\leq1\).  Dividing output row
\(i\) by \(\alpha_i\) gives
\[
 \left|\bigl(S_n^{-1}E^{\rm guard}_{\rm band}S_nq\bigr)_i\right|
 \ll_{B,W}\frac{y^2L^{A+3}}{n\alpha_i}
 \leq C_{B,W}\frac{y^2L^{A+3}}{n\alpha_0}=o(1).
\]
Indeed, once the cells are nonempty,
\(\alpha_0<\delta<\alpha_j\) for every \(j\geq1\): the low cell is
contained in \([t_0,\delta]\), every positive cell lies in
\((\delta,1]\), and nonemptiness follows for the fixed finite permitted
mesh from the prime number theorem.  Thus \(\alpha_0\) is the smallest
cell center.  The arithmetic quantities \(H_i,\alpha_i\) are prime
geometry and hence are unchanged by deleting candidates.  The gauge projections are uniformly
bounded in this sharp norm by
\eqref{eq:smooth-sharp-projection-convergence}.  This proves
\eqref{eq:marked-guard-relative-low-error}.  Finally
\(y^2L^{A+3}/n=n^{-5/9}L^{A+3}=o(1/\log L)=o(\alpha_0)\).
\end{proof}
 \section{The smooth-row bridge and exact finite-band fitting}
\label{sec:smooth-bridge}

In this section we complete the fractional construction left by
\Cref{sec:rough-selector}.  We use throughout
\[
 L=\log n,\qquad N=\frac nL,\qquad
 \theta=\frac29,\qquad y=n^\theta,\qquad U=\theta^{-1}=\frac92.
\]
The real cutoff \(W\), the head set
\(F_{\rm hd}=\{p:p\leq W\}\), and the head modulus \(P_{\rm hd}\) have
already been fixed.  All constants in this section may depend on \(W\),
but not on \(n\) or on a prime in \((W,y]\).

\subsection{The integer row ledger and the active physical target}
\label{subsec:smooth-ledger}

Recall that \(Y=T/D'\) is the target remaining after the anchor divisor
has been removed, and that \(B_0\) is the combined product of the fixed
lower factors and the base states of the precharged switch banks.  The
rough-stage residual is
\[
 r_p=v_p(Y)-v_p(B_0)-\sum_{a\in\mathcal A}x_av_p(a)
\]
as in \eqref{eq:rough-residual-convention}.  In the trivial complete
rough-signature row the fractional selector has the additive
decomposition
\begin{equation}
 x^{\rm sm}=x^{\rm top}+f^{\rm prot}+z^{\rm act}.
 \label{eq:smooth-additive-decomposition}
\end{equation}
Here \(x^{\rm top}\) is the constant-density component on the short top
interval \(H=(2n-Kh,2n]\), \(f^{\rm prot}\) is the broad layer of weight
\(\beta_{\rm prot}/L\) reserved for the tangent construction, and
\(z^{\rm act}\) is the adjustable broad excess.  The support of
\(z^{\rm act}\) is contained in
\[
                 J=(n,2n-Kh],
\]
and is therefore disjoint from the support of \(x^{\rm top}\).
The protected summand may share numerical coordinates with
\(z^{\rm act}\); in that case the final inclusion weight is their sum.
Only \(z^{\rm act}\) is differentiated below.  In particular, neither
\(x^{\rm top}\) nor \(f^{\rm prot}\) is ever folded into the bridge
probability measure.

We first record the integer-versus-real distinction which will also be
needed in the final rounding.  Freeze \(B_0\), all nonsmooth fractional
rows, the whole top component, the protected floor, and every other
smooth seed not used in the interpolation below.  Let
\(m_{\rm sm,fix}\) be the number of fixed factors of \(B_0\) in the
smooth signature row, and let \(m_{\rm sm,fr}\) be the total mass of all
frozen contributions in that row.  Put
\[
 m_{\rm sm,oth}:=m_{\rm sm,fr}-m_{\rm sm,fix}
       -\sum_a x_a^{\rm top}-\sum_a f_a^{\rm prot}.
\]
Let \(\widetilde q_0=\sum_m z_m^{\rm act}\) be the actual active mass
left by the guarded rough construction, before any interpolation in this
section.  Thus the actual post-guard mass in the smooth row is
\(m_{\rm sm,fr}+\widetilde q_0\).  Define, without an asymptotic
replacement,
\begin{equation}
 Q_{\rm sm}(0):=\operatorname*{argmin}_{k\in\mathbb Z}
       |k-(m_{\rm sm,fr}+\widetilde q_0)|,
 \qquad q_0:=Q_{\rm sm}(0)-m_{\rm sm,fr};
 \label{eq:smooth-nearest-actual-quota}
\end{equation}
break a possible tie by taking the smaller integer.  In particular,
\(|q_0-\widetilde q_0|\leq1/2\).  Choose the integer total smooth-row
quota \(Q_{\rm sm}(d)=Q_{\rm sm}(0)-d\), and define
\begin{align}
 q_{\rm sm}^{\rm flex}(d)
   &:=Q_{\rm sm}(d)-m_{\rm sm,fix}\in\mathbb Z,
       \label{eq:smooth-flexible-integer-quota}\\
 q^{\rm act}(d)
   &:=Q_{\rm sm}(d)-m_{\rm sm,fr}=q_0-d.
       \label{eq:smooth-active-real-mass}
\end{align}
Thus \(q^{\rm act}(d)\) is a real mass; it need not be an integer.
The exact ledger identity is
\begin{equation}
\boxed{\;
 Q_{\rm sm}(d)=m_{\rm sm,fix}
 +\sum_a x_a^{\rm top}+\sum_a f_a^{\rm prot}
 +q^{\rm act}(d)+m_{\rm sm,oth},
 \qquad q_{\rm sm}^{\rm flex}(d)\in\mathbb Z .
\;}
\label{eq:smooth-quota-ledger}
\end{equation}
Equivalently, the total weight of the flexible numerical coordinates in
the smooth row is \(q_{\rm sm}^{\rm flex}(d)\).  There is no separate
integrality requirement on any of the four real summands on the
right-hand side of \eqref{eq:smooth-quota-ledger}.

We now define the two frozen ledgers without identifying fractional mass
with a cardinality.  Let \(\mathscr F_0\) be the set of frozen integral
factors (so that its contribution includes \(B_0\)), and let
\(\mathscr X_0\) be the set of all frozen fractional coordinates, with
their actual weights denoted by \(x_a^0\).  The active coordinates
\(z^{\rm act}\) are excluded from \(\mathscr X_0\).  Put, exactly,
\begin{equation}
 m_0:=\#\mathscr F_0+\sum_{a\in\mathscr X_0}x_a^0,
 \qquad
 \Lambda_0:=\sum_{a\in\mathscr F_0}\log a
       +\sum_{a\in\mathscr X_0}x_a^0\log a.
 \label{eq:smooth-frozen-mass-log-definition}
\end{equation}
Thus \(m_0\) is the exact \emph{total frozen mass}, generally not an
integer, whereas \(\Lambda_0\) is its ordinary logarithmic contribution.
In particular \(\log B_0\) is included in \(\Lambda_0\), while
\(\log D'\) is not, since \(D'\) has already been divided from the
target.  The source of the mass estimate is explicit.  The charged fixed
tokens together with the flexible mass in every nontrivial row total
exactly \(t_R\).  Summing those identities and adjoining the actual
post-guard smooth raw mass gives
\[
 m_0+\widetilde q_0
 =\sum_{R\ne1}t_R+
   \{\hbox{actual post-guard smooth raw mass}\}
 =h+O(N/L).
\]
The nearest-integer initialization changes this by only \(O(1)\).
Together with the corresponding height estimates from the rough stage,
this gives
\begin{equation}
 m_0+q_0=h+O(N/L),\qquad
 \Lambda_0=m_0L+O(N),\qquad
 \log Y=hL+O(N).
\label{eq:smooth-baseline-height-ledger}
\end{equation}
The first relation uses the same exact convention for every retained
factor, bank base state, guard, and fractional row.  Consequently
\[
 A_0:=\log Y-\Lambda_0-q_0L=O(N).
\]
For \(q^{\rm act}(d)=q_0-d\), the height above \(n\) which the active
smooth component must supply is
\begin{equation}
 A(d):=\log Y-\Lambda_0-q^{\rm act}(d)L=A_0+dL.
\label{eq:smooth-active-height}
\end{equation}
Fix once and for all a number \(\mu\in(0,\log2)\), and let \(d\) be a
nearest integer to
\begin{equation}
 d_*:=\frac{\mu q_0-A_0}{L+\mu}.
\label{eq:smooth-height-center}
\end{equation}
The bounds above give \(d=O(N/L)\).  Since \(q_0\asymp N\), they also
give
\begin{equation}
 q^{\rm act}(d)\asymp N,\qquad
 \frac{A(d)}{q^{\rm act}(d)}
        =\mu+O(L/N).
\label{eq:smooth-physical-interiority}
\end{equation}
This centers the mean of the \emph{remaining active smooth mass}; it is
not a statement about the mean of all lower factors.  In particular,
retained exceptional factors and bank states have already been charged
in \(\Lambda_0\).

Here is the promised literal feasibility of the two mass changes.  Before
freezing the head cells, keep two disjoint positive-length subintervals of
the zero-head stratum
\[
 \{m\in J:P^+(m)\leq y,\ (m,P_{\rm hd})=1\},
\]
one below and one above \(e^\mu n\).  The rough active density is a fixed
positive multiple of \(1/L\) there, and the unused ceiling is bounded
below by a fixed constant.  Each subpool consequently has both removal
capacity and addition capacity \(\gg_W N\).  First change their combined
mass by \(q_0-\widetilde q_0=O(1)\), and then by \(-d=O(N/L)\), splitting
each change between the two pools.  Spreading these changes uniformly
alters an individual active coordinate by \(O_W(L^{-2})\).  All head
valuations remain zero, no
coordinate leaves \([0,1]\), and the two physical splitting parameters
remain a fixed distance from their endpoints.  Thus the mass after this
operation is exactly \(q^{\rm act}(d)=q_0-d\).  The later two-pool
interpolation changes the split, but not this mass.  This proves, rather
than assumes, that every integer \(d\) used in
\eqref{eq:smooth-height-center} is attainable.

\subsection{Exact head cells and physical interpolation}
\label{subsec:smooth-head-physical}

After the choice of \(d\), subtract from the target each frozen head
valuation and let \(A_p^{\rm act}\) be the valuation still assigned to
the active smooth component at \(p\leq W\).  This quantity has a useful
coordinatewise verification.  Put
\[
 \gamma_{p,W}:=
 \begin{cases}
 (c-C_0)/(3(p-1)),&p\in F_{\rm anc},\\
 c/(2(p-1)),&p\in F_{\rm hd}\setminus F_{\rm anc}.
 \end{cases}
\]
The nonsmooth flexible point, the protected layer, and the initial active
smooth point are all \(F_{\rm hd}\)-free.  Hence the only head charge
besides \(D'\) is the already combined base \(B_0\).  From
\eqref{eq:central-anchor-lemma-reserve}, its preserved form
\eqref{eq:bank-anchor-change-cost},
\eqref{eq:uniform-exceptional-charge}, and
\eqref{eq:base-guard-small-prime-profile}, we have simultaneously
\begin{equation}
 \begin{split}
 A_p^{\rm act}
 &=v_p(T)-v_p(D')-v_p(B_0)\\
 &\geq\left\{\gamma_{p,W}
       -C_W\frac{\delta_*}{\theta p}\right\}N
       -o_W(N)\geq 2\epsilon_WN
       \qquad(p\leq W),
 \end{split}
 \label{eq:smooth-head-reserve-minus-charge}
\end{equation}
after the prescribed choice of \(\delta_*\) and then sufficiently large
\(n\).  Also \(A_p^{\rm act}\leq v_p(T)\leq C_WN\).  Changing \(d\)
uses only the zero-head pool just displayed, so it changes every one of
these targets by zero (and, even without that choice, an
\(O_W(N/L)\) change would be smaller than the fixed linear margin).
Consequently
\begin{equation}
 \epsilon_WN\leq A_p^{\rm act}\leq C_WN
       \qquad(p\leq W).
\label{eq:smooth-head-interior}
\end{equation}
Choose a fixed integer \(E_W\) so large that
\begin{equation}
 \frac{(A_p^{\rm act})_{p\leq W}}{q^{\rm act}(d)}
 \in\operatorname{int}\operatorname{conv}
       \{0,E_W\mathbf e_p:p\leq W\}
\label{eq:smooth-head-simplex}
\end{equation}
with a fixed positive barycentric margin.  The margin follows directly
from \eqref{eq:smooth-head-reserve-minus-charge}, not from a limiting
interiority assertion.  Explicitly, the coefficient
of \(E_W\mathbf e_p\) is
\(A_p^{\rm act}/(E_Wq^{\rm act}(d))\), and the coefficient of \(0\) is
one minus their sum.

The vertices in \eqref{eq:smooth-head-simplex} are realized by exact
valuation/coprimality patterns: the zero vertex is
\(v_\ell(m)=0\) for every \(\ell\leq W\), and the \(p\)-vertex is
\[
 v_p(m)=E_W,\qquad v_\ell(m)=0
       \quad(\ell\leq W,\ \ell\ne p).
\]
These are finite divisibility inclusion--exclusion cells, not residue
classes.  By Proposition~\ref{prop:marked-structured-cells}, each such pattern has
\(\Theta_W(n)\) \(y\)-smooth members in every fixed positive-length
subinterval of \((n,2n]\).

Choose \(\eta_0>0\) so that
\[
 0<\mu-\eta_0<\mu+\eta_0<\log2
\]
and then fix constants
\begin{equation}
 1<a_-<b_-<e^{\mu-\eta_0},\qquad
 e^{\mu+\eta_0}<a_+<b_+<2.
 \label{eq:smooth-scaled-physical-constants}
\end{equation}
For every \(n\), define the scaled intervals
\begin{equation}
 I_-(n):=(a_-n,b_-n],\qquad I_+(n):=(a_+n,b_+n].
 \label{eq:smooth-scaled-physical-intervals}
\end{equation}
Since \((2-b_+)n>Kh\) for all sufficiently large \(n\), one has
\(I_+(n)\subset(n,2n-Kh]\); the corresponding containment for
\(I_-(n)\) is immediate.  In each head pattern reserve positive-density
smooth subpools in both \(I_-(n)\) and \(I_+(n)\).  Assign the barycentric masses from
\eqref{eq:smooth-head-simplex}, subtracting exactly the mass assigned
from the active zero-pattern stratum.  Thus no new \(\Theta(N)\) mass is
added.  Split each pattern mass between its two physical subpools.  The
two logarithmic means lie on opposite sides of \(\mu\); hence a compact
choice of the splitting coefficients makes the total active logarithm
equal to \eqref{eq:smooth-active-height} exactly.  The \(O(L)\) error
caused by rounding \(d_*\) is removed by changing one splitting
coefficient by \(O(L/N)=o(1)\).

Every assigned cell has \(\Theta_W(n)\) candidates and receives
\(O_W(N)\) mass.  Its individual weights are therefore \(O_W(1/L)\).
The barycentric and physical splitting margins leave a positive-density
zero-pattern active component with weights bounded below by a fixed
multiple of \(1/L\).  This component, rather than the frozen protected
floor, supplies all covariance lower bounds below.  The Stieltjes estimate
\eqref{eq:marked-tilted-physical-full} also shows that the physical and
head redistribution changes a marked valuation at \(p>W\) by at most
\begin{equation}
                         O_W\!\left(\frac{N}{pL}\right).
\label{eq:smooth-head-physical-marked-cost}
\end{equation}

We summarize this construction as a measure-realization lemma, so the
probability space used by the bridge is not left implicit.

\begin{lemma}[Baseline active-measure realization]
\label[lemma]{lem:smooth-baseline-active-measure}
Let \(\mathscr E=\{0\}\cup\{E_W\mathbf e_p:p\leq W\}\) be the finite
head-pattern simplex above, and let \(\sigma\in\{-,+\}\) indicate the
physical interval \(I_\sigma(n)\).  Put \(q_n:=q^{\rm act}(d)\).
After deleting all numerical guards,
there are disjoint structured cells \(C_{e,\sigma,n}\) and masses
\(q_{e,\sigma,n}>0\) such that
\begin{equation}
 \nu_n=\sum_{e\in\mathscr E}\sum_{\sigma\in\{-,+\}}
       \lambda_{e,\sigma}(n)\nu_{e,\sigma,n},
 \qquad
 \lambda_{e,\sigma}(n):=\frac{q_{e,\sigma,n}}{q_n},
 \label{eq:smooth-baseline-finite-mixture}
\end{equation}
is a probability measure, where \(\nu_{e,\sigma,n}\) is uniform on
\(C_{e,\sigma,n}\).  More precisely:
\begin{enumerate}[label=\textup{(\roman*)}]
\item the vector \((\lambda_{e,\sigma}(n))\) stays in a fixed compact
subset of the open simplex;
\item \(\#C_{e,\sigma,n}=\Theta_W(n)\), also after guard deletion;
\item the coordinate realization
\begin{equation}
 z_m^0=\frac{q_{e,\sigma,n}}{\#C_{e,\sigma,n}}
       \quad(m\in C_{e,\sigma,n})
 \label{eq:smooth-baseline-coordinate-realization}
\end{equation}
satisfies \(z_m^0=O_W(1/L)\), and on a fixed positive fraction of the
zero-head cells one has \(z_m^0\geq c_W/L\);
\item after addition of the frozen protected weight at a shared clean
coordinate, the total weight is still \(O_W(1/L)<1\) for large \(n\);
and
\item the total mass is exactly \(q_n\), all head moments are exactly
\((A_p^{\rm act})_{p\leq W}\), and the ordinary logarithmic moment is
exactly \(A(d)\).
\end{enumerate}
\end{lemma}

\begin{proof}
For a fixed head pattern and sign, let \(C_{e,\sigma,n}\) be the smooth
members of the corresponding scaled interval with that exact valuation
and coprimality pattern, after disjointification and guard deletion.
The marked cell estimate gives \(\Theta_W(n)\) members before deletion.
By \Cref{lem:rough-guard-census}, the whole excluded set has size
\(O_c(N)+O(yL)=o(n)\) (and its intersection with the active smooth row
is smaller).  Hence the same \(\Theta_W(n)\) estimate remains after
deletion.  The strict barycentric margin in
\eqref{eq:smooth-head-simplex} and the strict two-sided physical split
put every normalized assigned mass in a fixed compact subset of the open
simplex.  This proves (i), (ii), and
\eqref{eq:smooth-baseline-finite-mixture}.

Since \(q_n\asymp n/L\), formula
\eqref{eq:smooth-baseline-coordinate-realization} gives the upper bound
in (iii); the compact lower mass and the upper candidate-count bound give
the asserted zero-head lower bound.  These cells lie below
\(2n-Kh\), hence are disjoint from the top component, and all other
unavailable coordinates were deleted as guards.  The only allowed
overlap is with the protected layer, whose weight is itself \(O_W(1/L)\).
This proves (iv).  Finally the barycentric equations are exactly the
head equations, the two physical split parameters were chosen to solve
the ordinary logarithm, and their weights sum to \(q_n\).  These are
finite linear identities, proving (v).
\end{proof}

\subsection{The Poisson--Dickman covariance operator}
\label{subsec:smooth-PD-operator}

We next prove the analytic inverse used for the band fit.  Let \(\Pi\)
be the scale-invariant Poisson process on \((0,1]\), conditioned on
\(\sum_{t\in\Pi}t=U=9/2\).  Set
\[
 h(s)=\frac{\rho(U-s)}{\rho(U)},\qquad
 K(s,t)=\frac{\rho(U-s-t)}{\rho(U)}-h(s)h(t),
\]
with \(\rho(v)=0\) for \(v<0\).  The Palm identities
\eqref{eq:Palm-one-point}--\eqref{eq:Palm-two-point} imply that the
covariance of the additive statistic \(\sum_{t\in\Pi}f(t)\) is the
quadratic form of
\begin{equation}
 (Af)(s)=h(s)f(s)+\int_0^1K(s,t)f(t)\,\frac{dt}{t}.
\label{eq:smooth-PD-operator}
\end{equation}
Put \(F(x)=\rho(U-x)/\rho(U)\).  The Dickman delay equation shows that
\(F\in C^2[0,2]\): indeed \(U-x\in[5/2,9/2]\), a compact interval on
which the two differentiated delay equations are continuous.  We have
\(h(s)=F(s)\) and
\[
 K(s,t)=F(s+t)-F(s)F(t).
\]
Since \(F(0)=1\), a mixed first difference followed by Taylor's theorem
gives
\begin{equation}
                         |K(s,t)|\ll st
       \qquad(0\leq s,t\leq1).
\label{eq:smooth-kernel-product}
\end{equation}

\begin{lemma}[Poisson--Dickman quotient inverse]
\label[lemma]{lem:PD-quotient-inverse}
On \(\mathcal H=L^2((0,1],dt/t)\), the operator \(A\) is nonnegative,
Fredholm, and
\[
                         \ker A=\operatorname{span}\{t\}.
\]
Consequently there is a constant \(\kappa_0>0\) such that
\begin{equation}
 \langle f,Af\rangle_{\mathcal H}
 \geq\kappa_0\inf_{\lambda\in\mathbb R}
       \int_0^1|f(t)-\lambda t|^2\,\frac{dt}{t}.
\label{eq:PD-Poincare}
\end{equation}
On bounded functions one also has
\begin{equation}
 \inf_{\lambda\in\mathbb R}\|f-\lambda t\|_\infty
       \ll\|Af\|_\infty.
\label{eq:PD-ordinary-Linf-inverse}
\end{equation}
Moreover, on
\[
 X_t:=\{f:f(t)/t\in L^\infty(0,1]\},
 \qquad \|f\|_{X_t}:=\|f(t)/t\|_\infty,
\]
one has the weighted quotient estimate
\begin{equation}
 \inf_{\lambda\in\mathbb R}\|f-\lambda t\|_{X_t}
 \ll\left\|\frac{Af}{t}\right\|_\infty.
\label{eq:PD-weighted-inverse}
\end{equation}
\end{lemma}

\begin{proof}
Let
\[
 \mathcal D:=\{f\in L^\infty(0,1]:
   f=0\ \hbox{a.e. on }(0,\varepsilon)
   \hbox{ for some }\varepsilon>0\}.
\]
This is dense in \(\mathcal H=L^2((0,1],dt/t)\).  For
\(f\in\mathcal D\), the additive sum
\(S_f=\sum_{x\in\Pi}f(x)\) has only finitely many nonzero terms.
The one- and two-point Palm formulas therefore give, first for
\(f,g\in\mathcal D\),
\begin{align}
 \operatorname{Cov}_U(S_f,S_g)
 &=\int_0^1h(t)f(t)g(t)\,\frac{dt}{t}\notag\\
 &\quad+\int_0^1\!\int_0^1K(s,t)f(s)g(t)
       \,\frac{ds}{s}\frac{dt}{t}
 =\langle f,Ag\rangle_{\mathcal H}.
 \label{eq:PD-core-covariance}
\end{align}
Multiplication by \(h\) is boundedly invertible on \(\mathcal H\).  Also,
by \eqref{eq:smooth-kernel-product},
\[
 \int_0^1\!\int_0^1|K(s,t)|^2
       \,\frac{ds}{s}\frac{dt}{t}<\infty.
\]
Thus the integral part is Hilbert--Schmidt, the bilinear form
\eqref{eq:PD-core-covariance} extends continuously from
\(\mathcal D\) to \(\mathcal H\), and \(A\) is the resulting bounded
self-adjoint operator.  Nonnegativity on the dense core passes to the
extension.  Since \(A\) is an invertible multiplication operator plus a
compact operator, it is Fredholm of index zero.

If \(Af=0\), the integral equation and
\eqref{eq:smooth-kernel-product} give the explicit Cauchy--Schwarz bound
\begin{equation}
 |f(s)|\leq \frac1{\min h}
 \left(\int_0^1|K(s,t)|^2\frac{dt}{t}\right)^{1/2}
 \|f\|_{\mathcal H}
 \ll s\|f\|_{\mathcal H}
 \label{eq:PD-kernel-pointwise-control}
\end{equation}
for almost every \(s\).  Here \(\min_{[0,1]}h>0\), since
\(U-s\in[7/2,9/2]\).  Change the
representative on a null set so that this holds for every \(s\in(0,1]\).
Then
\[
 \sum_{x\in\Pi}|f(x)|\ll\sum_{x\in\Pi}x=U
 \qquad\mathbb P_U\hbox{-almost surely}.
\]
The same bound makes all one- and two-point second-moment integrals finite.
Truncating at \(x\geq\varepsilon\) and using dominated convergence extends
\eqref{eq:PD-core-covariance} to this particular \(f\).  Consequently
\[
 \operatorname{Var}_U(S_f)=\langle f,Af\rangle_{\mathcal H}=0.
\]
Thus \(S_f\) is a genuine almost-sure constant: for some \(c\),
\[
 Z_c(\pi):=\sum_{x\in\pi}f(x)-c=0
 \quad\text{for }\mathbb P_U\text{-almost every }\pi.
\]
We spell out the Palm disintegration which turns this almost-sure
statement into an additive identity.  Campbell--Mecke, disintegrated with
respect to the total mass, says for every nonnegative measurable
\(\Phi\)
\begin{align}
 &\mathbb E_U\sum_{r\in\Pi}
   \Phi(r,\Pi\setminus\{r\})
 =\frac1{\rho(U)}\int_0^1
   \rho(U-r)\mathbb E_{U-r}\Phi(r,\Pi)\,\frac{dr}{r},
 \label{eq:smooth-one-palm-disintegration}\\
 &\mathbb E_U\sum_{\substack{r,s\in\Pi\\r\ne s}}
   \Phi(r,s,\Pi\setminus\{r,s\})
 =\frac1{\rho(U)}\int_0^1\!\int_0^1
   \rho(U-r-s)\mathbb E_{U-r-s}\Phi(r,s,\Pi)
       \frac{dr}{r}\frac{ds}{s}.
 \label{eq:smooth-two-palm-disintegration}
\end{align}
Terms with a negative residual mass are interpreted as zero.  These
formulas follow first with all atoms at least \(\varepsilon\), where the
sums are finite, and then by monotone convergence.  Apply
\eqref{eq:smooth-two-palm-disintegration} to
\((Z_c(\pi)+f(r)+f(s))^2\) and
\eqref{eq:smooth-one-palm-disintegration} to
\((Z_c(\pi)+f(u))^2\).  Their left sides vanish.  Since
\(\rho(U-r-s)>0\) for \(r+s\leq1\), we obtain, for almost every such
pair \((r,s)\), under the \emph{same} residual law
\(\mathbb P_{U-r-s}\),
\begin{align*}
 Z_c(\Pi)+f(r)+f(s)&=0,\\
 Z_c(\Pi)+f(r+s)&=0.
\end{align*}
For the second equality we used the one-point formula at \(u=r+s\);
the pullback of its null set under \((r,s)\mapsto r+s\) is null by
Fubini.  Subtraction gives
\(f(r)+f(s)=f(r+s)\) for almost every \(r,s>0\) with \(r+s\leq1\).
For completeness, the measurable local Cauchy step can be seen by
convolution.  On each compact interval \(I\Subset(0,1)\), choose
\(\psi\in C_c^\infty(0,1-\sup I)\) with \(\int\psi=1\), and put
\[
 F_0(x):=\int\{f(x+y)-f(y)\}\psi(y)\,dy.
\]
The bound \(f(t)=O(t)\) gives local integrability, so translation
continuity in \(L^1_{\rm loc}\) makes \(F_0\) continuous.  Fubini and the
almost-everywhere Cauchy identity give \(F_0=f\) almost everywhere.
If two compact intervals overlap, their continuous representatives agree
almost everywhere on the overlap and hence, by continuity, agree
everywhere there.  They therefore glue to a continuous representative on
\((0,1)\).  The identity \(F_0(x+z)=F_0(x)+F_0(z)\) holds first almost
everywhere and then everywhere on each compact subtriangle by continuity.
A continuous local additive function is \(\lambda x\); the glued
representative shows that the same \(\lambda\) applies on all overlapping
intervals.  Thus \(f(t)=\lambda t\) almost everywhere.
Conversely this function is in the kernel because the
conditioned sum is constant.  Since zero is now an isolated eigenvalue of
a Fredholm self-adjoint operator, \eqref{eq:PD-Poincare} follows.

For completeness, put
\[
 \widetilde K(s,t):=\frac{K(s,t)}t,
\]
using at \(t=0\) the continuous value
\(F'(s)-F(s)F'(0)\).  On \(L^\infty(0,1]\),
\[
 Af=M_hf+\mathcal K_\infty f,\qquad
 (\mathcal K_\infty f)(s)
   =\int_0^1\widetilde K(s,t)f(t)\,dt.
\]
The image of the \(L^\infty\) unit ball under \(\mathcal K_\infty\) is
uniformly bounded and equicontinuous in \(C[0,1]\).  Arzel\`a--Ascoli
therefore makes \(\mathcal K_\infty\) compact as an operator into
\(L^\infty\).  Since \(M_h\) is invertible, \(A\) is Fredholm on
\(L^\infty\).  If a bounded \(f\) lies in its kernel, the integral
equation first gives \(f(t)=O(t)\), so the preceding \(L^2\) kernel
calculation applies and
\[
                         \ker_{L^\infty}A
              =\operatorname{span}\{t\}.
\]
The induced map
\[
 \overline A:L^\infty/\operatorname{span}\{t\}
       \longrightarrow\operatorname{Ran}A,\qquad [f]\longmapsto Af,
\]
is a bounded bijection between Banach spaces, because a Fredholm range is
closed.  The quotient norm here is
\(\|[f]\|=\inf_{\lambda\in\mathbb R}\|f-\lambda t\|_\infty\).
The open mapping theorem gives
\eqref{eq:PD-ordinary-Linf-inverse}.

For the weighted assertion write \(f(t)=tQ(t)\).  Then
\begin{equation}
 \frac{A(tQ)(s)}s
 =h(s)Q(s)+\int_0^1\frac{K(s,t)}sQ(t)\,dt.
\label{eq:PD-Xt-operator}
\end{equation}
The quotient kernel in \eqref{eq:PD-Xt-operator} is even more explicit:
\begin{equation}
 \frac{K(s,t)}s\longrightarrow F'(t)-F'(0)F(t)
 \quad(s\downarrow0),
 \label{eq:smooth-weighted-kernel-axis}
\end{equation}
uniformly in \(t\in[0,1]\).  Since \(F\in C^2[0,2]\), this defines a
continuous kernel on the whole compact square.  The corresponding
conjugated operator
\[
 (\mathcal BQ)(s):=\frac{A(tQ)(s)}s
 =h(s)Q(s)+\int_0^1\frac{K(s,t)}sQ(t)\,dt
\]
maps the \(L^\infty\) unit ball in its integral part into a uniformly
bounded, equicontinuous subset of \(C[0,1]\), and that part is compact by
Arzel\`a--Ascoli.  Thus \(\mathcal B\) is an invertible multiplication
operator plus a compact operator.  Its kernel consists of constants,
because \(Q\mapsto tQ\) is an isometry from \(L^\infty\) onto \(X_t\)
and \(tQ\in\ker A\) exactly when \(Q\) is constant.  The induced map from
\(L^\infty/\mathbb R\), equipped with the quotient norm
\(\inf_{\lambda\in\mathbb R}\|Q-\lambda\|_\infty\), to the closed range
of \(\mathcal B\) is a bounded bijection, and the open mapping theorem gives
\eqref{eq:PD-weighted-inverse}.
\end{proof}

\subsection{Finite bands, the low row, and arithmetic transfer}
\label{subsec:smooth-finite-bands}

Put
\[
 t_p:=\frac{\log p}{\log y},\qquad
 t_0:=\frac{\log W}{\log y}.
\]
Choose fixed \(0<\delta<1\) and a fixed regular relative partition
\[
 [t_0,\delta]=I_0,\qquad
 (\delta,1]=I_1\sqcup\cdots\sqcup I_s.
\]
Thus, after merging a possible terminal fragment with its neighbor,
\begin{equation}
 c_{\rm mesh}\eta\,\inf I_j
 \leq |I_j|\leq \eta\,\inf I_j
       \qquad(j\geq1)
\label{eq:smooth-regular-relative-mesh}
\end{equation}
for fixed \(c_{\rm mesh}>0\) and \(\eta>0\).  In the exponent convention
used in the tangent section, the low endpoint is
\(\delta_b=\theta\delta\), and
\[
 \theta(\delta+\eta)=\delta_b+\theta\eta.
\]
We write
\begin{align}
 \mathcal P_j&:=\{p:W<p\leq y,\ t_p\in I_j\},\notag\\
 \Omega_j(m)&:=\sum_{p\in\mathcal P_j}v_p(m),\notag\\
 H_j&:=\sum_{p\in\mathcal P_j}\frac1p,\qquad
 \alpha_j:=\frac1{H_j}
          \sum_{p\in\mathcal P_j}\frac{t_p}{p}.
\label{eq:smooth-band-data}
\end{align}
Let \(D=\operatorname{diag}(H_0,\ldots,H_s)\).

The active point constructed in
\Cref{subsec:smooth-head-physical} is supported on a fixed finite
disjoint union \(\mathcal S_n\) of positive-length physical intervals
with exact finite head patterns; any initial overlap is disjointified
before normalization.  Denote its weights by \(z_m^0\), and normalize
\[
 q_n=\sum_{m\in\mathcal S_n}z_m^0=q^{\rm act}(d)\asymp N,
 \qquad
 \nu_n(m):=\frac{z_m^0}{q_n}.
\]
Every \(z_m^0\) is \(O_W(1/L)\), and on a fixed positive-density
zero-pattern subpool it is bounded below by \(c_W/L\).  By
\Cref{lem:rough-guard-census}, only the promoted smooth anchors and bank
endpoints can meet these active smooth cells, so the relevant numerical
guards number \(O(yL)\).  By
\Cref{cor:marked-full-tilt-guard-stability}, deleting them changes every
normalized marked row by
\[
 O\!\left(\frac{y^2L^{O(1)}}n\right)=o(\alpha_0).
\]
This uses the actual guard geometry, rather than an assertion about an
arbitrary \(o(n)\) set.

For a compact vector of parameters, reweight only the active mass by
\begin{equation}
 z_m(\xi)=
 q_n\frac{z_m^0
   \exp\!\left\{L^{-1}\left(
       \sum_j\eta_j\Omega_j(m)+\zeta R(m)
             +\sum_\nu\kappa_\nu H_\nu^\circ(m)\right)\right\}}
 {\displaystyle\sum_{\ell\in\mathcal S_n}z_\ell^0
   \exp\!\left\{L^{-1}\left(
       \sum_j\eta_j\Omega_j(\ell)+\zeta R(\ell)
             +\sum_\nu\kappa_\nu H_\nu^\circ(\ell)\right)\right\}},
\label{eq:smooth-exponential-family}
\end{equation}
where \(R(m)=\log(m/n)\).  To define the head scores, choose one reference
pattern \(e_*\in\mathscr E\), put
\[
 H_e(m):=\1_{C_{e,-,n}\cup C_{e,+,n}}(m),\qquad
 H_e^\circ(m):=H_e(m)-\mathbb E_{\nu_n}H_e\quad(e\ne e_*),
\]
and enumerate these \(|\mathscr E|-1\) centered indicators as
\(H_1^\circ,\ldots,H_e^\circ\).  Changing the centering adds only a
constant to the exponential score and hence does not change the tilted
measure.  The denominator keeps the active mass exactly \(q_n\).  For any
two statistics \(F,G\),
\begin{equation}
 \frac{\partial}{\partial \xi_G}
       \sum_m z_m(\xi)F(m)
       =\frac{q_n}{L}\operatorname{Cov}_{\nu_{n,\xi}}(F,G).
\label{eq:smooth-covariance-Jacobian}
\end{equation}
Here and below a \emph{compact effective-tilt box} means that, after use
of the exact logarithmic identity, every coefficient of \(v_p(m)\) for
\(p>W\), and every remaining head or physical coefficient, has absolute
value at most a fixed \(B\).  This is the hypothesis of
\eqref{eq:marked-tilt-score}.  It does not impose a bounded value on an
unidentified redundant parameter.  In particular the whole low-band
contribution to the exponent is bounded by
\[
 \frac{B}{L}\sum_{p>W}v_p(m)\leq\frac{B}{\log W}+O_B(L^{-1}),
\]
so the fact that \(H_0\sim\log L\) creates no divergent count tilt.

We first state the quotient inverse, including the moving low row.

For a partition \(\mathcal P\), define the step injection and harmonic
averaging operators on \(\mathbb R^{s+1}\) by
\begin{align}
 (J_{\mathcal P}b)(t)&=b_j\quad(t\in I_j),
 & (E_{\mathcal P}f)_j&=\frac1{H_j^{(c)}}
       \int_{I_j}f(t)\,\frac{dt}{t},
 \label{eq:smooth-JE-definitions}\\
 B_{t_0,\mathcal P}&:=E_{\mathcal P}A^{(t_0)}J_{\mathcal P},
 & A^{(t_0)}f(s)&:=h(s)f(s)+
       \int_{t_0}^1K(s,t)f(t)\,\frac{dt}{t}.
 \label{eq:smooth-compressed-continuum-operator}
\end{align}
Here \(H_j^{(c)}=\int_{I_j}dt/t\),
\(D_c=\operatorname{diag}(H_j^{(c)})\), and
\(\alpha_j^{(c)}=E_{\mathcal P}(t)_j\).  Thus
\begin{equation}
 H_0^{(c)}=\log(\delta/t_0),\qquad
 \alpha_0^{(c)}=\frac{\delta-t_0}{H_0^{(c)}}.
 \label{eq:smooth-low-center}
\end{equation}
The continuum projection and gauge are
\begin{equation}
 P_{\alpha,c}b=b-\alpha^{(c)}
   \frac{(\alpha^{(c)})^TD_cb}{(\alpha^{(c)})^TD_c\alpha^{(c)}},
 \qquad \mathcal G_c:=\{b:(\alpha^{(c)})^TD_cb=0\}.
 \label{eq:smooth-continuum-gauge}
\end{equation}
The arithmetic projection and gauge, which are different finite-\(n\)
objects, are
\begin{equation}
 P_{\alpha,n}b=b-\alpha
   \frac{\alpha^TDb}{\alpha^TD\alpha},
 \qquad \mathcal G_n:=\{b:\alpha^TDb=0\}.
 \label{eq:smooth-arithmetic-gauge}
\end{equation}
From this point on \(\mathcal G\) without a subscript means
\(\mathcal G_n\); the continuum gauge is always written
\(\mathcal G_c\).
Because the two displayed denominators are bounded below, both
projections are uniformly bounded in \(\ell^\infty\).  We prove below,
for each fixed permitted mesh, that
\(\|P_{\alpha,n}-P_{\alpha,c}\|_{\infty\to\infty}=o_n(1)\).

Let \(Z=(R,H_1^\circ,\ldots,H_e^\circ)\), with one head indicator
deleted so that the constant score is absent, and put
\[
 \Gamma_{\xi,n}:=\operatorname{Cov}_{\nu_{n,\xi}}(Z,Z).
\]
When \(n\) is understood we write \(\Gamma_\xi=\Gamma_{\xi,n}\); in
particular
\(\Gamma_{0,n}:=\operatorname{Cov}_{\nu_n}(Z,Z)\) is the actual
finite-\(n\) baseline covariance, not a putative limiting covariance.
Define the genuine
head/physical Schur band matrix
\begin{equation}
 \mathcal C^Z_\xi:=q_n\left\{
  \operatorname{Cov}_\xi(\Omega,\Omega)
 -\operatorname{Cov}_\xi(\Omega,Z)\Gamma_\xi^{-1}
       \operatorname{Cov}_\xi(Z,\Omega)\right\}.
 \label{eq:smooth-joint-Schur-band-matrix}
\end{equation}
There is no ``row Schur complement'': normalization in
\eqref{eq:smooth-exponential-family} has already removed the constant
row, whose covariance would be zero.

We shall use the following finite-dimensional perturbation fact twice,
once for the arithmetic transfer and once for the nuisance Schur block.

\begin{lemma}[Stable projected Schur perturbation]
\label[lemma]{lem:stable-projected-Schur-perturbation}
Let \(D\) be positive diagonal, let \(P\) be the \(D\)-orthogonal
projection onto a gauge space \(\mathcal G\), and let \(B\) be symmetric.
Assume
\[
 b^TBb\geq\kappa b^TDb\quad(b\in\mathcal G),\qquad
 \left\|\left(PD^{-1}BP|_{\mathcal G}\right)^{-1}
       \right\|_{\infty\to\infty}\leq C.
\]
If a symmetric perturbation \(E\) satisfies
\begin{equation}
 |b^TEb|\leq\varepsilon b^TDb\quad(b\in\mathcal G),
 \qquad
 \|PD^{-1}EP\|_{\infty\to\infty}\leq\varepsilon,
 \label{eq:stable-Schur-perturbation-hypotheses}
\end{equation}
then, whenever \(\varepsilon<\min(\kappa/2,(2C)^{-1})\),
\(B-E\) has quadratic gap at least \(\kappa/2\) on \(\mathcal G\), and
\[
 \left\|\left(PD^{-1}(B-E)P|_{\mathcal G}\right)^{-1}
       \right\|_{\infty\to\infty}\leq2C.
\]
\end{lemma}

\begin{proof}
The quadratic assertion follows by subtracting the first bound in
\eqref{eq:stable-Schur-perturbation-hypotheses}.  On \(\mathcal G\),
factor the projected operator as
\[
 PD^{-1}(B-E)P
 =PD^{-1}BP\left\{I-
 (PD^{-1}BP|_{\mathcal G})^{-1}PD^{-1}EP\right\}.
\]
The second factor differs from the identity by an operator of
\(\ell^\infty\)-norm at most \(C\varepsilon<1/2\), so its Neumann
series has norm at most two.  This proves the inverse assertion.
\end{proof}

\medskip\noindent\emph{Uniformity convention for the next two lemmas.}\quad
The continuum constants used below depend only on \(U=9/2\), the
regularity constant \(c_{\rm mesh}\), and a fixed upper bound for
\(w:=\delta+\eta\); in particular they are chosen before \(W\), a tilt
box, or a particular mesh is instantiated.  The constants multiplying
the nonvanishing arithmetic errors \(W^{-1}\) are those already obtained
in the marked estimates and are independent of every later tilt-box
radius and mesh.  We first choose \(W\) large enough to absorb those
terms and then choose \(w\) below the continuum smallness threshold, as
required by the global order in
\eqref{eq:upper-construction-nested-quantifiers}.  After a radius \(B\)
and one permitted finite mesh \(\mathcal P\) have been fixed, the notation

\[
 r_{n;B,W,\mathcal P}=o_{n;B,W,\mathcal P}(1)
\]

means convergence uniformly for \(\xi\) in that \(B\)-box and uniformly
over the unit ball of the vector norm occurring in the estimate.  For a
quadratic-form estimate the unit ball means \(b^TDb\leq1\).  Its
threshold may depend on \((B,W,\delta,\eta,\mathcal P)\); no uniform
threshold over unbounded boxes or meshes varying with \(n\) is used.
Finally, the continuum comparison object is itself
\(B_{t_0,\mathcal P}\) with the same
\(t_0=(\log W)/(\log y)\) as the arithmetic operator at that \(n\).
Thus the argument compares two moving finite-dimensional objects at
each \(n\), rather than claiming convergence to a single fixed band
matrix.

\begin{lemma}[Moving-low-cell arithmetic quotient]
\label[lemma]{lem:arithmetic-band-quotient}
There are structural thresholds \(W_0<\infty\) and \(w_0>0\).
Choose \(W\geq W_0\) and then fixed \(\delta,\eta>0\) with
\(w:=\delta+\eta\leq w_0\).  There are constants \(\kappa>0\) and
\(C_W<\infty\), independent of the permitted regular mesh and of the
radius \(B\) of a subsequently fixed compact effective-tilt box, with
the following quantifiers: for every fixed \(B\) and every fixed
permitted mesh \(\mathcal P\), there is
\(n_0=n_0(B,W,\delta,\eta,\mathcal P)\) such that the conclusions below
hold for all \(n\geq n_0\) and every \(\xi\) in the \(B\)-box.  No common
threshold over arbitrarily fine meshes or over unbounded boxes is
asserted.  Then \(\Gamma_\xi\) is positive
definite and
\begin{align}
 b^T\mathcal C^Z_\xi b
 &\geq\kappa q_n\inf_{\lambda\in\mathbb R}
       \sum_jH_j|b_j-\lambda\alpha_j|^2,
 \label{eq:smooth-arithmetic-Poincare}\\
 \left\|
 \left(P_{\alpha,n}D^{-1}\frac{\mathcal C^Z_\xi}{q_n}P_{\alpha,n}
       \bigm|_{\alpha^TDb=0}\right)^{-1}
 \right\|_{\ell^\infty\to\ell^\infty}&\ll_W1.
 \label{eq:smooth-arithmetic-Linf-inverse}
\end{align}
The constants \(\kappa,C_W\) are uniform as \(t_0\downarrow0\) and over
all regular meshes satisfying
\eqref{eq:smooth-regular-relative-mesh}; only the threshold \(n_0\) may
depend on the particular fixed mesh and box.
\end{lemma}

\begin{proof}
\emph{Step 1: the continuum inverse and form gap, uniformly in the
moving low cell.}\quad
This step contains no primes and no tilt parameter.  In particular its
constants are independent of \(W,B,n\), and of the locations of the
positive-cell endpoints subject to
\eqref{eq:smooth-regular-relative-mesh}.
Write \(f=J_{\mathcal P}b\).  Since \(h(s)=1+O(s)\) and
\(|K(s,t)|\leq Cst\), direct harmonic averaging on
\(I_0=[t_0,\delta]\) gives
\begin{equation}
 (B_{t_0,\mathcal P}b)_0
 =\{1+O(\delta)\}b_0+
 O\!\left(\alpha_0^{(c)}\right)\|b\|_\infty;
 \label{eq:smooth-low-row}
\end{equation}
indeed
\[
 \frac1{H_0^{(c)}}\int_{I_0}s\,\frac{ds}{s}
   =\alpha_0^{(c)},\qquad
 \int_{t_0}^1t|f(t)|\,\frac{dt}{t}\leq\|b\|_\infty.
\]
The \(C^2\)-regularity of \(F\) also gives
\[
 |h'(s)|\leq C,\qquad |\partial_sK(s,t)|\leq Ct.
\]
Hence, on a positive relative cell \(I_j\), the oscillation of
\(A^{(t_0)}f\) is at most \(C\eta\|b\|_\infty\).  These are ordinary
supremum-norm estimates and their constants do not involve the growing
harmonic mass \(H_0^{(c)}\).

Suppose the projected inverse were not uniform.  There would be
\(t_{0,\nu}\downarrow0\), permitted partitions, and gauge-fixed vectors
\(b_\nu\) such that
\[
 \|b_\nu\|_\infty=1,\qquad
 \|P_{\alpha,c}B_{t_{0,\nu},\mathcal P_\nu}b_\nu\|_\infty=o(1).
\]
By \eqref{eq:smooth-continuum-gauge}, for bounded scalars
\(\beta_\nu=O(1)\),
\begin{equation}
 B_{t_{0,\nu},\mathcal P_\nu}b_\nu
 =\beta_\nu\alpha^{(c)}+o_{\ell^\infty}(1).
 \label{eq:smooth-unprojected-compressed-equation}
\end{equation}
More explicitly, applying \(I-P_{\alpha,c}\) gives
\[
 \beta_\nu=
 \frac{(\alpha^{(c)})^TD_c
       B_{t_{0,\nu},\mathcal P_\nu}b_\nu}
      {(\alpha^{(c)})^TD_c\alpha^{(c)}}+o(1).
\]
The denominator is bounded below by the cells meeting \([1/2,1]\), while
\(\sum_jH_j^{(c)}\alpha_j^{(c)}=1-t_0\) and the compressed operators are
uniformly bounded.  Hence \(\beta_\nu=O(1)\).  Formula
\eqref{eq:smooth-low-row}
now gives
\begin{equation}
 |b_{\nu,0}|=O(\delta)+o(1).
 \label{eq:smooth-low-value-control}
\end{equation}
On each positive cell, the oscillation estimate turns
\eqref{eq:smooth-unprojected-compressed-equation} into
\[
 A^{(t_{0,\nu})}J_{\mathcal P_\nu}b_\nu(s)
 =\beta_\nu s+O(\delta+\eta)+o(1).
\]
On the low cell both sides are \(O(\delta)+o(1)\), by
\eqref{eq:smooth-low-value-control} and
\eqref{eq:smooth-kernel-product}, so the same estimate holds there.
Extend the low value to \((0,t_{0,\nu})\) by
\(b_{\nu,0}t/t_{0,\nu}\).  For \(s\geq t_{0,\nu}\), the newly inserted
integral term is bounded by
\[
 \int_0^{t_{0,\nu}}|K(s,t)|\,|b_{\nu,0}|
       \frac{t}{t_{0,\nu}}\frac{dt}{t}
 \ll |b_{\nu,0}|s t_{0,\nu}=O(t_{0,\nu}).
\]
For \(0<s<t_{0,\nu}\), the diagonal term is
\(O(|b_{\nu,0}|)=O(\delta)+o(1)\), while the integral term is
\(O(s)\); this is already within the error allowed below.  Thus, for
the resulting bounded function \(\widetilde f_\nu\),
\begin{equation}
 A\widetilde f_\nu=\beta_\nu t+O(\delta+\eta)+o(1)
 \quad\hbox{in }L^\infty.
 \label{eq:smooth-full-reconstructed-equation}
\end{equation}
Pairing with the kernel vector \(t\) shows
\(\beta_\nu=O(\delta+\eta)+o(1)\).  Then
\eqref{eq:PD-ordinary-Linf-inverse} puts \(\widetilde f_\nu\) within
\(O(\delta+\eta)+o(1)\) of \(\lambda_\nu t\).  Finally,
\((\alpha^{(c)})^TD_cb_\nu=0\) is precisely
\(\int_{t_0}^1tJ_{\mathcal P}b_\nu\,dt/t=0\); the linear extension
changes this pairing by \(O(t_0)\).  Hence
\(\lambda_\nu=O(\delta+\eta)+o(1)\), contradicting
\(\|b_\nu\|_\infty=1\) once \(\delta+\eta\) is below a fixed
threshold.  For the quadratic estimate normalize
\(\|b\|_{D_c}=1\) and use the same linear extension on
\((0,t_0)\).  Its extra squared norm is
\[
 \int_0^{t_0}\left|b_0\frac{t}{t_0}\right|^2\frac{dt}{t}
 =\frac{|b_0|^2}{2}
 \leq\frac1{2H_0^{(c)}}=o(1),
\]
and its change in the \(t\)-pairing is at most
\(|b_0|t_0/2=o(1)\).  If the compressed quadratic gap failed, the
quadratic forms, and not only the norms, compare as follows.  With
\(f=J_{\mathcal P}b\) on \([t_0,1]\), harmonic averaging gives the
exact identity
\begin{equation}
 b^TD_cB_{t_0,\mathcal P}b
 =\int_{t_0}^1 f(s)(A^{(t_0)}f)(s)\,\frac{ds}{s}.
 \label{eq:smooth-compressed-form-identity}
\end{equation}
If \(e(s)=b_0s/t_0\) on \((0,t_0)\), then
\(|K(s,t)|\ll st\), Cauchy--Schwarz, and the preceding norm bound give
\begin{align}
 \left|\int_0^{t_0}h(s)e(s)^2\,\frac{ds}{s}\right|
     &\ll |b_0|^2=o(1),\notag\\
 \left|\int_0^{t_0}\!\int_{t_0}^1
       K(s,t)e(s)f(t)\,\frac{ds}{s}\frac{dt}{t}\right|
     &\ll |b_0|t_0\|f\|_{\mathcal H}=o(1),\notag\\
 \left|\int_0^{t_0}\!\int_0^{t_0}
       K(s,t)e(s)e(t)\,\frac{ds}{s}\frac{dt}{t}\right|
     &\ll |b_0|^2t_0^2=o(1).
 \label{eq:smooth-low-extension-form-error}
\end{align}
Consequently
\begin{equation}
 \langle\widetilde f,A\widetilde f\rangle_{\mathcal H}
       -b^TD_cB_{t_0,\mathcal P}b=o(1)
 \label{eq:smooth-full-compressed-form-comparison}
\end{equation}
under unit \(D_c\)-normalization.  If the compressed quadratic gap
failed, the extended functions would therefore have
\(L^2(dt/t)\)-norm \(1+o(1)\), asymptotically zero pairing with \(t\),
and \(\langle\widetilde f,A\widetilde f\rangle=o(1)\).  The
Poisson--Dickman Poincar\'e
inequality \eqref{eq:PD-Poincare} would put them \(o(1)\)-close to
\(\operatorname{span}\{t\}\), while the gauge would force the scalar
to be \(o(1)\), a contradiction.  This proves the continuum quadratic
gap uniformly as the low endpoint moves to zero.  We have therefore
proved, with structural constants \(C_c<\infty\) and \(\kappa_c>0\),
\begin{align}
 \left\|\left(P_{\alpha,c}B_{t_0,\mathcal P}P_{\alpha,c}
       \bigm|_{\mathcal G_c}\right)^{-1}\right\|_{\infty\to\infty}
 &\leq C_c,\label{eq:smooth-continuum-band-inverse}\\
 b^TD_cB_{t_0,\mathcal P}b
 &\geq\kappa_c\,b^TD_cb\qquad(b\in\mathcal G_c).
 \label{eq:smooth-continuum-band-gap}
\end{align}
Both constants are uniform for all sufficiently small \(t_0>0\), hence
uniform as \(t_0\downarrow0\), and for every permitted regular mesh once
\(w\leq w_0\).  This is precisely the range used because
\(t_0=(\log W)/(\log y)\to0\).  These are the continuum results; no
arithmetic object has yet been used.

\medskip\noindent
\emph{Step 2: arithmetic cell data and normalized operator
convergence.}\quad
We now give the discrete transfer rather than hiding it in a convergence
phrase.  Put \(I_p=1_{p\mid m}\) and
\(\Omega_i^{\rm sf}=\sum_{p\in\mathcal P_i}I_p\).  For a prime
coefficient vector \(a=(a_p)\), define the squarefree and full normalized
prime operators and band averaging by
\begin{align}
 (\mathsf A^{\rm sf}_{n,\xi}a)_p
   &:=p\operatorname{Cov}_{\nu_{n,\xi}}
       \left(I_p,\sum_q a_qI_q\right),\notag\\
 (\mathsf A_{n,\xi}a)_p
   &:=p\operatorname{Cov}_{\nu_{n,\xi}}
       \left(v_p,\sum_q a_qv_q\right),
 & (\mathsf E_na)_i&:=\frac1{H_i}
       \sum_{p\in\mathcal P_i}\frac{a_p}{p}.
 \label{eq:smooth-prime-band-operators}
\end{align}
Thus
\[
 D^{-1}\operatorname{Cov}_\xi(\Omega^{\rm sf},\Omega^{\rm sf})
 =\mathsf E_n\mathsf A^{\rm sf}_{n,\xi}J_n,
\]
where \((J_nb)_p=b_{j(p)}\).  The squarefree marked formulas give a diagonal
multiplier plus a compact kernel and an error \(E_{pq}\):
\begin{align}
 \operatorname{Var}(1_{p\mid m})
   &=\frac{h(t_p)}p+E_{pp},\notag\\
 \operatorname{Cov}(1_{p\mid m},1_{q\mid m})
   &=\frac{K(t_p,t_q)}{pq}+E_{pq}\quad(p\ne q),
 \label{eq:smooth-arithmetic-kernel}\\
 \sup_{W<p\le y}\sum_{W<q\le y}p|E_{pq}|
   &\le \frac{C_{\rm sf}}W+\epsilon^{\rm sf}_{B,W}(n).
 \label{eq:smooth-normalized-prime-row-error}
\end{align}
Here \(C_{\rm sf}\) is absolute and
\(\epsilon^{\rm sf}_{B,W}(n)\to0\) for each fixed compact tilt box.
The \(1/W\) term contains the squarefree diagonal correction
\(-h(t_p)^2/p^2\); equivalently one may include it in the exact diagonal
multiplier.  The leading functions are the same \(h,K\) throughout a compact
effective-tilt box: \Cref{lem:marked-tilted-squarefree} puts the entire
tilt dependence in \(E_{pq}\).  Thus the continuum kernel is independent
of the limiting fugacity.  The definition of the box above
and \(\sum_{p>W}v_p/L\leq1/\log W+O(1/L)\) are exactly what make the
constants in this assertion uniform.

The passage from the prime row norm to the band norm is the exact
calculation
\begin{equation}
 \left|\frac1{H_i}\sum_{p\in\mathcal P_i}\frac1p
       \sum_qpE_{pq}b_{j(q)}\right|
 \leq\frac1{H_i}\sum_{p\in\mathcal P_i}\frac1p
       \sum_qp|E_{pq}|\,|b_{j(q)}|.
 \label{eq:smooth-band-normalized-row-transfer}
\end{equation}
Thus the large value \(H_0\asymp\log L\) cancels rather than multiplying
the error.  Mertens' theorem gives, uniformly cell by cell,
\begin{equation}
 H_0=\log L+O_{W,\delta}(1),\qquad
 \alpha_0=\frac{\delta+o(1)}{H_0}.
 \label{eq:smooth-discrete-low-asymptotics}
\end{equation}

We record precisely which discrete cell data converge to their continuum
counterparts.  On every positive cell, quantitative Mertens summation
gives
\begin{equation}
 H_j-H_j^{(c)}=o_n(1),\qquad
 H_j\alpha_j-H_j^{(c)}\alpha_j^{(c)}=o_n(1)\quad(j\geq1).
 \label{eq:smooth-positive-cell-data-convergence}
\end{equation}
On the moving low cell the first absolute convergence need not hold and is
not used: the fixed lower cutoff can leave a bounded Mertens constant.
The correct statements are
\begin{equation}
 H_0=H_0^{(c)}+O_W(1),\quad
 \frac{H_0}{H_0^{(c)}}=1+o_n(1),\quad
 H_0\alpha_0=\delta-t_0+o_n(1)
                  =H_0^{(c)}\alpha_0^{(c)}+o_n(1).
 \label{eq:smooth-low-cell-data-convergence}
\end{equation}
Indeed the last relation is partial summation applied to
\((\log y)^{-1}\sum_{W<p\leq y^\delta}(\log p)/p\).
It follows, for the fixed finite mesh, that
\begin{align}
 \max_j|\alpha_j-\alpha_j^{(c)}|&=o_n(1),&
 \max_j\left|\frac{\alpha_j}{\alpha_j^{(c)}}-1\right|&=o_n(1),\notag\\
 \alpha^TD\alpha-(\alpha^{(c)})^TD_c\alpha^{(c)}&=o_n(1),&
 \|P_{\alpha,n}-P_{\alpha,c}\|_{\infty\to\infty}&=o_n(1).
 \label{eq:smooth-two-projection-convergence}
\end{align}
For the relative assertion, use
\eqref{eq:smooth-low-cell-data-convergence} on the low cell; on every
positive cell, \(\alpha_j^{(c)}\gg\delta\), so the absolute convergence
is already relative.  This is exactly the additional information needed
in the sharp norm.  Put
\[
 S_n:=\operatorname{diag}(\alpha_0,\ldots,\alpha_s),\qquad
 S_c:=\operatorname{diag}(\alpha_0^{(c)},\ldots,\alpha_s^{(c)}).
\]
Then
\begin{equation}
 \|S_nS_c^{-1}-I\|_{\infty\to\infty}=o_n(1),\qquad
 \|S_n^{-1}P_{\alpha,n}S_n
      -S_c^{-1}P_{\alpha,c}S_c\|_{\infty\to\infty}=o_n(1).
 \label{eq:smooth-sharp-projection-convergence}
\end{equation}
Indeed, if \(\omega_{n,j}=H_j\alpha_j^2\) and
\(\omega_{c,j}=H_j^{(c)}(\alpha_j^{(c)})^2\), then
\[
 S_n^{-1}P_{\alpha,n}S_n
   =I-\mathbf1\frac{\omega_n^T}{\mathbf1^T\omega_n},\qquad
 S_c^{-1}P_{\alpha,c}S_c
   =I-\mathbf1\frac{\omega_c^T}{\mathbf1^T\omega_c}.
\]
The positive-cell coordinates of \(\omega_n-\omega_c\) tend to zero by
\eqref{eq:smooth-positive-cell-data-convergence} and the relative
center convergence.  On the low cell, both coordinates are
\((\delta-t_0)^2/H_0^{(c)}+o(1)=o(1)\).  Their total masses converge to
the same positive limit by
\eqref{eq:smooth-two-projection-convergence}, proving the second operator
convergence.
Thus the continuum and arithmetic gauges are compared only through
quantities which genuinely converge; no assertion
\(\|D-D_c\|\to0\) is made at the moving low cell.

For clarity, here is the double prime-band quadrature as a normalized
operator calculation.  Put
\[
 (TJ_{\mathcal P}b)(s):=\int_{t_0}^1
       K(s,t)(J_{\mathcal P}b)(t)\,\frac{dt}{t}.
\]
Let \(e_k\) be the step vector supported on one cell.  Quantitative
Mertens summation and partial summation give, for every positive cell,
\[
 \sup_{0\leq s\leq1}\left|
  \sum_{t_q\in I_k}\frac{K(s,t_q)}q
  -\int_{I_k}K(s,t)\,\frac{dt}{t}\right|=o_n(1).
\]
For the moving low cell, split first at a fixed \(\tau>0\).  The same
estimate applies on \([\tau,\delta]\), whereas
\(|K(s,t)|\ll st\) makes both the prime sum and the integral over
\([t_0,\tau]\) at most \(Cs\tau\).  Letting first \(n\to\infty\)
and then \(\tau\downarrow0\) proves the same conclusion for \(I_0\),
uniformly in \(s\).  Since the fixed positive-cell mesh has finite
dimension, summing over the vectors \(e_k\) gives input
operator-norm convergence for \(\|b\|_\infty\leq1\).

The output averaging is a second application of the same argument.  Define
the actual doubly averaged kernel band operator by
\[
 (\widehat B_n^K b)_i:=\frac1{H_i}
   \sum_{p\in\mathcal P_i}\frac1p
   \sum_{W<q\leq y}\frac{K(t_p,t_q)}q b_{j(q)}.
\]
Equicontinuity in the first variable and Mertens summation now give
\begin{equation}
 \sup_{\|b\|_\infty\leq1}
 \left\|\widehat B_n^K b
       -E_{\mathcal P}TJ_{\mathcal P}b\right\|_\infty=o_n(1),
 \label{eq:smooth-collective-compact-quadrature}
\end{equation}
for each fixed permitted mesh.  On the low output row the unnormalized
kernel numerator is \(O(\delta)\), while
\(H_0\asymp H_0^{(c)}\asymp\log L\); hence the fixed-cutoff Mertens
constant is absorbed after normalization.  The diagonal multiplier is
handled at the same time:
\[
 \frac1{H_i}\sum_{p\in\mathcal P_i}\frac{h(t_p)}p
 -\frac1{H_i^{(c)}}\int_{I_i}h(t)\,\frac{dt}{t}=o_n(1).
\]
This includes \(I_0\), since \(h(t)=1+O(t)\) there.  Thus both prime
indices and the diagonal, rather than only the input prime index,
converge to the compressed continuum operator.  The operator constants
are mesh-uniform, but the convergence threshold may depend on the one
fixed mesh.  Equation \eqref{eq:smooth-band-normalized-row-transfer}
handles the remaining error \(E\).

Fix one permitted positive-cell mesh; only its low endpoint now moves
with \(n\).  Define the \emph{leading} arithmetic band operator, before
any marked, Bernoulli, prime-power, guard, or nuisance error is added, by
\begin{equation}
 (\widehat B_n^{(0)}b)_i:=\frac1{H_i}
 \sum_{p\in\mathcal P_i}\frac1p\left\{
 h(t_p)b_i+\sum_{\substack{W<q\leq y\\q\ne p}}
       \frac{K(t_p,t_q)}q\,b_{j(q)}\right\}.
 \label{eq:smooth-leading-arithmetic-band-operator}
\end{equation}
The omitted \(p=q\) kernel diagonal has operator norm \(o_n(1)\):
\[
 \max_i\frac1{H_i}\sum_{p\in\mathcal P_i}
       \frac{|K(t_p,t_p)|}{p^2}=o_n(1).
\]
Indeed \(K(t,t)=O(t^2)\); on \(I_0\) the numerator is
\(O_W((\log y)^{-2})\), and on a positive cell every prime tends to
infinity.  Consequently the preceding double quadrature, the multiplier
comparison, and \eqref{eq:smooth-two-projection-convergence} imply
\begin{equation}
 \left\|P_{\alpha,n}\widehat B_n^{(0)}P_{\alpha,n}
       -P_{\alpha,c}B_{t_0,\mathcal P}P_{\alpha,c}
 \right\|_{\infty\to\infty}=o_{n;W,\mathcal P}(1).
 \label{eq:smooth-leading-projected-operator-convergence}
\end{equation}
All operators in this display act on the common coordinate space; the
two ranges are still different and are identified next.

Put
\[
 \Phi_n:=P_{\alpha,n}|_{\mathcal G_c}:
       \mathcal G_c\longrightarrow\mathcal G_n,\qquad
 \varepsilon_n:=\|P_{\alpha,n}-P_{\alpha,c}\|_{\infty\to\infty}.
\]
For \(x\in\mathcal G_c\),
\(\|\Phi_nx-x\|_\infty\leq\varepsilon_n\|x\|_\infty\).  Hence, if
\(\varepsilon_n<1\), \(\Phi_n\) is injective; both gauge spaces have
codimension one, so it is onto, and
\begin{equation}
 \|\Phi_n\|\leq1+\varepsilon_n,\qquad
 \|\Phi_n^{-1}\|\leq(1-\varepsilon_n)^{-1}.
 \label{eq:smooth-ordinary-gauge-isomorphism}
\end{equation}
If \(y=\Phi_nx\in\mathcal G_n\), the same calculation gives the
ambient-space estimate
\[
 \|\Phi_n^{-1}y-y\|_\infty
 =\|x-P_{\alpha,n}x\|_\infty
 \leq\frac{\varepsilon_n}{1-\varepsilon_n}\|y\|_\infty.
\]
Pulling the left side of
\eqref{eq:smooth-leading-projected-operator-convergence} back by \(\Phi_n\)
therefore gives the following comparison.  Indeed, if
\(\widehat T_n=P_{\alpha,n}\widehat B_n^{(0)}P_{\alpha,n}\) and
\(T_c=P_{\alpha,c}B_{t_0,\mathcal P}P_{\alpha,c}\), then for
\(x\in\mathcal G_c\)
\[
 \widehat T_n\Phi_nx-\Phi_nT_cx
 =(\widehat T_n-T_c)\Phi_nx
   +T_c(\Phi_n-I)x+(I-P_{\alpha,n})T_cx.
\]
The last two terms are \(O(\varepsilon_n)\|x\|_\infty\), because
\(T_cx\in\mathcal G_c\) and the leading operators are uniformly bounded.
Together with \eqref{eq:smooth-ordinary-gauge-isomorphism}, this gives
\begin{equation}
 \left\|\Phi_n^{-1}
   \left(P_{\alpha,n}\widehat B_n^{(0)}P_{\alpha,n}
       \bigm|_{\mathcal G_n}\right)\Phi_n
  -P_{\alpha,c}B_{t_0,\mathcal P}P_{\alpha,c}
       \bigm|_{\mathcal G_c}\right\|_{\infty\to\infty}=o_n(1).
 \label{eq:smooth-pulled-ordinary-operator-convergence}
\end{equation}
The bounds in \eqref{eq:smooth-continuum-band-inverse} and a Neumann
series now establish the leading arithmetic inverse, with norm at most
\(2C_c\), for all sufficiently large \(n\).  Only \emph{after} this
inverse has been obtained do we use
\eqref{eq:smooth-normalized-prime-row-error},
\eqref{eq:smooth-band-normalized-row-transfer}, and
\eqref{eq:smooth-power-row-transfer}: choose \(W_0\) so that their
box-independent \(O(1/W)\) operator norm, after multiplication by the
uniformly bounded projections and gauge maps, is at most
\((8C_c)^{-1}\), and then, for the fixed box and mesh, increase \(n_0\)
until their vanishing remainders have the same bound.  A second Neumann
series gives the projected inverse for the actual un-Schur covariance.
This order removes any possibility of using an inverse to justify its
own leading approximation.

\medskip\noindent
\emph{Step 3: quadratic-form convergence on the moving unit sphere.}\quad
Let \(b^TDb=1\).  Then
\[
 |b_0|\leq H_0^{-1/2},\qquad
 \sum_{j\geq1}H_j|b_j|^2\leq1.
\]
The norm comparison is now literal.  On positive cells the dimension is
fixed, \(H_j\gg_{\delta,\eta,c_{\rm mesh}}1\), and
\(H_j-H_j^{(c)}=o_n(1)\).  On the low cell,
\[
 |H_0-H_0^{(c)}|\,|b_0|^2
       \ll_W H_0^{-1}=o_n(1).
\]
Thus
\[
 \|J_{\mathcal P}b\|_{L^2([t_0,1],dt/t)}^2=1+o_n(1).
\]
The extension on \((0,t_0)\) has squared norm
\(|b_0|^2/2=o_n(1)\), so the full reconstructed norm is \(1+o_n(1)\).
The same decomposition, with one factor \(t\), shows that an arithmetic
gauge vector has continuum \(t\)-pairing \(o_n(1)\).

We next compare the leading forms, splitting both prime indices into
\(I_0\) and the positive cells.  The positive--positive block converges
uniformly on the present unit sphere by the fixed finite-dimensional
two-index quadrature.  For a block with one low index,
\(|K(s,t)|\leq Cst\), Cauchy--Schwarz on the positive cells, and
\(\sum_{p\in\mathcal P_0}t_p/p=O(\delta)\) give
\[
 |\mathscr K_{0,+}(b)|
 \ll |b_0|\left(\sum_{p\in\mathcal P_0}\frac{t_p}{p}\right)
       \sum_{j\geq1}|b_j|
           \sum_{q\in\mathcal P_j}\frac{t_q}{q}
 \ll H_0^{-1/2}=o_n(1).
\]
Here the last sum is uniformly bounded by Cauchy--Schwarz:
\[
 \sum_{j\geq1}H_j\alpha_j|b_j|
 \leq\left(\sum_{j\geq1}H_j|b_j|^2\right)^{1/2}
      \left(\sum_{j\geq1}H_j\alpha_j^2\right)^{1/2}
 \ll1.
\]
The identical bound holds for the continuum low--positive block.  The
two low--low blocks are each
\[
 O\!\left(|b_0|^2
       \left(\sum_{p\in\mathcal P_0}\frac{t_p}{p}\right)^2\right)
   +O\!\left(|b_0|^2
       \left(\int_{I_0}dt\right)^2\right)
 =O(H_0^{-1})=o_n(1).
\]
For the multiplier part, write \(h=1+(h-1)\).  The constant term leaves
only \((H_0-H_0^{(c)})b_0^2=o(1)\) on the low row, and the remainder is
handled by weighted Mertens summation because
\((h(t)-1)/t\) extends continuously to \(t=0\).  The positive
multiplier rows converge by ordinary cellwise quadrature.  Finally, the
removed discrete \(p=q\) kernel diagonal is \(o_n(1)\) on the unit
sphere by the estimate preceding
\eqref{eq:smooth-leading-projected-operator-convergence}.  Combining
these four estimates with
\eqref{eq:smooth-compressed-form-identity}--
\eqref{eq:smooth-full-compressed-form-comparison} proves, uniformly for
\(b^TDb=1\),
\begin{equation}
 b^TD\widehat B_n^{(0)}b-
 \langle\widetilde f,A\widetilde f\rangle_{\mathcal H}=o_n(1),
 \label{eq:smooth-discrete-full-form-comparison}
\end{equation}
where \(\widetilde f\) is the linear low extension of
\(J_{\mathcal P}b\).  This explicitly proves both the norm and
quadratic-form convergence needed for the Poincar\'e gap; no equivalence
constant depending on \(H_0\) has been used.
Indeed, if in addition \(b\in\mathcal G_n\), then the norm and gauge
comparisons above give
\[
 \|\widetilde f\|_{\mathcal H}^2=1+o_n(1),\qquad
 \langle\widetilde f,t\rangle_{\mathcal H}=o_n(1).
\]
Since \(\|t\|_{\mathcal H}^2=1/2\), the exact Hilbert-space projection
formula yields
\[
 \inf_{\lambda\in\mathbb R}
 \|\widetilde f-\lambda t\|_{\mathcal H}^2
 =\|\widetilde f\|_{\mathcal H}^2
   -\frac{|\langle\widetilde f,t\rangle_{\mathcal H}|^2}
         {\|t\|_{\mathcal H}^2}
 =1+o_n(1).
\]
Thus \eqref{eq:PD-Poincare} and
\eqref{eq:smooth-discrete-full-form-comparison} give, for all
sufficiently large \(n\),
\begin{equation}
 b^TD\widehat B_n^{(0)}b\geq\frac{\kappa_0}{2}\,b^TDb
 \qquad(b\in\mathcal G_n).
 \label{eq:smooth-leading-arithmetic-form-gap}
\end{equation}
To transfer every remaining arithmetic error literally, let
\(R=(R_{ij})\) be any symmetric band covariance error satisfying
\[
 \sup_i\frac1{H_i}\sum_j|R_{ij}|\leq\varepsilon.
\]
The normalized operator \(D^{-1}R\) is self-adjoint for the
\(D\)-inner product, and
\begin{align*}
 |b^TRb|
 &\leq\frac12\sum_{i,j}|R_{ij}|(|b_i|^2+|b_j|^2)\\
 &=\sum_i|b_i|^2\sum_j|R_{ij}|\\
 &\leq\varepsilon\,b^TDb.
\end{align*}
Thus its \(D\)-Hilbert operator norm is at most \(\varepsilon\).
Applying this estimate to the squarefree, prime-power, and quadrature
errors transfers the continuum quadratic gap to the prime-band
covariance form.

Replacing \(1_{p\mid m}\) by \(v_p(m)\) adds
\begin{equation}
 \sup_p\sum_qp|E^{\rm pow}_{pq}|
 \le \frac{C_{\rm pow}}W+\epsilon_{B,W}(n)
 \label{eq:smooth-power-row-transfer}
\end{equation}
by \eqref{eq:marked-prime-power-aggregate}.  The constants
\(C_{\rm sf},C_{\rm pow}\) are known before the compact ODE box is
chosen.  This is the estimate invoked only after the leading inverse in
\eqref{eq:smooth-pulled-ordinary-operator-convergence} was established.
The already chosen \(W\geq W_0\) makes its \(1/W\) contribution smaller
than the reserved continuum inverse and form gaps; after a fixed box is
specified, only the threshold for \(n\) is increased to absorb the
vanishing remainder.

\medskip\noindent
\emph{Step 4: the finite nuisance block and its Schur perturbation.}\quad
It remains to justify the joint Schur complement in
\eqref{eq:smooth-joint-Schur-band-matrix}.  Every head pattern retains a
fixed mass in each of two separated positive-length physical subpools.
If a centered linear combination \(aR+\sum_\nu b_\nu H_\nu^\circ\)
vanished, varying \(R\) inside either physical subpool forces \(a=0\),
and the affinely independent simplex head patterns then force every
\(b_\nu=0\).  The fixed mass and separation margins make this quantitative:
\begin{equation}
 \Gamma_\xi\geq\gamma_W I
 \label{eq:smooth-joint-head-physical-gap}
\end{equation}
throughout every subsequently fixed compact tilt box, with
\(\gamma_W\) chosen before the radius of that box.  We prove this
uniformity directly; no convergence of the baseline mixture weights is
required.  By
\Cref{lem:smooth-baseline-active-measure}, there is a number
\(\lambda_W>0\) such that
\(\lambda_{e,\sigma}(n)\geq\lambda_W\) for every head pattern \(e\),
both signs \(\sigma\), and all sufficiently large \(n\).  Moreover
\[
 \frac{\inf I_+(n)}n-\frac{\sup I_-(n)}n=a_+-b_->0,
\]
so the corresponding ranges of \(R=\log(m/n)\) have a fixed positive
separation.

We first prove a finite-\(n\), uniform baseline bound.  For
\(F=aR+\sum_\nu b_\nu H_\nu^\circ\), compare, within each fixed head
pattern, its conditional means on the two physical cells.  The head term
is identical on those cells, whereas their \(R\)-means are separated by
at least \(\log a_+-\log b_->0\).  The elementary two-point variance
bound and the lower bound \(\lambda_W\) therefore give
\[
                         \operatorname{Var}_{\nu_n}(F)\geq c_{R,W}a^2.
\]
On the other hand the finitely many head-pattern vectors are affinely
independent after the constant indicator has been deleted.  Their masses
are all at least \(2\lambda_W\), so, uniformly in \(n\),
\[
 \operatorname{Var}_{\nu_n}\!\left(\sum_\nu b_\nu H_\nu^\circ\right)
       \geq c_{H,W}\sum_\nu b_\nu^2.
\]
Since \(R\) is uniformly bounded,
\[
 \operatorname{Var}_{\nu_n}\!\left(\sum_\nu b_\nu H_\nu^\circ\right)
 \leq2\operatorname{Var}_{\nu_n}(F)
      +2a^2\operatorname{Var}_{\nu_n}(R),
\]
and the preceding \(a^2\)-bound absorbs the last term, giving
\[
 \operatorname{Var}_{\nu_n}\!\left(\sum_\nu b_\nu H_\nu^\circ\right)
 \leq C_W\operatorname{Var}_{\nu_n}(F).
\]
The head-pattern lower bound now yields
\(\operatorname{Var}_{\nu_n}(F)\geq(c_{H,W}/C_W)\sum_\nu b_\nu^2\).
Combining this with the \(a^2\)-bound, one may
choose \(\gamma_W>0\), independently of \(n\) and of every later tilt
box, such that
\[
 \Gamma_{0,n}\geq2\gamma_W I
 \qquad(n\geq n_0(W)).
\]

It remains only to compare each tilted covariance with its own
finite-\(n\) baseline.  If
\(S_\xi=\sum_j\eta_j\Omega_j\) is the medium-prime part of the score,
then \(|S_\xi|\leq B\Omega\),
\(e^{|S_\xi|/L}\ll_{B,W}1\), and the one-mark estimates give
\(\mathbb E_{\nu_n}\Omega\ll_W\log L\).  Consequently
\[
 \sup_{\xi\ \text{ in the fixed }B\text{-box}}
 \mathbb E_{\nu_n}\left|e^{S_\xi/L}-1\right|
 \ll_{B,W}\frac{\log L}{L}=o(1).
\]
The physical and finite-head factors are
\(1+O_{B,W}(L^{-1})\), and normalization changes the same estimate by
only a constant factor.  Thus
\[
 \sup_{\xi\ \text{ in the fixed }B\text{-box}}
 \|\nu_{n,\xi}-\nu_n\|_{\rm TV}=o_{n;B,W}(1).
\]
Every coordinate of \(Z\) is uniformly bounded and its dimension depends
only on \(W\).  If \(d_W=\dim Z\) and
\(\max_\nu\|Z_\nu\|_\infty\leq M_W\), then for every Euclidean unit
vector \(a\),
\[
 \left|\operatorname{Var}_{\nu_{n,\xi}}(a^TZ)
       -\operatorname{Var}_{\nu_n}(a^TZ)\right|
 \leq 6d_WM_W^2\|\nu_{n,\xi}-\nu_n\|_{\rm TV}.
\]
Taking the supremum over \(a\) gives
\[
 \sup_{\xi\ \text{ in the fixed }B\text{-box}}
       \|\Gamma_{\xi,n}-\Gamma_{0,n}\|_{\mathrm{op}}
       =o_{n;B,W}(1).
\]
For each fixed \(B\), increasing only \(n_0(B,W)\) now proves
\eqref{eq:smooth-joint-head-physical-gap}.  Its constant
\(\gamma_W\) was chosen from the uniform finite-\(n\) geometry before
\(B\), so no unproved limiting mixture and no box-dependent main
constant is involved.
The constant score is absent: normalization has removed it, and one
centered head indicator was deleted before \(\Gamma_\xi\) was formed.
The head cross rows are
\(O_{B,W}(1/L)\) throughout the fixed box, and the physical cross row is
given by the Stieltjes
estimate \eqref{eq:marked-tilted-physical-full}.  Write
\[
 U_\xi:=\operatorname{Cov}_\xi(\Omega,Z),\qquad
 C_\xi:=\operatorname{Cov}_\xi(\Omega,\Omega).
\]
The actual block elimination is
\begin{equation}
 q_n^{-1}\mathcal C^Z_\xi
 =C_\xi-U_\xi\Gamma_\xi^{-1}U_\xi^T.
 \label{eq:smooth-displayed-joint-block-elimination}
\end{equation}
The marked head/physical bounds give, for \(\|b\|_\infty\le1\),
\[
 |(U_\xi^Tb)_\nu|
 \leq\sum_j|(U_\xi)_{j\nu}|\,|b_j|
 \ll_{B,W}\frac1L\sum_jH_j
 \ll_{B,W}\frac{\log L}{L}.
\]
Together with \(\|\Gamma_\xi^{-1}\|_{\rm op}\leq\gamma_W^{-1}\) this
gives, row by row,
\begin{equation}
 \left\|D^{-1}U_\xi\Gamma_\xi^{-1}U_\xi^Tb\right\|_\infty
 \ll_{B,W}\frac{\log L}{L^2},
 \label{eq:smooth-joint-Schur-row-error}
\end{equation}
because \(|(U_\xi)_{i\nu}|\ll_{B,W}H_i/L\) and
\(\sum_jH_j\ll\log L\).  This is the promised normalized row estimate,
with an explicit rate.

\medskip\noindent
\emph{Step 5: the exact physical direction and the full quotient
form.}\quad
The physical column is nevertheless responsible for the exact null
relation, and we record it rather than discarding it as a small error.
Let
\[
 T_y:=\sum_{W<p\le y}t_pv_p,\qquad
 S_\alpha:=\sum_j\alpha_j\Omega_j,\qquad
 S_g:=\sum_{W<p\le y}(\alpha_{j(p)}-t_p)v_p.
\]
For a statistic \(X\), put
\(X^{\circ,\xi}:=X-\mathbb E_{\nu_{n,\xi}}X\).  This notation is
distinct from the already defined baseline-centered head scores
\(H_e^\circ\); changing their centering subtracts only a constant and
therefore leaves the centered covariance span unchanged.  Every member
of every active structured cell is \(y\)-smooth, so its logarithm is the
sum of the \(p\leq y\) valuation terms.  Moreover, for each \(p\leq W\),
the function \(v_p\) is constant on every exact head-pattern cell and
hence is a linear combination of the retained head indicators and the
deleted constant indicator.  After centering, the constant disappears
and the whole \(p\leq W\) term lies in the centered head span.  Therefore,
identically at finite \(n\),
\begin{equation}
 T_y^{\circ,\xi}=
 \frac{R^{\circ,\xi}
       -\sum_{p\le W}(\log p)v_p^{\circ,\xi}}{\log y}
 \in\operatorname{span}Z^{\circ,\xi},\qquad
 S_\alpha^{\circ,\xi}=T_y^{\circ,\xi}+S_g^{\circ,\xi}.
 \label{eq:smooth-physical-alpha-block-relation}
\end{equation}
Let \(P_{Z^\perp}\) denote orthogonal projection in
\(L^2(\nu_{n,\xi})\) onto the orthogonal complement of the centered
coordinates of \(Z\).  Thus \(P_{Z^\perp}\) kills
\(T_y^{\circ,\xi}\)
exactly.  More generally, for \(b=q+\lambda\alpha\), with
\(q\in\mathcal G_n\), and for every \(\mu\in\mathbb R\),
\begin{equation}
 P_{Z^\perp}\{(b\cdot\Omega)^{\circ,\xi}\}
 =P_{Z^\perp}\left\{
   \left(\sum_{W<p\leq y}
       (b_{j(p)}-\mu t_p)v_p\right)^{\circ,\xi}\right\}.
 \label{eq:smooth-exact-arithmetic-null-relation}
\end{equation}
This is the precise finite-\(n\) post-Schur relation used below.  We next
record its arithmetic geometry.  Define
\begin{equation}
 V_n:=\sum_{W<p\leq y}\frac{(\alpha_{j(p)}-t_p)^2}{p},\qquad
 \mathfrak d_n(b)^2:=\inf_{\mu\in\mathbb R}
   \sum_{W<p\leq y}\frac{|b_{j(p)}-\mu t_p|^2}{p},
 \qquad A_{\alpha,n}:=\alpha^TD\alpha.
 \label{eq:smooth-arithmetic-physical-distance}
\end{equation}
Since
\(\sum_{p\in\mathcal P_j}(\alpha_j-t_p)/p=0\) in every band, and since
\(q^TD\alpha=0\), the following identities are exact.  For all
sufficiently large \(n\), every permitted band has positive harmonic
mass and \(\alpha_j>0\); hence \(A_{\alpha,n}>0\), while \(V_n\geq0\).
In particular the denominator in the third line below is nonzero:
\begin{align}
 \mathfrak d_n(q+\lambda\alpha)^2
 &=\inf_{\mu\in\mathbb R}\left\{
      \|q+(\lambda-\mu)\alpha\|_D^2+\mu^2V_n\right\}\notag\\
 &=\|q\|_D^2+\inf_{\mu\in\mathbb R}\left\{
      (\lambda-\mu)^2A_{\alpha,n}+\mu^2V_n\right\}\notag\\
 &=\|q\|_D^2+\lambda^2
      \frac{A_{\alpha,n}V_n}{A_{\alpha,n}+V_n}
 \geq\|q\|_D^2.
 \label{eq:smooth-quotient-after-physical-geometry}
\end{align}
We also record here the size of \(V_n\), independently of the later
regression argument.  Weighted Mertens summation and
\eqref{eq:smooth-two-projection-convergence} give
\[
 V_n=\sum_j\int_{I_j}|t-\alpha_j^{(c)}|^2\,\frac{dt}{t}
          +o_n\bigl((\delta+\eta)^2\bigr).
\]
On \(I_0=[t_0,\delta]\), the displayed integral is exactly
\[
 \frac{\delta^2-t_0^2}{2}
 -\frac{(\delta-t_0)^2}{\log(\delta/t_0)}
 \asymp\delta^2.
\]
On a positive cell \(I=[a,a+\ell]\), it is
\(\asymp\ell^3/a\); summing the regular relative mesh gives
\(\asymp\eta^2\).  Consequently, uniformly over permitted meshes,
\begin{equation}
                         V_n\asymp\delta^2+\eta^2
                              \asymp(\delta+\eta)^2.
 \label{eq:smooth-arithmetic-cell-variance-size}
\end{equation}

We now transfer this exact distance to the covariance form; this step is
needed because an arithmetic center is not a continuum center.  If
\(\mathfrak d_n(b)>0\), let \(\mu_*\) be its unique minimizer and put
\begin{equation}
 c_p:=b_{j(p)}-\mu_*t_p.
 \label{eq:smooth-minimizing-prime-coefficients}
\end{equation}
Differentiation of the finite sum gives the exact normal equation
\begin{equation}
                 \sum_{W<p\leq y}\frac{t_pc_p}{p}=0.
 \label{eq:smooth-arithmetic-physical-normal-equation}
\end{equation}
Normalize for the moment so that \(\mathfrak d_n(b)=1\).
Equation \eqref{eq:smooth-quotient-after-physical-geometry}, together
with \(A_{\alpha,n}\asymp1\) and
\(V_n\asymp(\delta+\eta)^2\), gives
\[
 \|q\|_D\leq1,\qquad |\lambda|+|\mu_*|\ll_{\delta,\eta}1.
\]
In particular \(|q_0|\leq H_0^{-1/2}\),
\(\alpha_0\ll_{\delta}H_0^{-1}\), and \(t_0\to0\).

Define the piecewise affine continuum coefficient
\[
 f_n(t):=b_j-\mu_*t\qquad(t\in I_j),
\]
and extend it on \((0,t_0)\) by
\(f_n(t_0)t/t_0\); denote the resulting function by
\(\widetilde f_n\).  The cellwise Mertens argument in
\eqref{eq:smooth-positive-cell-data-convergence}--
\eqref{eq:smooth-low-cell-data-convergence} gives
\begin{equation}
 \|\widetilde f_n\|_{\mathcal H}^2=1+o_n(1),\qquad
 \langle\widetilde f_n,t\rangle_{\mathcal H}=o_n(1).
 \label{eq:smooth-minimizer-continuum-norm-gauge}
\end{equation}
Here the second relation is
\eqref{eq:smooth-arithmetic-physical-normal-equation} after quadrature.
For clarity about the moving low cell, the only term that can retain the
fixed-cutoff Mertens constant is
\(|b_0|^2(H_0-H_0^{(c)})=O_W(|b_0|^2)=o_n(1)\); every term carrying a
factor \(t\) converges by weighted Mertens summation.  Also
\(|f_n(t_0)|=o_n(1)\), so the linear extension has \(o_n(1)\) squared
norm and \(o_n(1)\) \(t\)-pairing.  Thus no convergence \(D\to D_c\) has
been assumed.

For prime coefficients \(c=(c_p)\), write the leading discrete form as
\begin{equation}
 \mathscr Q_n^{(0)}(c):=
 \sum_{W<p\leq y}\frac{h(t_p)c_p^2}{p}
 +\sum_{\substack{W<p,r\leq y\\p\ne r}}
       \frac{K(t_p,t_r)c_pc_r}{pr}.
 \label{eq:smooth-leading-minimizer-form}
\end{equation}
The continuum double integral contains its diagonal as a measure-zero
set, whereas the finite prime double sum in the quadrature temporarily
contains \(p=r\).  The difference is
\begin{equation}
 \Delta_{\rm diag}(c):=
 \sum_{W<p\leq y}\frac{K(t_p,t_p)c_p^2}{p^2}=o_n(1)
 \qquad\bigl(\mathfrak d_n(b)=1\bigr).
 \label{eq:smooth-minimizer-kernel-diagonal}
\end{equation}
To verify this uniformly, on the positive cells \(p\geq y^\delta\), so
\[
 \sum_{t_p>\delta}\frac{|K(t_p,t_p)|c_p^2}{p^2}
 \ll y^{-\delta}\sum_{t_p>\delta}\frac{c_p^2}{p}=o(1).
\]
On \(I_0\), \(c_p=b_0-\mu_*t_p\), where
\(|b_0|=O(H_0^{-1/2})\) and \(\mu_*=O_{\delta,\eta}(1)\).  Since
\[
 \sum_{p>W}\frac{t_p^2}{p^2}
   =O_W((\log y)^{-2}),\qquad
 \sum_{p>W}\frac{t_p^4}{p^2}
   =O_W((\log y)^{-4}),
\]
and \(K(t,t)=O(t^2)\), the low-cell contribution is \(o(1)\) as well.
The two-index quadrature proved in
\eqref{eq:smooth-collective-compact-quadrature}--
\eqref{eq:smooth-discrete-full-form-comparison} applies also to
\(f_n=J_{\mathcal P}b-\mu_*t\): the step part is covered by the displayed
operator comparison, while the \(t\)-part follows from the same
Stieltjes summation, and its coefficient is bounded under the present
normalization.  Together with
\eqref{eq:smooth-minimizer-continuum-norm-gauge}, it gives
\[
 \mathscr Q_n^{(0)}(c)
 =\langle\widetilde f_n,A\widetilde f_n\rangle_{\mathcal H}+o_n(1).
\]
All \(o_n(1)\) terms in these comparisons are uniform on the normalized
class \(\mathfrak d_n(b)=1\): \(q\) is bounded in \(D\)-norm,
\(\lambda,\mu_*=O_{\delta,\eta}(1)\), the positive-cell dimension is
fixed, and the low coordinate tends to zero.  In particular, the
resulting threshold for \(n\) does not depend on \(b\).
The Poincar\'e inequality \eqref{eq:PD-Poincare} now applies: the two
relations in \eqref{eq:smooth-minimizer-continuum-norm-gauge} say that
the squared distance of \(\widetilde f_n\) from
\(\operatorname{span}\{t\}\) is \(1+o_n(1)\).  Hence, by homogeneity,
there is a mesh-uniform \(\kappa_1>0\) such that, for every \(b\) and all
sufficiently large \(n\),
\begin{equation}
                 \mathscr Q_n^{(0)}(c)
                 \geq\kappa_1\mathfrak d_n(b)^2
 \label{eq:smooth-arithmetic-physical-Poincare-transfer}
\end{equation}
when \(c\) is defined by
\eqref{eq:smooth-minimizing-prime-coefficients}.  If
\(\mathfrak d_n(b)=0\), then
\eqref{eq:smooth-quotient-after-physical-geometry} gives \(b=0\), so
the assertion is trivial.

It remains to verify that every arithmetic error is relative to this
same minimizing distance.  If \(E_{pr}=E_{rp}\) is any prime covariance
error with
\(\sup_p\sum_rp|E_{pr}|\leq\varepsilon\), symmetry and
\(2|c_pc_r|\leq c_p^2+c_r^2\) give the exact estimate
\begin{equation}
 \left|\sum_{p,r}E_{pr}c_pc_r\right|
 \leq\varepsilon\sum_p\frac{c_p^2}{p}
 =\varepsilon\mathfrak d_n(b)^2.
 \label{eq:smooth-minimizer-prime-row-error}
\end{equation}
The proof of \Cref{cor:marked-full-tilt-guard-stability}, before band
averaging, gives the corresponding prime-level guard estimate
\begin{equation}
 \sup_{W<p\leq y}\sum_{W<r\leq y}
       p|E^{\rm guard}_{pr}|
 \ll_{B,W}\frac{y^2L^{A+3}}n=o_n(1)
 \label{eq:smooth-guard-prime-row-error}
\end{equation}
for the fixed exponent \(A\) in that corollary.  Explicitly, a deleted
integer satisfies the two score envelopes
\eqref{eq:marked-guard-score-envelopes}.  Taking the signs of the
covariance differences as the coefficient vector \(b\) in
\eqref{eq:marked-guard-explicit-prime-row}, the exact conditional
covariance identity \eqref{eq:marked-conditional-deletion-covariance}
gives \(O_{B,W}(y^2L^{A+2}/n)\), which is stronger than the displayed
bound.  Hence
\eqref{eq:smooth-minimizer-prime-row-error} applies to guarded deletions
as well as to the analytic covariance remainders.
Therefore the squarefree and prime-power remainders in
\eqref{eq:smooth-normalized-prime-row-error} and
\eqref{eq:smooth-power-row-transfer}, including the guarded-deletion
remainder from \Cref{cor:marked-full-tilt-guard-stability}, contribute at
most
\[
 \left\{\frac{C_{\rm sf}+C_{\rm pow}}W
       +o_{n;B,W,\mathcal P}(1)\right\}\mathfrak d_n(b)^2.
\]
The main constants here are independent of the subsequently fixed tilt
box.
Equivalently, if \(\mathscr E_n(c)\) denotes the sum of these actual
covariance errors, then
\begin{equation}
 \operatorname{Var}_\xi\!\left(\sum_pc_pv_p\right)
 =\mathscr Q_n^{(0)}(c)+\mathscr E_n(c),\qquad
 |\mathscr E_n(c)|
 \leq\left\{\frac{C_{\rm sf}+C_{\rm pow}}W
       +o_{n;B,W,\mathcal P}(1)\right\}\mathfrak d_n(b)^2.
 \label{eq:smooth-actual-minimizer-covariance-ledger}
\end{equation}

Finally, the one-mark head estimates and
\eqref{eq:marked-tilted-physical-full} give
\[
 \left\|\operatorname{Cov}_\xi
   \left(Z,\sum_pc_pv_p\right)\right\|
 \ll_{B,W}\frac1L\sum_p\frac{|c_p|}{p}
 \ll_{B,W}\frac{\sqrt{\log L}}L\,\mathfrak d_n(b).
\]
Together with \eqref{eq:smooth-joint-head-physical-gap}, the nuisance
Schur subtraction is consequently
\[
 O_{B,W}\!\left(\frac{\log L}{L^2}\right)
                  \mathfrak d_n(b)^2.
\]
Indeed, the definition
\eqref{eq:smooth-joint-Schur-band-matrix} and the exact score identity
\eqref{eq:smooth-exact-arithmetic-null-relation} give
\[
 \frac1{q_n}b^T\mathcal C_\xi^Zb
 =\left\|P_{Z^\perp}\{(b\cdot\Omega)^{\circ,\xi}\}\right\|_2^2
 =\left\|P_{Z^\perp}
   \left\{\left(\sum_pc_pv_p\right)^{\circ,\xi}\right\}\right\|_2^2.
\]
First choose \(W\) so
that the box-independent \(1/W\) terms are smaller than half the
Poincar\'e gap, and then take \(n\) sufficiently large for the fixed box
and mesh.  Equations
\eqref{eq:smooth-arithmetic-physical-Poincare-transfer} and
\eqref{eq:smooth-quotient-after-physical-geometry} then yield
\begin{equation}
 \frac1{q_n}b^T\mathcal C_\xi^Zb
 \geq\kappa\,\mathfrak d_n(b)^2
 \geq\kappa\|q\|_D^2.
 \label{eq:smooth-full-quotient-from-minimizer}
\end{equation}
Every band vector has a unique decomposition
\(b=q+\lambda\alpha\), \(q\in\mathcal G_n\), and
\(\|q\|_D^2=\inf_\lambda\|b-\lambda\alpha\|_D^2\).  Thus the last display
is exactly \eqref{eq:smooth-arithmetic-Poincare}; no gap on
\(\mathcal G_n\) has been used as a substitute for the full quotient
inequality.

The inverse assertion is a separate projected statement.  The
gauge isomorphism
\eqref{eq:smooth-ordinary-gauge-isomorphism}, the pulled comparison
\eqref{eq:smooth-pulled-ordinary-operator-convergence}, and the two
successive Neumann steps following that display give the projected
\(\ell^\infty\)-inverse for the actual un-Schur band covariance
\(C_\xi\).  Its norm is bounded independently of the later box and
mesh; only the threshold for the vanishing arithmetic remainder depends
on them.  Apply
\Cref{lem:stable-projected-Schur-perturbation} with
\(P=P_{\alpha,n}\) and
\(E=U_\xi\Gamma_\xi^{-1}U_\xi^T\).
Equation \eqref{eq:smooth-joint-Schur-row-error} and the symmetric
band-row calculation preceding
\eqref{eq:smooth-power-row-transfer} give the required
operator and quadratic perturbation bounds
\(O_{B,W}(\log L/L^2)=o_{n;B,W}(1)\).  The lemma therefore preserves the
projected inverse, proving
\eqref{eq:smooth-arithmetic-Linf-inverse}.
This proves
\eqref{eq:smooth-arithmetic-Poincare}--
\eqref{eq:smooth-arithmetic-Linf-inverse} with the constant row removed
and the head/physical block jointly, rather than separately, inverted.
\end{proof}

\subsection{The weighted compensated score}
\label{subsec:smooth-compensated-score}

The quotient estimate alone is not sufficient: the scalar parameter
which remains after quotient fitting is of order \(1/(\delta+\eta)\).
We now prove the weighted estimate which shows that its \emph{effective}
prime coefficients remain bounded.

Put
\begin{equation}
 w:=\delta+\eta,\qquad
 g_p:=\alpha_{j(p)}-t_p,\qquad
 V:=V_n=\sum_{W<p\leq y}\frac{g_p^2}{p}.
\label{eq:smooth-raw-compensated-score}
\end{equation}
The regular lower bound in
\eqref{eq:smooth-regular-relative-mesh} is used in the next lemma.

For a band vector put
\begin{equation}
 \|b\|_\sharp:=\max_j\frac{|b_j|}{\alpha_j}.
 \label{eq:smooth-sharp-norm}
\end{equation}
This is a finite norm for every \(n\), but its low weight
\(\alpha_0\asymp\delta/H_0\) moves with \(n\).
For the continuum cells we use the distinct notation
\begin{equation}
 \|b\|_{\sharp,c}:=\max_j\frac{|b_j|}{\alpha_j^{(c)}}.
 \label{eq:smooth-continuum-sharp-norm}
\end{equation}

\begin{lemma}[Weighted discrete quotient inverse]
\label[lemma]{lem:smooth-weighted-discrete-inverse}
Assume that \(W\geq W_0\) and
\(w=\delta+\eta\leq w_0\) have been chosen as in
\Cref{lem:arithmetic-band-quotient}.
Let
\[
 \mathsf B^Z_\xi:=D^{-1}\mathcal C^Z_\xi/q_n,
 \qquad \mathcal G:=\{b:\alpha^TDb=0\}.
\]
If \(u\in\mathcal G\), \(|u_j|\leq Cw\alpha_j\), and the gauge-fixed
solution \(q\in\mathcal G\) is defined by
\begin{equation}
 P_{\alpha,n}\mathsf B^Z_\xi P_{\alpha,n}q=u,
 \label{eq:smooth-weighted-discrete-system}
\end{equation}
then
\begin{equation}
                         |q_j|\leq C_{\sharp}w\alpha_j
                         \quad(0\leq j\leq s).
 \label{eq:smooth-weighted-discrete-conclusion}
\end{equation}
Here \(C_{\sharp}\) depends on the displayed right-side constant \(C\)
and the continuum operator, but not on \(W\), the compact tilt radius,
or the permitted fine mesh once \(W\geq W_0\).
This is uniform in the moving cutoff, every permitted regular mesh, all
sufficiently large \(n\), and every fixed compact effective-tilt box, in
the precise sense that the constant is independent of the mesh and box
radius, while the threshold for \(n\) may depend on
\((B,W,\mathcal P)\).
\end{lemma}

\begin{proof}
\emph{Step 1: a sharp continuum inverse.}\quad
We first prove the continuum assertion, using only the continuum gauge
\(\mathcal G_c\), centers \(\alpha^{(c)}\), and harmonic matrix \(D_c\).
For a band vector define the
linear reconstruction
\begin{equation}
 (J_{\mathcal P}^\sharp q)(t):=\frac{t}{\alpha_j^{(c)}}q_j
       \quad(t\in I_j).
 \label{eq:smooth-linear-cell-reconstruction}
\end{equation}
Thus
\begin{equation}
 \|J_{\mathcal P}^\sharp q\|_{X_t}=\|q\|_{\sharp,c},
 \label{eq:smooth-sharp-reconstruction-isometry}
\end{equation}
including on the moving low cell.  If a uniform continuum inverse did
not exist, there would be a sequence of moving cutoffs and permitted
meshes, numbers \(\varepsilon_\nu\downarrow0\), and
\(q_\nu\in\mathcal G_{c,\nu}\) such that
\[
 \|q_\nu\|_{\sharp,c}=1,\qquad
 \max_j\frac{|[P_{\alpha,c}B_{t_0,\mathcal P}q_\nu]_j|}
                  {\alpha_j^{(c)}}\leq\varepsilon_\nu.
\]
As in \eqref{eq:smooth-unprojected-compressed-equation}, there is a
bounded scalar \(\beta_\nu\) such that
\begin{equation}
 (B_{t_0,\mathcal P}q_\nu)_j
 =\beta_\nu\alpha_j^{(c)}+r_{\nu,j},\qquad
 |r_{\nu,j}|\leq\varepsilon_\nu\alpha_j^{(c)}.
 \label{eq:smooth-weighted-unprojected-equation}
\end{equation}
Here one may take
\[
 r_\nu=P_{\alpha,c}B_{t_0,\mathcal P}q_\nu,\qquad
 \beta_\nu=
 \frac{(\alpha^{(c)})^TD_cB_{t_0,\mathcal P}q_\nu}
      {(\alpha^{(c)})^TD_c\alpha^{(c)}}.
\]
The denominator is bounded below uniformly by the cells meeting
\([1/2,1]\).  Also
\(\|q_\nu\|_\infty\leq\|q_\nu\|_{\sharp,c}=1\),
\(\sum_jH_j^{(c)}\alpha_j^{(c)}=1-t_0\), and the compressed continuum
operators are uniformly bounded.  Hence \(|\beta_\nu|\leq C\), with \(C\)
independent of the moving cutoff and mesh.
For the local estimates below write \(q=q_\nu\),
\(\beta=\beta_\nu\), and suppress \(\nu\) from the cell data.  Put
\(x_j=q_j/\alpha_j^{(c)}\) and, on \(I_j\), set
\[
 d(t):=(J_{\mathcal P}^\sharp-J_{\mathcal P})q(t)
      =x_j(t-\alpha_j^{(c)}).
\]
Then \(\|x\|_\infty=\|q\|_{\sharp,c}\) and, exactly on every cell,
\begin{equation}
 \int_{I_j}d(t)\,\frac{dt}{t}=0.
 \label{eq:smooth-sharp-cell-centering}
\end{equation}
For the multiplier part the constant value
\(h(\alpha_j^{(c)})\) therefore cancels.  On a positive relative cell,
Lipschitz continuity of \(h\), the mesh bounds, and
\(\alpha_j^{(c)}\asymp\inf I_j\) give
\[
 \frac1{H_j^{(c)}}\left|
 \int_{I_j}h(t)d(t)\,\frac{dt}{t}\right|
 \leq C\eta\alpha_j^{(c)}\|q\|_{\sharp,c}.
\]
On \(I_0\), use \(h(t)=1+O(t)\) and the two explicit estimates
\[
 \frac1{H_0^{(c)}}\int_{I_0}(t-\alpha_0^{(c)})\,\frac{dt}{t}=0,
 \qquad
 \frac1{H_0^{(c)}}\int_{I_0}|t-\alpha_0^{(c)}|\,\frac{dt}{t}
 \leq2\alpha_0^{(c)}.
\]
They give an \(O(\delta\alpha_0^{(c)})\|q\|_{\sharp,c}\) low multiplier
error.  For the integral part, \(|K(s,t)|\leq Cst\) gives
\[
 \left|\frac1{H_i^{(c)}}\int_{I_i}\int_{t_0}^1
       K(s,t)d(t)\,\frac{dt}{t}\frac{ds}{s}\right|
 \leq C\alpha_i^{(c)}\|q\|_{\sharp,c}
       \sum_k\int_{I_k}|t-\alpha_k^{(c)}|\,dt.
\]
The low summand is \(O(\delta^2)\).  On the regular positive cells the
sum is \(O(\eta)\): each summand is
\(O(\eta^2(\inf I_k)^2)\), and comparison with the corresponding
geometric Riemann sum gives \(O(\eta)\).  Hence, uniformly over all
permitted meshes,
\begin{equation}
 \max_j\frac{\left|
   [E_{\mathcal P}A^{(t_0)}(J_{\mathcal P}^\sharp-J_{\mathcal P})q]_j
                 \right|}{\alpha_j^{(c)}}
 \le Cw\|q\|_{\sharp,c}.
 \label{eq:smooth-step-linear-reconstruction-error}
\end{equation}
If
\(G(s)=A^{(t_0)}J_{\mathcal P}^\sharp q(s)/s\), then
\eqref{eq:smooth-weighted-kernel-axis}, \(h\in C^1\), and the fact that
\(J_{\mathcal P}^\sharp q(t)/t=x_j\) is constant on each cell give
\[
 \operatorname{osc}_{I_0}G\leq C\delta\|q\|_{\sharp,c},\qquad
\operatorname{osc}_{I_j}G\leq C\eta\|q\|_{\sharp,c}\quad(j\geq1).
\]
Moreover, since
\(H_j^{(c)}\alpha_j^{(c)}=|I_j|\), for every cell
\begin{equation}
 (E_{\mathcal P}A^{(t_0)}J_{\mathcal P}^\sharp q)_j
 =\alpha_j^{(c)}\frac1{|I_j|}\int_{I_j}G(s)\,ds.
 \label{eq:smooth-sharp-average-to-pointwise}
\end{equation}
In particular, on the low cell,
\[
 (E_{\mathcal P}A^{(t_0)}J_{\mathcal P}^\sharp q)_0
 =\alpha_0^{(c)}\frac1{\delta-t_0}
       \int_{t_0}^{\delta}G(s)\,ds.
\]
Consequently \eqref{eq:smooth-weighted-unprojected-equation} implies
\begin{equation}
 \left\|\frac{A^{(t_0)}J_{\mathcal P}^\sharp q}{t}-\beta\right\|_\infty
 \ll w+o(1).
 \label{eq:smooth-weighted-reconstructed-equation}
\end{equation}
Extend the low linear formula to \((0,t_0)\), and denote the resulting
function by \(\widetilde f\).  Since
\(|K(s,t)/s|\ll t\), the omitted input changes the weighted output by
\(O(t_0^2)\).  For the extended function, its quotient by \(t\) is
constant on \((0,\delta]\).  The formulas for \(h\) and \(K(s,t)/s\),
with \(F\in C^2[0,2]\), therefore give
\[
 \operatorname{osc}_{0\leq s\leq\delta}
 \frac{A\widetilde f(s)}s=O(\delta).
\]
At \(s=0\), the quotient kernel is interpreted by
\eqref{eq:smooth-weighted-kernel-axis}.  Hence
\eqref{eq:smooth-weighted-reconstructed-equation} holds on all of
\((0,1]\) for the full operator, with the same \(O(w)+o(1)\) right
side.  Put \(e(t)=A\widetilde f(t)/t-\beta\).  Self-adjointness and
\(At=0\) give the exact identity
\[
 0=\langle t,A\widetilde f\rangle_{\mathcal H}
  =\frac{\beta}{2}+\int_0^1t\,e(t)\,dt.
\]
Hence \(|\beta|\leq\|e\|_\infty=O(w)+o(1)\).  The \(X_t\) inverse
\eqref{eq:PD-weighted-inverse} then puts the reconstruction within
\(O(w)+o(1)\) of \(\lambda t\).  Since
\((\alpha^{(c)})^TD_cq=0\), this means exactly
\[
 0=\sum_jH_j^{(c)}(\alpha_j^{(c)})^2
       \frac{q_j}{\alpha_j^{(c)}}.
\]
The total weight in this sum is bounded below by a structural positive
constant: the cells contained in \([1/2,1]\) have centers at least
\(1/2\) and total harmonic mass bounded below.  Averaging
\(q_j/\alpha_j^{(c)}=\lambda+O(w)+o(1)\) against these positive weights
therefore gives \(\lambda=O(w)+o(1)\).  Once the structural threshold
\(w_0\) is sufficiently small, this contradicts
\(\|q\|_{\sharp,c}=1\).  Equivalently, with
\begin{equation}
 \Pi_c:=S_c^{-1}P_{\alpha,c}S_c,\qquad
 \mathcal R_c:=\operatorname{Ran}\Pi_c,\qquad
 \mathcal T_c:=\Pi_cS_c^{-1}B_{t_0,\mathcal P}S_c\Pi_c,
 \label{eq:smooth-continuum-sharp-operator}
\end{equation}
there is a mesh- and \(t_0\)-uniform constant \(C_\sharp^{(0)}\) such
that
\begin{equation}
 \left\|(\mathcal T_c|_{\mathcal R_c})^{-1}\right\|_{\infty\to\infty}
 \le C_\sharp^{(0)}.
 \label{eq:smooth-continuum-sharp-inverse}
\end{equation}

\medskip\noindent
\emph{Step 2: arithmetic cell data and sharp \(h,K\) quadrature.}\quad
The arithmetic transfer must be relative.  If
\(\|b\|_\sharp\leq1\), then \(|b_{j(q)}|\leq\alpha_{j(q)}\ll t_q\) on
positive cells, while on the low cell
\(|b_0|\leq\alpha_0\).  The pointwise form of
\eqref{eq:marked-tilted-squarefree} and
\eqref{eq:marked-one-two-covariance} is, for \(p\ne q\),
\begin{equation}
 p|E_{pq}^{\rm sf}|\ll_{B,W}\frac1{qL}.
 \label{eq:smooth-sharp-pointwise-squarefree-error}
\end{equation}
The diagonal has the same \(L^{-1}\) error, in addition to the
Bernoulli correction treated below.  Therefore the exact band average
\eqref{eq:smooth-band-normalized-row-transfer} gives
\begin{equation}
 \begin{split}
 |(\mathsf E_nE^{\rm sf}J_nb)_0|
   &\ll_{B,W} L^{-1}
       \ll_{B,W}\alpha_0\frac{\log L}{\delta L},\\
 |(\mathsf E_nE^{\rm sf}J_nb)_i|
   &\ll_{B,W}\alpha_i\frac1{\delta L}\quad(i\ge1),\\
 \max_i\frac{|(\mathsf E_nE^{\rm sf}J_nb)_i|}{\alpha_i}
   &\ll_{B,W}\frac{\log L}{\delta L}.
 \end{split}
 \label{eq:smooth-sharp-squarefree-transfer}
\end{equation}
For the low row, the first inequality is the explicit calculation
\[
 \frac1{H_0}\sum_{p\in\mathcal P_0}\frac1p
 \sum_q O_{B,W}\!\left(\frac{|b_{j(q)}|}{qL}\right)
 \ll_{B,W}\frac1L,
\]
because
\(\sum_q|b_{j(q)}|/q\ll
H_0\alpha_0+\sum_{j\ge1}H_j\alpha_j\ll1\).
Thus the numerator is genuinely \(o(\alpha_0)\), since
\(\alpha_0\asymp\delta/\log L\).
The active cells have the numerical guards deleted.  By
\eqref{eq:marked-guard-relative-low-error}, that deletion contributes an
operator \(G_n\) satisfying, for \(\|b\|_\sharp\leq1\),
\begin{equation}
 \max_i\frac{|(G_nb)_i|}{\alpha_i}=o_{n;B,W,\mathcal P}(1),
 \label{eq:smooth-sharp-guard-deletion}
\end{equation}
because its band output is \(o_n(\alpha_0)\),
\(\|b\|_\infty\leq1\), and \(\alpha_i\gg\alpha_0\) on every positive
cell.  Thus guard deletion is part of the vanishing sharp remainder.

We next state the sharp quadrature which replaces an unspecified
additive discretization error.  Let \(\widehat B_n^{h,K}\) be the band
operator with no marked error, defined by
\begin{equation}
 (\widehat B_n^{h,K}b)_i:=\frac1{H_i}
 \sum_{p\in\mathcal P_i}\frac1p\left\{
     h(t_p)b_i+\sum_{\substack{W<q\le y\\q\ne p}}
          \frac{K(t_p,t_q)}q b_{j(q)}\right\}.
 \label{eq:smooth-discrete-hK-operator}
\end{equation}
Then, for every fixed permitted mesh,
\begin{equation}
 \left\|S_n^{-1}\widehat B_n^{h,K}S_n
       -S_c^{-1}B_{t_0,\mathcal P}S_c
 \right\|_{\infty\to\infty}
       =o_{n;W,\mathcal P}(1).
 \label{eq:smooth-sharp-main-quadrature}
\end{equation}
Here is the low-cell calculation needed for this relative statement.
Since \(F\in C^2[0,2]\) and \(F(0)=1\), the function
\[
 \varkappa(s,t):=\frac{K(s,t)}{st}
\]
has a continuous extension to \([0,1]^2\).  The \((i,k)\)-entry of the
conjugated discrete kernel, with the diagonal temporarily inserted, is
\begin{equation}
 \frac{\alpha_k}{H_i\alpha_i}
 \sum_{p\in\mathcal P_i}\frac{t_p}{p}
 \sum_{q\in\mathcal P_k}\frac{t_q}{q}
       \varkappa(t_p,t_q).
 \label{eq:smooth-sharp-kernel-entry}
\end{equation}
To see every normalization, introduce the probability measures
\begin{align*}
 \rho_{j,n}
 &:=\frac1{H_j\alpha_j}
       \sum_{p\in\mathcal P_j}\frac{t_p}{p}\,\delta_{t_p},&
 \rho_{j,c}
 &:=\frac1{H_j^{(c)}\alpha_j^{(c)}}
       \1_{I_j}(t)\,dt,
\end{align*}
and retain
\(\omega_{n,j}=H_j\alpha_j^2\),
\(\omega_{c,j}=H_j^{(c)}(\alpha_j^{(c)})^2\).
Because \(H_j\alpha_j=\sum_{p\in\mathcal P_j}t_p/p\) and
\(H_j^{(c)}\alpha_j^{(c)}=|I_j|\), these are indeed probability
measures.  Formula \eqref{eq:smooth-sharp-kernel-entry} is exactly
\begin{equation}
 \omega_{n,k}\iint\varkappa(s,t)\,
       d\rho_{i,n}(s)d\rho_{k,n}(t),
 \label{eq:smooth-sharp-kernel-probability-entry}
\end{equation}
and its continuum counterpart is the same expression with \(n\)
replaced by \(c\).

The moving-low convergence follows from the uniform weighted prime
count
\begin{equation}
 \sup_{t_0\leq x\leq1}\left|
 \sum_{W<p\leq y^x}\frac{t_p}{p}-(x-t_0)\right|
 =O_W((\log y)^{-1}).
 \label{eq:smooth-weighted-prime-distribution}
\end{equation}
This is
\(\sum_{p\leq X}(\log p)/p=\log X+O(1)\), divided by
\(\log y\).  Stieltjes summation and uniform continuity of
\(\varkappa\) therefore give
\(\rho_{j,n}\Rightarrow\rho_{j,c}\), uniformly against all sections of
\(\varkappa\), on every positive cell and also on
\(I_0=[t_0,\delta]\).  On the low cell,
\[
 \omega_{n,0}=(H_0\alpha_0)\alpha_0=o_n(1),\qquad
 \omega_{c,0}=(\delta-t_0)\alpha_0^{(c)}=o_n(1),
\]
whereas \(\omega_{n,k}-\omega_{c,k}=o_n(1)\) on every positive cell.
Thus \eqref{eq:smooth-sharp-kernel-probability-entry} converges in all
four low/positive orientations: a low input carries the vanishing
factor \(\omega_{0,\bullet}\), and a low output is averaged against a
probability measure.  In particular neither \(H_0\) nor
\(\alpha_0^{-1}\) is lost.
Removing the inserted diagonal changes the conjugated operator by at
most
\[
 \max_i\frac1{H_i}\sum_{p\in\mathcal P_i}
       \frac{|K(t_p,t_p)|}{p^2}=o_{n;W,\mathcal P}(1).
\]
Indeed \(|K(t,t)|\ll t^2\).  On a positive cell its primes tend to
infinity, while on the low cell
\(t_p^2=(\log p)^2/(\log y)^2\) and
\(\sum_{p>W}(\log p)^2/p^2<\infty\).
For the diagonal multiplier on \(I_0\), write
\(h(t)=1+t\widetilde h(t)\), where \(\widetilde h\) is continuous.
The constant part has average exactly one in both the arithmetic and
continuum rows; \eqref{eq:smooth-weighted-prime-distribution} handles
the remaining weighted average, together with
\(H_0/H_0^{(c)}\to1\).  Ordinary Mertens quadrature applies on every
positive cell.  The dimension is fixed after the mesh is chosen, so the
entrywise convergence is exactly
\eqref{eq:smooth-sharp-main-quadrature}.  Together with
\eqref{eq:smooth-sharp-projection-convergence}, this also verifies the
sharp convergence of the arithmetic and continuum gauges; no
convergence of \(D\) to \(D_c\) is used.

\medskip\noindent
\emph{Step 3: the remaining sharp arithmetic errors and the nuisance
Schur block.}\quad
The Bernoulli diagonal correction, kept out of
\(E^{\rm sf}_{pq}\) in the preceding off-diagonal calculation, satisfies
\begin{equation}
 |E_i^{\rm sf,diag}|
 \ll\frac1{H_i}\sum_{p\in\mathcal P_i}
       \frac{|b_{j(p)}|}{p^2}
 \le\frac{\alpha_i}{W}.
 \label{eq:smooth-sharp-squarefree-diagonal}
\end{equation}
For the prime-power part we use the product-weighted aggregate bounds,
rather than an unavailable unweighted comparison of
\(H_0^{-1}\sum_{p\in\mathcal P_0}p^{-2}\) with \(\alpha_0\).
For \(\|b\|_\sharp\le1\), harmonic centering and the regular mesh give
\[
 \sum_q\frac{t_q|b_{j(q)}|}{q}\ll1,\qquad
 \sum_q\frac{t_q|b_{j(q)}|}{q^2}\ll\frac1W.
\]
Indeed, on a positive cell \(\alpha_j\asymp t_q\), while on the low
cell
\[
 |b_0|\sum_{q\in\mathcal P_0}\frac{t_q}{q}
 \le H_0\alpha_0^2\ll1.
\]
The unweighted vanishing remainders in the three aggregate estimates are
also harmless at the moving low cell.  More precisely, for the fixed
regular mesh,
\begin{equation}
 \sum_q\frac{|b_{j(q)}|}{q^2}\ll_W\alpha_0+o(\alpha_0),
 \qquad
 \frac1{H_i\alpha_i}\sum_{p\in\mathcal P_i}\frac1{p^2}
 =O_{W,\delta,\eta}(1),
 \label{eq:smooth-sharp-unweighted-power-bands}
\end{equation}
uniformly in \(i\).  For the first estimate, the low-cell contribution
is \(|b_0|\sum_{p>W}p^{-2}\ll_W\alpha_0\), while every positive cell
starts at a power of \(y\) and is \(o(\alpha_0)\); the second estimate is
the same reciprocal calculation after division by the cell's harmonic
mass.  Also \(\sum_q|b_{j(q)}|/q\ll1\), as in the squarefree calculation
above.  Thus the unweighted error coefficient in
\eqref{eq:marked-pair-error-explicit-rate} remains \(o(\alpha_i)\) in
each orientation.
Denote by \(E_i^{JI},E_i^{IJ},E_i^{JJ}\), and
\(E_i^{\rm diag}\) the four prime-power contributions to the normalized
band output.  We write every orientation rather than invoke symmetry.
The box-independent main terms satisfy
\begin{align}
 |E_i^{JI}|_{\rm main}
 &\ll \frac1{H_i}\sum_{p\in\mathcal P_i}\frac{t_p}{p^2}
       \sum_q\frac{t_q|b_{j(q)}|}{q}
 \leq\frac{\alpha_i}{W}
       \sum_q\frac{t_q|b_{j(q)}|}{q}
 \ll\frac{\alpha_i}{W},\notag\\
 |E_i^{IJ}|_{\rm main}
 &\ll \frac1{H_i}\sum_{p\in\mathcal P_i}\frac{t_p}{p}
       \sum_q\frac{t_q|b_{j(q)}|}{q^2}
 \ll\frac{\alpha_i}{W},\notag\\
 |E_i^{JJ}|_{\rm main}
 &\ll \frac1{H_i}\sum_{p\in\mathcal P_i}\frac{t_p}{p^2}
       \sum_q\frac{t_q|b_{j(q)}|}{q^2}
 \ll\frac{\alpha_i}{W^2},\notag\\
 |E_i^{\rm diag}|_{\rm main}
 &\ll \frac{|b_i|}{H_i}\sum_{p\in\mathcal P_i}\frac1{p^2}
 \leq\frac{\alpha_i}{W}.
 \label{eq:smooth-sharp-four-power-main-directions}
\end{align}
Let \(\epsilon^{(0)}_{B,W}(n)\) be the single pair remainder in
\eqref{eq:marked-pair-error-explicit-rate}.  The four unweighted
remainders, divided by \(\alpha_i\), are bounded respectively by
\begin{align}
 &\frac{\epsilon^{(0)}_{B,W}(n)}{H_i\alpha_i}
       \sum_{p\in\mathcal P_i}\frac1{p^2}
       \sum_q\frac{|b_{j(q)}|}{q},\notag\\
 &\frac{\epsilon^{(0)}_{B,W}(n)}{\alpha_i}
       \sum_q\frac{|b_{j(q)}|}{q^2},\notag\\
 &\frac{\epsilon^{(0)}_{B,W}(n)}{H_i\alpha_i}
       \sum_{p\in\mathcal P_i}\frac1{p^2}
       \sum_q\frac{|b_{j(q)}|}{q^2},\notag\\
 &\frac{\epsilon^{(0)}_{B,W}(n)|b_i|}{H_i\alpha_i}
       \sum_{p\in\mathcal P_i}\frac1{p^2}.
 \label{eq:smooth-sharp-four-power-error-directions}
\end{align}
Each is \(o_{n;B,W,\mathcal P}(1)\) uniformly in \(i\), by
\eqref{eq:smooth-sharp-unweighted-power-bands},
\(\sum_q|b_{j(q)}|/q\ll1\), and the fact that every positive-cell
\(\alpha_i\) is bounded below by a fixed multiple of \(\delta\), whereas
the low-cell value is \(\alpha_0\).  We also keep the literal endpoints
separate.  For \(\|b\|_\sharp\leq1\), the endpoint ledger
\eqref{eq:marked-four-tail-endpoint-ledger} gives, orientation by
orientation,
\begin{align}
 \frac{|E_i^{JI,\mathrm{end}}|}{\alpha_i}
 &\ll_{B,W}\frac Ln\frac{\#\mathcal P_i}{H_i\alpha_i}
       \sum_q|b_{j(q)}|
 \ll_{B,W,\delta,\eta}\frac{y^2L}{n},\notag\\
 \frac{|E_i^{IJ,\mathrm{end}}|}{\alpha_i}
 &\ll_{B,W}\frac Ln\frac{\#\mathcal P_i}{H_i\alpha_i}
       \sum_q|b_{j(q)}|
 \ll_{B,W,\delta,\eta}\frac{y^2L}{n},\notag\\
 \frac{|E_i^{JJ,\mathrm{end}}|}{\alpha_i}
 &\ll_{B,W}\frac {L^2}n\frac{\#\mathcal P_i}{H_i\alpha_i}
       \sum_q|b_{j(q)}|
 \ll_{B,W,\delta,\eta}\frac{y^2L^2}{n},\notag\\
 \frac{|E_i^{\mathrm{diag,end}}|}{\alpha_i}
 &\ll_{B,W}\frac {L^2}n\frac{\#\mathcal P_i}{H_i}
       \frac{|b_i|}{\alpha_i}
 \ll_{B,W,\delta,\eta}\frac{yL^2}{n}.
 \label{eq:smooth-sharp-four-endpoint-directions}
\end{align}
Indeed \(\#\mathcal P_i\leq y\),
\(\sum_q|b_{j(q)}|\leq y\), and \(|b_i|\leq\alpha_i\).  Moreover,
\(H_i\alpha_i=\sum_{p\in\mathcal P_i}t_p/p\) tends to \(\delta\) on
the low cell and to \(|I_i|\) on each positive cell, while \(H_i\) is
bounded below on every permitted cell.  Hence the four displays imply
the common relative bound
\begin{equation}
 \max_i\frac{|E_i^{\rm endpoint}|}{\alpha_i}
 \ll_{B,W,\mathcal P}\frac{y^2L^2}{n}
       =o_{n;B,W,\mathcal P}(1).
 \label{eq:smooth-sharp-four-power-endpoints}
\end{equation}
Thus all box dependence in the four directions is contained in one
remainder
\(\epsilon^{\sharp}_{B,W,\mathcal P}(n)\to0\) for each fixed mesh, and
each direction separately is relative at the moving low row.
Consequently
\begin{equation}
 |E_i^{\rm pow}|
 \le\alpha_i\left\{\frac{C_{\rm pow}^{\sharp}}W
       +\epsilon^{\sharp}_{B,W,\mathcal P}(n)\right\},
 \label{eq:smooth-sharp-power-transfer}
\end{equation}
where \(C_{\rm pow}^{\sharp}\) is independent of the compact tilt box
and of the subsequently chosen fine mesh.
Finally \eqref{eq:smooth-joint-head-physical-gap}, the marked
head/physical rows, and the exact removal in
\eqref{eq:smooth-physical-alpha-block-relation} give, for
\(\|b\|_\sharp\leq1\),
\[
 |(U_\xi)_{i\nu}|\ll_{B,W}\frac{H_i}{L},\qquad
 \sum_jH_j|b_j|\leq\sum_jH_j\alpha_j\ll1,
 \qquad \|\Gamma_\xi^{-1}\|\ll_W1.
\]
Consequently
\begin{equation}
 |E_i^{\rm head/phys}|
 \ll_{B,W}\frac1{L^2}
 \leq\alpha_i\,O_{B,W}\!\left(\frac{\log L}{\delta L^2}\right)
 =o_{n;B,W}(\alpha_i).
 \label{eq:smooth-sharp-head-physical-transfer}
\end{equation}
For reference, the relative estimates just proved can be read in the
following ledger; the displayed quantity is always the supremum over
\(\|b\|_\sharp\leq1\):
\[
\begin{array}{c|c}
 \text{source}&
 \displaystyle\max_i |(E_nb)_i|/\alpha_i\\ \hline
 \text{leading \(h,K\) quadrature}
   &o_{n;W,\mathcal P}(1)\\
 \text{squarefree marked off-diagonal}
   &O_{B,W}(\log L/(\delta L))\\
 \text{guard deletion}
   &o_{n;B,W,\mathcal P}(1)\\
 \text{squarefree Bernoulli diagonal}
   &O(1/W)\\
 \text{prime powers}
   &C_{\rm pow}^{\sharp}/W+
       \epsilon^\sharp_{B,W,\mathcal P}(n)\\
 \text{head/physical Schur block}
   &O_{B,W}(\log L/(\delta L^2))\\
 \text{arithmetic/continuum projection}
   &o_{n;W,\mathcal P}(1).
\end{array}
\]
Thus every vanishing row error is \(o(\alpha_i)\), including the low
row \(i=0\); no additive \(o(1)\) is divided by \(\alpha_0\) after the
fact.  It follows that the sharp discrete operator is a perturbation, in
the sharp norm,
whose nonvanishing part is
\((C_{\rm sf}+C_{\rm pow}^{\sharp})/W\), with every dependence on the
later compact box and fixed mesh in a remainder tending to zero with
\(n\).  Equivalently, before either gauge projection is applied,
\begin{equation}
 \left\|S_n^{-1}\mathsf B^Z_\xi S_n
       -S_c^{-1}B_{t_0,\mathcal P}S_c\right\|_{\infty\to\infty}
 \leq\frac{C_{\rm sf}+C_{\rm pow}^{\sharp}}W
       +o_{n;B,W,\mathcal P}(1).
 \label{eq:smooth-unprojected-sharp-operator-comparison}
\end{equation}

\medskip\noindent
\emph{Step 4: identification of the varying sharp gauges and transfer
of the inverse.}\quad
We finish by identifying the two quotient spaces explicitly.  Put
\begin{align}
 \Pi_n&:=S_n^{-1}P_{\alpha,n}S_n,
 &\mathcal R_n&:=\operatorname{Ran}\Pi_n,\notag\\
 \mathcal T_{n,\xi}
 &:=\Pi_nS_n^{-1}\mathsf B^Z_\xi S_n\Pi_n,
 &\mathcal T_c&:=\Pi_cS_c^{-1}B_{t_0,\mathcal P}S_c\Pi_c.
 \label{eq:smooth-two-sharp-projected-operators}
\end{align}
Equations \eqref{eq:smooth-sharp-projection-convergence} and
\eqref{eq:smooth-unprojected-sharp-operator-comparison}, together with
the uniform boundedness of the continuum operator, give the explicit
projected comparison
\begin{equation}
 \|\mathcal T_{n,\xi}-\mathcal T_c\|_{\infty\to\infty}
 \le \frac{C_{\rm sf}+C_{\rm pow}^{\sharp}}W
       +o_{n;B,W,\mathcal P}(1).
 \label{eq:smooth-two-sharp-operator-comparison}
\end{equation}
All operators in this display are extended to the full coordinate
space by the displayed projections.  Put
\[
 \varepsilon_n:=\|\Pi_n-\Pi_c\|_{\infty\to\infty}=o_n(1).
\]
The exact formulas preceding
\eqref{eq:smooth-sharp-projection-convergence} show that both projections
have kernel \(\operatorname{span}\{\mathbf1\}\) and hence rank \(s\).
For \(x\in\mathcal R_c\), the restriction
\[
 U_n:=\Pi_n|_{\mathcal R_c}:\mathcal R_c\longrightarrow\mathcal R_n
\]
satisfies
\[
 \|U_nx-x\|_\infty
 =\|(\Pi_n-\Pi_c)x\|_\infty
 \leq\varepsilon_n\|x\|_\infty.
\]
Thus \(U_n\) is injective when \(\varepsilon_n<1\), and equality of the
two ranks makes it onto.  Moreover,
\begin{equation}
 \|U_n\|\leq1+\varepsilon_n,\qquad
 \|U_n^{-1}\|\leq(1-\varepsilon_n)^{-1}.
 \label{eq:smooth-sharp-range-isomorphism}
\end{equation}
If \(y=U_nx\in\mathcal R_n\), then, under the ambient coordinate
identification,
\[
 \|U_n^{-1}y-y\|_\infty=\|x-\Pi_nx\|_\infty
 \leq\frac{\varepsilon_n}{1-\varepsilon_n}\|y\|_\infty.
\]
For \(x\in\mathcal R_c\), the exact expansion inside \(U_n^{-1}\) is
\[
 \mathcal T_{n,\xi}U_nx-U_n\mathcal T_cx
 =(\mathcal T_{n,\xi}-\mathcal T_c)U_nx
   +\mathcal T_c(U_n-I)x
   +(I-\Pi_n)\mathcal T_cx.
\]
The last two terms are each at most
\(\varepsilon_n\|\mathcal T_c\|\,\|x\|_\infty\), up to the harmless
factor \(1+\varepsilon_n\) in the first one.  Equations
\eqref{eq:smooth-sharp-range-isomorphism} and
\eqref{eq:smooth-two-sharp-operator-comparison} therefore give
the pulled-back comparison
\begin{equation}
 \left\|U_n^{-1}(\mathcal T_{n,\xi}|_{\mathcal R_n})U_n
       -\mathcal T_c|_{\mathcal R_c}\right\|_{\infty\to\infty}
 \leq\frac{C_{\rm sf}+C_{\rm pow}^{\sharp}}W
       +o_{n;B,W,\mathcal P}(1).
 \label{eq:smooth-pulled-sharp-operator-comparison}
\end{equation}
We enlarge the structural threshold \(W_0\) once and for all, before the
actual \(W\) is selected, so that the
box-independent \(1/W\) term is smaller than
\((4C_\sharp^{(0)})^{-1}\) for every \(W\geq W_0\).  Only after this
choice are the compact box and mesh fixed and \(n\) taken large enough
to absorb the remainder.  The Neumann
series and \eqref{eq:smooth-continuum-sharp-inverse} now give
\[
 \left\|(\mathcal T_{n,\xi}|_{\mathcal R_n})^{-1}\right\|_{\infty\to\infty}
 \le 2C_\sharp^{(0)}+o(1).
\]
Finally put \(x=S_n^{-1}q\) and \(v=S_n^{-1}u\).  Because
\(q,u\in\mathcal G_n\), the exact projection formula gives
\(x,v\in\mathcal R_n\), and
\eqref{eq:smooth-weighted-discrete-system} becomes
\(\mathcal T_{n,\xi}x=v\).  The hypothesis gives
\(\|v\|_\infty\leq Cw\), while
\(\|x\|_\infty=\|q\|_\sharp\).  Applying the preceding inverse bound
therefore yields \(\|q\|_\sharp\leq C_\sharp w\), which is
\eqref{eq:smooth-weighted-discrete-conclusion}.  This also makes the
order of the choices \(W\), box, mesh, and \(n\) explicit.
\end{proof}

\begin{lemma}[Augmented weighted compensated-score lemma]
\label[lemma]{lem:augmented-bridge}
Uniformly on every fixed compact effective-tilt box, one has
\begin{equation}
 \sum_{W<p\leq y}\frac{|g_p|}{p}\ll w,\qquad
                         V\asymp w^2.
\label{eq:smooth-g-size}
\end{equation}
Let \(Z\) and \(\mathcal C^Z_\xi\) be as in
\eqref{eq:smooth-joint-Schur-band-matrix}.  First Schur-project the raw
score \(S_g=\sum_pg_pv_p\) off \(Z\), and then regress it on the
gauge-fixed quotient band space \(\mathcal G\).  Explicitly, put
\begin{align}
 b_i^Z&:=\operatorname{Cov}_\xi(\Omega_i,S_g)
 -\operatorname{Cov}_\xi(\Omega_i,Z)\Gamma_\xi^{-1}
       \operatorname{Cov}_\xi(Z,S_g),\notag\\
 P_{\alpha,n}D^{-1}\frac{\mathcal C^Z_\xi}{q_n}P_{\alpha,n}q^{\rm reg}
   &=P_{\alpha,n}D^{-1}b^Z,\qquad q^{\rm reg}\in\mathcal G,
 \label{eq:smooth-block-explicit-regression}\\
 a_Z&:=\Gamma_\xi^{-1}
   \operatorname{Cov}_\xi\left(
      Z,S_g-\sum_jq_j^{\rm reg}\Omega_j\right),\notag\\
 C_g&:=S_g-\sum_jq_j^{\rm reg}\Omega_j-a_Z^TZ,\qquad
 c_p:=g_p-q_{j(p)}^{\rm reg}.
 \label{eq:smooth-post-Schur-score}
\end{align}
Thus the Schur operation, quotient regression, and gauge occur in the
displayed order.  The number \(c_p\) is the valuation coefficient before
the explicit finite nuisance regression \(a_Z^TZ\); that nuisance block
is kept separate below.
In the display below, the vertical bar denotes the Schur-complement
variance after orthogonal regression on the named nuisance scores; it is
not conditioning on a probability-zero event.
Then
\begin{align}
 \|c\|_\infty&\le C_{\rm cmp}w,&
 \sum_p\frac{|c_p|}{p}&\le C_{\rm cmp}w,&
 \sum_p\frac{c_p^2}{p}&\le C_{\rm cmp}w^2,
\label{eq:smooth-compensated-coefficients}\\
 \operatorname{Var}_{\nu_{n,\xi}}
      \left(\sum_pc_pv_p(m)\,\middle|\,\mathcal G,Z\right)
     &\asymp_W w^2,
\label{eq:smooth-compensated-Schur-variance}\\
 \left|\operatorname{Cov}_{\nu_{n,\xi}}
      \left(v_p(m),C_g(m)\right)\right|
     &\ll_W\frac wp
       \qquad(W<p\leq y).
\label{eq:smooth-compensated-marked-row}
\end{align}
The contribution of prime powers to
\eqref{eq:smooth-compensated-Schur-variance} is
\begin{equation}
 \frac{C_{\rm rel}(C_{\rm cmp})}{W}w^2
       +\epsilon^{\rm rel}_{B,W,C_{\rm cmp}}(n)w^2,
\label{eq:smooth-compensated-power-relative}
\end{equation}
rather than an additive error independent of \(w\).
Here \(C_{\rm cmp}\) is the continuum regression constant, fixed before
\(W\) and independent of the compact tilt radius; all finite-head and
box dependence is in thresholds or in the displayed vanishing remainder,
which is defined explicitly in
\eqref{eq:marked-relative-remainder-definition}.
\end{lemma}

\begin{proof}
Mertens' theorem and
\eqref{eq:smooth-regular-relative-mesh} give the first assertion with
uniform constants.  In the continuum low cell,
\begin{equation}
 \int_{t_0}^{\delta}(t-\alpha_0^{(c)})^2\frac{dt}{t}
 =\frac{\delta^2-t_0^2}{2}
       -\frac{(\delta-t_0)^2}{\log(\delta/t_0)}
 =\frac{\delta^2}{2}+o(1).
 \label{eq:smooth-low-V-constant}
\end{equation}
For a positive cell \(I=[a,a+\ell]\), direct integration and the
two-sided regular-mesh bounds give
\begin{equation}
 \int_I(t-\alpha_I)^2\frac{dt}{t}
       \asymp\frac{\ell^3}{a},\qquad
 \int_I|t-\alpha_I|\frac{dt}{t}\asymp\frac{\ell^2}{a}.
 \label{eq:smooth-positive-cell-moments}
\end{equation}
The merged terminal cell satisfies the same bounds.  Summing the
geometric mesh and applying Mertens gives
\[
 \sum_p|g_p|/p\ll\delta+\eta,\qquad
 V\asymp\delta^2+\eta^2\asymp w^2.
\]

We next justify the sharp right-side estimate for the \emph{non-step}
coefficient \(g_p=\alpha_{j(p)}-t_p\).  This does not follow merely by
applying the preceding step-input operator estimate, since on the moving
low cell one can have \(|g_p|\gg\alpha_0\).  All prime sums in this
paragraph range over \(W<p\leq y\).  Put
\begin{equation}
 B_1(g):=\sum_p\frac{|g_p|}{p},\qquad
 G_1(g):=\sum_p\frac{t_p|g_p|}{p},\qquad
 B_2(g):=\sum_p\frac{|g_p|}{p^2},\qquad
 G_2(g):=\sum_p\frac{t_p|g_p|}{p^2}.
 \label{eq:smooth-raw-g-moment-ledger-definition}
\end{equation}
The cell moment calculations in
\eqref{eq:smooth-low-V-constant}--
\eqref{eq:smooth-positive-cell-moments}, followed by weighted Mertens
summation, give
\begin{align}
 \|g\|_\infty&\ll w,&
 B_1(g)&\ll w,&G_1(g)&\ll w,\notag\\
 G_2(g)&\leq W^{-1}G_1(g)\ll w/W,
 \label{eq:smooth-raw-g-first-moment-ledger}\\
 \frac1{H_i}\sum_{p\in\mathcal P_i}
       \frac{t_p|g_p|}{p}&\ll w\alpha_i
       \qquad(0\leq i\leq s).
 \label{eq:smooth-raw-g-local-product-moment}
\end{align}
For the two unweighted reciprocal-square quantities needed below one has,
for the fixed regular mesh,
\begin{align}
 B_2(g)&=O_W(\alpha_0)+o_n(\alpha_0),
 \label{eq:smooth-raw-g-unweighted-square-sum}\\
 d_i(g):=\frac1{H_i}\sum_{p\in\mathcal P_i}\frac{|g_p|}{p^2}
       &=o_{n;W,\delta,\eta,\mathcal P}(w\alpha_i)
       \quad\hbox{uniformly in }i.
 \label{eq:smooth-raw-g-local-diagonal}
\end{align}
Here is the low-row calculation, where these assertions are not a
consequence of a step bound.  Since \(|g_p|\leq\alpha_0+t_p\) on
\(\mathcal P_0\),
\begin{equation}
 d_0(g)\leq
 \frac{\alpha_0}{H_0}\sum_{p>W}\frac1{p^2}
 +\frac1{H_0\log y}\sum_{p>W}\frac{\log p}{p^2}
 =o_n(w\alpha_0).
 \label{eq:smooth-raw-g-low-diagonal-calculation}
\end{equation}
Indeed, after division by \(w\alpha_0\), the two terms are respectively
\(O_W((wH_0)^{-1})\) and
\(O_W((w\delta\log y)^{-1})\).  The same estimate before division by
\(H_0\) proves the low-cell contribution to
\eqref{eq:smooth-raw-g-unweighted-square-sum}.  Every positive cell
starts at \(y^\delta\), so its contribution to either display is
exponentially small relative to \(\alpha_0\).  This proves
\eqref{eq:smooth-raw-g-unweighted-square-sum}--
\eqref{eq:smooth-raw-g-local-diagonal} without losing a factor
\(\alpha_0^{-1}\).

We now contract every covariance term against \(g\).  The leading
squarefree band output is
\[
 R_i^{h,K}(g):=\frac1{H_i}\left\{
  \sum_{p\in\mathcal P_i}\frac{h(t_p)g_p}{p}
  +\sum_{p\in\mathcal P_i}\sum_{q\ne p}
       \frac{K(t_p,t_q)g_q}{pq}\right\}.
\]
Exact harmonic centering gives
\(\sum_{p\in\mathcal P_i}g_p/p=0\).  Since
\(h(t)-1=O(t)\), \(|K(s,t)|\ll st\), and
\(H_i^{-1}\sum_{p\in\mathcal P_i}t_p/p=\alpha_i\),
\eqref{eq:smooth-raw-g-first-moment-ledger}--
\eqref{eq:smooth-raw-g-local-product-moment} give
\begin{equation}
 \left|R_i^{h,K}(g)\right|
 \ll \frac1{H_i}\sum_{p\in\mathcal P_i}
          \frac{t_p|g_p|}{p}
    +\alpha_iG_1(g)
 \ll w\alpha_i.
 \label{eq:smooth-raw-g-leading-band-row}
\end{equation}
In particular this includes the low-row estimate that was formerly
implicit.  The squarefree marked error
\eqref{eq:smooth-sharp-pointwise-squarefree-error} contracts to
\begin{equation}
 |R_i^{\rm sf,err}(g)|
 \ll_{B,W}\frac{B_1(g)}L\ll_{B,W}\frac wL,
 \qquad
 \max_i\frac{|R_i^{\rm sf,err}(g)|}{w\alpha_i}
 \ll_{B,W}\frac{\log L}{\delta L}=o(1).
 \label{eq:smooth-raw-g-squarefree-error}
\end{equation}
The Bernoulli diagonal is bounded by \(d_i(g)\), hence is
\(o(w\alpha_i)\) by
\eqref{eq:smooth-raw-g-local-diagonal}.  Guard deletion is also
relative: apply \eqref{eq:marked-guard-relative-low-error} to
\(g/\|g\|_\infty\) and use \(\|g\|_\infty\ll w\).  Its output is
therefore \(o(w\alpha_0)=o(w\alpha_i)\) on every row.

For completeness, decompose \(V_p=I_p+J_p\) and use
\eqref{eq:marked-product-JI}--
\eqref{eq:marked-product-diagonal}.  The box-independent main terms in
the four orientations satisfy the following bounds, where
\(R_i^{JI},R_i^{IJ},R_i^{JJ},R_i^{\rm diag}\) denote their respective
contributions to
\((D^{-1}\operatorname{Cov}_\xi(\Omega,S_g))_i\):
\begin{align}
 |R_i^{JI}(g)|_{\rm main}
 &\ll \frac1{H_i}\sum_{p\in\mathcal P_i}\frac{t_p}{p^2}\,G_1(g)
      \ll\frac{w\alpha_i}{W},\notag\\
 |R_i^{IJ}(g)|_{\rm main}
 &\ll \frac1{H_i}\sum_{p\in\mathcal P_i}\frac{t_p}{p}\,G_2(g)
      \ll\frac{w\alpha_i}{W},\notag\\
 |R_i^{JJ}(g)|_{\rm main}
 &\ll \frac1{H_i}\sum_{p\in\mathcal P_i}\frac{t_p}{p^2}\,G_2(g)
      \ll\frac{w\alpha_i}{W^2},\notag\\
 |R_i^{\rm diag}(g)|_{\rm main}
 &\ll d_i(g)=o_n(w\alpha_i).
 \label{eq:smooth-raw-g-four-power-main-directions}
\end{align}
Let \(\epsilon_0=\epsilon^{(0)}_{B,W}(n)\) be the pair remainder in
\eqref{eq:marked-pair-error-explicit-rate}.  After division by
\(w\alpha_i\), its four contractions are bounded, respectively, by
\begin{equation}
 \epsilon_0\frac1{H_i\alpha_i}
       \sum_{p\in\mathcal P_i}\frac1{p^2}\frac{B_1(g)}w,
 \quad
 \epsilon_0\frac{B_2(g)}{w\alpha_i},
 \quad
 \epsilon_0\frac1{H_i\alpha_i}
       \sum_{p\in\mathcal P_i}\frac1{p^2}\frac{B_2(g)}w,
 \quad
 \epsilon_0\frac{d_i(g)}{w\alpha_i}.
 \label{eq:smooth-raw-g-four-power-error-directions}
\end{equation}
The first expression tends to zero by
\eqref{eq:smooth-sharp-unweighted-power-bands} and
\eqref{eq:smooth-raw-g-first-moment-ledger}; the third and fourth do so
by \eqref{eq:smooth-raw-g-unweighted-square-sum}--
\eqref{eq:smooth-raw-g-local-diagonal}.  For the second expression, the
low row is \(O_W(\epsilon_0/w)\), while on every positive row
\(\alpha_i\gg\delta\); it too tends to zero after the already chosen
positive \(w\), \(W\), box, and mesh have been fixed.  This is threshold
dependence only and creates no box-dependent main constant.  Finally,
the literal endpoint terms obey
\begin{equation}
 \max_i\frac{|R_i^{\rm endpoint}(g)|}{w\alpha_i}
 \ll_{B,W,\delta,\eta,\mathcal P}\frac{y^2L^2}{n}=o(1).
 \label{eq:smooth-raw-g-four-power-endpoints}
\end{equation}
Indeed \(\#\mathcal P_i\leq y\),
\(\sum_q|g_q|\ll wy\), and \(H_i\alpha_i\) tends to
\(\delta\) on the low cell and is bounded below on every fixed positive
cell.  Thus all \(JI,IJ,JJ\), diagonal, tail, endpoint, squarefree, and
guard terms have now been contracted against the actual non-step score.

The finite head and physical Schur subtraction is relative as well.  We
record its raw nuisance vector here so it is not inferred from the later
post-regression calculation.  Before guard deletion,
\eqref{eq:marked-full-tilt-component-valuation-comparison} and centered
component-weight cancellation give
\[
 |\operatorname{Cov}_\xi(H_\nu^\circ,S_g)|
 \ll_{B,W}\frac1L\sum_p\frac{|g_p|}{p}
 \ll_{B,W}\frac wL.
\]
The physical coordinate has the same bound by
\eqref{eq:marked-tilted-physical-full}.  Since
\(\|S_g\|_\infty\ll_WwL\),
\eqref{eq:marked-guard-deleted-mass} and
\eqref{eq:marked-conditional-deletion-covariance} show that guard
deletion changes either covariance by
\(O_{B,W}(wyL^{A+1}/n)=o_{n;B,W}(w/L)\).  Hence
\begin{equation}
 \|\operatorname{Cov}_\xi(Z,S_g)\|
 \ll_{B,W}\frac wL+o_{n;B,W}(w/L).
 \label{eq:smooth-raw-g-nuisance-vector}
\end{equation}
Together with
\(\|(U_\xi)_{i,*}\|\ll_{B,W}H_i/L\),
\eqref{eq:smooth-joint-head-physical-gap}, and
\(\|\Gamma_\xi^{-1}\|\ll_W1\), this gives
\begin{equation}
 \frac1{H_i}\left|
  \bigl(U_\xi\Gamma_\xi^{-1}
       \operatorname{Cov}_\xi(Z,S_g)\bigr)_i\right|
 \ll_{B,W}\frac w{L^2}+o_n(w/L^2)
 =o_{n;B,W}(w\alpha_i).
 \label{eq:smooth-raw-g-nuisance-Schur-row}
\end{equation}

Combining
\eqref{eq:smooth-raw-g-leading-band-row}--
\eqref{eq:smooth-raw-g-nuisance-Schur-row}, we have proved the explicit
sharp ledger
\begin{equation}
 \max_i\frac{|(D^{-1}b^Z)_i|}{w\alpha_i}
 \leq C_{\rm rhs}+\frac{C_{\rm pow}^{\rm raw}}W
       +r_{n;B,W,\mathcal P},
 \qquad r_{n;B,W,\mathcal P}\longrightarrow0.
 \label{eq:smooth-raw-g-complete-sharp-ledger}
\end{equation}
Here \(C_{\rm rhs},C_{\rm pow}^{\rm raw}\) depend only on the structural Dickman
and regular-mesh bounds; in particular they are independent of the later
box radius and of the permitted mesh.  All dependence on the fixed box
and mesh is confined to the vanishing remainder and its \(n\)-threshold.
Moreover, \eqref{eq:smooth-sharp-projection-convergence} shows that
\(S_n^{-1}P_{\alpha,n}S_n\) is uniformly bounded, so projecting the
right side in \eqref{eq:smooth-block-explicit-regression} preserves this
sharp estimate.  Increase only the corresponding \(n\)-threshold until
\(r_{n;B,W,\mathcal P}\leq1\).  The right-side constant passed to the
weighted inverse can therefore be chosen from \(C_{\rm rhs}\), the already fixed
structural \(1/W_0\) bound, and one spare unit, independently of the
actual \(W\), box radius, or permitted mesh.
Applying
\Cref{lem:smooth-weighted-discrete-inverse} gives
\begin{equation}
 |q_0^{\rm reg}|\le C_{\sharp}w\alpha_0,\qquad
 |q_j^{\rm reg}|\le C_{\sharp}w\alpha_j\quad(j\geq1).
\label{eq:smooth-regression-weighted-bound}
\end{equation}
We next prove, rather than assume, the sharper nuisance estimate needed
for the lower Schur bound.  Put
\[
 P_{\rm med}:=S_g-q^{\rm reg}\!\cdot\Omega=\sum_pc_pV_p,
 \qquad c_p=g_p-q_{j(p)}^{\rm reg}.
\]
Equations \eqref{eq:smooth-g-size} and
\eqref{eq:smooth-regression-weighted-bound}, together with
\(\sum_jH_j\alpha_j=O(1)\), give
\begin{equation}
 \|c\|_\infty\ll w,\qquad
 \sum_p\frac{|c_p|}{p}\leq
 \sum_p\frac{|g_p|}{p}
 +\sum_jH_j|q_j^{\rm reg}|\ll w.
 \label{eq:smooth-postband-coefficient-first-moment}
\end{equation}
Moreover, the quadratic coefficient bound used below is the direct
calculation
\begin{equation}
 \sum_p\frac{c_p^2}{p}
 \leq2V+2\sum_jH_j(q_j^{\rm reg})^2
 \leq2V+2C_\sharp^2w^2\sum_jH_j\alpha_j^2
 \ll w^2.
 \label{eq:smooth-postband-coefficient-second-moment}
\end{equation}

First omit the numerical guards and decompose the full tilted mixture
into its finitely many exact component cells \(C_a\), with weights
\(\lambda_a\).  By
\eqref{eq:marked-full-tilt-component-valuation-comparison}, for a fixed
reference component one may write
\[
 \mathbb E_{C_a,\xi}V_p=m_p+e_a(p),
 \qquad |e_a(p)|\ll_{B,W}\frac1{pL},
\]
uniformly in \(a,p\).  Each head indicator \(H_\nu^\circ\) is constant on
\(C_a\); write that value as \(h_{\nu,a}\) and its mixture mean as
\(\bar h_\nu\).  The law of total covariance and
\(\sum_a\lambda_a(h_{\nu,a}-\bar h_\nu)=0\) give the exact cancellation
\begin{align}
 \operatorname{Cov}_\xi(H_\nu^\circ,P_{\rm med})
 &=\sum_a\lambda_a(h_{\nu,a}-\bar h_\nu)
       \sum_pc_p\{m_p+e_a(p)\}\notag\\
 &=\sum_a\lambda_a(h_{\nu,a}-\bar h_\nu)
       \sum_pc_pe_a(p)
 \ll_{B,W}\frac1L\sum_p\frac{|c_p|}{p}
 \ll_{B,W}\frac wL.
 \label{eq:smooth-explicit-head-postband-cancellation}
\end{align}
The physical coordinate satisfies the same bound directly by
\eqref{eq:marked-tilted-physical-full}:
\begin{equation}
 |\operatorname{Cov}_\xi(R,P_{\rm med})|
 \leq\sum_p|c_p|\,|\operatorname{Cov}_\xi(R,V_p)|
 \ll_{B,W}\frac wL.
 \label{eq:smooth-explicit-physical-postband-cancellation}
\end{equation}

It remains to check that numerical guards do not destroy these
cancellations.  Every coordinate of \(Z\) is bounded by a constant
depending only on \(W\), while
\[
 \|P_{\rm med}\|_\infty\leq\|c\|_\infty
       \sum_{W<p\leq y}V_p\ll_WwL.
\]
Equations \eqref{eq:marked-guard-deleted-mass} and
\eqref{eq:marked-conditional-deletion-covariance} therefore show that
guard deletion changes each covariance in
\eqref{eq:smooth-explicit-head-postband-cancellation}--
\eqref{eq:smooth-explicit-physical-postband-cancellation} by
\[
 O_{B,W}\!\left(\frac{wyL^{A+1}}n\right)
 =o_{n;B,W}(w/L).
\]
Since \(\dim Z\) is fixed after \(W\), we have proved the vector estimate
\begin{equation}
 \left\|\operatorname{Cov}_\xi(Z,P_{\rm med})\right\|
 \ll_{B,W}\frac wL+o_{n;B,W}(w/L).
 \label{eq:smooth-explicit-postband-nuisance-vector}
\end{equation}
Together with \(\|\Gamma_\xi^{-1}\|\ll_W1\), this gives
\[
 \|a_Z\|=\|\Gamma_\xi^{-1}\operatorname{Cov}_\xi(Z,P_{\rm med})\|
 \ll_{B,W}\frac wL+o_n(w/L)\leq C_Ww
\]
after increasing only the \(n\)-threshold for the already fixed box.
Thus the nonvanishing comparison constant may still be chosen
independently of the box radius.  Equations \eqref{eq:smooth-g-size} and
\eqref{eq:smooth-regression-weighted-bound} imply all three bounds in
\eqref{eq:smooth-compensated-coefficients}.  The constants used for the
medium-prime vector are the Mertens main-term constants and the continuum
weighted inverse underlying \(C_{\sharp}\).  The preceding relative
display shows that finite nuisance regression changes the medium vector
only by a vanishing relative term; its nonvanishing finite-dimensional
part is recorded separately in \(a_Z\).  Hence \(C_{\rm cmp}\) is fixed
from the continuum regression before \(W\), the tilt-box radius, or the
final mesh is chosen.

For the lower Schur bound we again keep the continuum and arithmetic
objects separate.  Define
\begin{equation}
 V_c:=\int_{t_0}^1
       |\alpha_{\mathcal P}^{(c)}(t)-t|^2\,\frac{dt}{t},
 \qquad \alpha_{\mathcal P}^{(c)}(t):=\alpha_j^{(c)}\quad(t\in I_j).
 \label{eq:smooth-continuum-cell-variance}
\end{equation}
The moment calculations above and weighted Mertens summation give
\begin{equation}
 V=V_c+o_{n;W,\mathcal P}(w^2),\qquad V_c\asymp w^2.
 \label{eq:smooth-two-cell-variances}
\end{equation}
To transfer the actual fitted vector \(q^{\rm reg}\in\mathcal G_n\),
put
\begin{equation}
 x_n:=S_n^{-1}q^{\rm reg},\qquad
 x_c:=\Pi_cx_n,\qquad q^{(c)}:=S_cx_c.
 \label{eq:smooth-regression-continuum-identification}
\end{equation}
Then \(q^{(c)}\in\mathcal G_c\), and
\eqref{eq:smooth-sharp-projection-convergence} together with
\eqref{eq:smooth-regression-weighted-bound} gives
\begin{equation}
 \max_j\left|
   \frac{q_j^{(c)}}{\alpha_j^{(c)}}
       -\frac{q_j^{\rm reg}}{\alpha_j}\right|
       =o_{n;W,\mathcal P}(w).
 \label{eq:smooth-regression-two-gauges-relative}
\end{equation}

Let \(q_{\mathcal P}^{(c)}(t)=q_j^{(c)}\) on \(I_j\).  Harmonic
centering of \(t-\alpha_j^{(c)}\) on each cell gives the exact
continuum identity
\begin{align}
 &\inf_{\lambda\in\mathbb R}
 \int_{t_0}^1\left|\alpha_{\mathcal P}^{(c)}(t)
       -q_{\mathcal P}^{(c)}(t)-\lambda t\right|^2\frac{dt}{t}
 \notag\\
 &\qquad=\inf_{\lambda\in\mathbb R}
 \left\{\|(1-\lambda)\alpha^{(c)}-q^{(c)}\|_{D_c}^2
       +\lambda^2V_c\right\}
 \geq \frac{A_cV_c}{A_c+V_c}\gg V_c,
 \label{eq:smooth-Schur-geometric-lower}
\end{align}
where
\[
 A_c:=(\alpha^{(c)})^TD_c\alpha^{(c)}\asymp1,
 \qquad (\alpha^{(c)})^TD_cq^{(c)}=0.
\]
The second relation is the exact continuum gauge.  Thus no arithmetic
orthogonality is inserted into a continuum integral.

The exact finite-\(n\) identity
\begin{equation}
 \sum_{W<p\leq y}t_pv_p(m)
 =\frac{L+R(m)-\sum_{p\leq W}v_p(m)\log p}{\log y}
\label{eq:smooth-exact-log-identity}
\end{equation}
shows that, after row centering, the physical and head columns subtract
the function \(t\) exactly at finite \(n\).  Hence the distance from
the kernel direction in \eqref{eq:smooth-Schur-geometric-lower} is the
one relevant after the joint nuisance Schur complement.

We spell out the discrete-to-continuum quadratic transfer for the actual
non-step coefficient.  On \(I_j\cap[t_0,1]\), put
\[
 c_{\mathcal P}^{(c)}(t):=\alpha_j^{(c)}-t-q_j^{(c)},
\]
and extend its value at \(t_0\) linearly to \((0,t_0)\), exactly as in
\eqref{eq:smooth-full-reconstructed-equation}; denote the resulting
function on \((0,1]\) by
\(\widetilde c_{\mathcal P}^{(c)}\).  Cellwise Mertens
quadrature, \eqref{eq:smooth-two-cell-variances}, and
\eqref{eq:smooth-regression-two-gauges-relative} first give
\begin{equation}
 \begin{aligned}
 \sum_{W<p\le y}\frac{
   |(\alpha_{j(p)}-t_p-q_{j(p)}^{\rm reg})
       -c_{\mathcal P}^{(c)}(t_p)|}{p}
       &=o_{n;W,\mathcal P}(w),\\
 \sum_{W<p\le y}\frac{
   |(\alpha_{j(p)}-t_p-q_{j(p)}^{\rm reg})
       -c_{\mathcal P}^{(c)}(t_p)|^2}{p}
       &=o_{n;W,\mathcal P}(w^2).
 \end{aligned}
 \label{eq:smooth-actual-continuum-coefficient-comparison}
\end{equation}
Here \(W,\delta,\eta\), and hence \(w>0\), have already been fixed
before \(n\to\infty\).  Thus the relative center convergence in
\eqref{eq:smooth-two-projection-convergence} is legitimately
\(o_n(w)\) at this stage; no estimate uniform as \(w\downarrow0\) is
being asserted.  The first line controls the marked off-diagonal error,
and the second, together with boundedness of the continuum covariance
operator and Cauchy--Schwarz, changes its quadratic form by \(o(w^2)\).
The two-index marked formula then gives, uniformly on each fixed box,
\begin{equation}
 \operatorname{Var}_\xi\!\left(\sum_{W<p\le y}c_pI_p\right)
 =\langle \widetilde c_{\mathcal P}^{(c)},
       A\widetilde c_{\mathcal P}^{(c)}\rangle_{\mathcal H}
   +O_{B,W}(w^2/L)+O(w^2/W)
       +o_{n;B,W,\mathcal P}(w^2).
 \label{eq:smooth-nonstep-quadratic-transfer}
\end{equation}
Here the \(O(w^2/W)\) term is precisely the squarefree diagonal
correction; it is retained rather than hidden in the little-oh term.
At the low boundary,
\[
 |c_{\mathcal P}^{(c)}(t_0)|
 \leq \alpha_0^{(c)}+t_0+|q_0^{(c)}|
 =O(\alpha_0^{(c)}+t_0)=o_n(w).
\]
Consequently the low extension changes both the squared norm and the
quadratic form by \(o_n(w^2)\), by the calculation in
\eqref{eq:smooth-low-extension-form-error}.  Therefore
\eqref{eq:PD-Poincare}, \eqref{eq:smooth-Schur-geometric-lower}, and
\eqref{eq:smooth-two-cell-variances} give a fixed positive multiple of
\(w^2\) before the following relative perturbations are subtracted.
The finite nuisance Schur loss is, by
\eqref{eq:smooth-explicit-postband-nuisance-vector} and
\eqref{eq:smooth-joint-head-physical-gap},
\[
 \operatorname{Cov}_\xi\!\left(Z,
       S_g-q^{\rm reg}\!\cdot\Omega\right)^T\Gamma_\xi^{-1}
 \operatorname{Cov}_\xi\!\left(Z,
       S_g-q^{\rm reg}\!\cdot\Omega\right)
 =O_{B,W}(w^2/L^2)+o_{n;B,W,\mathcal P}(w^2).
\]
Thus the continuum Poincar\'e gap is transferred at the relative
\(w^2\) scale before the prime-power replacement is made.  The upper
bound in \eqref{eq:smooth-compensated-Schur-variance} follows directly
from the unprojected score and \eqref{eq:smooth-g-size}.

At the \(w^2\) scale every transfer is relative.  From the marked
squarefree errors and the coefficient bounds already proved,
\begin{align}
 E_{\rm sqfree}
 &\ll_{B,W}\frac1L\left\{
       \left(\sum_p\frac{|c_p|}{p}\right)^2
       +\sum_p\frac{c_p^2}{p}\right\}
   =O_{B,W}(w^2/L),\notag\\
 E_{\rm sf,diag}
 &\ll\sum_p\frac{c_p^2}{p^2}
   \leq\frac1W\sum_p\frac{c_p^2}{p}
   =O(w^2/W),\notag\\
 E_{\rm disc}
 &\leq \varepsilon_{W,\delta,\eta}(n)
       \left\{\left(\sum_p\frac{|c_p|}{p}\right)^2
       +\sum_p\frac{c_p^2}{p}\right\}
   =o_n(w^2),
 \label{eq:smooth-relative-squarefree-errors}
\end{align}
where quantitative Mertens summation gives
\(\varepsilon_{W,\delta,\eta}(n)\ll
1/\log L+\log L/L+e^{-c\sqrt L}\).  The endpoint terms in the four-mark formula
are included in this \(\varepsilon\); the cofactor margin
\(n/y^4=n^{1/9}\) makes their total \(o_n(w^2)\) for the fixed mesh.
For the nuisance block,
\eqref{eq:smooth-explicit-postband-nuisance-vector}, the gap
\(\Gamma_\xi^{-1}=O_W(1)\), and the exact physical removal in
\eqref{eq:smooth-physical-alpha-block-relation} give
\begin{equation}
                         E_{\rm head/phys}=o_n(w^2).
 \label{eq:smooth-relative-head-physical-error}
\end{equation}
These displays include respectively the off-diagonal squarefree,
squarefree-diagonal, discretization, and head/physical errors; none is
being replaced by a generic additive \(o_n(1)\).

It remains to justify the relative prime-power statement.  Write
\(I_p=1_{p\mid m}\) and \(J_p=v_p-I_p\).
The product-weighted aggregate estimate
\eqref{eq:marked-product-JI}, with its transposed and two-extra-power
analogues, gives, for \(p\ne q\),
\begin{equation}
 |\operatorname{Cov}(J_p,I_q)|
 \le C_{\rm pow}\frac{t_pt_q}{p^2q}
       +\frac{\epsilon_{B,W}(n)}{p^2q}
       +O_{B,W}(L/n),
\label{eq:smooth-JI-bound}
\end{equation}
with product-weighted analogues for the diagonal, transpose, and
\(\operatorname{Cov}(J_p,J_q)\).  Hence
\begin{align}
 &\left|\operatorname{Var}\left(\sum_pc_pv_p\right)
       -\operatorname{Var}\left(\sum_pc_pI_p\right)\right|
 \notag\\
 &\quad\le
 \frac{C_{\rm rel}(C_{\rm cmp})}{W}w^2
       +\epsilon^{\rm rel}_{B,W,C_{\rm cmp}}(n)w^2.
\label{eq:smooth-power-variance-calculation}
\end{align}
This is \eqref{eq:marked-compensated-power-relative}, applied with
\eqref{eq:smooth-compensated-coefficients}.  Together with the
\(O(w^2/W)\) squarefree diagonal in
\eqref{eq:smooth-relative-squarefree-errors} and
\eqref{eq:smooth-relative-head-physical-error}, this proves
\eqref{eq:smooth-compensated-power-relative}.  For clarity, the complete
relative-error ledger is therefore
\begin{equation}
 \begin{array}{c|c}
 \text{contribution}&\text{bound}\ 
 \\ \hline
 E_{\rm sf,off}&O_W(w^2/L)\\
 E_{\rm sf,diag}&O(w^2/W)\\
 E_{\rm pow}&C_{\rm rel}(C_{\rm cmp})w^2/W
                  +\epsilon^{\rm rel}_{B,W,C_{\rm cmp}}(n)w^2\\
 E_{\rm disc}&o_{n;W,\mathcal P}(w^2)\\
 E_{\rm head}&o_{n;B,W,\mathcal P}(w^2)\\
 E_{\rm phys}&o_{n;B,W,\mathcal P}(w^2).
 \end{array}
 \label{eq:smooth-complete-relative-error-ledger}
\end{equation}
The last two entries are listed separately: the head term follows from
the finite-dimensional gap \eqref{eq:smooth-joint-head-physical-gap}
and the marked head rows, while the physical leading term is removed
exactly by \eqref{eq:smooth-physical-alpha-block-relation} and its
remaining Stieltjes error is \(o_n(w^2)\).  The constant
\(C_{\rm cmp}\), and hence \(C_{\rm rel}(C_{\rm cmp})\), was fixed from
the continuum regression before \(W\), the mesh, or the tilt box; in
particular it is independent of all three.  Choose \(W\) using the
box-independent constant \(C_{\rm rel}(C_{\rm cmp})\), and only after
the compact box is known take \(n\) large enough to absorb
\(\epsilon^{\rm rel}_{B,W,C_{\rm cmp}}(n)\); the lower Schur bound is
then preserved.

Finally we verify the marked row without suppressing an orientation,
diagonal, or endpoint term.  Fix \(W<p\leq y\); all covariances in the
following calculation are under \(\nu_{n,\xi}\).  Since
\(C_g=\sum_qc_qV_q-a_Z^TZ\), first consider the prime part.  The
squarefree diagonal is \(O(|c_p|/p)=O(w/p)\).  For \(q\ne p\),
\eqref{eq:smooth-arithmetic-kernel} and
\eqref{eq:smooth-normalized-prime-row-error} give
\begin{align}
 \sum_{q\ne p}|c_q|\,
       |\operatorname{Cov}(I_p,I_q)|
 &\ll \frac{t_p}{p}\sum_q\frac{t_q|c_q|}{q}
   +\frac{\|c\|_\infty}{p}
       \left\{\frac{C_{\rm sf}}W+\epsilon^{\rm sf}_{B,W}(n)\right\}
 \ll_{B,W}\frac wp.
 \label{eq:smooth-compensated-squarefree-marked-row}
\end{align}
Here and below the threshold for the vanishing remainder may depend on
the already fixed box.  The coefficient estimates
\eqref{eq:smooth-compensated-coefficients} also imply
\begin{equation}
 \sum_q|c_q|\ll wy,\qquad
 \sum_q\frac{t_q|c_q|}{q}\ll w,\qquad
 \sum_q\frac{t_q|c_q|}{q^2}\ll\frac wW.
 \label{eq:smooth-compensated-marked-row-coefficient-ledger}
\end{equation}
Using \(V=I+J\), the four estimates
\eqref{eq:marked-product-JI}--\eqref{eq:marked-product-diagonal} now give
the following separate contractions:
\begin{align}
 \sum_{q\ne p}|c_q|\,|\operatorname{Cov}(J_p,I_q)|
 &\ll \frac w{p^2}+\epsilon_{B,W}(n)\frac w{p^2}
       +\frac{wyL}{n},\notag\\
 \sum_{q\ne p}|c_q|\,|\operatorname{Cov}(I_p,J_q)|
 &\ll \frac w{pW}+\epsilon_{B,W}(n)\frac w{pW}
       +\frac{wyL}{n},\notag\\
 \sum_{q\ne p}|c_q|\,|\operatorname{Cov}(J_p,J_q)|
 &\ll \frac w{p^2W}+\epsilon_{B,W}(n)\frac w{p^2W}
       +\frac{wyL^2}{n},\notag\\
 |c_p|\,|\operatorname{Var}(V_p)-\operatorname{Var}(I_p)|
 &\ll \frac w{p^2}+\epsilon_{B,W}(n)\frac w{p^2}
       +\frac{wL^2}{n}.
 \label{eq:smooth-compensated-power-marked-row-ledger}
\end{align}
For the last line, \(J_p>0\) implies \(I_p=1\), whence
\(|\operatorname{Var}(I_p+J_p)-\operatorname{Var}(I_p)|
 \leq3\mathbb E J_p^2\); this is exactly the diagonal estimate in
\eqref{eq:marked-product-diagonal}.  All endpoint terms in
\eqref{eq:smooth-compensated-power-marked-row-ledger} are uniformly
\(o(w/p)\), since
\begin{equation}
 \sup_{p\leq y}\frac{pyL^2}{n}
 \leq\frac{y^2L^2}{n}=n^{-5/9}L^2=o(1).
 \label{eq:smooth-compensated-marked-row-endpoints}
\end{equation}
Thus the full prime part of \(C_g\) contributes \(O_W(w/p)\), uniformly
on the fixed box once its \(n\)-threshold is chosen.  We also spell out
the nuisance marked family used here and later.  Before guard deletion,
\eqref{eq:marked-full-tilt-component-valuation-comparison} and the same
centered component-weight cancellation as in
\eqref{eq:smooth-explicit-head-postband-cancellation} give
\[
 |\operatorname{Cov}(V_p,H_\nu^\circ)|\ll_{B,W}\frac1{pL}
\]
for every head coordinate, while
\eqref{eq:marked-tilted-physical-full} gives this estimate for the
physical coordinate.  Guard deletion changes either covariance by at
most \(O_{B,W}(\delta_\xi L)\), by
\eqref{eq:marked-conditional-deletion-covariance}.  Uniformly for
\(p\leq y\),
\[
 \delta_\xi L\,(pL)\ll_{B,W}\frac{y^2L^{A+2}}n=o(1),
\]
so, in the fixed-dimensional Euclidean norm,
\begin{equation}
 \|\operatorname{Cov}_\xi(V_p,Z)\|\ll_{B,W}\frac1{pL}.
 \label{eq:smooth-explicit-nuisance-marked-family}
\end{equation}
Since \(\|a_Z\|\ll_Ww\), it follows that
\[
 |\operatorname{Cov}(V_p,a_Z^TZ)|\ll_{B,W}\frac w{pL}
 \ll_W\frac wp
\]
after increasing the threshold for the fixed box.
Adding this nuisance row to
\eqref{eq:smooth-compensated-squarefree-marked-row} and
\eqref{eq:smooth-compensated-power-marked-row-ledger} proves
\eqref{eq:smooth-compensated-marked-row}.
\end{proof}

\subsection{Exact nonlinear fitting}
\label{subsec:smooth-nonlinear-fit}

At the physically and head-fitted point constructed above, all frozen
charges are included in the residual.  Complete rough-signature equations
give \(r_p=0\) for \(p>y\), and the head fit gives \(r_p=0\) for
\(p\leq W\).  The ordinary logarithm was fitted exactly, so
\[
 0=\log Y-\log B_0-\sum_ax_a\log a
   =\sum_pr_p\log p.
\]
Consequently the following compatibility is an exact finite-\(n\)
identity:
\begin{equation}
                         \sum_{W<p\leq y}t_pr_p=0.
\label{eq:smooth-exact-physical-compatibility}
\end{equation}

Let
\[
 \Delta_j:=\sum_{p\in\mathcal P_j}r_p.
\]
After absorbing the marked cost
\eqref{eq:smooth-head-physical-marked-cost} into the constant, the
strict-rate output \eqref{eq:strict-rate-rough-residual} gives
\begin{equation}
 |\Delta_j|\ll_W\frac{NH_j}{L}.
\label{eq:smooth-band-target-envelope}
\end{equation}
The scalar target requires the sharper estimate below, rather than a
generic \(O(N/L)\) bound.
By \eqref{eq:smooth-exact-physical-compatibility},
\begin{align}
 \sum_j\alpha_j\Delta_j
 &=\sum_{W<p\leq y}(\alpha_{j(p)}-t_p)r_p,\notag\\
 \left|\sum_j\alpha_j\Delta_j\right|
 &\ll_W\frac NL\sum_{W<p\leq y}
       \frac{|\alpha_{j(p)}-t_p|}{p}
 \ll_W\frac{Nw}{L}.
\label{eq:smooth-compensated-target}
\end{align}

\begin{proposition}[Physically centered fixed-partition fit]
\label[proposition]{prop:smooth-fixed-partition-fit}
For the regular partition above, the active weights can be changed while
keeping their total mass, every head valuation, and the ordinary
logarithm exact, so that
\begin{equation}
 \sum_{p\in\mathcal P_j}r_p=0
       \qquad(0\leq j\leq s).
\label{eq:smooth-exact-band-balance}
\end{equation}
At the resulting point,
\begin{equation}
 |r_p|\leq C_{1,W}\frac{N}{pL}
       \qquad(W<p\leq y),
\label{eq:smooth-postbridge-fixed-rate}
\end{equation}
where \(C_{1,W}\) is independent of every sufficiently fine regular
partition chosen after \(W\).  The top component and protected layer in
\eqref{eq:smooth-additive-decomposition} remain literally unchanged, all
total coordinate weights lie in \([0,1]\), and the integer equation
\eqref{eq:smooth-flexible-integer-quota} is preserved.
\end{proposition}

\begin{proof}
We define one finite-dimensional map and solve all moments
simultaneously.  For \(u\in\mathcal G\), put
\(\mathscr B(u)=\sum_ju_j\Omega_j\).  Use the fixed statistic
\(S_g=\sum_pg_pv_p\) and the nuisance vector \(Z\) from
\eqref{eq:smooth-joint-Schur-band-matrix}.  A gauge-fixed parameter is
\(\xi=(u,a,\lambda)\in\mathcal G\times\mathbb R^{e+1}\times\mathbb R\),
with exponential score
\[
                         \mathscr B(u)+a^TZ+\lambda S_g.
\]
The gauge coordinates are concrete.  If \(e_0,\ldots,e_s\) are the
standard band vectors, use, for \(1\leq j\leq s\),
\begin{equation}
 q^{(j)}:=e_j-\frac{H_j\alpha_j}{H_0\alpha_0}e_0.
 \label{eq:smooth-concrete-gauge-basis}
\end{equation}
Since \(H_0\alpha_0=\delta+o(1)\), these vectors are well defined and
form a basis of \(\mathcal G=\{u:\alpha^TDu=0\}\).  Primal coordinates
mean coefficients in this basis; dual coordinates mean evaluation on
the same basis.  On the nuisance block we use the ordinary Euclidean
norm.  The intrinsic quotient-dual norm is specified below and is the
one used for all mesh-uniform estimates.
This is merely a nonsingular change of coordinates in
\eqref{eq:smooth-exponential-family}, by
\eqref{eq:smooth-exact-log-identity}.  For a statistic \(F\), put
\[
                         m_F(\xi)=\sum_m z_m(\xi)F(m).
\]
To avoid concealing a coordinate factor, regard the quotient moment as
an element of the finite dual
\(\mathcal G^*\): its value at \(v\in\mathcal G\) is the moment of
\(\mathscr B(v)\).  The normalized moment map is
\begin{equation}
 \mathcal M(\xi):=\frac L{q_n}\left(
 \big[v\mapsto m_{\mathscr B(v)}(\xi)-m_{\mathscr B(v)}(0)\big],
 m_Z(\xi)-m_Z(0),
 m_{S_g}(\xi)-m_{S_g}(0)\right).
 \label{eq:smooth-normalized-moment-map}
\end{equation}
Its prescribed endpoint is
\begin{equation}
\tau:=\frac L{q_n}\left(
 \big[v\mapsto v^T\Delta\big],\ 0,\ \alpha^T\Delta\right).
 \label{eq:smooth-normalized-moment-target}
\end{equation}
To specify the quotient norm without a coordinate convention, put
\[
 d_D:=\frac L{q_n}D^{-1}\Delta.
\]
Then
\[
 \tau_{\mathcal G}(v)=\langle v,d_D\rangle_D,\qquad
 \|d_D\|_\infty\ll_W1,
\]
by \eqref{eq:smooth-band-target-envelope}; the actual gauge
representative is \(P_{\alpha,n}d_D\).  We define
\(\|\tau_{\mathcal G}\|_\infty:=\|P_{\alpha,n}d_D\|_\infty\).
The projection is uniformly bounded by
\eqref{eq:smooth-continuum-gauge}.
This target is exact: if the \(Z\)-moments do not change, then
\eqref{eq:smooth-exact-physical-compatibility} says that changing the
band moments by \(\Delta\) changes the \(S_g\)-moment by
\(\alpha^T\Delta\).  The envelopes
\eqref{eq:smooth-band-target-envelope} and
\eqref{eq:smooth-compensated-target} give
\begin{equation}
 \|\tau_{\mathcal G}\|_\infty+\|\tau_Z\|\ll_W1,\qquad
                         |\tau_g|\ll_Ww.
 \label{eq:smooth-target-coordinate-sizes}
\end{equation}

We display the full inverse used at an arbitrary point of the path.
Let \(\mathscr Q\) denote the joint vector of the quotient band scores
\((\mathscr B(q^{(j)}))_{1\leq j\leq s}\) and \(Z\), in the concrete
coordinates \eqref{eq:smooth-concrete-gauge-basis}, and
write
\begin{equation}
 A_\xi:=\operatorname{Cov}_\xi(\mathscr Q,\mathscr Q),\quad
 b_\xi:=\operatorname{Cov}_\xi(\mathscr Q,S_g),\quad
 d_\xi:=\operatorname{Var}_\xi(S_g),\quad
 \sigma_\xi:=d_\xi-b_\xi^TA_\xi^{-1}b_\xi.
 \label{eq:smooth-full-covariance-blocks}
\end{equation}
In these primal/dual gauge coordinates, the Jacobian of \(\mathcal M\)
is exactly the covariance block
\begin{equation}
 D\mathcal M(\xi)=
 \begin{pmatrix}A_\xi&b_\xi\\ b_\xi^T&d_\xi\end{pmatrix}.
 \label{eq:smooth-full-moment-Jacobian}
\end{equation}
For a right side \((v_{\mathscr Q},v_g)\), its inverse is not a
sequential nonlinear fit but the single block formula
\begin{align}
 \dot\lambda&=\sigma_\xi^{-1}
       \{v_g-b_\xi^TA_\xi^{-1}v_{\mathscr Q}\},\notag\\
 \dot{\mathscr Q}&=A_\xi^{-1}
       \{v_{\mathscr Q}-b_\xi\dot\lambda\}.
 \label{eq:smooth-full-block-inverse}
\end{align}
Thus the quotient and nuisance coordinates are continuously refitted
whenever the slow coordinate moves.

We now identify both terms involving \(A_\xi^{-1}\) with the operators
estimated in the preceding lemmas.  This is an exact finite-dimensional
calculation.  Let \(\mathbf Q:\mathbb R^s\to\mathcal G\) be the matrix
whose columns are the vectors \(q^{(j)}\) in
\eqref{eq:smooth-concrete-gauge-basis}; thus a concrete quotient
coordinate \(x\) represents the actual band vector \(q=\mathbf Qx\).
Set
\begin{align*}
 C_\xi&:=\operatorname{Cov}_\xi(\Omega,\Omega),&
 U_\xi&:=\operatorname{Cov}_\xi(\Omega,Z),&
 \Gamma_\xi&:=\operatorname{Cov}_\xi(Z,Z),\\
 h_\xi&:=\operatorname{Cov}_\xi(\Omega,S_g),&
 z_\xi&:=\operatorname{Cov}_\xi(Z,S_g),&
 \overline C_\xi&:=C_\xi-U_\xi\Gamma_\xi^{-1}U_\xi^T.
\end{align*}
Thus \(\overline C_\xi=q_n^{-1}\mathcal C_\xi^Z\).  Literally in the
concrete primal/dual coordinates of
\eqref{eq:smooth-concrete-gauge-basis},
\begin{equation}
 A_\xi=
 \begin{pmatrix}
  \mathbf Q^TC_\xi\mathbf Q&\mathbf Q^TU_\xi\\
  U_\xi^T\mathbf Q&\Gamma_\xi
 \end{pmatrix},
 \qquad
 b_\xi=\binom{\mathbf Q^Th_\xi}{z_\xi}.
 \label{eq:smooth-concrete-full-blocks}
\end{equation}

Let a joint right side be \((r_{\mathcal G},r_Z)\), and choose any
\(r^\flat\in\mathbb R^{s+1}\) representing its quotient functional:
\(r_{\mathcal G}(v)=v^Tr^\flat\) for \(v\in\mathcal G\).  Directly
eliminating the nuisance variable in
\(A_\xi(x,a)=(r_{\mathcal G},r_Z)\), with \(q=\mathbf Qx\), gives the
equivalent pair of equations
\begin{align}
 a&=\Gamma_\xi^{-1}(r_Z-U_\xi^Tq),\notag\\
 v^T\overline C_\xi q
   &=v^T(r^\flat-U_\xi\Gamma_\xi^{-1}r_Z)
                         \qquad(v\in\mathcal G).
 \label{eq:smooth-exact-joint-block-elimination}
\end{align}
Since the Euclidean annihilator of \(\mathcal G\) is
\(\operatorname{span}\{D\alpha\}\), the second line is, in turn,
equivalent to the intrinsic projected equation
\begin{equation}
 P_{\alpha,n}D^{-1}\overline C_\xi P_{\alpha,n}q
 =P_{\alpha,n}D^{-1}
       (r^\flat-U_\xi\Gamma_\xi^{-1}r_Z),
 \qquad q\in\mathcal G.
 \label{eq:smooth-exact-projected-block-elimination}
\end{equation}
This equivalence uses the actual band values of \(q\); no norm of the
coordinate matrix \(\mathbf Q\) or its inverse is introduced.

Take first \((r^\flat,r_Z)=(h_\xi,z_\xi)\).  Equations
\eqref{eq:smooth-exact-joint-block-elimination}--
\eqref{eq:smooth-exact-projected-block-elimination} are then exactly
\eqref{eq:smooth-block-explicit-regression} and
\eqref{eq:smooth-post-Schur-score}.  If
\(q^{\rm reg}=\mathbf Qx^{\rm reg}\), they prove, without an implicit
coordinate identification,
\begin{equation}
 A_\xi^{-1}b_\xi=(x^{\rm reg},a_Z),\qquad
 \sigma_\xi=\operatorname{Var}_\xi
   (S_g-q^{\rm reg}\!\cdot\Omega-a_Z^TZ)
 =\operatorname{Var}_\xi(C_g).
 \label{eq:smooth-exact-full-regression-identification}
\end{equation}
The variance identity follows also by expanding the square and using the
two normal equations in
\eqref{eq:smooth-exact-joint-block-elimination}; thus it is an ordinary
finite Schur complement, not an additional probabilistic assumption.

For the actual fast target, \(v_{\mathscr Q}=(\tau_{\mathcal G},0)\),
choose
\[
 r^\flat=\frac L{q_n}\Delta=Dd_D.
\]
Write \(A_\xi^{-1}v_{\mathscr Q}=(x^{(0)},a^{(0)})\) and put
\(u^{(0)}=\mathbf Qx^{(0)}\in\mathcal G\).  The same exact elimination
now reads
\begin{equation}
 P_{\alpha,n}D^{-1}\overline C_\xi P_{\alpha,n}u^{(0)}
       =P_{\alpha,n}d_D,\qquad
 a^{(0)}=-\Gamma_\xi^{-1}U_\xi^Tu^{(0)}.
 \label{eq:smooth-fast-block-elimination}
\end{equation}
By \eqref{eq:smooth-arithmetic-Linf-inverse} and
\(\|P_{\alpha,n}d_D\|_\infty\ll_W1\),
\begin{equation}
 \|u^{(0)}\|_\infty\ll_W1.
 \label{eq:smooth-fast-quotient-bound}
\end{equation}
Moreover \(|(U_\xi)_{i,*}|\ll_{B,W}H_i/L\),
\(\sum_iH_i\ll\log L\), and \(\|\Gamma_\xi^{-1}\|\ll_W1\).  Therefore
\[
 \|a^{(0)}\|\ll_{B,W}\frac{\log L}{L}=o_{n;B,W}(1).
\]
In particular the fast nuisance coefficient is bounded by an absolute
multiple depending only on \(W\), after increasing the threshold for the
already fixed box; its bound does not enter the choice of that box.

We may now apply \Cref{lem:augmented-bridge} through the exact
identification \eqref{eq:smooth-exact-full-regression-identification}.
It gives
\begin{equation}
 \sigma_\xi\asymp_Ww^2,\qquad
 |q_j^{\rm reg}|\ll w\alpha_j,\qquad
 \|a_Z\|\ll_Ww,\qquad
 \|c\|_\infty=\max_p|g_p-q_{j(p)}^{\rm reg}|\ll w.
 \label{eq:smooth-full-regression-effective-bounds}
\end{equation}
The only nontrivial target pairing in the numerator of
\eqref{eq:smooth-full-block-inverse} is now an exact dual pairing:
\begin{align}
 |b_\xi^TA_\xi^{-1}v_{\mathscr Q}|
 &=|(A_\xi^{-1}b_\xi)^Tv_{\mathscr Q}|
  =\left|\frac L{q_n}\sum_jq_j^{\rm reg}\Delta_j\right|\notag\\
 &\ll_Ww\sum_jH_j\alpha_j\ll_Ww.
 \label{eq:smooth-exact-slow-target-pairing}
\end{align}
Here the nuisance target is zero,
\eqref{eq:smooth-band-target-envelope} bounds \(\Delta_j\), and
\(\sum_jH_j\alpha_j=\sum_{W<p\leq y}t_p/p=O(1)\).
Together with \(|v_g|\ll_Ww\), the first line of
\eqref{eq:smooth-full-block-inverse} therefore gives
\(|\dot\lambda|\ll_Ww^{-1}\).  Its second line, interpreted through
\eqref{eq:smooth-exact-full-regression-identification} and
\eqref{eq:smooth-fast-block-elimination}, is the exact pair
\[
 \dot u=u^{(0)}-\dot\lambda q^{\rm reg},\qquad
 \dot a=a^{(0)}-\dot\lambda a_Z.
\]
Consequently, prime by prime,
\begin{equation}
 \dot u_{j(p)}+\dot\lambda g_p
 =u^{(0)}_{j(p)}+\dot\lambda
       (g_p-q_{j(p)}^{\rm reg})
 =u^{(0)}_{j(p)}+\dot\lambda c_p.
 \label{eq:smooth-exact-effective-prime-coefficient}
\end{equation}
Equations \eqref{eq:smooth-fast-quotient-bound},
\eqref{eq:smooth-full-regression-effective-bounds}, and the bounds for
\(\dot\lambda,a^{(0)}\) now prove
\begin{equation}
 |\dot\lambda|\ll_Ww^{-1},\qquad
 \max_{W<p\le y}|\dot u_{j(p)}+\dot\lambda g_p|
       +\|\dot a\|+w|\dot\lambda|\leq C_*.
 \label{eq:smooth-effective-block-inverse-bound}
\end{equation}
The leading continuum blocks give a constant \(C_{{\rm inv},W}\) independent of
the tilt box and the regular mesh.  Before this constant or the ODE box
is introduced, \(W\) has already made the box-independent terms
\(C_{\rm pow}/W\), \(C_{\rm pow}^{\sharp}/W\), and
\(C_{\rm rel}(C_{\rm cmp})/W\) smaller than the fixed continuum gaps.
All dependence on a later fixed tilt box is in
\(\epsilon_{B,W}(n)\), its explicitly rescaled relative version
\(\epsilon^{\rm rel}_{B,W,C_{\rm cmp}}(n)\), or the analogous fixed-mesh
remainders, all of which tend to zero with \(n\).  Fix
\(C_*=2C_{{\rm inv},W}\).  After the box below is prechosen, only the threshold
for \(n\) is increased so that all these \(B_*,W\)-dependent vanishing
errors use the spare factor two.  Thus
\eqref{eq:smooth-effective-block-inverse-bound} holds with this
noncircular \(C_*\).

We now remove the possible bootstrap circularity.  Define the effective
norm
\[
 \|\xi\|_{\rm eff}:=
 \max_{W<p\le y}|u_{j(p)}+\lambda g_p|
       +\|a\|+w|\lambda|.
\]
Choose in advance \(B_*=4C_*\), after \(W\) but before the ODE.  On the prechosen box
\(\|\xi\|_{\rm eff}\leq B_*\), the logarithm of the Radon--Nikodym
derivative is bounded by
\[
 \frac{B_*}{L}\sum_{p>W}v_p(m)+O_W(B_*/L)
 \leq\frac{B_*}{\log W}+o_W(1).
\]
This gives a fixed dominating density.  More sharply, the tilted marked
lemmas show that the leading \(h,K\) blocks are unchanged and that the
box-dependent local errors are \(O_{B_*,W}(1/L)\).  After increasing the
fixed-mesh threshold for \(n\), all covariance gaps and row estimates
used in \eqref{eq:smooth-effective-block-inverse-bound} remain within
the reserved factor two.  In particular its right side is at most
\(2C_{{\rm inv},W}=C_*<B_*/2\) throughout this box.  The box was chosen before
the path, so this conclusion is not being used to define its own domain.

The constants \(C_*\) and \(B_*\) are uniform over all permitted regular
meshes.  They may therefore be fixed after \(W\) and before a particular
mesh is instantiated; after the mesh and the box are fixed, only the
threshold for \(n\) is increased.  In particular, the choice of \(W\)
never depends on \(B_*\).  Fix such a mesh and such an \(n\).  In the
gauge coordinates the parameter
space is finite-dimensional and \(\mathcal M\) is \(C^\infty\).  On the
closed effective box
\[
 \mathcal K_*:=\{\xi:\|\xi\|_{\rm eff}\leq B_*\}
\]
the preceding covariance gaps make \(D\mathcal M(\xi)\) invertible and
\[
 \|D\mathcal M(\xi)^{-1}\tau\|_{\rm eff}\leq C_*
 \qquad(\xi\in\mathcal K_*).
\]
Solve the straight-target ODE
\begin{equation}
 \xi'(t)=D\mathcal M(\xi(t))^{-1}\tau,\qquad
                         \xi(0)=0.
 \label{eq:smooth-moment-map-ODE}
\end{equation}
Let \([0,T_{\max})\) be its maximal solution interval while the path
remains in the interior of \(\mathcal K_*\).  Local existence follows
from \eqref{eq:smooth-full-moment-Jacobian}.  More explicitly,
\(D\mathcal M\) is smooth and invertible on a neighborhood of the
closed box, so
\(\xi\mapsto D\mathcal M(\xi)^{-1}\tau\) is locally Lipschitz there;
the finite-dimensional Picard--Lindelöf theorem gives both local
existence and uniqueness.  For
\(t<\min(T_{\max},1)\), integration gives
\[
 \|\xi(t)\|_{\rm eff}\leq C_*t\leq C_*=B_*/4.
\]
If \(T_{\max}\leq1\), the same derivative bound makes \(\xi(t)\)
Cauchy as \(t\uparrow T_{\max}\).  Its limit lies in the strictly
smaller box \(\|\xi\|_{\rm eff}\leq B_*/4\), where the Jacobian remains
invertible.  The local existence theorem therefore restarts the
solution past \(T_{\max}\), a contradiction.  Hence \(T_{\max}>1\).
The chain rule gives
\[
 \frac d{dt}\mathcal M(\xi(t))=\tau,
 \qquad \mathcal M(\xi(1))=\tau.
\]
To check the converse coordinate
implication, let \(\Delta^{\rm out}\) be the resulting vector of band-moment
increments.  Equality against every \(v\in\mathcal G\) implies
\[
 \Delta^{\rm out}-\Delta=cD\alpha
\]
for some scalar \(c\), because the ordinary annihilator of
\(\mathcal G\) is \(\operatorname{span}(D\alpha)\).  Preservation of
the \(Z\)-moments preserves the physical and head moments, so the exact
logarithmic identity turns the slow equation into
\(\alpha^T\Delta^{\rm out}=\alpha^T\Delta\).  Hence
\(c\,\alpha^TD\alpha=0\), and \(c=0\).  Thus
\(\mathcal M(\xi(1))=\tau\) is precisely all band, head, physical, and
mass equations.  Existence, rather than strict convexity, has now
been proved; strict convexity only supplies uniqueness in the fixed
gauge.

For later use we record the primewise row bound for every block of this
inverse.  Along \eqref{eq:smooth-moment-map-ODE},
\begin{align}
 \left|\operatorname{Cov}_\xi(v_p,\dot{\mathscr Q}_{\rm quot})\right|
   &\ll_Wp^{-1},\notag\\
 \left|\operatorname{Cov}_\xi(v_p,\dot a^TZ)\right|
   &\ll_W(pL)^{-1},\notag\\
 \left|\dot\lambda\operatorname{Cov}_\xi(v_p,C_g)\right|
   &\ll_Ww^{-1}\frac wp\ll_Wp^{-1}.
 \label{eq:smooth-each-block-prime-row}
\end{align}
For the first line one must include the leading row, not only its
arithmetic error.  If \(u_j\) are the bounded effective quotient
coefficients, then
\[
 \left|\operatorname{Cov}_\xi\left(
 v_p,\sum_ju_j\Omega_j\right)\right|
 \le\frac{C_W}{p}
   +C_W\frac{t_p}{p}\sum_{W<q\le y}\frac{t_q}{q}
   +o(1/p)
 \ll_W\frac1p.
\]
This uses the diagonal multiplier, the product kernel
\eqref{eq:smooth-kernel-product}, and the box-uniform full-valuation
transfer \eqref{eq:smooth-power-row-transfer}.
The second line follows from the full head/physical marked family
\eqref{eq:smooth-explicit-nuisance-marked-family} (whose physical
coordinate is \eqref{eq:marked-tilted-physical-full}), and the third from
\eqref{eq:smooth-compensated-marked-row}.  Multiplication by
\(q_n/L\asymp N/L\) and integration over \(0\leq t\leq1\) show that the
bridge changes every marked moment by
\(O_W(N/(pL))\).  Combined with
\eqref{eq:strict-rate-rough-residual}, this proves
\eqref{eq:smooth-postbridge-fixed-rate}.

Only \(z^{\rm act}\) was reweighted.  The effective box bounds its density
distortion between two fixed positive constants.  Hence, for a constant
\(C_{B_*,W}\) independent of \(m,n\),
\begin{equation}
 0\leq z_m(\xi(1))\leq C_{B_*,W}z_m^0
       \leq\frac{C'_{B_*,W}}L.
 \label{eq:smooth-final-active-coordinate-bound}
\end{equation}
On a shared clean coordinate the only frozen summand is the protected
layer, also \(O_W(1/L)\); the top support is disjoint.  Therefore the
complete coordinate weight is between zero and
\((C'_{B_*,W}+C''_W)/L<1\) for all sufficiently large \(n\).  The
positive zero-head pool persists.
The disjoint top support and the frozen additive protected summand are
untouched, proving the remaining margin assertions.
\end{proof}

\begin{corollary}[Mesh-uniform marked-rate constant]
\label[corollary]{cor:smooth-mesh-uniform-marked-rate}
After \(W\) is fixed there exist \(\eta_0>0\) and a finite constant
\(C_{\rm tan}\), chosen before the particular fine mesh and before
\(w=\delta+\eta\), with the following property.  For every permitted
regular mesh with \(w<\eta_0\), once that mesh is fixed there is a
threshold \(n_0=n_0(W,\mathcal P,B_*)\) such that, for all
\(n\geq n_0\), the post-bridge residual satisfies
\begin{equation}
                         |r_p|\leq C_{\rm tan}\frac{N}{pL}
                         \qquad(W<p\leq y).
 \label{eq:smooth-uniform-Ctan-quantifier}
\end{equation}
The threshold is not asserted to be uniform over an infinite family of
meshes, but the displayed constant is.
\end{corollary}

\begin{proof}
Choose \(\eta_0\) below the continuum inverse thresholds in
\Cref{lem:arithmetic-band-quotient} and
\Cref{lem:smooth-weighted-discrete-inverse}.  The quotient and nuisance rows
in \eqref{eq:smooth-each-block-prime-row} are \(O_W(1/p)\), with
constants independent of the permitted mesh.  The slow row is uniform
for the nontrivial reason
\begin{equation}
 |\dot\lambda|\,
 |\operatorname{Cov}_\xi(v_p,C_g)|
 \ll_W\frac1w\frac wp\ll_W\frac1p,
 \label{eq:smooth-uniform-slow-row-cancellation}
\end{equation}
so its constant is also independent of \(w\).  Integrating the three
rows for \(0\leq t\leq1\) and multiplying by
\(q_n/L\asymp N/L\) gives a mesh- and \(w\)-independent addition to the
rough-stage marked-rate constant.  Enlarge the previously fixed
\(C_{\rm tan}\) by this amount.  All box- and mesh-dependent errors tend
to zero after the mesh is fixed and are absorbed by increasing only
\(n_0\).  This proves the stated quantifiers.
\end{proof}

\subsection{Choice of constants and the bridge output}
\label{subsec:smooth-bridge-output}

The order of constants is important.  First fix the tangent comparison
ratios \(r_0<3/2\) and \(\rho>1\) with \(\rho^3<r_0\), together with the
anchor and bank data.  Next choose one sufficiently large, fixed,
nonprime \(W\) containing the anchor support, put every prime at most
\(W\) in the finite head, and make the box-independent quantities
\[
 \frac{C_{\rm sf}+C_{\rm pow}+C_{\rm pow}^{\sharp}
          +C_{\rm rel}(C_{\rm cmp})}{W}
\]
smaller than the reserved continuum covariance and inverse gaps.  Never
enlarge \(W\) afterward.  The compact ODE box is chosen later; its
dependence appears only in remainders tending to zero with \(n\).  With this
head fixed, choose the rough-selector constants and the protected/active
split, then choose \(\delta_*<1/18\) and construct the guarded selector.
Record its surviving slack \(\sigma\) and fixed-rate constant
\(C_{\rm tan}\).  Only then choose the fixed regular band parameters
\(\delta,\eta\) small enough for the bridge and tangent inequalities.
Finally fix that mesh and let \(n\to\infty\), with its threshold allowed
to depend on all preceding fixed choices.  This order uses no uniformity
for a growing head block.

Combining \Cref{subsec:smooth-ledger,subsec:smooth-head-physical} and
\Cref{prop:smooth-fixed-partition-fit}, the fractional point passed to the
tangent stage has all of the following exact properties:
\begin{enumerate}[label=\textup{(\roman*)}]
\item every nontrivial complete rough-signature row has its prescribed
integer quota, and the flexible smooth row has the integer quota
\(q_{\rm sm}^{\rm flex}(d)\);
\item \(r_p=0\) for \(p\leq W\) and for \(p>y\);
\item the physical logarithm is exact, so
\eqref{eq:smooth-exact-physical-compatibility} holds;
\item every band satisfies \eqref{eq:smooth-exact-band-balance}, and every
medium-prime residual satisfies
\eqref{eq:smooth-postbridge-fixed-rate}; and
\item on clean smooth endpoints the still-frozen protected layer and the
active ceiling provide two-sided slack \(\sigma/L\); on clean
nonsmooth endpoints the two-sided slack is instead the broad-selector
margin left by \eqref{eq:rough-row-correction-slack}.
\end{enumerate}
No integer deletion bracket and no unverified endpoint sign is used in
this construction.
 \section{The finite-band tangent absorber}
\label{sec:tangent-absorber}

We now cancel the remaining residual at the primes \(W<p\leq y\).
The inputs from \Cref{sec:smooth-bridge} are
\begin{equation}
 \sum_{p\in\mathcal P_j}r_p=0,\qquad
 |r_p|\leq C_{\rm tan}\frac{N}{pL},
 \qquad C_{\rm tan}:=C_{1,W},
\label{eq:tangent-input}
\end{equation}
for every fixed exponent band, together with exact complete rough rows,
exact head valuations and exact ordinary logarithm.  There is also a fixed
\(\sigma>0\) such that every clean broad endpoint in an active,
nondedicated row has two-sided slack \(\sigma/L\).  In the smooth row its
lower margin comes from the frozen protected summand (and its upper margin
from the surviving active ceiling); in a nontrivial rough row both margins
are the broad-selector margins left after the \(o(1/L)\) row correction.
The residual convention is always target minus current.  All valuations in
\eqref{eq:tangent-input} are full valuations, including prime powers.

\subsection{Uniform clean common-multiplier lists}
\label{subsec:tangent-common-lists}

Fix \(1<r_0<3/2\).  For distinct primes
\[
                 W<v\leq u\leq y,\qquad u/v\leq r_0,
\]
put
\[
 I_{uv}:=\left(\frac nv,\frac{2n-Kh}{u}\right].
\]
For \(a\in I_{uv}\), both \(va\) and \(ua\) lie in the broad interval
\(J=(n,2n-Kh]\).  Let \(\mathscr D_{\rm row}\) be the set of fully
dedicated nonexceptional rows fixed in
\eqref{eq:rough-dedicated-row-data}.  Under the endpoint-deletion convention
used here this set is empty; we nevertheless retain the condition below so
that the list definition itself records the required active-row test.
Use the exhaustive numerical guard set
\(\Gamma_{\rm num}\) defined in \eqref{eq:numerical-guard-set}.  It
consists of the modified external anchors, the fixed exceptional upper
factors, every withheld bank donor, and both possible states of every bank
path.  Later bridge and tangent coordinates are flexible and hence are
not an additional class of guards.  Define
\begin{equation}
 \mathcal L^+_{uv}:=
 \left\{a\in I_{uv}\cap\mathbb Z:
  (a,P_{\rm hd})=1,\quad
  X_{R_y(a)}\geq X_0,\quad
  R_y(a)\notin\mathscr D_{\rm row},\quad
  ua,va\notin\Gamma_{\rm num}\right\}.
\label{eq:tangent-clean-list}
\end{equation}
Since \(u,v\leq y\),
\begin{equation}
                 R_y(ua)=R_y(a)=R_y(va).
\label{eq:tangent-row-identity}
\end{equation}
Thus a switch between \(ua\) and \(va\) stays in one complete rough
signature row.

\begin{lemma}[Uniform common list]
\label[lemma]{lem:tangent-common-list}
With the exceptional cutoff fixed once by \eqref{eq:rough-cutoff}, there
is a fixed \(\kappa>0\) such that
\begin{equation}
                         |\mathcal L^+_{uv}|
             \geq\kappa\frac nu
\label{eq:tangent-list-lower}
\end{equation}
uniformly for every permitted pair \(u,v\) and every sufficiently large
\(n\).
\end{lemma}

\begin{proof}
Put
\(\delta_{\rm hd}=\varphi(P_{\rm hd})/P_{\rm hd}\).
The interval length is
\[
 |I_{uv}|=\frac nu
       \left(2-\frac uv-\frac{Kh}{n}\right)+O(1).
\]
Counting the fixed reduced residue classes modulo \(P_{\rm hd}\) gives
\begin{equation}
 \#\{a\in I_{uv}:(a,P_{\rm hd})=1\}
 =\left\{\delta_{\rm hd}\left(2-\frac uv\right)+o(1)\right\}
       \frac nu.
\label{eq:tangent-raw-list}
\end{equation}
The error is uniform because \(n/u\geq n/y=n^{7/9}\).

We next remove the exceptional rows.  Write \(a=Rb\), where
\(R=R_y(a)\) and \(P^+(b)\leq y\).  If \(X_R=2n/R<X_0=n^{\delta_*}\),
then
\begin{equation}
                 R>\frac{2n}{X_0},\qquad
                 b<\frac{2X_0}{u}.
\label{eq:tangent-exceptional-factorization}
\end{equation}
For an upper bound, discard the additional lower cutoff
\(R>2n/X_0\).  For fixed \(b\), the full physical interval forced by
\(a\in I_{uv}\) has location and length both
\(\asymp_{r_0}n/(ub)\); whenever it can contain an exceptional multiplier,
\eqref{eq:tangent-exceptional-factorization} makes this length
\(\gg_{r_0}n/X_0\).  The interval Selberg sieve
\eqref{eq:Selberg-specialized} applies uniformly: its lcm remainder level
is \(y^4=n^{8/9}\), while
\[
                 n/X_0=n^{1-\delta_*},
                 \qquad \delta_*<1/18.
\]
For a fixed sieve constant \(C_{\rm sv}>0\), it follows that the number
\(E_{uv}\) of exceptional multipliers is
\begin{align}
 E_{uv}
 &\ll\frac{n}{u\log y}
       \sum_{b\leq2X_0/u}\frac1b
       +\frac{X_0}{u}\frac{y^4}{(\log y)^2}\notag\\
 &\leq
 \left\{C_{\rm sv}\frac{\delta_*}{\theta}+o(1)\right\}
       \frac nu.
\label{eq:tangent-exceptional-loss}
\end{align}
Indeed the second term is \(o(n/u)\), since
\[
 \frac{X_0y^4}{n(\log y)^2}
       =n^{\delta_*+8/9-1+o(1)}=o(1).
\]
If \(u>2X_0\), this sum is empty.  Notice that the loss is a small fixed
fraction, not \(o(n/u)\).

It remains to check the actual guards.  We give the complete census of
those guards that can meet a lower endpoint \(ua\) or \(va\).
Every modified prefix anchor is \(Pq\), where \(P>y\) is its unique
rough marker and every prime factor of \(q\) is at most
\(2R+1\leq W\); hence it contains no medium label.  A row-zero anchor is
a prime greater than \(n\), so it contains no medium label either.  A
promoted anchor is \(2^{k_p}p^{e_p}\).  If \(p>y\), its only smooth prime
is \(2\); if \(p\leq y\), it is one of at most \(\pi(y)\) promoted smooth
anchors.  Consequently only \(O(\pi(y))\) anchors can be divisible by a
prime in \((W,y]\).

There are \(O(yL)\) component factors in the fully reserved bank by
\eqref{eq:bank-total-size}; taking both states changes only the implied
constant.  These are the only remaining lower guards.  Indeed,
\(G_{\rm fix}\subset(2n,M]\) by
\eqref{eq:rough-fixed-list-closed}, and every member of
\(E_{\rm donor}\) lies above \(2n\), so neither class can equal a
lower endpoint.  Thus the number of endpoint-relevant guards is
\[
 O(\pi(y))+O(yL)=O(yL).
\]
For a fixed label \(u\), the equation \(ua=g\) determines at most the one
integer multiplier \(a=g/u\); the same statement holds for \(v\).
Therefore the guard deletion is, uniformly in the permitted pair,
\begin{equation}
                 O(yL)=o(n/u),
\label{eq:tangent-guard-loss}
\end{equation}
because \(u\leq y\) and \(y^2L/n=o(1)\).
The dedicated-row condition costs zero because
\(\mathscr D_{\rm row}=\varnothing\).  More generally, the quantitative
data in \eqref{eq:rough-dedicated-row-data} would bound its loss by
\(O\{yL(y/u+1)\}=o(n/u)\).  Thus the definition cannot call on a row
whose flexible quota was set to zero.  Terminal exceptional rows were
already removed in \eqref{eq:tangent-exceptional-loss}.  Forced exceptional upper factors
lie above \(2n\) and cannot occur in a lower list.

Combining \eqref{eq:tangent-raw-list},
\eqref{eq:tangent-exceptional-loss}, and
\eqref{eq:tangent-guard-loss}, the one-time inequality
\eqref{eq:rough-cutoff} gives
\[
 C_{\rm sv}\frac{\delta_*}{\theta}
 <\frac14\delta_{\rm hd}(2-r_0),
\qquad
 \delta_{\rm hd}(2-r_0)
       -C_{\rm sv}\frac{\delta_*}{\theta}
 >\frac34\delta_{\rm hd}(2-r_0)>0.
\]
Thus, without changing \(\delta_*\), one may take
\[
 \kappa=\frac12\left\{\delta_{\rm hd}(2-r_0)
       -C_{\rm sv}\frac{\delta_*}{\theta}\right\}>0.
\]
\end{proof}

We record explicitly why every endpoint now retained in
\eqref{eq:tangent-clean-list} has usable slack.  If \(R_y(a)=1\), the
frozen protected smooth summand supplies the lower floor and the fitted
active summand stays below its broad ceiling.  If \(R_y(a)>1\), the row is
nondedicated by definition of the list, and the broad raw floor and ceiling
survive the row correction because that correction is \(o(1/L)\) per clean
coordinate.  After decreasing the fixed constant \(\sigma\), if needed,
both cases give
\begin{equation}
                 \frac{\sigma}{L}\leq x_{ua},x_{va}
                 \leq1-\frac{\sigma}{L}.
\label{eq:tangent-endpoint-slack}
\end{equation}
Thus one pair \(u,v\) has at least
\begin{equation}
                         \kappa\sigma\frac{N}{u}
\label{eq:tangent-pair-capacity}
\end{equation}
real two-sided capacity.  This floor is used for tangents only now; it
was not a source of bridge covariance.

\subsection{Finite-band earthmover estimates}
\label{subsec:tangent-earthmover}

Use the exponent bands corresponding to
\eqref{eq:smooth-regular-relative-mesh}:
\begin{equation}
 \mathcal B_0=(W,n^{\delta_b}],\qquad
 \mathcal B_j=(n^{a_{j-1}},n^{a_j}]
       \quad(1\leq j\leq s),
\label{eq:tangent-exponent-bands}
\end{equation}
where \(a_0=\delta_b\), \(a_s=\theta=2/9\), and
\[
 c_{\rm mesh}\eta a_{j-1}\leq a_j-a_{j-1}
       \leq\eta a_{j-1}.
\]
Set
\begin{equation}
                         w_{\rm tan}:=\delta_b+\theta\eta.
\label{eq:tangent-width}
\end{equation}
These are the same bands as in the bridge, merely written in the
\(\log n\)-exponent convention.

Choose a fixed \(\rho>1\) with \(\rho^3<r_0\).  Partition each band,
separately, into multiplicative \(\rho\)-cells, merging a possible
terminal fragment with its neighbor.  Any labels paired within one cell
or in adjacent cells then have ratio at most \(\rho^3<r_0\).
The prime number theorem, after \(W\) is chosen sufficiently large
depending on \(\rho\), gives
\begin{equation}
 \#\{p\text{ in a cell at scale }Q\}
                         \asymp_\rho\frac{Q}{\log Q}.
\label{eq:tangent-cell-prime-count}
\end{equation}
This is a fixed-ratio interval, so no short-interval prime theorem is
being used.

\begin{lemma}[Finite-band transport]
\label[lemma]{lem:tangent-finite-band-transport}
In each exponent band \([a,b]\), there is a nonnegative directed flow
\((f_{uv})\) on its prime labels such that
\begin{align}
 \sum_v(f_{pv}-f_{vp})&=r_p,
       \label{eq:tangent-flow-divergence}\\
 f_{uv}>0&\ \Longrightarrow\
       \max(u,v)/\min(u,v)\leq\rho^3<r_0.
       \label{eq:tangent-flow-locality}
\end{align}
Writing \(f_p:=\sum_v(f_{pv}+f_{vp})\), the union of these flows over all
bands may be chosen to satisfy
\begin{align}
 \#\{(u,v):f_{uv}>0\}&=O_\rho(\pi(y)),
       \label{eq:tangent-sparse-support}\\
 \sum_{u,v}f_{uv}
   &\ll_\rho C_{\rm tan}Nw_{\rm tan}+o(N),
       \label{eq:tangent-total-traffic}\\
 \frac{pf_p}{N}
   &\ll_\rho C_{\rm tan}
       \left\{t\log\frac bt+\frac1L\right\}
       \quad(p\asymp_\rho n^t),
       \label{eq:tangent-label-load}\\
 \sup_p\frac{pf_p}{N}
   &\ll_\rho C_{\rm tan}w_{\rm tan}+o(1).
       \label{eq:tangent-max-label-load}
\end{align}
All constants are uniform over the finitely many bands in the permitted
mesh.
\end{lemma}

\begin{proof}
Fix one band and write its consecutive prime cells as
\(\mathcal Q_1,\ldots,\mathcal Q_m\).  Put
\[
 s_i:=\sum_{p\in\mathcal Q_i}r_p,\qquad
 F_i:=\sum_{h\leq i}s_h\quad(1\leq i<m),\qquad
 F_0=F_m=0.
\]
The last equality is exactly the band-balance assertion in
\eqref{eq:tangent-input}.  If the boundary between
\(\mathcal Q_i\) and \(\mathcal Q_{i+1}\) is \(Q=n^t\), band balance also
gives
\[
 |F_i|
 =\left|\sum_{Q<p\leq n^b}r_p\right|.
\]
Using the pointwise bound in \eqref{eq:tangent-input} and subtracting the
two instances of \eqref{eq:Mertens-harmonic}, at their actual endpoints,
we obtain uniformly
\begin{equation}
 |F_i|\ll C_{\rm tan}\frac NL
       \left\{\log\frac bt+\frac1{tL}\right\}.
\label{eq:tangent-cut-load}
\end{equation}
The term \(1/(tL)\) is the endpoint error
\(O(1/\log Q)\); this remains valid in the moving first band, where
\(t\geq\log W/L\).  Thus no fixed-positive lower exponent is being
assumed.

We next give the promised flow algorithm.  Across boundary \(i\), send
the signed load \(F_i\) from \(\mathcal Q_i\) to
\(\mathcal Q_{i+1}\): if \(F_i<0\), reverse the direction.  On either
side distribute its absolute value uniformly among all primes of that
cell.  The two resulting lists of nonnegative port masses have the same
total \(|F_i|\); the usual greedy matching of their first still-positive
entries realizes this transport with at most
\(|\mathcal Q_i|+|\mathcal Q_{i+1}|-1\) positive edges.

Let \(d_p\) be the boundary-flow divergence at \(p\), with outgoing minus
incoming convention.  Since the signed boundary load on the right of
cell \(i\) is \(F_i\), while that on its left is \(F_{i-1}\),
\[
 \sum_{p\in\mathcal Q_i}d_p=F_i-F_{i-1}=s_i.
\]
Consequently \(q_p:=r_p-d_p\) has sum zero in each cell.  Greedily match
the positive and negative \(q_p\)'s inside that cell, directing flow from
the former to the latter.  This uses at most
\(|\mathcal Q_i|-1\) positive edges and gives divergence exactly \(q_p\).
Adding the boundary and internal flows proves
\eqref{eq:tangent-flow-divergence}.  Endpoints of an internal edge lie
in one cell, and endpoints of a boundary edge lie in adjacent cells.
The possible merged terminal cell was allowed for in the exponent
\(\rho^3\), so \eqref{eq:tangent-flow-locality} follows.

This construction also proves the quantitative claims without an
unproved distribution step.  Put \(m_i:=|\mathcal Q_i|\).  The total
boundary traffic incident to \(p\in\mathcal Q_i\) is at most
\[
 b_p:=\frac{|F_{i-1}|+|F_i|}{m_i}.
\]
The internal traffic incident to \(p\) is \(|q_p|\), and
\(|d_p|\leq b_p\); hence
\begin{equation}
                         f_p\leq |r_p|+2b_p.
\label{eq:tangent-port-load-identity}
\end{equation}
For \(p\asymp_\rho Q=n^t\),
\eqref{eq:tangent-cell-prime-count} says
\(m_i\asymp_\rho p/(tL)\).  Substitution of
\eqref{eq:tangent-cut-load} into
\eqref{eq:tangent-port-load-identity}, together with
\(|r_p|\leq C_{\rm tan}N/(pL)\), proves
\eqref{eq:tangent-label-load}.

There are only a bounded number of greedy matchings incident to a cell:
one internal matching and at most two boundary matchings.  Summing their
edge counts proves \eqref{eq:tangent-sparse-support}.  Moreover, the
cross-boundary traffic is \(\sum_i|F_i|\), while the internal traffic is
\[
 \frac12\sum_p|q_p|
 \leq\frac12\sum_p|r_p|+\sum_i|F_i|.
\]
Therefore
\begin{equation}
 \sum_{u,v}f_{uv}
 \leq\frac12\sum_{W<p\leq y}|r_p|+2\sum_{\rm boundaries}|F_i|.
\label{eq:tangent-traffic-ledger}
\end{equation}
The first term is
\begin{equation}
 \frac12\sum_{W<p\leq y}|r_p|
       \ll C_{\rm tan}\frac{N\log L}{L}=o(N)
\label{eq:tangent-direct-traffic}
\end{equation}
by \eqref{eq:Mertens-harmonic}.  Consecutive cell boundaries have
exponent spacing \(\asymp_\rho1/L\), so
\eqref{eq:tangent-cut-load} and an upper Riemann sum give
\begin{align}
 \sum_{\rm boundaries}|F_i|
 &\ll_\rho C_{\rm tan}N\left\{
 \int_0^{\delta_b}\log\frac{\delta_b}{t}\,dt
 +
 \sum_{j=1}^s\int_{a_{j-1}}^{a_j}
                 \log\frac{a_j}{t}\,dt\right\}+o(N)\notag\\
 &\ll_\rho C_{\rm tan}N(\delta_b+\theta\eta)+o(N).
\label{eq:tangent-path-traffic}
\end{align}
Here the accumulated \(1/(tL)\) endpoint terms are
\(O(N\log L/L)=o(N)\).  Also
\(\int_0^{\delta_b}\log(\delta_b/t)\,dt=\delta_b\), while
\[
 \int_a^b\log(b/t)\,dt
 \leq\frac{(b-a)^2}{a}
 \leq\eta(b-a)
\]
for every regular positive band; summing gives \(O(\theta\eta)\).
Equations \eqref{eq:tangent-traffic-ledger}--\eqref{eq:tangent-path-traffic} prove
\eqref{eq:tangent-total-traffic}.  Finally,
\(t\log(b/t)\leq b-t\); it is at most \(\delta_b\) in the first band and
at most \(\theta\eta\) in a positive regular band.  This proves
\eqref{eq:tangent-max-label-load} and completes the construction.
\end{proof}

At a boundary \(Q=n^t\), \eqref{eq:tangent-cell-prime-count} and
\eqref{eq:tangent-pair-capacity} give aggregate port capacity
\begin{equation}
 \asymp_\rho\kappa\sigma\frac{N}{\log Q}
       =\asymp_\rho\kappa\sigma\frac{N}{tL}.
\label{eq:tangent-boundary-capacity}
\end{equation}
Comparing this with \eqref{eq:tangent-cut-load} shows explicitly that
the maximum boundary utilization is
\begin{equation}
 \ll_\rho
 \begin{cases}
 C_{\rm tan}\delta_b/(\kappa\sigma),&\mathcal B_0,\\
 C_{\rm tan}\theta\eta/(\kappa\sigma),&\mathcal B_j,\ j\geq1.
 \end{cases}
\label{eq:tangent-utilization}
\end{equation}

\subsection{Requests and simultaneous collision avoidance}
\label{subsec:tangent-LLL}

For a directed support edge \(e=(s(e),t(e))\), put
\[
 f_e:=f_{s(e)t(e)},\qquad
 U_e:=\max\{s(e),t(e)\},\qquad V_e:=\min\{s(e),t(e)\}.
\]
The common list is indexed by the unordered pair
\((U_e,V_e)\), while the source--sink orientation is retained as part of
the request.  For \(f>0\), put
\begin{equation}
 k(f):=\max\left\{1,\left\lceil\frac{4Lf}{\sigma}\right\rceil\right\}.
\label{eq:tangent-request-splitting}
\end{equation}
Split an edge \(e\) with \(f_e>0\) into \(k(f_e)\) equal requests.  The
definition gives
\[
                  k(f)\geq\frac{4Lf}{\sigma},
 \qquad \frac{f}{k(f)}\leq\frac{\sigma}{4L},
 \qquad f\leq\frac{\sigma k(f)}{4L}.
\]
Thus every request has size at most \(\sigma/(4L)\), and
\begin{equation}
                         k(f)\leq\frac{4Lf}{\sigma}+1.
\label{eq:tangent-request-count-edge}
\end{equation}
The second inequality is essential: a small-piece assertion alone would
not control the number of independent random choices.

Let \(k_p\) be the number of requests incident to the prime label \(p\),
and let \(K_{\rm req}\) be the total number of requests.  By
\eqref{eq:tangent-sparse-support},
\begin{align}
 k_p&\leq\frac{4L}{\sigma}f_p+\deg_{\rm supp}(p),\notag\\
 K_{\rm req}
 &\leq\frac{4L}{\sigma}\sum_{u,v}f_{uv}
       +O(\pi(y)).
\label{eq:tangent-request-counts}
\end{align}
Equations \eqref{eq:tangent-total-traffic},
\eqref{eq:tangent-label-load}, and
\eqref{eq:tangent-max-label-load} imply
\begin{align}
 \frac{pk_p}{n}
 &\ll_\rho\frac{C_{\rm tan}}{\sigma}w_{\rm tan}+o(1),
       \notag\\
 \frac{K_{\rm req}}n
 &\ll_\rho\frac{C_{\rm tan}}{\sigma}w_{\rm tan}+o(1).
\label{eq:tangent-normalized-request-loads}
\end{align}
Indeed, the direct term contributes
\(O(C_{\rm tan}\log L/(\sigma L))=o(1)\).  The \(+1\) terms in
\eqref{eq:tangent-request-count-edge} are harmless because
\[
 \frac{\pi(y)}n=o(1),\qquad
 \frac{y\pi(y)}n=n^{2\theta-1+o(1)}=o(1).
\]
In particular, if
\begin{equation}
 \frac{C_{\rm tan}}{\kappa\sigma}w_{\rm tan}
\label{eq:tangent-port-smallness}
\end{equation}
is sufficiently small, every prime port retains a fixed capacity margin.
Under the stronger smallness condition imposed in
\eqref{eq:tangent-final-smallness}, the implicit constant in
\eqref{eq:tangent-normalized-request-loads} may also be absorbed so that,
uniformly for every directed support edge \(e\),
\begin{equation}
k(f_e)\leq k_{U_e}\leq\frac{\kappa}{4}\frac n{U_e}
       \leq\frac14|\mathcal L^+_{U_eV_e}|.
\label{eq:tangent-requests-fit-list}
\end{equation}
Moreover, this request bound implies the real edge-capacity bound
\begin{equation}
 f_e\leq\frac{\sigma k(f_e)}{4L}
       \leq\frac{\kappa\sigma}{16}\frac{n}{U_eL}
       =\frac{\kappa\sigma}{16}\frac N{U_e}.
\label{eq:tangent-real-edge-load}
\end{equation}
On the other hand, the actual two-sided capacity of the common list is
\[
 |\mathcal L^+_{U_eV_e}|\frac{\sigma}{L}
       \geq\kappa\frac n{U_e}\frac{\sigma}{L}
       =\kappa\sigma\frac N{U_e}.
\]
Hence each real edge uses at most one sixteenth of its literal available
capacity.  Even before the random experiment, it also has fewer requests
than list entries.  The later local lemma enforces the stronger global
endpoint disjointness.

Let \(\mathscr R\) be the finite request set.  For a request
\(\mathfrak r\) on labels
\(\{U_{\mathfrak r},V_{\mathfrak r}\}\), with
\(U_{\mathfrak r}\geq V_{\mathfrak r}\), let \(Z_{\mathfrak r}\) be
uniform on \(\mathcal L^+_{U_{\mathfrak r}V_{\mathfrak r}}\).  Its
directed source and sink are also retained from the flow edge that
created it.  Choose all
the variables \(Z_{\mathfrak r}\) independently.  Its two numerical
endpoints are
\[
 E_{\mathfrak r}:=
 \{U_{\mathfrak r}Z_{\mathfrak r},
   V_{\mathfrak r}Z_{\mathfrak r}\}.
\]
They are distinct within one request because
\(U_{\mathfrak r}\ne V_{\mathfrak r}\).
For every unordered pair of distinct requests define the bad event
\begin{equation}
 \mathscr E_{\mathfrak r,\mathfrak s}
 :=\{E_{\mathfrak r}\cap E_{\mathfrak s}\ne\varnothing\}.
\label{eq:tangent-bad-event}
\end{equation}
This definition includes all four cross-request endpoint equalities.

Suppose first that the two label pairs are disjoint.  For fixed endpoint
labels \(\alpha\) and \(\beta\), the equation
\(\alpha a=\beta b\) has, because the two primes are distinct,
\[
                         a=\beta m,\qquad b=\alpha m.
\]
The endpoint range gives \(m=O(n/(\alpha\beta))\); the harmless literal
\(+1\) is absorbed because \(\alpha\beta\leq y^2=o(n)\).  Moreover the
larger label in the first request is at most \(r_0\alpha\), and similarly
for the second.  Lemma~\ref{lem:tangent-common-list} therefore gives a
product of list sizes at least
\(\kappa^2n^2/(r_0^2\alpha\beta)\).  Taking the union over the four
endpoint equations yields
\begin{equation}
 \mathbb P(\mathscr E_{\mathfrak r,\mathfrak s})
       \ll_{r_0}\frac1{\kappa^2n}.
\label{eq:tangent-collision-disjoint-labels}
\end{equation}

Now suppose that the requests share a label \(p\).  After writing their
pairs as \(\{p,r\}\) and \(\{p,q\}\), the four equalities are
\[
 pa=pb,\qquad pa=qb,\qquad ra=pb,\qquad ra=qb.
\]
For \(pa=pb\), there are at most \(O_{r_0}(n/p)\) possible common
multipliers, while both lists have size
\(\gg_{r_0}\kappa n/p\).  Its probability is therefore
\(O_{r_0}(p/(\kappa^2n))\).  Each equality with distinct endpoint labels
has \(O(n/(\alpha\beta))\) raw solutions as above and probability
\(O_{r_0}(1/(\kappa^2n))\).  This also covers identical label pairs:
there are then two same-label equalities, with the same bound.  Hence
\begin{equation}
 \mathbb P(\mathscr E_{\mathfrak r,\mathfrak s})
       \ll_{r_0}\frac{p}{\kappa^2n}.
\label{eq:tangent-collision-shared-label}
\end{equation}
These cases exhaust the event in \eqref{eq:tangent-bad-event}.

Fix a request with labels \(u,v\).  There are at most \(k_u\) other
requests sharing \(u\), at most \(k_v\) sharing \(v\), and at most
\(K_{\rm req}\) remaining requests.  Thus
\begin{align*}
 \sum_{\mathfrak s\ne\mathfrak r}
   \mathbb P(\mathscr E_{\mathfrak r,\mathfrak s})
 &\ll_{r_0}\frac1{\kappa^2}
    \left(\frac{uk_u+vk_v}{n}+\frac{K_{\rm req}}n\right).
\end{align*}
Equations \eqref{eq:tangent-normalized-request-loads},
\eqref{eq:tangent-collision-disjoint-labels}, and
\eqref{eq:tangent-collision-shared-label} therefore give
\begin{equation}
 B_{\rm coll}:=
 \sup_{\mathfrak r}
 \sum_{\mathfrak s\ne\mathfrak r}
       \mathbb P(\mathscr E_{\mathfrak r,\mathfrak s})
 \ll_{r_0,\rho}
 \frac{C_{\rm tan}}{\kappa^2\sigma}w_{\rm tan}+o(1).
\label{eq:tangent-per-request-collision-sum}
\end{equation}
Thus every bad event, including every repeated-label case, is present in
the displayed probability ledger.  Choose the fixed band parameters so
that \(B_{\rm coll}\leq1/8\).

For completeness, we now specify the dependency graph and verify the
factor-four bookkeeping.  Join two bad events when their unordered
request pairs intersect.  An event is a function of its two request
variables, and the underlying variables are mutually independent;
therefore it is independent of the sigma-algebra generated by all
nonneighbors.  For an event \(E\), its neighborhood lies in the union of
the bad events involving either of its two requests, and hence
\[
 \sum_{E'\sim E}\mathbb P(E')\leq2B_{\rm coll}.
\]
Uniformly, \(\mathbb P(E)=O_{r_0}(y/(\kappa^2n))=o(1)\), so
\(x_E:=2\mathbb P(E)\) belongs to \([0,1)\) for sufficiently large \(n\).
Also
\[
 \sum_{E'\sim E}x_{E'}\leq4B_{\rm coll}\leq\frac12.
\]
Using \(\prod_i(1-z_i)\geq1-\sum_i z_i\) for nonnegative \(z_i\) with
sum at most one,
\[
 x_E\prod_{E'\sim E}(1-x_{E'})
 \geq2\mathbb P(E)
       \left(1-\sum_{E'\sim E}x_{E'}\right)
 \geq\mathbb P(E).
\]
The asymmetric local lemma from \Cref{subsec:LLL} therefore gives a
simultaneous choice with no endpoint collision.  Since every list was
guarded in advance, these endpoints are also distinct from every anchor,
fixed factor, bank state, and withheld donor.

\subsection{Exact flow algebra and preserved invariants}
\label{subsec:tangent-invariants}

Lemma~\ref{lem:tangent-finite-band-transport} constructs the directed
flow, and the request splitting preserves its divergence exactly.  If
\(s(e)\) and \(t(e)\) denote respectively the source and sink label of a
request and \(\tau_e>0\) its size, then
\begin{equation}
 \sum_e\tau_e
       \bigl(\mathbf e_{s(e)}-\mathbf e_{t(e)}\bigr)=r.
\label{eq:tangent-flow-invariant}
\end{equation}
This equation fixes the sign convention: the selector change is \(+r\).
Equivalently, graph boundary is source minus sink and equals \(r\).

For the multiplier chosen for a request, use the signed coordinate move
\[
                 \tau_e\bigl(\mathbf1_{s(e)a(e)}
                         -\mathbf1_{t(e)a(e)}\bigr).
\]
Its full valuation change is
\begin{equation}
 v(s(e)a(e))-v(t(e)a(e))
             =\mathbf e_{s(e)}-\mathbf e_{t(e)}.
\label{eq:tangent-exact-valuation-switch}
\end{equation}
This remains true even if the multiplier contains \(s(e)\), \(t(e)\),
or higher powers of either, because the complete vector \(v(a(e))\)
cancels.

The moves have the following exact invariants.
\begin{enumerate}[label=\textup{(\roman*)}]
\item By \eqref{eq:tangent-row-identity}, each move stays inside one
complete rough-signature row; cardinality and every valuation above
\(y\) are unchanged.
\item The multiplier is \(P_{\rm hd}\)-free and \(s(e),t(e)>W\), so all
head valuations are unchanged.
\item Both labels lie in one exponent band, so every exact band total is
unchanged.
\item Equation \eqref{eq:tangent-flow-invariant} and
\eqref{eq:tangent-exact-valuation-switch} make every medium-prime
valuation exact.
\end{enumerate}
The preliminary selector was exact in ordinary logarithm.  Since its head
and rough coordinates were already exact, the residual convention gives
\begin{equation}
                         \sum_{W<p\leq y}r_p\log p=0.
\label{eq:tangent-log-compatibility}
\end{equation}
Indeed, the target and the current selector have the same valuations
outside \((W,y]\), so the logarithm of their product ratio is exactly the
left side.  Taking the prime-log scalar product of
\eqref{eq:tangent-flow-invariant} now gives the complete calculation
\begin{equation}
 \Delta\log
 =\sum_e\tau_e\{\log s(e)-\log t(e)\}
 =\sum_{W<p\leq y}r_p\log p
 =0.
\label{eq:tangent-exact-log-invariant}
\end{equation}
Thus no separate logarithmic balance is required inside an exponent
band.

All chosen numerical endpoints are distinct, and every request has size
at most \(\sigma/(4L)\).  Thus each coordinate is changed at most once
and \eqref{eq:tangent-endpoint-slack} gives the explicit post-move margin
\begin{equation}
 \frac{3\sigma}{4L}\leq x_a^{\rm new}
       \leq1-\frac{3\sigma}{4L}
 \qquad\text{for every endpoint used by the tangent}.
\label{eq:tangent-postmove-margin}
\end{equation}
Every unused coordinate keeps its previous feasible value, so the whole
vector remains in \([0,1]\).  At this point the protected layer has been
released: the object passed to rounding is the single total vector
\((x_a)\), not a separate floor plus an independently roundable selector.

\begin{proposition}[Exact post-tangent certificate]
\label[proposition]{prop:post-tangent-certificate}
Let \(\mathcal A_R\) denote the remaining flexible candidates of actual
complete signature \(R=R_y(a)\); thus \(R=1\), not \(0\), is the smooth
row.  After the tangent
construction,
\begin{equation}
\boxed{
\begin{gathered}
 0\leq x_a\leq1\quad(a\in\mathcal A),\\
 \sum_{a\in\mathcal A_R}x_a=q_R\in\mathbb Z
       \quad(R\ne1),\qquad
 \sum_{a\in\mathcal A_1}x_a
       =q_{\rm sm}^{\rm flex}(d)\in\mathbb Z,\\
 v(B_0)+\sum_{a\in\mathcal A}x_av(a)=v(Y)
       \qquad\text{at every prime}.
\end{gathered}}
\label{eq:post-tangent-certificate}
\end{equation}
\end{proposition}

\begin{proof}
The row identities are unchanged by the tangent and were integer
identities before it.  Primes \(p\leq W\) were exact after the head fit,
primes \(W<p\leq y\) are exact by
\eqref{eq:tangent-flow-invariant}, and primes \(p>y\) were exact by the
complete rough-signature equations.  This proves the valuation identity
simultaneously at all primes.  Coordinate feasibility was proved above.
\end{proof}

Since every valuation on the left of
\eqref{eq:post-tangent-certificate} is nonnegative, the same identity
also re-verifies \(B_0\mid Y\) for the \emph{combined} charged base.  It
is not obtained by multiplying separate divisibility claims.  The
certificate \eqref{eq:post-tangent-certificate} is the precise input to
the deterministic rounding and bank absorption in the next section.

\subsection{Final parameter order}
\label{subsec:tangent-parameter-order}

We summarize the noncircular choice.  Fix \(r_0<3/2\), then fix
\(\rho>1\) with \(\rho^3<r_0\).  Choose the single finite head cutoff
\(W\) large enough for the head cells, marked estimates, uniform bridge
perturbation, guard estimates, and the fixed-ratio cell prime counts.
For the bridge perturbation this choice uses only the box-independent
leading constants in
\eqref{eq:marked-box-uniform-power-row},
\eqref{eq:marked-compensated-power-relative}, and
\eqref{eq:smooth-sharp-power-transfer}; it does not use the radius of
the later ODE box.
Use the one-time cutoff \(\delta_*\) already fixed in
\eqref{eq:rough-cutoff}, construct the guarded rough selector and the
head/physical baseline, and apply \Cref{lem:tangent-common-list}; this
fixes \(\kappa\) and the surviving endpoint margin \(\sigma\).  The
uniform estimates of \Cref{prop:smooth-fixed-partition-fit} now supply a
mesh threshold and a constant \(C_{\rm tan}\) valid for every sufficiently
fine regular partition; at this point no particular partition has yet been
instantiated.  Choose a fixed
\(\varepsilon_{\rm tan}=\varepsilon_{\rm tan}
(r_0,\rho,W,\kappa,\sigma)>0\) small enough to absorb the fixed constants
in \eqref{eq:tangent-utilization},
\eqref{eq:tangent-normalized-request-loads}, and
\eqref{eq:tangent-per-request-collision-sum}.  Then choose fixed
\(\delta_b,\eta>0\) inside the uniform mesh range and so small that
\begin{equation}
 C_{\rm tan}(\delta_b+\theta\eta)
       \leq\varepsilon_{\rm tan}\kappa^2\sigma.
\label{eq:tangent-final-smallness}
\end{equation}
Instantiate the bridge on this partition and then run the tangent.  The
displayed condition implies both the port condition
\eqref{eq:tangent-port-smallness} and the local-lemma condition
\eqref{eq:tangent-per-request-collision-sum}; request sizes themselves are
bounded by the previously fixed slack in
\eqref{eq:tangent-request-splitting}.  Finally take \(n\) sufficiently
large to absorb every accompanying \(o(1)\) term, including all
\(B_*,W\)-dependent power-transfer remainders, and then let
\(n\to\infty\).
Neither \(W\) nor any already recorded \(W\)-dependent constant is changed
afterward.
 \section{Deterministic rounding and final assembly}
\label{sec:rounding-assembly}

We now convert the exact fractional selector into an exact subset of legal
factors.  We first recall the chronology of the charge, because it rules
out a possible circularity.  Before the rough selector was solved, the
fixed residual factors and the chosen zero states of the universal bank
were combined into
\begin{equation}
 B_0=
 \prod_{a\in G_{\rm fix}}a
 \prod_{a\in G_{\rm bank}^0}a
 \label{eq:assembly-base-product}
\end{equation}
as in \eqref{eq:rough-charged-base}.  The nonnegative exceptional estimate
\eqref{eq:uniform-exceptional-charge}, the bank and guard estimate
\[
 v_p(B_0/B_{\rm exc})=o_W\left(\frac{N}{pL}\right)
 \qquad(p\le y),
\]
and the uniform tail supply
\[
 v_p(T_n)=\frac{(c+o(1))N}{p-1}
\]
were compared coordinatewise.  At primes in \(F_{\rm anc}\), this
comparison used the strict reserve remaining after \(D\) was changed to
\(D'\); away from \(F_{\rm anc}\), the divisor \(D'\) has no valuation.
After \(\delta_*\) was fixed sufficiently small, these inequalities held
simultaneously at every \(p\le y\).  Above \(y\), every fixed or bank-base
token was injected into a different withheld tail token with the same
complete signature.  Thus, before any fractional solve,
\begin{equation}
                         D'B_0\mid T_n,\qquad
                         B_0\mid Y:=T_n/D'.
 \label{eq:assembly-pre-solve-divisibility}
\end{equation}
This is the combined certificate
\eqref{eq:combined-precharge-divisibility}; it is not obtained by
multiplying separate divisors of \(T_n\).

Let \(\mathcal A=\coprod_S\mathcal A_S\) be the flexible lower candidates
remaining after all declared guards.  The tangent stage supplies one total
weight \(x_a\in[0,1]\) for each \(a\in\mathcal A\) and the exact
post-tangent certificate
\begin{align}
 \sum_{a\in\mathcal A_S}x_a&=q_S\in\mathbb Z
       &&\text{for every flexible complete-signature row }S,
 \label{eq:assembly-integer-row-certificate}\\
 v(B_0)+\sum_{a\in\mathcal A}x_av(a)&=v(Y)
       &&\text{at every prime}.
 \label{eq:assembly-valuation-certificate}
\end{align}
These are precisely the equations in
\eqref{eq:post-tangent-certificate}.  In the trivial signature row, the
integer in \eqref{eq:assembly-integer-row-certificate} is
\(q_{\rm sm}^{\rm flex}(d)\).  It is the total flexible smooth-row quota
from \eqref{eq:smooth-quota-ledger}.  The bridge-active quantity
\(q^{\rm act}(d)\) is a real mass and is not used as a rounding equation.
Thus the distinction between the integer total row and the real active
normalization survives intact into the present argument.

\subsection{Column-sparse floating rounding}
\label{subsec:floating-rounding}

The complete column family used below is
\begin{equation}
 \mathcal C:=\{C_{p,j}:p\ {\rm prime},\ j\geq1,\ p^j\leq M\},
 \qquad C_{p,j}(a):=\mathbf1_{\{p^j\mid a\}}.
 \label{eq:rounding-prime-power-column}
\end{equation}
It is finite: necessarily \(p\leq M\) and
\(1\leq j\leq\lfloor\log_2M\rfloor\).
Because \(M<3n\) for large \(n\), a legal factor belongs to at most
\begin{equation}
 \sum_pv_p(a)=\Omega(a)\le\log_2a\le d_n
 \label{eq:rounding-column-degree}
\end{equation}
of these columns.

\begin{lemma}[Floating rounding]
\label[lemma]{lem:floating-rounding}
Suppose a finite set \(\mathcal A\) is partitioned into rows
\(\mathcal A_S\), \(0\le x_a\le1\), and
\(\sum_{a\in\mathcal A_S}x_a\) is an integer in every row.  If every
coordinate belongs to at most \(d\) members of a finite family of
zero-one columns, there are
\(X_a\in\{0,1\}\) such that
\begin{align}
 \sum_{a\in\mathcal A_S}X_a
   &=\sum_{a\in\mathcal A_S}x_a
       &&\text{for every }S,
 \label{eq:rounding-row-preservation}\\
 \left|\sum_{a\in\mathcal A}(X_a-x_a)C(a)\right|
   &\le4d
       &&\text{for every column }C.
 \label{eq:rounding-column-discrepancy}
\end{align}
\end{lemma}

\begin{proof}
Freeze every coordinate as soon as it becomes integral.  At an
intermediate stage let \(F\) be the set of strictly fractional
coordinates and put \(m=|F|\).  A row meeting \(F\) contains at least two
members of \(F\): after the frozen integral coordinates in that row are
removed, the sum of the remaining coordinates is an integer, which one
strictly fractional coordinate cannot supply.  Hence at most \(m/2\) row
equations are active.

Retain every active row equation and every column having more than \(4d\)
members of \(F\).  There are at most \(dm\) coordinate--column incidences,
so fewer than \(m/4\) column equations are retained.  The retained system
is a finite real linear system and has rank less than \(3m/4<m\).  It
therefore has a nonzero null vector \(z=(z_a)_{a\in F}\).  Its complete
feasible movement interval is
\[
 I_z:=\{t\in\mathbb R:
            0\leq x_a+tz_a\leq1\ \text{for every }a\in F\}.
\]
Because every \(x_a\), \(a\in F\), is strictly between zero and one,
\(0\) is an interior point of \(I_z\).  Since \(z\ne0\), this is a
bounded closed interval with two finite endpoints.  Move to either
endpoint.  Every retained equation is preserved and at least one
additional coordinate reaches \(0\) or \(1\); it is then frozen.
Consequently the procedure stops after at most \(|\mathcal A|\)
iterations with an integral vector.

Row equations are never dropped, which proves
\eqref{eq:rounding-row-preservation}.  At the first stage when a column
is not retained, it has at most \(4d\) fractional coordinates.  Its value
was preserved at every earlier stage, and only those then-fractional
coordinates can subsequently change.  Each changes by at most one, so
its total change from the original vector is at most \(4d\).  A column
retained through the last stage is preserved exactly.  This proves
\eqref{eq:rounding-column-discrepancy}.
\end{proof}

Apply the lemma to the complete-signature rows, with the prime-power
columns \eqref{eq:rounding-prime-power-column} and \(d=d_n\).  Define the
rounded-minus-fractional valuation error
\begin{equation}
                  e_p:=\sum_{a\in\mathcal A}(X_a-x_a)v_p(a).
 \label{eq:rounding-error-sign}
\end{equation}
It is an integer.  Indeed, by
\eqref{eq:assembly-valuation-certificate},
\(\sum_ax_av_p(a)=v_p(Y)-v_p(B_0)\) is an integer, while the rounded sum
is also integral.  Since
\[
                         v_p(a)=\sum_{j\ge1}C_{p,j}(a),
\]
\eqref{eq:rounding-column-discrepancy} gives
\begin{equation}
 |e_p|\le4d_n\left\lceil\frac{\log(3n)}{\log p}\right\rceil
         =\beta_p.
 \label{eq:rounding-error-box}
\end{equation}

The complete-row equations do essential additional work.  If \(P>y\),
then \(v_P(a)=S_P\) is constant on \(\mathcal A_S\); hence
\begin{equation}
 e_P=\sum_S S_P
      \sum_{a\in\mathcal A_S}(X_a-x_a)=0
 \label{eq:rounding-no-rough-error}
\end{equation}
by \eqref{eq:rounding-row-preservation}.  Thus \(e\) is supported on
\(p\le y\) and lies in the full universal box reserved in
Section~\ref{sec:rounding-bank}.

\subsection{The guarded exactification lemma}
\label{subsec:guarded-exactification}

For clarity, we isolate the exact hypotheses used in the last discrete
step.

\begin{proposition}[Guarded integral exactification]
\label[proposition]{prop:guarded-integral-exactification}
Assume the following.
\begin{enumerate}
 \item The flexible candidates are partitioned by their complete
       signatures, all their row sums are the integers in
       \eqref{eq:assembly-integer-row-certificate}, and their fractional
       weights lie in \([0,1]\).
 \item The fixed factors, flexible candidates, both states of every bank
       path, donors, and external anchors satisfy the joint guarding and
       distinctness conditions of Section~\ref{sec:rounding-bank}.
 \item The charged base is exactly \eqref{eq:assembly-base-product}, and
       the valuation certificate
       \eqref{eq:assembly-valuation-certificate} holds after that charge.
 \item Every path satisfies the all-factor row identity
       \eqref{eq:bank-all-factor-invariance}; every base-state row token
       replaced a distinct withheld donor token before the row quotas were
       frozen; and the resulting quotas are nonnegative.
 \item The bank is the two-sided \(\beta\)-universal reserve of
       Lemma~\ref{lem:precharged-universal-bank}.
\end{enumerate}
Then there are \(X_a\in\{0,1\}\) and state indices
\(\varepsilon_g\in\{0,1\}\), \(g\in\mathcal B\), such that, on putting
\[
 G_{\rm bank}^{\rm fin}:=
       \bigsqcup_{g\in\mathcal B}G_g^{\varepsilon_g},
\]
the three sets \(G_{\rm fix}\), \(G_{\rm bank}^{\rm fin}\), and
\(\{a\in\mathcal A:X_a=1\}\) are pairwise disjoint, the flexible row
quotas are preserved, and the product of their union is exactly \(Y\).
\end{proposition}

\begin{proof}
Use Lemma~\ref{lem:floating-rounding}.  Equations
\eqref{eq:rounding-error-box}--\eqref{eq:rounding-no-rough-error} put the
integer error \(e\) in the universal box.  Apply
Lemma~\ref{lem:precharged-universal-bank} to \(-e\).  It supplies a
subcollection of whole paths to toggle.  Set \(\varepsilon_g=1\) exactly
for those paths and \(\varepsilon_g=0\) otherwise.  The operation is the
state replacement
\begin{equation}
 G_{\rm bank}^0=\bigsqcup_{g\in\mathcal B}G_g^0
 \quad\rightsquigarrow\quad
 G_{\rm bank}^{\rm fin}
       =\bigsqcup_{g\in\mathcal B}G_g^{\varepsilon_g};
\label{eq:bank-state-replacement}
\end{equation}
no bank factor is multiplied onto the already charged base.  An
untoggled path remains literally in its base state, whereas a toggled
path is replaced once, simultaneously in all of its components.

By universality and the definition of a path change,
\begin{equation}
 v(G_{\rm bank}^{\rm fin})-v(G_{\rm bank}^0)
 =\sum_{\varepsilon_g=1}
       \{v(G_g^1)-v(G_g^0)\}=-e.
\label{eq:bank-final-valuation-change}
\end{equation}
Likewise, for every complete signature \(S\),
\[
 m_S(G_{\rm bank}^{\rm fin})-m_S(G_{\rm bank}^0)
 =\sum_{\varepsilon_g=1}
       \{m_S(G_g^1)-m_S(G_g^0)\}=0
\]
by \eqref{eq:bank-all-factor-invariance}.  Thus every fixed row token
charged in the base remains one fixed row token after exactification, and
the flexible row equations preserved by floating rounding remain valid.

The rounded flexible change is \(+e\), while
\eqref{eq:bank-final-valuation-change} is \(-e\).  Starting with
\eqref{eq:assembly-valuation-certificate}, their sum leaves every prime
valuation equal to \(v(Y)\).  Finally, both possible states of every path
were jointly guarded before the choice of the toggle subplan.  Therefore
an arbitrary subplan, including the one just selected, is collision-free
internally and disjoint from \(G_{\rm fix}\), the anchors, and all
flexible candidates.  This proves both the product identity and the
asserted distinctness.
\end{proof}

No coordinatewise relative-interior hypothesis is used in this proof.
Such interiority was useful upstream in constructing the fractional point,
but floating rounding and the universal bank require only
\(0\le x_a\le1\), integer row sums, and the exact charged certificate.

\subsection{The final residual set}
\label{subsec:final-residual-set}

\begin{theorem}[Upper-bound construction]
\label[theorem]{thm:upper-bound-construction}
Fix \(c>C_0\) and put
\[
 M=2n+\left\lceil c\frac n{\log n}\right\rceil.
\]
For every sufficiently large \(n\), there is a set of distinct integers
\[
                         \mathcal S\subset(n,M]
\]
such that
\[
                         \prod_{a\in\mathcal S}a=Q(n,M).
\]
\end{theorem}

\begin{proof}
We first discharge, in order, the five hypotheses of
Proposition~\ref{prop:guarded-integral-exactification}.
\begin{enumerate}[label=\textup{(\roman*)}]
\item Proposition~\ref{prop:post-tangent-certificate} gives
      \(0\leq x_a\leq1\) and every integer complete-signature row sum.
\item The joint guarding required there is supplied by
      Lemmas~\ref{lem:bank-donor-injection} and
      \ref{lem:rough-guard-census}, together with the anchor avoidance in
      Lemma~\ref{lem:bank-anchor-modification}.  In particular, donor
      occurrences are not separately reinserted.
\item The charged base is exactly
      \eqref{eq:assembly-base-product}; the simultaneous combined charge
      is \eqref{eq:combined-precharge-divisibility}; and the valuation
      identity after that charge is
      \eqref{eq:assembly-valuation-certificate}.
\item The donor-backed row-token construction in
      Section~\ref{subsec:bank-precharge} gives nonnegative quotas
      \eqref{eq:postcharge-rough-quota}, and every complete path obeys the
      all-factor identity \eqref{eq:bank-all-factor-invariance}.
\item The required two-sided box universality is precisely
      Lemma~\ref{lem:precharged-universal-bank}.
\end{enumerate}
Thus there is no additional assumption at the final interface.

Take \(X\), the state indices \((\varepsilon_g)\), and
\(G_{\rm bank}^{\rm fin}\) from
Proposition~\ref{prop:guarded-integral-exactification}, and define
\begin{equation}
 \mathcal G_{\rm res}:=
 G_{\rm fix}\ \cup\ G_{\rm bank}^{\rm fin}\
 \cup\{a\in\mathcal A:X_a=1\}.
\label{eq:final-residual-set}
\end{equation}
The union is disjoint by that proposition.  Since
\[
 v(B_0)=v(G_{\rm fix})+v(G_{\rm bank}^0)
\]
and the rounded-minus-fractional convention
\eqref{eq:rounding-error-sign} gives
\[
 \sum_{a\in\mathcal A}X_av(a)
 =\sum_{a\in\mathcal A}x_av(a)+e,
\]
the state-replacement identity
\eqref{eq:bank-final-valuation-change} and the exact post-tangent
certificate give, at every prime,
\begin{align}
 v\left(\prod_{a\in\mathcal G_{\rm res}}a\right)
 &=v(B_0)+\sum_{a\in\mathcal A}x_av(a)+e-e \notag\\
 &=v(Y).
 \label{eq:final-residual-valuations}
\end{align}
Equality at every prime proves the integer identity
\begin{equation}
                    \prod_{a\in\mathcal G_{\rm res}}a
                    =Y=\frac{T_n}{D'}.
\label{eq:final-residual-product}
\end{equation}

The full central-anchor construction, with the guarded modifications of
Section~\ref{subsec:bank-anchor-guards}, gives a set
\(\mathcal H'\subset(n,2n]\), disjoint from
\(\mathcal G_{\rm res}\), and
\[
                         \prod_{a\in\mathcal H'}a=C_nD'
\]
by \eqref{eq:bank-modified-anchor-identity}.  Every factor in
\(\mathcal G_{\rm res}\) is legal and lies in \((n,M]\): this holds for
\(G_{\rm fix}\) and the two bank states by their construction and for
the flexible set by the definition of \(\mathcal A\).  Hence
\begin{equation}
                    \mathcal S:=\mathcal H'\cup\mathcal G_{\rm res}
 \label{eq:final-complement-subset}
\end{equation}
consists of distinct integers in \((n,M]\) and satisfies
\begin{equation}
 \prod_{a\in\mathcal S}a
  =(C_nD')\frac{T_n}{D'}
  =C_nT_n
  =Q(n,M).
 \label{eq:final-complement-product}
\end{equation}
\end{proof}

To return to the original formulation, take the complement of the set in
Theorem~\ref{thm:upper-bound-construction}.  Since the
product of all integers in \((n,M]\) is \(M!/n!\),
\begin{equation}
 \prod_{\substack{n<a\le M\\a\notin\mathcal S}}a
   =\frac{M!/n!}{M!/(n!)^2}=n!.
 \label{eq:final-factorial-product}
\end{equation}
This is a product of distinct integers greater than \(n\), all at most
\(M=2n+\lceil cn/\log n\rceil\).  Consequently
\begin{equation}
 f(n)\le2n+\left\lceil c\frac n{\log n}\right\rceil
 \qquad(c>C_0)
 \label{eq:final-upper-bound}
\end{equation}
for every sufficiently large \(n\) depending on \(c\).  Since \(c>C_0\)
was arbitrary,
\begin{equation}
 \limsup_{n\to\infty}
 \frac{(f(n)-2n)\log n}{n}\le C_0.
 \label{eq:final-limsup}
\end{equation}
Lemma~\ref{lem:thirteen-layer-lower-bound} gives the reverse liminf.
Together they prove Theorem~\ref{thm:main}.
 \appendix
\section{Numerical verification}
\label{app:numerical-verification}

The companion numerical verifier is the single human-readable file
\texttt{numerical\_verifier.py}, distributed with this paper.  It contains
the complete 211-entry rational array used in
Lemma~\ref{lem:finite-allocation-certificate} and certifies precisely the
finite claims made there.  More specifically, it verifies
\begin{itemize}
 \item positivity, the allowed cofactor range, and every exact row sum;
 \item the prime factorization of every cofactor and every finite prime-load
       capacity inequality;
 \item the finite--tail overlap inequalities for \(201<p\le401\), together
       with the exact scalar inequality used for the Nagura tail.
\end{itemize}
The program uses only Python's standard library and \texttt{Fraction}
arithmetic.  It performs no numerical search and makes no floating-point
comparison.  It therefore checks the finite certificate exactly; it does
not replace the analytic tail argument in
Lemma~\ref{lem:nagura-allocation-tail} or any later analytic step of the
proof.
 
\bibliographystyle{alpha}
\newcommand{\etalchar}[1]{$^{#1}$}
\begin{thebibliography}{ACR{\etalchar{+}}26}

\bibitem[ABT99]{ArratiaBarbourTavare1999}
Richard Arratia, A.~D. Barbour, and Simon Tavar{\'e}.
\newblock The {Poisson--Dirichlet} distribution and the scale-invariant poisson
  process.
\newblock {\em Combinatorics, Probability and Computing}, 8(5):407--416, 1999.

\bibitem[ACR{\etalchar{+}}26]{AlexeevEtAl2026}
Boris Alexeev, Evan Conway, Matthieu Rosenfeld, Andrew~V. Sutherland, Terence
  Tao, Markus Uhr, and Kevin Ventullo.
\newblock Decomposing a factorial into large factors.
\newblock arXiv preprint, \url{https://arxiv.org/abs/2503.20170}, 2026.
\newblock Version 4, 3 April 2026.

\bibitem[AS16]{AlonSpencer2016}
Noga Alon and Joel~H. Spencer.
\newblock {\em The Probabilistic Method}.
\newblock Wiley Series in Discrete Mathematics and Optimization. John Wiley \&
  Sons, Hoboken, NJ, 4 edition, 2016.

\bibitem[EGS82]{ErdosGuySelfridge1982}
Paul Erd{\H{o}}s, Richard~K. Guy, and John~L. Selfridge.
\newblock Another property of 239 and some related questions.
\newblock {\em Congressus Numerantium}, 34:243--257, 1982.

\bibitem[HT93]{HildebrandTenenbaum1993}
Adolf Hildebrand and G{\'e}rald Tenenbaum.
\newblock Integers without large prime factors.
\newblock {\em Journal de Th{\'e}orie des Nombres de Bordeaux}, 5(2):411--484,
  1993.

\bibitem[LP17]{LastPenrose2017}
G{\"u}nter Last and Mathew Penrose.
\newblock {\em Lectures on the Poisson Process}.
\newblock Institute of Mathematical Statistics Textbooks. Cambridge University
  Press, Cambridge, 2017.

\bibitem[Mau26]{Mausberg2026}
Samuel Mausberg.
\newblock A thirteen-layer lower bound for {Erd\H{o}s} problem \#390.
\newblock Note linked from the Erd\H{o}s Problems discussion thread,
  \url{https://www.erdosproblems.com/forum/thread/390}, 2026.
\newblock Dated 2 May 2026.

\bibitem[Mer74]{Mertens1874}
Franz Mertens.
\newblock Ein beitrag zur analytischen zahlentheorie.
\newblock {\em Journal f{\"u}r die reine und angewandte Mathematik}, 78:46--62,
  1874.

\bibitem[MV07]{MontgomeryVaughan2007}
Hugh~L. Montgomery and Robert~C. Vaughan.
\newblock {\em Multiplicative Number Theory I: Classical Theory}, volume~97 of
  {\em Cambridge Studies in Advanced Mathematics}.
\newblock Cambridge University Press, Cambridge, 2007.

\bibitem[Nag52]{Nagura1952}
Jitsuro Nagura.
\newblock On the interval containing at least one prime number.
\newblock {\em Proceedings of the Japan Academy}, 28(4):177--181, 1952.

\bibitem[Sai89]{Saias1989}
{\'E}ric Saias.
\newblock Sur le nombre des entiers sans grand facteur premier.
\newblock {\em Journal of Number Theory}, 32(1):78--99, 1989.

\bibitem[Tao25]{Tao390Comment}
Terence Tao.
\newblock Comment on {Erd\H{o}s} problem \#390.
\newblock Erd\H{o}s Problems discussion thread,
  \url{https://www.erdosproblems.com/forum/thread/390}, 2025.
\newblock Posted 30 August 2025.

\end{thebibliography}
 

\end{document}
